% Generated mechanically from the frozen sites-cohomology.tex target.
% Do not edit this generated file; edit the bound locale source instead.

% 好，从这里开始。
%

\setcounter{chapter}{20}
\chapter{站点上的上同调}



\phantomsection
\label{sites-cohomology-section-phantom}


\section{引言}
\label{sites-cohomology-section-introduction}

\noindent
本文讨论层上同调的若干主题。
我们考察站点上的层会发生什么，
不过许多时候只是把拓扑情形中的讨论、
构造和证明在这里重述一遍。
基本参考文献为 \cite{SGA4}、\cite{Godement} 和 \cite{Iversen}。




\section{层上同调}
\label{sites-cohomology-section-cohomology-sheaves}

\noindent
设 $\mathcal{C}$ 为站点，见《站点》定义 \ref{sites-definition-site}。
设 $\mathcal{F}$ 为 $\mathcal{C}$ 上的阿贝尔层。
我们知道 $\mathcal{C}$ 上的阿贝尔层范畴有足够多的内射对象，
见《内射对象》定理 \ref{injectives-theorem-sheaves-injectives}。
因此可以选取一个内射分解
$\mathcal{F}[0] \to \mathcal{I}^\bullet$.
对站点 $\mathcal{C}$ 的任意对象 $U$，定义
\begin{equation}
\label{sites-cohomology-equation-cohomology-object-site}
H^i(U, \mathcal{F}) = H^i(\Gamma(U, \mathcal{I}^\bullet))
\end{equation}
为{\it 阿贝尔层 $\mathcal{F}$ 在对象 $U$ 上的第 $i$ 个上同调群}。
换言之，它们是函子 $\mathcal{F} \mapsto \mathcal{F}(U)$ 的右导出函子。
函子族 $H^i(U, -)$ 构成万有 $\delta$-函子
$\textit{Ab}(\mathcal{C}) \to \textit{Ab}$.

\medskip\noindent
站点 $\mathcal{C}$ 有时没有终对象。在这种情形下，
我们把 $\mathcal{C}$ 上集合预层 $\mathcal{F}$ 的{\it 整体截面}定义为集合
\begin{equation}
\label{sites-cohomology-equation-global-sections}
\Gamma(\mathcal{C}, \mathcal{F}) =
\Mor_{\textit{PSh}(\mathcal{C})}(e, \mathcal{F})
\end{equation}
其中 $e$ 是 $\mathcal{C}$ 上预层范畴的终对象。
在这种情形下，给定 $\mathcal{C}$ 上的阿贝尔层 $\mathcal{F}$，
如下定义{\it $\mathcal{F}$ 在 $\mathcal{C}$ 上的第 $i$ 个上同调群}：
\begin{equation}
\label{sites-cohomology-equation-cohomology}
H^i(\mathcal{C}, \mathcal{F}) = H^i(\Gamma(\mathcal{C}, \mathcal{I}^\bullet))
\end{equation}
换言之，它是整体截面函子的第 $i$ 个右导出函子。
函子族 $H^i(\mathcal{C}, -)$ 构成万有 $\delta$-函子
$\textit{Ab}(\mathcal{C}) \to \textit{Ab}$.

\medskip\noindent
设 $f : \Sh(\mathcal{C}) \to \Sh(\mathcal{D})$ 为拓扑斯的态射，
见《站点》定义 \ref{sites-definition-topos}。
对上述 $\mathcal{F}[0] \to \mathcal{I}^\bullet$，定义
\begin{equation}
\label{sites-cohomology-equation-higher-direct-image}
R^if_*\mathcal{F} = H^i(f_*\mathcal{I}^\bullet)
\end{equation}
为{\it $\mathcal{F}$ 的第 $i$ 个高阶直接像}。
它们是 $f_*$ 的右导出函子。
函子族 $R^if_*$ 构成从
$\textit{Ab}(\mathcal{C}) \to \textit{Ab}(\mathcal{D})$ 的万有 $\delta$-函子。

\medskip\noindent
设 $(\mathcal{C}, \mathcal{O})$ 为环化站点，
见《站点上的模》定义 \ref{sites-modules-definition-ringed-site}。
设 $\mathcal{F}$ 为 $\mathcal{O}$-模。
我们知道 $\mathcal{O}$-模范畴有足够多的内射对象，
见《内射对象》定理 \ref{injectives-theorem-sheaves-modules-injectives}。
因此可以选取一个内射分解
$\mathcal{F}[0] \to \mathcal{I}^\bullet$.
对站点 $\mathcal{C}$ 的任意对象 $U$，定义
\begin{equation}
\label{sites-cohomology-equation-cohomology-object-site-modules}
H^i(U, \mathcal{F}) = H^i(\Gamma(U, \mathcal{I}^\bullet))
\end{equation}
为{\it $\mathcal{F}$ 在 $U$ 上的第 $i$ 个上同调群}。
函子族 $H^i(U, -)$ 构成万有 $\delta$-函子
$\textit{Mod}(\mathcal{O}) \to \text{Mod}_{\mathcal{O}(U)}$。类似地，
\begin{equation}
\label{sites-cohomology-equation-cohomology-modules}
H^i(\mathcal{C}, \mathcal{F}) = H^i(\Gamma(\mathcal{C}, \mathcal{I}^\bullet))
\end{equation}
它是{\it $\mathcal{F}$ 在 $\mathcal{C}$ 上的第 $i$ 个上同调群}。
函子族 $H^i(\mathcal{C}, -)$ 构成万有 $\delta$-函子
$\textit{Mod}(\mathcal{C}) \to \text{Mod}_{\Gamma(\mathcal{C}, \mathcal{O})}$.

\medskip\noindent
设 $f : (\Sh(\mathcal{C}), \mathcal{O}) \to (\Sh(\mathcal{D}), \mathcal{O}')$
为环化拓扑斯的态射，
见《站点上的模》定义 \ref{sites-modules-definition-ringed-topos}。
对上述 $\mathcal{F}[0] \to \mathcal{I}^\bullet$，定义
\begin{equation}
\label{sites-cohomology-equation-higher-direct-image-modules}
R^if_*\mathcal{F} = H^i(f_*\mathcal{I}^\bullet)
\end{equation}
为{\it $\mathcal{F}$ 的第 $i$ 个高阶直接像}。
它们是 $f_*$ 的右导出函子。
函子族 $R^if_*$ 构成从
$\textit{Mod}(\mathcal{O}) \to \textit{Mod}(\mathcal{O}')$ 的万有 $\delta$-函子。








\section{导出函子}
\label{sites-cohomology-section-derived-functors}

\noindent
我们简要说明一种利用分解函子处理右导出函子的方法。
设 $(\mathcal{C}, \mathcal{O})$ 为环化站点。本章记
$$
K(\mathcal{O}) = K(\textit{Mod}(\mathcal{O}))
\quad
\text{和}
\quad
D(\mathcal{O}) = D(\textit{Mod}(\mathcal{O}))
$$
；对于《导出范畴》定义 \ref{derived-definition-complexes-notation}
和定义 \ref{derived-definition-unbounded-derived-category} 中引入的三角范畴，
其有界版本也采用类似记号。
由《导出范畴》注记 \ref{derived-remark-big-abelian-category}，
存在分解函子
$$
j = j_{(\mathcal{C}, \mathcal{O})} :
K^{+}(\textit{Mod}(\mathcal{O}))
\longrightarrow
K^{+}(\mathcal{I})
$$
其中 $\mathcal{I}$ 是 $\textit{Mod}(\mathcal{O})$ 中由内射
$\mathcal{O}$-模组成的严格满加性子范畴。
对任意阿贝尔范畴 $\mathcal{B}$ 以及任意左正合函子
$F : \textit{Mod}(\mathcal{O}) \to \mathcal{B}$，
记 $RF$ 为《导出范畴》第 \ref{derived-section-right-derived-functor} 节中
用上述分解函子 $j$ 构造的右导出函子：
\begin{equation}
\label{sites-cohomology-equation-RF}
RF = F \circ j' : D^{+}(\mathcal{O}) \longrightarrow D^{+}(\mathcal{B})
\end{equation}
记号见《导出范畴》引理 \ref{derived-lemma-right-derived-functor}。
注意，视情形而定，可以把 $RF$ 看作定义在
$\textit{Mod}(\mathcal{O})$、$\text{Comp}^{+}(\textit{Mod}(\mathcal{O}))$ 或
$K^{+}(\mathcal{O})$ 上。按照《导出范畴》定义
\ref{derived-definition-higher-derived-functors}，得到第 $i$ 个右导出函子
\begin{equation}
\label{sites-cohomology-equation-RFi}
R^iF = H^i \circ RF : \textit{Mod}(\mathcal{O}) \longrightarrow \mathcal{B}
\end{equation}
于是 $R^0F = F$，且 $\{R^iF, \delta\}_{i \geq 0}$ 是万有
$\delta$-函子；见《导出范畴》引理
\ref{derived-lemma-higher-derived-functors}。

\medskip\noindent
下面给出该构造的两个特例。给定环 $R$，记
$K(R) = K(\text{Mod}_R)$ 和 $D(R) = D(\text{Mod}_R)$，
其有界版本也类似。对 $\mathcal{C}$ 的任意对象 $U$，有左正合函子
$
\Gamma(U, -) :
\textit{Mod}(\mathcal{O})
\longrightarrow
\text{Mod}_{\mathcal{O}(U)}
$
由上述讨论，它给出
$$
R\Gamma(U, -) :
D^{+}(\mathcal{O})
\longrightarrow
D^{+}(\mathcal{O}(U))
$$
注意，$H^i(U, -) = R^i\Gamma(U, -)$ 与上面的
(\ref{sites-cohomology-equation-cohomology-object-site-modules}) 相容。类似地，有
$$
R\Gamma(\mathcal{C}, -) :
D^{+}(\mathcal{O})
\longrightarrow
D^{+}(\Gamma(\mathcal{C}, \mathcal{O}))
$$
它与 (\ref{sites-cohomology-equation-cohomology-modules}) 相容。若
$f : (\Sh(\mathcal{C}), \mathcal{O}) \to (\Sh(\mathcal{D}), \mathcal{O}')$
是环化拓扑斯的态射，则得到左正合函子
$f_* : \textit{Mod}(\mathcal{O}) \to \textit{Mod}(\mathcal{O}')$
，它给出{\it 导出推前}
$$
Rf_* : D^{+}(\mathcal{O}) \to D^+(\mathcal{O}')
$$
按照 (\ref{sites-cohomology-equation-higher-direct-image-modules})，
$Rf_*\mathcal{F}^\bullet$ 的第 $i$ 个上同调层记作
$R^if_*\mathcal{F}^\bullet$，称为第 $i$ 个{\it 高阶直接像}。
上面所示的函子都是导出范畴之间的正合函子。







\section{第一上同调与挠子}
\label{sites-cohomology-section-h1-torsors}

\begin{definition}
\label{sites-cohomology-definition-torsor}
设 $\mathcal{C}$ 为站点，$\mathcal{G}$ 为 $\mathcal{C}$ 上的
（不必交换的）群层。{\it 伪挠子}，更确切地说，
{\it 伪 $\mathcal{G}$-挠子}，是 $\mathcal{C}$ 上的集合层
$\mathcal{F}$，配备作用
$\mathcal{G} \times \mathcal{F} \to \mathcal{F}$，使得
\begin{enumerate}
\item 只要 $\mathcal{F}(U)$ 非空，作用
$\mathcal{G}(U) \times \mathcal{F}(U) \to \mathcal{F}(U)$
就是单传递的。
\end{enumerate}
{\it 伪 $\mathcal{G}$-挠子的态射}
$\mathcal{F} \to \mathcal{F}'$，是与 $\mathcal{G}$-作用相容的集合层态射。
{\it 挠子}，更确切地说，{\it $\mathcal{G}$-挠子}，
是还满足下列条件的伪 $\mathcal{G}$-挠子：
\begin{enumerate}
\item[(2)] 对每个 $U \in \Ob(\mathcal{C})$，存在 $U$ 的覆盖
$\{U_i \to U\}_{i \in I}$，使得对所有 $i \in I$，
$\mathcal{F}(U_i)$ 均非空。
\end{enumerate}
{\it $\mathcal{G}$-挠子的态射}就是伪 $\mathcal{G}$-挠子的态射。
{\it 平凡 $\mathcal{G}$-挠子}是配备显然左 $\mathcal{G}$-作用的层
$\mathcal{G}$。
\end{definition}

\noindent
显然，挠子的态射自动是同构。

\begin{lemma}
\label{sites-cohomology-lemma-trivial-torsor}
设 $\mathcal{C}$ 为站点，$\mathcal{G}$ 为 $\mathcal{C}$ 上的
（不必交换的）群层。$\mathcal{G}$-挠子 $\mathcal{F}$ 平凡，
当且仅当
$\Gamma(\mathcal{C}, \mathcal{F}) \not = \emptyset$.
\end{lemma}

\begin{proof}
略。
\end{proof}

\begin{lemma}
\label{sites-cohomology-lemma-torsors-h1}
设 $\mathcal{C}$ 为站点，$\mathcal{H}$ 为 $\mathcal{C}$ 上的阿贝尔层。
$\mathcal{H}$-挠子的同构类集合与 $H^1(\mathcal{C}, \mathcal{H})$
之间存在典范双射。
\end{lemma}

\begin{proof}
设 $\mathcal{F}$ 为 $\mathcal{H}$-挠子。考虑 $\mathcal{F}$ 上的自由阿贝尔层
$\mathbf{Z}[\mathcal{F}]$。它是如下规则的层化：对
$U \in \Ob(\mathcal{C})$，赋予所有有限形式和
$\sum n_i[s_i]$ 的集合，其中 $n_i \in \mathbf{Z}$ 且
$s_i \in \mathcal{F}(U)$。有自然映射
$$
\sigma : \mathbf{Z}[\mathcal{F}] \longrightarrow \underline{\mathbf{Z}}
$$
，它把局部截面 $\sum n_i[s_i]$ 映到 $\sum n_i$。
$\sigma$ 的核由形如 $[s] - [s']$ 的截面生成。有典范映射
$a : \Ker(\sigma) \to \mathcal{H}$
，它把 $[s] - [s']$ 映到 $h$，其中 $h$ 是满足
$h \cdot s' = s$ 的 $\mathcal{H}$ 的局部截面。考虑推出图
$$
\xymatrix{
0 \ar[r] &
\Ker(\sigma) \ar[r] \ar[d]^a &
\mathbf{Z}[\mathcal{F}] \ar[r] \ar[d] &
\underline{\mathbf{Z}} \ar[r] \ar[d] &
0 \\
0 \ar[r] &
\mathcal{H} \ar[r] &
\mathcal{E} \ar[r] &
\underline{\mathbf{Z}} \ar[r] &
0
}
$$
这里 $\mathcal{E}$ 是由推出得到的扩张。对下方短正合列所对应的
上同调长正合列，把边界算子作用于
$1 \in H^0(\mathcal{C}, \underline{\mathbf{Z}})$，得到元素
$\xi = \xi_\mathcal{F} \in H^1(\mathcal{C}, \mathcal{H})$
。

\medskip\noindent
反过来，给定 $\xi \in H^1(\mathcal{C}, \mathcal{H})$，
可如下为 $\xi$ 配置一个挠子。选取 $\mathcal{H}$ 到内射阿贝尔层
$\mathcal{I}$ 中的嵌入 $\mathcal{H} \to \mathcal{I}$。
令 $\mathcal{Q} = \mathcal{I}/\mathcal{H}$，于是有短正合列
$$
\xymatrix{
0 \ar[r] &
\mathcal{H} \ar[r] &
\mathcal{I} \ar[r] &
\mathcal{Q} \ar[r] &
0
}
$$
元素 $\xi$ 是某个整体截面
$q \in H^0(\mathcal{C}, \mathcal{Q})$
的像，因为 $H^1(\mathcal{C}, \mathcal{I}) = 0$（见《导出范畴》引理
\ref{derived-lemma-higher-derived-functors}）。
令 $\mathcal{F} \subset \mathcal{I}$ 为由那些在层 $\mathcal{Q}$ 中
映到 $q$ 的截面组成的（集合）子层。容易验证，
$\mathcal{F}$ 是 $\mathcal{H}$-挠子。

\medskip\noindent
上述两个构造互逆的验证留给读者。
\end{proof}






\section{第一上同调与扩张}
\label{sites-cohomology-section-h1-extensions}

\begin{lemma}
\label{sites-cohomology-lemma-h1-extensions}
设 $(\mathcal{C}, \mathcal{O})$ 为环化站点，
$\mathcal{F}$ 为 $\mathcal{C}$ 上的 $\mathcal{O}$-模层。存在典范双射
$$
\Ext^1_{\textit{Mod}(\mathcal{O})}(\mathcal{O}, \mathcal{F})
\longrightarrow
H^1(\mathcal{C}, \mathcal{F})
$$
，它把扩张
$$
0 \to \mathcal{F} \to \mathcal{E} \to \mathcal{O} \to 0
$$
映到 $1 \in \Gamma(\mathcal{C}, \mathcal{O})$ 在
$H^1(\mathcal{C}, \mathcal{F})$ 中的像。
\end{lemma}

\begin{proof}
我们构造引理中映射的逆。设 $\xi \in H^1(\mathcal{C}, \mathcal{F})$。
选取单射 $\mathcal{F} \subset \mathcal{I}$，其中 $\mathcal{I}$ 在
$\textit{Mod}(\mathcal{O})$ 中内射。令
$\mathcal{Q} = \mathcal{I}/\mathcal{F}$。由上同调长正合列可知，
$\xi$ 是某个截面
$\tilde \xi \in \Gamma(\mathcal{C}, \mathcal{Q}) =
\Hom_\mathcal{O}(\mathcal{O}, \mathcal{Q})$.
的像。现在只需作拉回
$$
\xymatrix{
0 \ar[r] &
\mathcal{F} \ar[r] \ar@{=}[d] &
\mathcal{E} \ar[r] \ar[d] &
\mathcal{O} \ar[r] \ar[d]^{\tilde \xi} &
0 \\
0 \ar[r] &
\mathcal{F} \ar[r] &
\mathcal{I} \ar[r] &
\mathcal{Q} \ar[r] &
0
}
$$
；见《同调代数》第 \ref{homology-section-extensions} 节。
\end{proof}

\noindent
下面的引理将由更一般的引理
\ref{sites-cohomology-lemma-cohomology-modules-abelian-agree} 取代。

\begin{lemma}
\label{sites-cohomology-lemma-h1-mod-ab-agree}
设 $(\mathcal{C}, \mathcal{O})$ 为环化站点，
$\mathcal{F}$ 为 $\mathcal{C}$ 上的 $\mathcal{O}$-模层。
以 $\mathcal{F}_{ab}$ 表示其底阿贝尔群层。则有函子性同构
$$
H^1(\mathcal{C}, \mathcal{F}_{ab})
=
H^1(\mathcal{C}, \mathcal{F})
$$
，其中左端是在 $\textit{Ab}(\mathcal{C})$ 中计算的上同调，
右端是在 $\textit{Mod}(\mathcal{O})$ 中计算的上同调。
\end{lemma}

\begin{proof}
以 $\underline{\mathbf{Z}}$ 表示常值层 $\mathbf{Z}$。由于
$\textit{Ab}(\mathcal{C}) = \textit{Mod}(\underline{\mathbf{Z}})$
，可以两次应用引理 \ref{sites-cohomology-lemma-h1-extensions}，
因而只需证明
$$
\Ext^1_{\textit{Mod}(\mathcal{O})}(\mathcal{O}, \mathcal{F})
=
\Ext^1_{\textit{Mod}(\underline{\mathbf{Z}})}(
\underline{\mathbf{Z}}, \mathcal{F}_{ab}).
$$
设 $0 \to \mathcal{F} \to \mathcal{E} \to \mathcal{O} \to 0$
是 $\textit{Mod}(\mathcal{O})$ 中的扩张。利用显然的阿贝尔层态射
$1 : \underline{\mathbf{Z}} \to \mathcal{O}$
作拉回，得到如下扩张 $\mathcal{E}_{ab}$：
$$
\xymatrix{
0 \ar[r] &
\mathcal{F}_{ab} \ar[r] \ar@{=}[d] &
\mathcal{E}_{ab} \ar[r] \ar[d] &
\underline{\mathbf{Z}} \ar[r] \ar[d]^{1} &
0 \\
0 \ar[r] &
\mathcal{F} \ar[r] &
\mathcal{E} \ar[r] &
\mathcal{O} \ar[r] &
0
}
$$
反向构造稍微有趣一些。设
$0 \to \mathcal{F}_{ab} \to \mathcal{E}_{ab} \to \underline{\mathbf{Z}} \to 0$
是 $\textit{Mod}(\underline{\mathbf{Z}})$ 中的扩张。
由于 $\underline{\mathbf{Z}}$ 是平坦的 $\underline{\mathbf{Z}}$-模，
可知序列
$$
0 \to \mathcal{F}_{ab} \otimes_{\underline{\mathbf{Z}}} \mathcal{O}
\to \mathcal{E}_{ab} \otimes_{\underline{\mathbf{Z}}} \mathcal{O}
\to \underline{\mathbf{Z}} \otimes_{\underline{\mathbf{Z}}} \mathcal{O}
\to 0
$$
正合；见《站点上的模》引理 \ref{sites-modules-lemma-flat-tor-zero}。
当然，
$\underline{\mathbf{Z}} \otimes_{\underline{\mathbf{Z}}} \mathcal{O}
= \mathcal{O}$.
因此可以沿（$\mathcal{O}$-线性的）乘法映射
$\mu : \mathcal{F} \otimes_{\underline{\mathbf{Z}}} \mathcal{O}
\to \mathcal{F}$ 作推出，得到如下 $\mathcal{O}$ 由
$\mathcal{F}$ 的扩张：
$$
\xymatrix{
0 \ar[r] &
\mathcal{F}_{ab} \otimes_{\underline{\mathbf{Z}}} \mathcal{O}
\ar[r] \ar[d]^\mu &
\mathcal{E}_{ab} \otimes_{\underline{\mathbf{Z}}} \mathcal{O}
\ar[r] \ar[d] &
\mathcal{O} \ar[r] \ar@{=}[d] &
0 \\
0 \ar[r] &
\mathcal{F} \ar[r] &
\mathcal{E} \ar[r] &
\mathcal{O} \ar[r] &
0
}
$$
这就是所需扩张。这两个构造互逆的验证从略。
\end{proof}





\section{第一上同调与可逆层}
\label{sites-cohomology-section-invertible-sheaves}

\noindent
环化站点的 Picard 群定义于《站点上的模》第
\ref{sites-modules-section-invertible} 节。

\begin{lemma}
\label{sites-cohomology-lemma-h1-invertible}
设 $(\mathcal{C}, \mathcal{O})$ 为局部环化站点。存在阿贝尔群的典范同构
$$
H^1(\mathcal{C}, \mathcal{O}^*) = \Pic(\mathcal{O}).
$$
。
\end{lemma}

\begin{proof}
设 $\mathcal{L}$ 为可逆 $\mathcal{O}$-模。考虑由下列规则定义的预层
$\mathcal{L}^*$：
$$
U \longmapsto \{s \in \mathcal{L}(U)
\text{使得 } \mathcal{O}_U \xrightarrow{s \cdot -} \mathcal{L}_U
\text{ 为同构}\}
$$
这个预层满足层条件。此外，若 $f \in \mathcal{O}^*(U)$ 且
$s \in \mathcal{L}^*(U)$，则显然 $fs \in \mathcal{L}^*(U)$。
同理，若 $s, s' \in \mathcal{L}^*(U)$，则存在唯一的
$f \in \mathcal{O}^*(U)$ 使得 $fs = s'$。再者，由《站点上的模》引理
\ref{sites-modules-lemma-invertible-is-locally-free-rank-1}，
层 $\mathcal{L}^*$ 局部有截面。换言之，
$\mathcal{L}^*$ 是 $\mathcal{O}^*$-挠子。因此得到映射
$$
\begin{matrix}
\text{在 }(\mathcal{C}, \mathcal{O})\text{ 上的可逆层} \\
\text{的同构类集合}
\end{matrix}
\longrightarrow
\begin{matrix}
\mathcal{O}^*\text{-挠子} \\
\text{的同构类集合}
\end{matrix}
$$
略去验证这是阿贝尔群同态。由引理 \ref{sites-cohomology-lemma-torsors-h1}，
右端与 $H^1(\mathcal{C}, \mathcal{O}^*)$ 典范双射。
因此只需证明这个映射既单又满。

\medskip\noindent
单射性。若挠子 $\mathcal{L}^*$ 平凡，则由引理
\ref{sites-cohomology-lemma-trivial-torsor}，$\mathcal{L}^*$ 有整体截面。
这恰好意味着 $\mathcal{L} \cong \mathcal{O}$ 是
$\Pic(\mathcal{O})$ 中的单位元。

\medskip\noindent
满射性。设 $\mathcal{F}$ 为 $\mathcal{O}^*$-挠子。考虑集合预层
$$
\mathcal{L}_1 : U \longmapsto
(\mathcal{F}(U) \times \mathcal{O}(U))/\mathcal{O}^*(U)
$$
其中 $f \in \mathcal{O}^*(U)$ 在 $(s, g)$ 上的作用为
$(fs, f^{-1}g)$。规定
$(s, g) + (s', g') = (s, g + (s'/s)g')$，其中 $s'/s$ 是满足
$fs = s'$ 的 $\mathcal{O}^*$ 的局部截面 $f$；并对
$\mathcal{O}$ 的局部截面 $h$ 规定 $h(s, g) = (s, hg)$，
则 $\mathcal{L}_1$ 成为 $\mathcal{O}$-模预层。
其层化 $\mathcal{L} = \mathcal{L}_1^\#$ 是可逆
$\mathcal{O}$-模，且相关的 $\mathcal{O}^*$-挠子
$\mathcal{L}^*$ 与 $\mathcal{F}$ 同构；验证从略。
\end{proof}









\section{上同调的局部性}
\label{sites-cohomology-section-locality}

\noindent
下面的引理说明，定义层 $\mathcal{F}$ 在站点的一个对象上的上同调时没有歧义。

\begin{lemma}
\label{sites-cohomology-lemma-cohomology-of-open}
设 $(\mathcal{C}, \mathcal{O})$ 为环化站点，
$U$ 为 $\mathcal{C}$ 的对象。
\begin{enumerate}
\item 若 $\mathcal{I}$ 是内射 $\mathcal{O}$-模，
则 $\mathcal{I}|_U$ 是内射 $\mathcal{O}_U$-模。
\item 对任意 $\mathcal{O}$-模层 $\mathcal{F}$，有
$H^p(U, \mathcal{F}) = H^p(\mathcal{C}/U, \mathcal{F}|_U)$.
\end{enumerate}
\end{lemma}

\begin{proof}
回忆，限制到 $U$ 的函子 $j_U^{-1}$ 是零延拓函子 $j_{U!}$ 的右伴随，
见《站点上的模》第 \ref{sites-modules-section-localize} 节。
此外，$j_{U!}$ 正合。因此 (1) 由《同调代数》引理
\ref{homology-lemma-adjoint-preserve-injectives} 得到。

\medskip\noindent
按照定义，$H^p(U, \mathcal{F}) = H^p(\mathcal{I}^\bullet(U))$，
其中 $\mathcal{F} \to \mathcal{I}^\bullet$ 是
$\textit{Mod}(\mathcal{O})$ 中的内射分解。
由上可知，$\mathcal{F}|_U \to \mathcal{I}^\bullet|_U$
是 $\textit{Mod}(\mathcal{O}_U)$ 中的内射分解。
故 $H^p(U, \mathcal{F}|_U)$ 等于
$H^p(\mathcal{I}^\bullet|_U(U))$。
当然，对 $\mathcal{C}$ 上的任意层 $\mathcal{F}$，都有
$\mathcal{F}(U) = \mathcal{F}|_U(U)$。这就给出 (2) 中的等式。
\end{proof}

\noindent
下面的引理将用于考察改变一个部分宇宙时会发生什么，
也用于比较小 étale 站点与大 étale 站点的上同调。

\begin{lemma}
\label{sites-cohomology-lemma-cohomology-bigger-site}
设 $\mathcal{C}$ 和 $\mathcal{D}$ 为站点，
$u : \mathcal{C} \to \mathcal{D}$ 为函子。
假设 $u$ 满足《站点》引理 \ref{sites-lemma-bigger-site} 的假设。
设 $g : \Sh(\mathcal{C}) \to \Sh(\mathcal{D})$
为相关的拓扑斯态射。对 $\mathcal{D}$ 上的任意阿贝尔层
$\mathcal{F}$，有同构
$$
R\Gamma(\mathcal{C}, g^{-1}\mathcal{F}) = R\Gamma(\mathcal{D}, \mathcal{F}),
$$
特别地，
$H^p(\mathcal{C}, g^{-1}\mathcal{F}) = H^p(\mathcal{D}, \mathcal{F})$
；并且对任意 $U \in \Ob(\mathcal{C})$，有同构
$$
R\Gamma(U, g^{-1}\mathcal{F}) = R\Gamma(u(U), \mathcal{F}),
$$
特别地，
$H^p(U, g^{-1}\mathcal{F}) = H^p(u(U), \mathcal{F})$。
所有这些同构关于 $\mathcal{F}$ 都是函子性的。
\end{lemma}

\begin{proof}
由假设 (e)，显然
$\Gamma(\mathcal{C}, g^{-1}\mathcal{F}) = \Gamma(\mathcal{D}, \mathcal{F})$
，所以只需证明 $g^{-1}$ 把内射阿贝尔层变为内射阿贝尔层。
照例应用《同调代数》引理
\ref{homology-lemma-adjoint-preserve-injectives}。
采用《站点》引理 \ref{sites-lemma-bigger-site} 的记号，
$g^{-1}$ 的左伴随是 $g_! = f^{-1}$，而它是正合函子。
因此该引理确实适用。
\end{proof}

\noindent
设 $(\mathcal{C}, \mathcal{O})$ 为环化站点，
$\mathcal{F}$ 为 $\mathcal{O}$-模层。
设 $\varphi : U \to V$ 为 $\mathcal{O}$ 的态射。
则有典范的{\it 限制映射}
\begin{equation}
\label{sites-cohomology-equation-restriction-mapping}
H^n(V, \mathcal{F})
\longrightarrow
H^n(U, \mathcal{F}), \quad
\xi \longmapsto \xi|_U
\end{equation}
，它关于 $\mathcal{F}$ 是函子性的。事实上，任取内射分解
$\mathcal{F} \to \mathcal{I}^\bullet$。层 $\mathcal{I}^p$ 的限制映射
给出复形态射
$$
\Gamma(V, \mathcal{I}^\bullet)
\longrightarrow
\Gamma(U, \mathcal{I}^\bullet)
$$
左端是表示 $R\Gamma(V, \mathcal{F})$ 的复形，右端是表示
$R\Gamma(U, \mathcal{F})$ 的复形。应用函子 $H^n$，得到上同调群上的映射。
如上所示，我们用记号 $\xi \mapsto \xi|_U$ 表示这个映射。
因此规则 $U \mapsto H^n(U, \mathcal{F})$ 是 $\mathcal{O}$-模预层。
这个预层通常记作 $\underline{H}^n(\mathcal{F})$。
引理 \ref{sites-cohomology-lemma-include} 将给出它的另一种解释。

\medskip\noindent
下面的引理说明，转到一个覆盖后可以消去高阶上同调类。

\begin{lemma}
\label{sites-cohomology-lemma-kill-cohomology-class-on-covering}
设 $(\mathcal{C}, \mathcal{O})$ 为环化站点，
$\mathcal{F}$ 为 $\mathcal{O}$-模层，$U$ 为 $\mathcal{C}$ 的对象。
设 $n > 0$ 且 $\xi \in H^n(U, \mathcal{F})$。
则存在 $\mathcal{C}$ 的覆盖 $\{U_i \to U\}$，
使得对所有 $i \in I$ 都有 $\xi|_{U_i} = 0$。
\end{lemma}

\begin{proof}
设 $\mathcal{F} \to \mathcal{I}^\bullet$ 为内射分解。则
$$
H^n(U, \mathcal{F}) =
\frac{\Ker(\mathcal{I}^n(U) \to \mathcal{I}^{n + 1}(U))}
{\Im(\mathcal{I}^{n - 1}(U) \to \mathcal{I}^n(U))}.
$$
选取表示上述商中上同调类的元素
$\tilde \xi \in \mathcal{I}^n(U)$。由于 $\mathcal{I}^\bullet$
是 $\mathcal{F}$ 的内射分解且 $n > 0$，复形
$\mathcal{I}^\bullet$ 在次数 $n$ 正合。因此
$\Im(\mathcal{I}^{n - 1} \to \mathcal{I}^n) =
\Ker(\mathcal{I}^n \to \mathcal{I}^{n + 1})$ 作为层成立。
由于 $\tilde \xi$ 是核层在 $U$ 上的截面，
故存在站点的覆盖 $\{U_i \to U\}$，使得
$\tilde \xi|_{U_i}$ 是某个截面
$\xi_i \in \mathcal{I}^{n - 1}(U_i)$ 在 $d$ 下的像。
按照限制 $\xi|_{U_i}$ 的定义，它对应于
$\tilde \xi|_{U_i}$ 的类，结论随即得到。
\end{proof}

\begin{lemma}
\label{sites-cohomology-lemma-higher-direct-images}
设 $f : (\mathcal{C}, \mathcal{O}_\mathcal{C}) \to
(\mathcal{D}, \mathcal{O}_\mathcal{D})$ 为环化站点的态射，
对应于连续函子 $u : \mathcal{D} \to \mathcal{C}$。
对任意 $\mathcal{F} \in \Ob(\textit{Mod}(\mathcal{O}_\mathcal{C}))$，
层 $R^if_*\mathcal{F}$ 是预层
$$
V \longmapsto H^i(u(V), \mathcal{F})
$$
的相关层。
\end{lemma}

\begin{proof}
设 $\mathcal{F} \to \mathcal{I}^\bullet$ 为内射分解。
按照定义，$R^if_*\mathcal{F}$ 是复形
$$
f_*\mathcal{I}^0 \to f_*\mathcal{I}^1 \to f_*\mathcal{I}^2 \to \ldots
$$
的第 $i$ 个上同调层。按照 $\mathcal{O}_\mathcal{D}$-模上的
阿贝尔范畴结构之定义，这个上同调层是预层
$$
V
\longmapsto
\frac{\Ker(f_*\mathcal{I}^i(V) \to f_*\mathcal{I}^{i + 1}(V))}
{\Im(f_*\mathcal{I}^{i - 1}(V) \to f_*\mathcal{I}^i(V))}
$$
的相关层，而它显然等于
$$
\frac{\Ker(\mathcal{I}^i(u(V)) \to \mathcal{I}^{i + 1}(u(V)))}
{\Im(\mathcal{I}^{i - 1}(u(V)) \to \mathcal{I}^i(u(V)))}
$$
，后者等于 $H^i(u(V), \mathcal{F})$，证毕。
\end{proof}






\section{{\v C}ech 复形与 {\v C}ech 上同调}
\label{sites-cohomology-section-cech}

\noindent
设 $\mathcal{C}$ 为范畴，$\mathcal{U} = \{U_i \to U\}_{i \in I}$
为具有固定靶的态射族，见《站点》定义
\ref{sites-definition-family-morphisms-fixed-target}。
假设所有纤维积 $U_{i_0} \times_U \ldots \times_U U_{i_p}$
都存在于 $\mathcal{C}$ 中。设 $\mathcal{F}$ 为 $\mathcal{C}$ 上的阿贝尔预层。令
$$
\check{\mathcal{C}}^p(\mathcal{U}, \mathcal{F})
=
\prod\nolimits_{(i_0, \ldots, i_p) \in I^{p + 1}}
\mathcal{F}(U_{i_0} \times_U \ldots \times_U U_{i_p}).
$$
这是阿贝尔群。对
$s \in \check{\mathcal{C}}^p(\mathcal{U}, \mathcal{F})$，
以 $s_{i_0\ldots i_p}$ 表示它在因子
$\mathcal{F}(U_{i_0} \times_U \ldots \times_U U_{i_p})$ 中的值。定义
$$
d : \check{\mathcal{C}}^p(\mathcal{U}, \mathcal{F})
\longrightarrow
\check{\mathcal{C}}^{p + 1}(\mathcal{U}, \mathcal{F})
$$
如下：
\begin{equation}
\label{sites-cohomology-equation-d-cech}
d(s)_{i_0\ldots i_{p + 1}} =
\sum\nolimits_{j = 0}^{p + 1}
(-1)^j s_{i_0\ldots \hat i_j \ldots i_{p + 1}}
|_{U_{i_0} \times_U \ldots \times_U U_{i_{p + 1}}}
\end{equation}
其中限制取自投影映射
$$
U_{i_0} \times_U \ldots \times_U U_{i_{p + 1}} \longrightarrow
U_{i_0} \times_U \ldots \times_U \widehat{U_{i_j}} \times_U
\ldots \times_U U_{i_{p + 1}}.
$$
不难看出 $d \circ d = 0$。换言之，
$\check{\mathcal{C}}^\bullet(\mathcal{U}, \mathcal{F})$ 是复形。

\begin{definition}
\label{sites-cohomology-definition-cech-complex}
设 $\mathcal{C}$ 为范畴，$\mathcal{U} = \{U_i \to U\}_{i \in I}$
为具有固定靶的态射族，使得所有纤维积
$U_{i_0} \times_U \ldots \times_U U_{i_p}$ 都存在于 $\mathcal{C}$ 中。
设 $\mathcal{F}$ 为 $\mathcal{C}$ 上的阿贝尔预层。
复形 $\check{\mathcal{C}}^\bullet(\mathcal{U}, \mathcal{F})$
称为与 $\mathcal{F}$ 和族 $\mathcal{U}$ 相关的{\it {\v C}ech 复形}。
它的上同调群
$H^i(\check{\mathcal{C}}^\bullet(\mathcal{U}, \mathcal{F}))$
称为 $\mathcal{F}$ 关于 $\mathcal{U}$ 的{\it {\v C}ech 上同调群}，
记作 $\check H^i(\mathcal{U}, \mathcal{F})$。
\end{definition}

\noindent
注意，站点 $\mathcal{C}$ 的任意覆盖 $\{U_i \to U\}$
都是这个定义适用的具有固定靶的态射族。

\begin{lemma}
\label{sites-cohomology-lemma-cech-h0}
设 $\mathcal{C}$ 为站点，$\mathcal{F}$ 为 $\mathcal{C}$ 上的阿贝尔预层。
下列条件等价：
\begin{enumerate}
\item $\mathcal{F}$ 是 $\mathcal{C}$ 上的阿贝尔层；
\item 对站点 $\mathcal{C}$ 的每个覆盖
$\mathcal{U} = \{U_i \to U\}_{i \in I}$，自然映射
$$
\mathcal{F}(U) \to \check{H}^0(\mathcal{U}, \mathcal{F})
$$
（见《站点》第 \ref{sites-section-sheafification} 节）是双射。
\end{enumerate}
\end{lemma}

\begin{proof}
这是因为层条件恰好是：对 $\mathcal{C}$ 的每个覆盖，
$\mathcal{F}(U) \to \check{H}^0(\mathcal{U}, \mathcal{F})$
都是双射。
\end{proof}

\noindent
设 $\mathcal{C}$ 为范畴，$\mathcal{U} = \{U_i \to U\}_{i\in I}$
为 $\mathcal{C}$ 中具有固定靶的态射族，并假设所有纤维积
$U_{i_0} \times_U \ldots \times_U U_{i_p}$ 都存在于 $\mathcal{C}$ 中。
设 $\mathcal{V} = \{V_j \to V\}_{j\in J}$ 为另一个这样的族。
设 $f : U \to V$、$\alpha : I \to J$ 和
$f_i : U_i \to V_{\alpha(i)}$ 构成具有固定靶的态射族之间的态射，
见《站点》第 \ref{sites-section-refinements} 节。
在这种情形下，得到 {\v C}ech 复形的映射
\begin{equation}
\label{sites-cohomology-equation-map-cech-complexes}
\varphi : \check{\mathcal{C}}^\bullet(\mathcal{V}, \mathcal{F})
\longrightarrow
\check{\mathcal{C}}^\bullet(\mathcal{U}, \mathcal{F})
\end{equation}
，它在次数 $p$ 由下式给出：
$$
\varphi(s)_{i_0 \ldots i_p} =
(f_{i_0} \times \ldots \times f_{i_p})^*s_{\alpha(i_0) \ldots \alpha(i_p)}
$$


\section{{\v C}ech 上同调作为预层上的函子}
\label{sites-cohomology-section-cech-functor}

\noindent
警告：本节只讨论范畴上的阿贝尔预层。
在层和层范畴的语境中，以下结果完全不成立！

\medskip\noindent
设 $\mathcal{C}$ 为范畴，$\mathcal{U} = \{U_i \to U\}_{i \in I}$
为具有固定靶的态射族，使得所有纤维积
$U_{i_0} \times_U \ldots \times_U U_{i_p}$ 都存在于 $\mathcal{C}$ 中。
设 $\mathcal{F}$ 为 $\mathcal{C}$ 上的阿贝尔预层。构造
$$
\mathcal{F} \longmapsto \check{\mathcal{C}}^\bullet(\mathcal{U}, \mathcal{F})
$$
关于 $\mathcal{F}$ 是函子性的。事实上，它是函子
\begin{equation}
\label{sites-cohomology-equation-cech-functor}
\check{\mathcal{C}}^\bullet(\mathcal{U}, -) :
\textit{PAb}(\mathcal{C})
\longrightarrow
\text{Comp}^{+}(\textit{Ab})
\end{equation}
。记号见《导出范畴》定义 \ref{derived-definition-complexes-notation}。
回忆，阿贝尔范畴中下有界复形所成的范畴仍是阿贝尔范畴，
见《同调代数》引理 \ref{homology-lemma-cat-cochain-abelian}。

\begin{lemma}
\label{sites-cohomology-lemma-cech-exact-presheaves}
方程 (\ref{sites-cohomology-equation-cech-functor}) 给出的函子是正合函子
（见《同调代数》引理 \ref{homology-lemma-exact-functor}）。
\end{lemma}

\begin{proof}
对 $\mathcal{C}$ 的任意对象 $W$，函子
$\mathcal{F} \mapsto \mathcal{F}(W)$ 是从
$\textit{PAb}(\mathcal{C})$ 到 $\textit{Ab}$ 的加性正合函子。
复形的各项 $\check{\mathcal{C}}^p(\mathcal{U}, \mathcal{F})$
是这些正合函子的乘积，因而也是正合的。
此外，复形列正合当且仅当它在每个给定次数上的项所成之列正合。
引理得证。
\end{proof}

\begin{lemma}
\label{sites-cohomology-lemma-cech-cohomology-delta-functor-presheaves}
设 $\mathcal{C}$ 为范畴，
$\mathcal{U} = \{U_i \to U\}_{i \in I}$ 为具有固定靶的态射族，
使得所有纤维积 $U_{i_0} \times_U \ldots \times_U U_{i_p}$
都存在于 $\mathcal{C}$ 中。函子
$\mathcal{F} \mapsto \check{H}^n(\mathcal{U}, \mathcal{F})$
构成从阿贝尔范畴 $\textit{PAb}(\mathcal{C})$ 到
$\mathbf{Z}$-模范畴的 $\delta$-函子
（见《同调代数》定义 \ref{homology-definition-cohomological-delta-functor}）。
\end{lemma}

\begin{proof}
由引理 \ref{sites-cohomology-lemma-cech-exact-presheaves}，阿贝尔预层的短正合列
$0 \to \mathcal{F}_1 \to \mathcal{F}_2 \to \mathcal{F}_3 \to 0$
被变为 $\mathbf{Z}$-模复形的短正合列。因此可应用《同调代数》引理
\ref{homology-lemma-long-exact-sequence-cochain}，得到边界映射
$\delta_{\mathcal{F}_1 \to \mathcal{F}_2 \to \mathcal{F}_3} :
\check{H}^n(\mathcal{U}, \mathcal{F}_3) \to
\check{H}^{n + 1}(\mathcal{U}, \mathcal{F}_1)$
以及相应的长正合列。这些映射与预层短正合列之间的映射相容，
其验证从略。
\end{proof}

\begin{lemma}
\label{sites-cohomology-lemma-cech-map-into}
设 $\mathcal{C}$ 为范畴，$\mathcal{U} = \{U_i \to U\}_{i \in I}$
为具有固定靶的态射族，使得所有纤维积
$U_{i_0} \times_U \ldots \times_U U_{i_p}$ 都存在于 $\mathcal{C}$ 中。
考虑阿贝尔预层的链复形 $\mathbf{Z}_{\mathcal{U}, \bullet}$：
$$
\ldots
\to
\bigoplus_{i_0i_1i_2} \mathbf{Z}_{U_{i_0} \times_U U_{i_1} \times_U U_{i_2}}
\to
\bigoplus_{i_0i_1} \mathbf{Z}_{U_{i_0} \times_U U_{i_1}}
\to
\bigoplus_{i_0} \mathbf{Z}_{U_{i_0}}
\to 0 \to \ldots
$$
其中最后一个非零项置于次数 $0$，且映射
$$
\mathbf{Z}_{U_{i_0} \times_U \ldots \times_u U_{i_{p + 1}}}
\longrightarrow
\mathbf{Z}_{U_{i_0} \times_U
\ldots \widehat{U_{i_j}} \ldots \times_U U_{i_{p + 1}}}
$$
由典范映射的 $(-1)^j$ 倍给出。则存在同构
$$
\Hom_{\textit{PAb}(\mathcal{C})}(\mathbf{Z}_{\mathcal{U}, \bullet}, \mathcal{F})
=
\check{\mathcal{C}}^\bullet(\mathcal{U}, \mathcal{F})
$$
，它关于 $\mathcal{F} \in \Ob(\textit{PAb}(\mathcal{C}))$ 是函子性的。
\end{lemma}

\begin{proof}
这直接来自下列事实：
\begin{align*}
\Hom_{\textit{PAb}(\mathcal{C})}(
\bigoplus_{i_0 \ldots i_p}
\mathbf{Z}_{U_{i_0} \times_U \ldots \times_U U_{i_p}},
\mathcal{F})
& =
\prod_{i_0 \ldots i_p}
\Hom_{\textit{PAb}(\mathcal{C})}(
\mathbf{Z}_{U_{i_0} \times_U \ldots \times_U U_{i_p}},
\mathcal{F}) \\
& =
\prod_{i_0 \ldots i_p}
\mathcal{F}(U_{i_0} \times_U \ldots \times_U U_{i_p})
\end{align*}
见《站点上的模》引理 \ref{sites-modules-lemma-obvious-adjointness}。
\end{proof}

\begin{lemma}
\label{sites-cohomology-lemma-homology-complex}
设 $\mathcal{C}$ 为范畴，
$\mathcal{U} = \{f_i : U_i \to U\}_{i \in I}$ 为具有固定靶的态射族，
使得所有纤维积 $U_{i_0} \times_U \ldots \times_U U_{i_p}$
都存在于 $\mathcal{C}$ 中。上面引理 \ref{sites-cohomology-lemma-cech-map-into}
中的预层链复形 $\mathbf{Z}_{\mathcal{U}, \bullet}$ 在正次数正合；
也就是说，对 $i > 0$，同调预层
$H_i(\mathbf{Z}_{\mathcal{U}, \bullet})$ 为零。
\end{lemma}

\begin{proof}
设 $V$ 为 $\mathcal{C}$ 的对象。我们需要证明阿贝尔群链复形
$\mathbf{Z}_{\mathcal{U}, \bullet}(V)$ 在次数 $> 0$ 正合。
这个复形是
$$
\xymatrix{
\ldots \ar[d] \\
\bigoplus_{i_0i_1i_2}
\mathbf{Z}[
\Mor_\mathcal{C}(V, U_{i_0} \times_U U_{i_1} \times_U U_{i_2})
]
\ar[d] \\
\bigoplus_{i_0i_1}
\mathbf{Z}[
\Mor_\mathcal{C}(V, U_{i_0} \times_U U_{i_1})
]
\ar[d] \\
\bigoplus_{i_0}
\mathbf{Z}[
\Mor_\mathcal{C}(V, U_{i_0})
] \ar[d] \\
0
}
$$
对任意态射 $\varphi : V \to U$，记
$\Mor_\varphi(V, U_i) = \{\varphi_i : V \to U_i \mid
f_i \circ \varphi_i = \varphi\}$。对
$\Mor_\varphi(V, U_{i_0} \times_U \ldots \times_U U_{i_p})$
采用类似记号。注意，与纤维积
$U_{i_0} \times_U \ldots \times_U U_{i_p}$ 之间的各个投影映射复合，
保持这些态射集合。因此上面的复形等于复形
$$
\xymatrix{
\ldots \ar[d] \\
\bigoplus_\varphi
\bigoplus_{i_0i_1i_2}
\mathbf{Z}[
\Mor_\varphi(V, U_{i_0} \times_U U_{i_1} \times_U U_{i_2})
]
\ar[d] \\
\bigoplus_\varphi
\bigoplus_{i_0i_1}
\mathbf{Z}[
\Mor_\varphi(V, U_{i_0} \times_U U_{i_1})
]
\ar[d] \\
\bigoplus_\varphi
\bigoplus_{i_0}
\mathbf{Z}[
\Mor_\varphi(V, U_{i_0})
] \ar[d] \\
0
}
$$
其次，注意有
$$
\Mor_\varphi(V, U_{i_0} \times_U \ldots \times_U U_{i_p})
=
\Mor_\varphi(V, U_{i_0}) \times \ldots
\times \Mor_\varphi(V, U_{i_p})
$$
利用这一点以及 $\mathbf{Z}[A] \oplus \mathbf{Z}[B] =
\mathbf{Z}[A \amalg B]$，可见该复形变为
$$
\xymatrix{
\ldots \ar[d] \\
\bigoplus_\varphi
\mathbf{Z}\left[
\coprod_{i_0i_1i_2}
\Mor_\varphi(V, U_{i_0}) \times \Mor_\varphi(V, U_{i_1}) \times
\Mor_\varphi(V, U_{i_2})
\right]
\ar[d] \\
\bigoplus_\varphi
\mathbf{Z}\left[
\coprod_{i_0i_1}
\Mor_\varphi(V, U_{i_0}) \times \Mor_\varphi(V, U_{i_1})
\right]
\ar[d] \\
\bigoplus_\varphi
\mathbf{Z}\left[
\coprod_{i_0}
\Mor_\varphi(V, U_{i_0})
\right] \ar[d] \\
0
}
$$
最后，令 $S_\varphi = \coprod_{i \in I} \Mor_\varphi(V, U_i)$，
得到
$$
\bigoplus\nolimits_\varphi \left(\ldots \to
\mathbf{Z}[S_\varphi \times S_\varphi \times S_\varphi] \to
\mathbf{Z}[S_\varphi \times S_\varphi] \to
\mathbf{Z}[S_\varphi] \to 0 \to \ldots
\right)
$$
这样，任务得到了简化。只需证明，对任意非空集合 $S$，
自由阿贝尔群的（增广）复形
$$
\ldots \to
\mathbf{Z}[S \times S \times S] \to
\mathbf{Z}[S \times S] \to
\mathbf{Z}[S] \xrightarrow{\Sigma} \mathbf{Z} \to 0 \to \ldots
$$
在所有次数都正合。为此固定元素 $s \in S$，并采用同伦
$$
n_{(s_0, \ldots, s_p)} \longmapsto n_{(s, s_0, \ldots, s_p)}
$$
，其中记号的含义显然。
\end{proof}

\begin{lemma}
\label{sites-cohomology-lemma-complex-tensored-still-exact}
\begin{slogan}
整数预层的 {\v C}ech 复形是整数常值预层的平坦分解。
\end{slogan}
设 $\mathcal{C}$ 为范畴，
$\mathcal{U} = \{f_i : U_i \to U\}_{i \in I}$ 为具有固定靶的态射族，
使得所有纤维积 $U_{i_0} \times_U \ldots \times_U U_{i_p}$
都存在于 $\mathcal{C}$ 中。设 $\mathcal{O}$ 为 $\mathcal{C}$ 上的环预层。
链复形
$$
\mathbf{Z}_{\mathcal{U}, \bullet}
\otimes_{p, \mathbf{Z}}
\mathcal{O}
$$
在正次数正合。这里 $\mathbf{Z}_{\mathcal{U}, \bullet}$
是引理 \ref{sites-cohomology-lemma-cech-map-into} 的链复形，张量积取在
取值为 $\mathbf{Z}$ 的常值环预层上。
\end{lemma}

\begin{proof}
设 $V$ 为 $\mathcal{C}$ 的对象。在引理
\ref{sites-cohomology-lemma-homology-complex} 的证明中，我们已经看到，
作为复形，$\mathbf{Z}_{\mathcal{U}, \bullet}(V)$ 同构于若干复形的直和，
而这些复形都同伦等价于置于次数零的 $\mathbf{Z}$。因此，
$\mathbf{Z}_{\mathcal{U}, \bullet}(V) \otimes_\mathbf{Z} \mathcal{O}(V)$
作为复形也同构于若干复形的直和，而这些复形都同伦等价于
置于次数零的 $\mathcal{O}(V)$。
也可以应用《站点上的模》引理
\ref{sites-modules-lemma-flat-resolution-of-flat}：
它适用是因为各预层 $\mathbf{Z}_{\mathcal{U}, i}$ 都平坦，
且引理 \ref{sites-cohomology-lemma-homology-complex} 的证明还说明
$H_0(\mathbf{Z}_{\mathcal{U}, \bullet})$ 是平坦预层。
\end{proof}

\begin{lemma}
\label{sites-cohomology-lemma-cech-cohomology-derived-presheaves}
设 $\mathcal{C}$ 为范畴，
$\mathcal{U} = \{f_i : U_i \to U\}_{i \in I}$ 为具有固定靶的态射族，
使得所有纤维积 $U_{i_0} \times_U \ldots \times_U U_{i_p}$
都存在于 $\mathcal{C}$ 中。作为 $\delta$-函子，{\v C}ech 上同调函子
$\check{H}^p(\mathcal{U}, -)$ 与函子
$$
\check{H}^0(\mathcal{U}, -) :
\textit{PAb}(\mathcal{C})
\longrightarrow
\textit{Ab}.
$$
的右导出函子典范同构。
此外，存在函子性的拟同构
$$
\check{\mathcal{C}}^\bullet(\mathcal{U}, \mathcal{F})
\longrightarrow
R\check{H}^0(\mathcal{U}, \mathcal{F})
$$
，其中右端表示左正合函子 $\check{H}^0(\mathcal{U}, -)$ 的导出函子
$$
R\check{H}^0(\mathcal{U}, -) :
D^{+}(\textit{PAb}(\mathcal{C}))
\longrightarrow
D^{+}(\mathbf{Z})
$$
\end{lemma}

\begin{proof}
注意，阿贝尔预层范畴有足够多的内射对象，见《内射对象》命题
\ref{injectives-proposition-presheaves-injectives}。
还要注意，$\check{H}^0(\mathcal{U}, -)$ 是从阿贝尔预层范畴到
$\mathbf{Z}$-模范畴的左正合函子。因此导出函子和右导出函子都存在，
见《导出范畴》第 \ref{derived-section-right-derived-functor} 节。

\medskip\noindent
设 $\mathcal{I}$ 为内射阿贝尔预层。在这种情形下，函子
$\Hom_{\textit{PAb}(\mathcal{C})}(-, \mathcal{I})$
在 $\textit{PAb}(\mathcal{C})$ 上正合。由引理
\ref{sites-cohomology-lemma-cech-map-into}，有
$$
\Hom_{\textit{PAb}(\mathcal{C})}(
\mathbf{Z}_{\mathcal{U}, \bullet}, \mathcal{I})
=
\check{\mathcal{C}}^\bullet(\mathcal{U}, \mathcal{I}).
$$
由引理 \ref{sites-cohomology-lemma-homology-complex}，
$\mathbf{Z}_{\mathcal{U}, \bullet}$ 在正次数正合。
再由上述向 $\mathcal{I}$ 取 Hom 的正合性可知，对所有 $i > 0$，
$\check{H}^i(\mathcal{U}, \mathcal{I}) = 0$。
因此 $\delta$-函子 $(\check{H}^n, \delta)$
（见引理 \ref{sites-cohomology-lemma-cech-cohomology-delta-functor-presheaves}）
满足《同调代数》引理 \ref{homology-lemma-efface-implies-universal}
的假设，从而是万有 $\delta$-函子。

\medskip\noindent
由《导出范畴》引理 \ref{derived-lemma-higher-derived-functors}，
序列 $R^i\check{H}^0(\mathcal{U}, -)$ 也构成万有 $\delta$-函子。
根据万有 $\delta$-函子的唯一性（见《同调代数》引理
\ref{homology-lemma-uniqueness-universal-delta-functor}），得到
$R^i\check{H}^0(\mathcal{U}, -) = \check{H}^i(\mathcal{U}, -)$.
这对大多数应用已经足够；读者可以跳过证明的剩余部分。

\medskip\noindent
设 $\mathcal{F}$ 为 $\mathcal{C}$ 上的任意阿贝尔预层。
在范畴 $\textit{PAb}(\mathcal{C})$ 中选取内射分解
$\mathcal{F} \to \mathcal{I}^\bullet$。考虑双复形
$\check{\mathcal{C}}^\bullet(\mathcal{U}, \mathcal{I}^\bullet)$
，其项为 $\check{\mathcal{C}}^p(\mathcal{U}, \mathcal{I}^q)$。
再考虑与这个双复形相关的全复形
$\text{Tot}(\check{\mathcal{C}}^\bullet(\mathcal{U}, \mathcal{I}^\bullet))$
，见《同调代数》第 \ref{homology-section-double-complexes} 节。
有复形态射
$$
\check{\mathcal{C}}^\bullet(\mathcal{U}, \mathcal{F})
\longrightarrow
\text{Tot}(\check{\mathcal{C}}^\bullet(\mathcal{U}, \mathcal{I}^\bullet))
$$
，它来自映射
$\check{\mathcal{C}}^p(\mathcal{U}, \mathcal{F})
\to \check{\mathcal{C}}^p(\mathcal{U}, \mathcal{I}^0)$
；另有复形态射
$$
\check{H}^0(\mathcal{U}, \mathcal{I}^\bullet)
\longrightarrow
\text{Tot}(\check{\mathcal{C}}^\bullet(\mathcal{U}, \mathcal{I}^\bullet))
$$
，它来自映射
$\check{H}^0(\mathcal{U}, \mathcal{I}^q) \to
\check{\mathcal{C}}^0(\mathcal{U}, \mathcal{I}^q)$.
应用《同调代数》引理
\ref{homology-lemma-double-complex-gives-resolution}，
这两个映射都是拟同构。事实上，由于作为函子的 {\v C}ech 复形正合
（引理 \ref{sites-cohomology-lemma-cech-exact-presheaves}），双复形的各列在正次数正合；
又由于我们刚刚看到内射预层 $\mathcal{I}^q$ 的高阶
{\v C}ech 上同调群为零，双复形的各行也在正次数正合。
拟同构在 $D^{+}(\mathbf{Z})$ 中可逆，因此得到引理中最后显示的态射。
该态射的函子性验证从略。
\end{proof}





\section{{\v C}ech 上同调与上同调}
\label{sites-cohomology-section-cech-cohomology-cohomology}

\noindent
上同调与 {\v C}ech 上同调之间的关系来自这样一个事实：
内射阿贝尔层的高阶 {\v C}ech 上同调为零。
为说明这一点，先注意内射阿贝尔层也是内射阿贝尔预层，
然后应用前一节关于 {\v C}ech 上同调的结果。

\begin{lemma}
\label{sites-cohomology-lemma-injective-abelian-sheaf-injective-presheaf}
设 $\mathcal{C}$ 为站点。内射阿贝尔层作为范畴
$\textit{PAb}(\mathcal{C})$ 中的对象也是内射的。
\end{lemma}

\begin{proof}
对范畴 $\mathcal{A} = \textit{Ab}(\mathcal{C})$、
$\mathcal{B} = \textit{PAb}(\mathcal{C})$，以及包含函子和层化函子，
应用《同调代数》引理 \ref{homology-lemma-adjoint-preserve-injectives}。
该引理的所有假设均满足，见《站点上的模》第
\ref{sites-modules-section-abelian-sheaves} 节。
\end{proof}

\begin{lemma}
\label{sites-cohomology-lemma-injective-trivial-cech}
设 $\mathcal{C}$ 为站点，
$\mathcal{U} = \{U_i \to U\}_{i \in I}$ 为 $\mathcal{C}$ 的覆盖。
设 $\mathcal{I}$ 为内射阿贝尔层，即
$\textit{Ab}(\mathcal{C})$ 中的内射对象。则
$$
\check{H}^p(\mathcal{U}, \mathcal{I}) =
\left\{
\begin{matrix}
\mathcal{I}(U) & \text{若} & p = 0 \\
0 & \text{若} & p > 0
\end{matrix}
\right.
$$
\end{lemma}

\begin{proof}
由引理 \ref{sites-cohomology-lemma-injective-abelian-sheaf-injective-presheaf}，
$\mathcal{I}$ 是 $\textit{PAb}(\mathcal{C})$ 中的内射对象。
因此可应用引理 \ref{sites-cohomology-lemma-cech-cohomology-derived-presheaves}
（或其证明），得到高阶 {\v C}ech 上同调群的消失。
次数零的情形见引理 \ref{sites-cohomology-lemma-cech-h0}。
\end{proof}

\begin{lemma}
\label{sites-cohomology-lemma-cech-cohomology}
设 $\mathcal{C}$ 为站点，
$\mathcal{U} = \{U_i \to U\}_{i \in I}$ 为 $\mathcal{C}$ 的覆盖。
存在函子
$$
\check{\mathcal{C}}^\bullet(\mathcal{U}, -)
\longrightarrow
R\Gamma(U, -)
$$
之间的变换，其中函子的定义域和值域为
$\textit{Ab}(\mathcal{C}) \to D^{+}(\mathbf{Z})$。
特别地，当 $\mathcal{F}$ 取遍 $\textit{Ab}(\mathcal{C})$ 时，
它给出函子变换
$\check{H}^p(\mathcal{U}, \mathcal{F}) \to H^p(U, \mathcal{F})$。
\end{lemma}

\begin{proof}
设 $\mathcal{F}$ 为阿贝尔层。选取内射分解
$\mathcal{F} \to \mathcal{I}^\bullet$。考虑双复形
$\check{\mathcal{C}}^\bullet(\mathcal{U}, \mathcal{I}^\bullet)$
，其项为 $\check{\mathcal{C}}^p(\mathcal{U}, \mathcal{I}^q)$。
再考虑相关的全复形
$\text{Tot}(\check{\mathcal{C}}^\bullet(\mathcal{U}, \mathcal{I}^\bullet))$,
见《同调代数》定义 \ref{homology-definition-associated-simple-complex}。
有复形态射
$$
\alpha :
\Gamma(U, \mathcal{I}^\bullet)
\longrightarrow
\text{Tot}(\check{\mathcal{C}}^\bullet(\mathcal{U}, \mathcal{I}^\bullet))
$$
，它来自映射
$\mathcal{I}^q(U) \to \check{H}^0(\mathcal{U}, \mathcal{I}^q)$
；另有复形态射
$$
\beta :
\check{\mathcal{C}}^\bullet(\mathcal{U}, \mathcal{F})
\longrightarrow
\text{Tot}(\check{\mathcal{C}}^\bullet(\mathcal{U}, \mathcal{I}^\bullet))
$$
，它来自映射 $\mathcal{F} \to \mathcal{I}^0$。
应用《同调代数》引理
\ref{homology-lemma-double-complex-gives-resolution}，可知
$\alpha$ 是拟同构。事实上，引理 \ref{sites-cohomology-lemma-injective-trivial-cech}
说明双复形 $\check{\mathcal{C}}^\bullet(\mathcal{U}, \mathcal{I}^\bullet)$
的第 $q$ 行是 $\Gamma(U, \mathcal{I}^q)$ 的分解。
故 $\alpha$ 在 $D^{+}(\mathbf{Z})$ 中可逆，引理中的变换是
$\beta$ 后接 $\alpha$ 的逆所得的复合。其函子性验证从略。
\end{proof}

\begin{lemma}
\label{sites-cohomology-lemma-cech-h1}
设 $\mathcal{C}$ 为站点，$\mathcal{G}$ 为 $\mathcal{C}$ 上的阿贝尔层，
$\mathcal{U} = \{U_i \to U\}_{i \in I}$ 为 $\mathcal{C}$ 的覆盖。映射
$$
\check{H}^1(\mathcal{U}, \mathcal{G})
\longrightarrow
H^1(U, \mathcal{G})
$$
为单射；并且经由引理 \ref{sites-cohomology-lemma-torsors-h1} 的双射，
它把 $\check{H}^1(\mathcal{U}, \mathcal{G})$ 等同于如下集合：
那些在每个 $U_i$ 上限制为平凡挠子的
$\mathcal{G}|_U$-挠子的同构类集合。
\end{lemma}

\begin{proof}
为此构造一个逆映射。设 $\mathcal{F}$ 为 $\mathcal{C}/U$ 上的
$\mathcal{G}|_U$-挠子，其在 $\mathcal{C}/U_i$ 上的限制平凡。
由引理 \ref{sites-cohomology-lemma-trivial-torsor}，这意味着存在截面
$s_i \in \mathcal{F}(U_i)$。在 $U_{i_0} \times_U U_{i_1}$ 上，
存在 $\mathcal{G}$ 的唯一截面 $s_{i_0i_1}$，使得
$s_{i_0i_1} \cdot s_{i_0}|_{U_{i_0} \times_U U_{i_1}} =
s_{i_1}|_{U_{i_0} \times_U U_{i_1}}$。直接计算可知，
$s_{i_0i_1}$ 是 {\v C}ech 余圈，而且它的类有良好定义
（即不依赖于截面 $s_i$ 的选择）。逆映射把 $\mathcal{F}$ 的同构类
映到余圈 $(s_{i_0i_1})$ 的上同调类。略去验证该映射确实为逆。
\end{proof}

\begin{lemma}
\label{sites-cohomology-lemma-include}
设 $\mathcal{C}$ 为站点。考虑函子
$i : \textit{Ab}(\mathcal{C}) \to \textit{PAb}(\mathcal{C})$.
它是左正合函子，其右导出函子为
$$
R^pi(\mathcal{F}) = \underline{H}^p(\mathcal{F}) :
U \longmapsto H^p(U, \mathcal{F})
$$
；见第 \ref{sites-cohomology-section-locality} 节的讨论。
\end{lemma}

\begin{proof}
显然 $i$ 左正合。选取内射分解
$\mathcal{F} \to \mathcal{I}^\bullet$。按照定义，$R^pi$ 是复形
$\mathcal{I}^\bullet$ 的第 $p$ 个上同调{\it 预层}。
换言之，$R^pi(\mathcal{F})$ 在 $\mathcal{C}$ 的对象 $U$ 上的截面为
$$
\frac{\Ker(\mathcal{I}^n(U) \to \mathcal{I}^{n + 1}(U))}
{\Im(\mathcal{I}^{n - 1}(U) \to \mathcal{I}^n(U))}.
$$
，这正是 $H^p(U, \mathcal{F})$ 的定义。
\end{proof}

\begin{lemma}
\label{sites-cohomology-lemma-cech-spectral-sequence}
设 $\mathcal{C}$ 为站点，
$\mathcal{U} = \{U_i \to U\}_{i \in I}$ 为 $\mathcal{C}$ 的覆盖。
对任意阿贝尔层 $\mathcal{F}$，存在谱序列
$(E_r, d_r)_{r \geq 0}$，满足
$$
E_2^{p, q} = \check{H}^p(\mathcal{U}, \underline{H}^q(\mathcal{F}))
$$
并收敛到 $H^{p + q}(U, \mathcal{F})$。
这个谱序列关于 $\mathcal{F}$ 是函子性的。
\end{lemma}

\begin{proof}
这是下列函子的 Grothendieck 谱序列（见《导出范畴》引理
\ref{derived-lemma-grothendieck-spectral-sequence}）：
$$
i :  \textit{Ab}(\mathcal{C}) \to \textit{PAb}(\mathcal{C})
\quad\text{和}\quad
\check{H}^0(\mathcal{U}, - ) : \textit{PAb}(\mathcal{C})
\to \textit{Ab}.
$$
由引理 \ref{sites-cohomology-lemma-cech-h0}，有
$\check{H}^0(\mathcal{U}, i(\mathcal{F})) = \mathcal{F}(U)$。
由引理 \ref{sites-cohomology-lemma-injective-trivial-cech}，$i(\mathcal{I})$
是 {\v C}ech 无环的。又由引理
\ref{sites-cohomology-lemma-cech-cohomology-derived-presheaves}，作为
$\textit{PAb}(\mathcal{C})$ 上的函子，有
$\check{H}^p(\mathcal{U}, -) = R^p\check{H}^0(\mathcal{U}, -)$。
合并这些事实即得引理。
\end{proof}

\begin{lemma}
\label{sites-cohomology-lemma-cech-spectral-sequence-application}
设 $\mathcal{C}$ 为站点，
$\mathcal{U} = \{U_i \to U\}_{i \in I}$ 为覆盖，
$\mathcal{F} \in \Ob(\textit{Ab}(\mathcal{C}))$。
假设对所有 $i > 0$、所有 $p \geq 0$ 以及所有
$i_0, \ldots, i_p \in I$，都有
$H^i(U_{i_0} \times_U \ldots \times_U U_{i_p}, \mathcal{F}) = 0$。
则 $\check{H}^p(\mathcal{U}, \mathcal{F}) = H^p(U, \mathcal{F})$。
\end{lemma}

\begin{proof}
应用引理 \ref{sites-cohomology-lemma-cech-spectral-sequence} 的谱序列。
假设意味着，只要 $q \not = 0$，就有 $E_2^{p, q} = 0$。
因此谱序列在 $E_2$ 退化，结论随即得到。
\end{proof}

\begin{lemma}
\label{sites-cohomology-lemma-ses-cech-h1}
设 $\mathcal{C}$ 为站点。设
$$
0 \to \mathcal{F} \to \mathcal{G} \to \mathcal{H} \to 0
$$
为 $\mathcal{C}$ 上阿贝尔层的短正合列。设 $U$ 为 $\mathcal{C}$ 的对象。
若存在 $U$ 的覆盖所成的共尾系，使得对其中每个覆盖 $\mathcal{U}$，
都有 $\check{H}^1(\mathcal{U}, \mathcal{F}) = 0$，
则映射 $\mathcal{G}(U) \to \mathcal{H}(U)$ 为满射。
\end{lemma}

\begin{proof}
取元素 $s \in \mathcal{H}(U)$。选取覆盖
$\mathcal{U} = \{U_i \to U\}_{i \in I}$，使得
(a) $\check{H}^1(\mathcal{U}, \mathcal{F}) = 0$，且
(b) $s|_{U_i}$ 是某个截面 $s_i \in \mathcal{G}(U_i)$ 的像。
显然可以找到满足 (b) 的覆盖；再由引理的假设，
可以找到同时满足 (a) 和 (b) 的覆盖。考虑截面
$$
s_{i_0i_1} =
s_{i_1}|_{U_{i_0} \times_U U_{i_1}} - s_{i_0}|_{U_{i_0} \times_U U_{i_1}}.
$$
由于 $s_i$ 提升 $s$，可知
$s_{i_0i_1} \in \mathcal{F}(U_{i_0} \times_U U_{i_1})$.
由 $\check{H}^1(\mathcal{U}, \mathcal{F})$ 的消失，
可以找到截面 $t_i \in \mathcal{F}(U_i)$，使得
$$
s_{i_0i_1} =
t_{i_1}|_{U_{i_0} \times_U U_{i_1}} - t_{i_0}|_{U_{i_0} \times_U U_{i_1}}.
$$
于是截面 $s_i - t_i$ 显然满足层条件，并粘合为
$\mathcal{G}$ 在 $U$ 上的一个截面，它映到 $s$。证毕。
\end{proof}

\begin{lemma}
\label{sites-cohomology-lemma-cech-vanish-collection}
（《层上同调》引理 \ref{cohomology-lemma-cech-vanish} 的变式。）
设 $\mathcal{C}$ 为站点，$\text{Cov}_\mathcal{C}$ 为
$\mathcal{C}$ 的覆盖所成的集合（见《站点》定义
\ref{sites-definition-site}）。设
$\mathcal{B} \subset \Ob(\mathcal{C})$，且
$\text{Cov} \subset \text{Cov}_\mathcal{C}$
为子集。设 $\mathcal{F}$ 为 $\mathcal{C}$ 上的阿贝尔层。假设
\begin{enumerate}
\item 对每个 $\mathcal{U} \in \text{Cov}$，若
$\mathcal{U} = \{U_i \to U\}_{i \in I}$，则
$U, U_i \in \mathcal{B}$，并且每个
$U_{i_0} \times_U \ldots \times_U U_{i_p} \in \mathcal{B}$.
\item 对每个 $U \in \mathcal{B}$，$\text{Cov}$ 中出现的
$U$ 的覆盖构成 $U$ 的覆盖所成的共尾系。
\item 对每个 $\mathcal{U} \in \text{Cov}$ 和所有 $p > 0$，都有
$\check{H}^p(\mathcal{U}, \mathcal{F}) = 0$。
\end{enumerate}
则对所有 $p > 0$ 以及任意 $U \in \mathcal{B}$，都有
$H^p(U, \mathcal{F}) = 0$。
\end{lemma}

\begin{proof}
设 $\mathcal{F}$ 和 $\text{Cov}$ 如引理所述。我们把这一点表述为：
“对任意 $\mathcal{U} \in \text{Cov}$，$\mathcal{F}$ 的高阶
{\v C}ech 上同调消失”。选取 $\mathcal{F}$ 到内射阿贝尔层
$\mathcal{I}$ 中的嵌入 $\mathcal{F} \to \mathcal{I}$。
由引理 \ref{sites-cohomology-lemma-injective-trivial-cech}，对任意
$\mathcal{U} \in \text{Cov}$，$\mathcal{I}$ 的高阶
{\v C}ech 上同调消失。令 $\mathcal{Q} = \mathcal{I}/\mathcal{F}$，
于是有短正合列
$$
0 \to \mathcal{F} \to \mathcal{I} \to \mathcal{Q} \to 0.
$$
由引理 \ref{sites-cohomology-lemma-ses-cech-h1} 和假设 (2)，
这个序列给出正合列
$$
0 \to \mathcal{F}(U) \to \mathcal{I}(U) \to \mathcal{Q}(U) \to 0.
$$
，它对每个 $U \in \mathcal{B}$ 都成立。因此对任意
$\mathcal{U} \in \text{Cov}$，得到 {\v C}ech 复形的短正合列
$$
0 \to
\check{\mathcal{C}}^\bullet(\mathcal{U}, \mathcal{F}) \to
\check{\mathcal{C}}^\bullet(\mathcal{U}, \mathcal{I}) \to
\check{\mathcal{C}}^\bullet(\mathcal{U}, \mathcal{Q}) \to 0
$$
，因为由假设 (1)，{\v C}ech 复形的每一项都是
$\mathcal{B}$ 中各元素上的取值之乘积。特别地，对任意覆盖
$\mathcal{U} \in \text{Cov}$，都有 {\v C}ech 上同调群的长正合列。
这说明 $\mathcal{Q}$ 也是阿贝尔层，并且对所有
$\mathcal{U} \in \text{Cov}$，其高阶 {\v C}ech 上同调消失。

\medskip\noindent
接着考察上同调长正合列
$$
\xymatrix{
0 \ar[r] &
H^0(U, \mathcal{F}) \ar[r] &
H^0(U, \mathcal{I}) \ar[r] &
H^0(U, \mathcal{Q}) \ar[lld] \\
&
H^1(U, \mathcal{F}) \ar[r] &
H^1(U, \mathcal{I}) \ar[r] &
H^1(U, \mathcal{Q}) \ar[lld] \\
&
\ldots & \ldots & \ldots \\
}
$$
，其中 $U \in \mathcal{B}$ 任意。由于 $\mathcal{I}$ 内射，
对 $n > 0$ 有 $H^n(U, \mathcal{I}) = 0$（见《导出范畴》引理
\ref{derived-lemma-higher-derived-functors}）。由上可知，
$H^0(U, \mathcal{I}) \to H^0(U, \mathcal{Q})$ 为满射，
故 $H^1(U, \mathcal{F}) = 0$。
由于 $\mathcal{F}$ 是任意一个对所有 $\mathcal{U} \in \text{Cov}$
高阶 {\v C}ech 上同调消失的阿贝尔层，而 $\mathcal{Q}$
也是这种层，所以同样有 $H^1(U, \mathcal{Q}) = 0$。
由长正合列，这又推出 $H^2(U, \mathcal{F}) = 0$。
如此继续即可。
\end{proof}














\section{第二上同调与 gerbe}
\label{sites-cohomology-section-gerbes}

\noindent
设 $p : \mathcal{S} \to \mathcal{C}$ 是站点上的一个 gerbe，
其所有自同构群均为交换群。在这种情形下，自同构层的第一与第二上同调群
（Stacks，引理 \ref{stacks-lemma-gerbe-abelian-auts}）
控制对象的存在性。

\medskip\noindent
以下引理将由我们今后补充的、关于这一关系的更完整讨论所取代。

\begin{lemma}
\label{sites-cohomology-lemma-existence}
设 $\mathcal{C}$ 为站点。设 $p : \mathcal{S} \to \mathcal{C}$
为此站点上的一个 gerbe，且其自同构层均为阿贝尔层。
令 $\mathcal{G}$ 为 Stacks，引理
\ref{stacks-lemma-gerbe-abelian-auts} 中构造的阿贝尔群层。
设 $U$ 为 $\mathcal{C}$ 的一个对象，并满足
\begin{enumerate}
\item $U$ 在 $\mathcal{C}$ 中存在一个由覆盖 $\{U_i \to U\}$
组成的共尾系，使得 $H^1(U_i, \mathcal{G}) = 0$ 且
$H^1(U_i \times_U U_j, \mathcal{G}) = 0$
对所有 $i, j$ 成立；
\item $H^2(U, \mathcal{G}) = 0$.
\end{enumerate}
则存在一个位于 $U$ 上的 $\mathcal{S}$ 中对象。
\end{lemma}

\begin{proof}
由 Stacks，定义 \ref{stacks-definition-gerbe}，
存在覆盖 $\mathcal{U} = \{U_i \to U\}$，并且对每个 $i$，
存在位于 $U_i$ 上的 $\mathcal{S}$ 中对象 $x_i$。
记 $U_{ij} = U_i \times_U U_j$。由 (1)，加细该覆盖后，
可以假设 $H^1(U_i, \mathcal{G}) = 0$ 且
$H^1(U_{ij}, \mathcal{G}) = 0$.
考虑 $\mathcal{C}/U_{ij}$ 上的层
$$
\mathcal{F}_{ij} =
\mathit{Isom}(x_i|_{U_{ij}}, x_j|_{U_{ij}})
$$
。由于
$\mathcal{G}|_{U_{ij}} = \mathit{Aut}(x_i|_{U_{ij}})$，可见存在由预复合给出的作用
$$
\mathcal{G}|_{U_{ij}} \times \mathcal{F}_{ij} \to \mathcal{F}_{ij}
$$
。显然，$\mathcal{F}_{ij}$ 是一个伪
$\mathcal{G}|_{U_{ij}}$-挠子；事实上它是挠子，因为 gerbe 的任意两个对象
局部同构。由我们对覆盖的选取以及引理 \ref{sites-cohomology-lemma-torsors-h1}，
这些挠子均平凡（从而由引理 \ref{sites-cohomology-lemma-trivial-torsor}，它们均有整体截面）。
换言之，可以选取同构
$$
\varphi_{ij} : x_i|_{U_{ij}} \longrightarrow x_j|_{U_{ij}}
$$
。为了得到位于 $U$ 上的对象 $x$，我们将适当调整这些
$\varphi_{ij}$ 的选取以得到下降数据（由于
$p : \mathcal{S} \to \mathcal{C}$ 是叠，该下降数据必为有效）。
成为下降数据的障碍在于上闭链条件可能不成立。令
$U_{ijk} = U_i \times_U U_j \times_U U_k$。于是可以考虑
$$
g_{ijk} = \varphi_{ik}^{-1}|_{U_{ijk}} \circ \varphi_{jk}|_{U_{ijk}} \circ
\varphi_{ij}|_{U_{ijk}}
$$
，这是 $U_{ijk}$ 上 $x_i$ 的一个自同构。因此可以并且确实将
$g_{ijk}$ 视为 $U_{ijk}$ 上 $\mathcal{G}$ 的一个截面。
计算（从略）表明，$(g_{i_0i_1i_2})$ 是 $\mathcal{G}$ 关于覆盖
$\mathcal{U}$ 的 {\v C}ech 复形
${\check C}^\bullet(\mathcal{U}, \mathcal{G})$ 中的一个 $2$-上闭链。
由引理 \ref{sites-cohomology-lemma-cech-spectral-sequence} 的谱序列以及
$H^1(U_i, \mathcal{G}) = 0$ 对所有 $i$ 成立，可知
${\check H}^2(\mathcal{U}, \mathcal{G}) \to H^2(U, \mathcal{G})$
为单射。因此，由假设 $H^2(U, \mathcal{G}) = 0$，
$(g_{i_0i_1i_2})$ 是上边缘。
从而可以找到截面 $g_{ij} \in \mathcal{G}(U_{ij})$，使得
$g_{ik}^{-1}|_{U_{ijk}} g_{jk}|_{U_{ijk}}  g_{ij}|_{U_{ijk}} = g_{ijk}$
对所有 $i,j,k$ 成立。
将 $\varphi_{ij}$ 替换为 $\varphi_{ij}g_{ij}^{-1}$ 后，
可见 $\varphi_{ij}$ 给出了位于 $U_i$ 上的对象 $x_i$ 的下降数据，
证明完毕。
\end{proof}












\section{模的上同调}
\label{sites-cohomology-section-cohomology-modules}

\noindent
前面对阿贝尔层上同调所述的一切，也适用于模的上同调，因为二者一致。

\begin{lemma}
\label{sites-cohomology-lemma-injective-module-injective-presheaf}
设 $(\mathcal{C}, \mathcal{O})$ 为环化站点。
一个内射模层作为范畴 $\textit{PMod}(\mathcal{O})$ 中的对象仍是内射的。
\end{lemma}

\begin{proof}
将《同调代数》引理 \ref{homology-lemma-adjoint-preserve-injectives}
应用于范畴 $\mathcal{A} = \textit{Mod}(\mathcal{O})$、
$\mathcal{B} = \textit{PMod}(\mathcal{O})$、嵌入函子及层化函子。
该引理的所有假设均满足，见《站点上的模》第
\ref{sites-modules-section-sheafification-presheaves-modules} 节。
\end{proof}

\begin{lemma}
\label{sites-cohomology-lemma-include-modules}
设 $(\mathcal{C}, \mathcal{O})$ 为环化站点。考虑函子
$i : \textit{Mod}(\mathcal{C}) \to \textit{PMod}(\mathcal{C})$.
这是左正合函子，其右导出函子由下式给出：
$$
R^pi(\mathcal{F}) = \underline{H}^p(\mathcal{F}) :
U \longmapsto H^p(U, \mathcal{F})
$$
；见第 \ref{sites-cohomology-section-locality} 节中的讨论。
\end{lemma}

\begin{proof}
显然 $i$ 左正合。在 $\textit{Mod}(\mathcal{O})$ 中选取内射分解
$\mathcal{F} \to \mathcal{I}^\bullet$。
根据定义，$R^pi$ 是复形 $\mathcal{I}^\bullet$ 的第 $p$ 个上同调
{\it 预层}。换言之，$R^pi(\mathcal{F})$ 在 $\mathcal{C}$ 的对象
$U$ 上的截面由下式给出：
$$
\frac{\Ker(\mathcal{I}^n(U) \to \mathcal{I}^{n + 1}(U))}
{\Im(\mathcal{I}^{n - 1}(U) \to \mathcal{I}^n(U))}.
$$
这正是 $H^p(U, \mathcal{F})$ 的定义。
\end{proof}

\begin{lemma}
\label{sites-cohomology-lemma-injective-module-trivial-cech}
设 $(\mathcal{C}, \mathcal{O})$ 为环化站点，
$\mathcal{U} = \{U_i \to U\}_{i \in I}$ 为 $\mathcal{C}$ 中的覆盖。
设 $\mathcal{I}$ 为内射 $\mathcal{O}$-模，即
$\textit{Mod}(\mathcal{O})$ 的内射对象。则
$$
\check{H}^p(\mathcal{U}, \mathcal{I}) =
\left\{
\begin{matrix}
\mathcal{I}(U) & \text{若} & p = 0 \\
0 & \text{若} & p > 0
\end{matrix}
\right.
$$
\end{lemma}

\begin{proof}
引理 \ref{sites-cohomology-lemma-cech-map-into} 给出下列一串等式中的第一个等式：
\begin{align*}
\check{\mathcal{C}}^\bullet(\mathcal{U}, \mathcal{I})
& =
\Mor_{\textit{PAb}(\mathcal{C})}(
\mathbf{Z}_{\mathcal{U}, \bullet}, \mathcal{I}) \\
& =
\Mor_{\textit{PMod}(\mathbf{Z})}(
\mathbf{Z}_{\mathcal{U}, \bullet}, \mathcal{I}) \\
& =
\Mor_{\textit{PMod}(\mathcal{O})}(
\mathbf{Z}_{\mathcal{U}, \bullet} \otimes_{p, \mathbf{Z}} \mathcal{O},
\mathcal{I})
\end{align*}
第三个等式来自《站点上的模》引理
\ref{sites-modules-lemma-adjointness-tensor-restrict-presheaves}。
由引理 \ref{sites-cohomology-lemma-injective-module-injective-presheaf}，
$\mathcal{I}$ 是 $\textit{PMod}(\mathcal{O})$ 中的内射对象。
因此 $\Hom_{\textit{PMod}(\mathcal{O})}(-, \mathcal{I})$ 是正合函子。
由引理 \ref{sites-cohomology-lemma-complex-tensored-still-exact}，高阶 {\v C}ech
上同调群均消失。零阶情形见引理 \ref{sites-cohomology-lemma-cech-h0}。
\end{proof}

\begin{lemma}
\label{sites-cohomology-lemma-cohomology-modules-abelian-agree}
设 $\mathcal{C}$ 为站点，$\mathcal{O}$ 为 $\mathcal{C}$ 上的环层。
设 $\mathcal{F}$ 为 $\mathcal{O}$-模，并以 $\mathcal{F}_{ab}$
记其底层阿贝尔群层。则
$$
H^i(\mathcal{C}, \mathcal{F}_{ab})
=
H^i(\mathcal{C}, \mathcal{F})
$$
，且对 $\mathcal{C}$ 的任意对象 $U$，还有
$$
H^i(U, \mathcal{F}_{ab})
=
H^i(U, \mathcal{F}).
$$
这里左边是在 $\textit{Ab}(\mathcal{C})$ 中计算的上同调，
右边是在 $\textit{Mod}(\mathcal{O})$ 中计算的上同调。
\end{lemma}

\begin{proof}
由《导出范畴》引理 \ref{derived-lemma-higher-derived-functors}，
$\delta$-函子 $(\mathcal{F} \mapsto H^p(U, \mathcal{F}))_{p \geq 0}$
是万有的。函子
$\textit{Mod}(\mathcal{O}) \to \textit{Ab}(\mathcal{C})$,
$\mathcal{F} \mapsto \mathcal{F}_{ab}$ 正合。因此
$(\mathcal{F} \mapsto H^p(U, \mathcal{F}_{ab}))_{p \geq 0}$
也是 $\delta$-函子。若能证明
$(\mathcal{F} \mapsto H^p(U, \mathcal{F}_{ab}))_{p \geq 0}$
也万有，则由万有 $\delta$-函子的唯一性可得引理的第二个断言；
见《同调代数》引理
\ref{homology-lemma-uniqueness-universal-delta-functor}。
由于 $\textit{Mod}(\mathcal{O})$ 有充分内射对象，只需证明：
对 $\textit{Mod}(\mathcal{O})$ 中任意内射对象 $\mathcal{I}$，
$H^i(U, \mathcal{I}_{ab}) = 0$；见《同调代数》引理
\ref{homology-lemma-efface-implies-universal}。

\medskip\noindent
设 $\mathcal{I}$ 为 $\textit{Mod}(\mathcal{O})$ 的内射对象。
在引理 \ref{sites-cohomology-lemma-cech-vanish-collection} 中取
$\mathcal{F} = \mathcal{I}$、$\mathcal{B} = \mathcal{C}$ 且
$\text{Cov} = \text{Cov}_\mathcal{C}$。
由引理 \ref{sites-cohomology-lemma-injective-module-trivial-cech}，该引理的假设 (3) 成立。
因此对 $\mathcal{C}$ 的每个对象 $U$，都有
$H^i(U, \mathcal{I}_{ab}) = 0$。

\medskip\noindent
若 $\mathcal{C}$ 有终对象，则这也推出第一个等式。否则，
由《站点》引理 \ref{sites-lemma-topos-good-site}，环化拓扑斯
$(\Sh(\mathcal{C}), \mathcal{O})$ 等价于一个底层站点确有终对象的
环化拓扑斯。因此引理成立。
\end{proof}

\begin{lemma}
\label{sites-cohomology-lemma-cohomology-products}
设 $\mathcal{C}$ 为站点，$I$ 为集合。对 $i \in I$，设
$\mathcal{F}_i$ 为 $\mathcal{C}$ 上的阿贝尔层。设
$U \in \Ob(\mathcal{C})$。典范映射
$$
H^p(U, \prod\nolimits_{i \in I} \mathcal{F}_i)
\longrightarrow
\prod\nolimits_{i \in I} H^p(U, \mathcal{F}_i)
$$
在 $p = 0$ 时为同构，在 $p = 1$ 时为单射。
\end{lemma}

\begin{proof}
对于 $p = 0$，断言成立，因为层的乘积等于其底层预层的乘积；
见《站点》引理 \ref{sites-lemma-limit-sheaf}。
下面证明 $p = 1$ 的情形。令 $\mathcal{F} = \prod \mathcal{F}_i$。
设 $\xi \in H^1(U, \mathcal{F})$ 在
$\prod H^1(U, \mathcal{F}_i)$ 中映为零。由上同调的局部性，
见引理 \ref{sites-cohomology-lemma-kill-cohomology-class-on-covering}，存在覆盖
$\mathcal{U} = \{U_j \to U\}$，使得对所有 $j$ 都有
$\xi|_{U_j} = 0$。由引理 \ref{sites-cohomology-lemma-cech-h1}，这意味着
$\xi$ 来自某个元素
$\check \xi \in \check H^1(\mathcal{U}, \mathcal{F})$.
由于对所有 $i$，映射
$\check H^1(\mathcal{U}, \mathcal{F}_i) \to H^1(U, \mathcal{F}_i)$
均为单射（由引理 \ref{sites-cohomology-lemma-cech-h1}），并且 $\xi$ 在
$\prod H^1(U, \mathcal{F}_i)$ 中的像为零，故其像
$\check \xi_i$ 在 $\check H^1(\mathcal{U}, \mathcal{F}_i)$ 中为零。
另一方面，由于 $\mathcal{F} = \prod \mathcal{F}_i$，
$\check{\mathcal{C}}^\bullet(\mathcal{U}, \mathcal{F})$ 是诸复形
$\check{\mathcal{C}}^\bullet(\mathcal{U}, \mathcal{F}_i)$ 的乘积。
因此由《同调代数》引理
\ref{homology-lemma-product-abelian-groups-exact}，可得所需结论
$\check \xi = 0$。
\end{proof}

\begin{lemma}
\label{sites-cohomology-lemma-restriction-along-monomorphism-surjective}
设 $(\mathcal{C}, \mathcal{O})$ 为环化站点，
$a : U' \to U$ 为 $\mathcal{C}$ 中的单态射。则对任意内射
$\mathcal{O}$-模 $\mathcal{I}$，限制映射
$\mathcal{I}(U) \to \mathcal{I}(U')$ 为满射。
\end{lemma}

\begin{proof}
令 $j : \mathcal{C}/U \to \mathcal{C}$ 与
$j' : \mathcal{C}/U' \to \mathcal{C}$ 为局部化态射
（《站点上的模》第 \ref{sites-modules-section-localize} 节）。
由于 $j_!$ 是限制函子的左伴随，对任意 $\mathcal{O}$-模层
$\mathcal{F}$，有
$$
\Hom_\mathcal{O}(j_!\mathcal{O}_U, \mathcal{F})
=
\Hom_{\mathcal{O}_U}(\mathcal{O}_U, \mathcal{F}|_U)
=
\mathcal{F}(U)
$$
。类似地，层 $j'_!\mathcal{O}_{U'}$ 表示函子
$\mathcal{F} \mapsto \mathcal{F}(U')$。此外，下面构造一个典范的
$\mathcal{O}$-模态射
$$
j'_!\mathcal{O}_{U'} \longrightarrow j_!\mathcal{O}_U
$$
；由 Yoneda 引理（《范畴》引理 \ref{categories-lemma-yoneda}），
它对应于限制映射 $\mathcal{F}(U) \to \mathcal{F}(U')$。
只需证明所显示的模态射是单射；见《同调代数》引理
\ref{homology-lemma-characterize-injectives}。

\medskip\noindent
为构造此映射，只需在如下两个预层之间构造映射：它们分别把
$\mathcal{C}$ 的对象 $V$ 送到 $\mathcal{O}(V)$-模
$$
\bigoplus\nolimits_{\varphi' \in \Mor_\mathcal{C}(V, U')} \mathcal{O}(V)
\quad\text{与}\quad
\bigoplus\nolimits_{\varphi \in \Mor_\mathcal{C}(V, U)} \mathcal{O}(V)
$$
；见《站点上的模》引理
\ref{sites-modules-lemma-extension-by-zero}。
取这样一个映射：它通过 $\mathcal{O}(V)$ 上的恒等映射，
把对应于 $\varphi'$ 的直和项映到对应于
$\varphi = a \circ \varphi'$ 的直和项。由于 $a$ 是单态射，
此映射为单射。层化正合，故结论成立。
\end{proof}






\section{完全无环层}
\label{sites-cohomology-section-limp}

\noindent
设 $(\mathcal{C}, \mathcal{O})$ 为环化站点。
设 $K$ 为 $\mathcal{C}$ 上的集合预层（这里特意使用正体大写字母，
以便与阿贝尔层区分）。给定 $\mathcal{O}$-模层 $\mathcal{F}$，令
$$
\mathcal{F}(K) =
\Mor_{\textit{PSh}(\mathcal{C})}(K, \mathcal{F}) =
\Mor_{\Sh(\mathcal{C})}(K^\#, \mathcal{F})
$$
。函子 $\mathcal{F} \mapsto \mathcal{F}(K)$ 是左正合函子
$\textit{Mod}(\mathcal{O}) \to \textit{Ab}$，故有右导出函子。
我们以 $H^p(K, \mathcal{F})$ 记这些右导出函子，从而
$H^0(K, \mathcal{F}) = \mathcal{F}(K)$。

\medskip\noindent
有如下几点观察：
\begin{enumerate}
\item 由于 $\mathcal{F}(K) = \mathcal{F}(K^\#)$，有
$H^p(K, \mathcal{F}) = H^p(K^\#, \mathcal{F})$.
上述定义允许 $K$ 为预层，仅是记号上的便利。
\item 假设对 $\mathcal{C}$ 的某个对象 $U$，有
$K = h_U$ 或 $K = h_U^\#$。则
$H^p(K, \mathcal{F}) = H^p(U, \mathcal{F})$，因为
$\Mor_{\Sh(\mathcal{C})}(h_U^\#, \mathcal{F}) = \mathcal{F}(U)$；
见《站点》第 \ref{sites-section-representable-sheaves} 节。
\item 若 $\mathcal{O} = \mathbf{Z}$（常值层），则上同调群给出函子
$H^p(K, - ) : \textit{Ab}(\mathcal{C}) \to \textit{Ab}$
，因为在此情形下
$\textit{Mod}(\mathcal{O}) = \textit{Ab}(\mathcal{C})$。
\end{enumerate}
用这种语言可以重述一些已经证明的结果。

\begin{lemma}
\label{sites-cohomology-lemma-compute-cohomology-on-sheaf-sets}
设 $(\mathcal{C}, \mathcal{O})$ 为环化站点，
$K$ 为 $\mathcal{C}$ 上的集合预层。设 $\mathcal{F}$ 为
$\mathcal{O}$-模，并以 $\mathcal{F}_{ab}$ 记其底层阿贝尔群层。
则 $H^p(K, \mathcal{F}) = H^p(K, \mathcal{F}_{ab})$。
\end{lemma}

\begin{proof}
可以用 $K$ 的层化替换 $K$，并假设 $K$ 是层。
注意，$H^p(K, \mathcal{F})$ 与 $H^p(K, \mathcal{F}_{ab})$
都只依赖于拓扑斯，而不依赖于底层站点。因此由《站点》引理
\ref{sites-lemma-topos-good-site}，可以用一个``更大''站点替换
$\mathcal{C}$，使得对 $\mathcal{C}$ 的某个对象 $U$，有
$K = h_U$。此时结论由引理
\ref{sites-cohomology-lemma-cohomology-modules-abelian-agree} 得出。
\end{proof}

\begin{lemma}
\label{sites-cohomology-lemma-cech-to-cohomology-sheaf-sets}
设 $\mathcal{C}$ 为站点，$K' \to K$ 为 $\mathcal{C}$ 上集合预层的
态射，且其层化为满射。令
$K'_p = K' \times_K \ldots \times_K K'$（$p + 1$ 个因子）。
对每个阿贝尔层 $\mathcal{F}$，存在谱序列，其
$E_1^{p, q} = H^q(K'_p, \mathcal{F})$，并收敛到
$H^{p + q}(K, \mathcal{F})$。
\end{lemma}

\begin{proof}
由于层化正合，可知 $(K_p')^\#$ 等于
$(K')^\# \times_{K^\#} \ldots \times_{K^\#} (K')^\#$
（$p + 1$ 个因子）。因此可以用各自的层化替换 $K$ 与 $K'$，
并假设 $K \to K'$ 是层的满射。按照《站点》引理
\ref{sites-lemma-topos-good-site}，用一个``更大''站点替换
$\mathcal{C}$ 后，可以假设 $K,K'$ 是 $\mathcal{C}$ 的对象，且
$\mathcal{U} = \{K' \to K\}$ 是覆盖。于是得到引理
\ref{sites-cohomology-lemma-cech-spectral-sequence} 的 {\v C}ech 到上同调谱序列，
其 $E_1$ 页如引理陈述所示。
\end{proof}

\begin{lemma}
\label{sites-cohomology-lemma-cohomology-on-sheaf-sets}
设 $\mathcal{C}$ 为站点，$K$ 为 $\mathcal{C}$ 上的集合层。
考虑拓扑斯态射
$j : \Sh(\mathcal{C}/K) \to \Sh(\mathcal{C})$，见《站点》引理
\ref{sites-lemma-localize-topos-site}。则 $j^{-1}$ 保持内射对象，且
$H^p(K, \mathcal{F}) = H^p(\mathcal{C}/K, j^{-1}\mathcal{F})$
对 $\mathcal{C}$ 上的任意阿贝尔层 $\mathcal{F}$ 成立。
\end{lemma}

\begin{proof}
由《站点》引理 \ref{sites-lemma-localize-topos} 与
\ref{sites-lemma-localize-topos-site}，拓扑斯态射 $j$
等价于一个局部化。因此结论由引理
\ref{sites-cohomology-lemma-cohomology-of-open} 得出。
\end{proof}

\noindent
结合引理 \ref{sites-cohomology-lemma-compute-cohomology-on-sheaf-sets} 可知，
对环化站点上的模层而言，以下定义也是``正确的''定义。

\begin{definition}
\label{sites-cohomology-definition-limp}
设 $\mathcal{C}$ 为站点。若对每个集合层 $K$ 及所有 $p \geq 1$，
都有 $H^p(K, \mathcal{F}) = 0$，则称阿贝尔层 $\mathcal{F}$ 为
{\it 完全无环的}\footnote{尽管这一术语用于
\cite[Vbis, Proposition 1.3.10]{SGA4}，它很可能不是标准术语。
在 \cite[V, Definition 4.1]{SGA4} 中，这一性质称为 ``flasque''，
但我们不能采用该术语，因为它会与我们对拓扑空间上 flasque 层的定义冲突。
如果你有更好的建议，请发送电子邮件至
\href{mailto:stacks.project@gmail.com}{stacks.project@gmail.com}
。}
\end{definition}

\noindent
显然，完全无环是内蕴性质，即在拓扑斯等价下保持。
完全无环层在站点的所有对象上都有消失的高阶上同调，
但一般而言，完全无环的条件严格更强。
下面给出完全无环层的一个有时很有用的刻画。

\begin{lemma}
\label{sites-cohomology-lemma-characterize-limp}
设 $\mathcal{C}$ 为站点，$\mathcal{F}$ 为阿贝尔层。若
\begin{enumerate}
\item 对 $p > 0$ 及 $U \in \Ob(\mathcal{C})$，有
$H^p(U, \mathcal{F}) = 0$；
\item 对集合层的每个满射 $K' \to K$，增广 {\v C}ech 复形
$$
0 \to H^0(K, \mathcal{F}) \to H^0(K', \mathcal{F}) \to
H^0(K' \times_K K', \mathcal{F}) \to \ldots
$$
正合，
\end{enumerate}
则 $\mathcal{F}$ 完全无环（反之亦然）。
\end{lemma}

\begin{proof}
由假设 (1)，对所有 $p > 0$ 及 $\mathcal{C}$ 的所有对象 $U$，有
$H^p(h_U^\#, g^{-1}\mathcal{I}) = 0$。注意，若
$K = \coprod K_i$ 是 $\mathcal{C}$ 上集合层的余积，则
$H^p(K, g^{-1}\mathcal{I}) = \prod H^p(K_i, g^{-1}\mathcal{I})$。
对任意集合层 $K$，存在满射
$$
K' = \coprod h_{U_i}^\# \longrightarrow K
$$
；见《站点》引理 \ref{sites-lemma-sheaf-coequalizer-representable}。
因此可得：(*) 对每个集合层 $K$，存在集合层的满射
$K' \to K$，使得对 $p > 0$ 有 $H^p(K', \mathcal{F}) = 0$。
我们断言，(*) 与条件 (2) 推出 $\mathcal{F}$ 完全无环。
注意，条件 (*) 与 (2) 只依赖于 $\mathcal{F}$ 作为拓扑斯
$\Sh(\mathcal{C})$ 中的对象，而不依赖于底层站点。
（在证明的其余部分将不再使用性质 (1)。）

\medskip\noindent
我们对 $n \geq 0$ 归纳证明 (*) 与 (2) 推出下列归纳假设 $IH_n$：
对所有 $0 < p \leq n$ 及所有集合层 $K$，有
$H^p(K, \mathcal{F}) = 0$。显然 $IH_0$ 成立。假设 $IH_n$ 成立。
选取集合层 $K$，并选取满射 $K' \to K$，使得对所有 $p > 0$，
$H^p(K', \mathcal{F}) = 0$。有谱序列
$$
E_1^{p, q} = H^q(K'_p, \mathcal{F})
$$
收敛到 $H^{p + q}(K, \mathcal{F})$；见引理
\ref{sites-cohomology-lemma-cech-to-cohomology-sheaf-sets}。
由 $IH_n$，对 $0 < q \leq n$ 有 $E_1^{p, q} = 0$；由假设 (2)，
对 $p > 0$ 有 $E_2^{p, 0} = 0$。最后，由 $K'$ 的选取，
$H^q(K', \mathcal{F}) = 0$，故对 $q > 0$ 有
$E_1^{0, q} = 0$。因此 $p + q = n + 1$ 的所有项
$E_2^{p, q}$ 均为零，从而 $H^{n + 1}(K, \mathcal{F}) = 0$。
\end{proof}







\section{Leray 谱序列}
\label{sites-cohomology-section-leray}

\noindent
证明 Leray 谱序列存在性的关键是下列引理。

\begin{lemma}
\label{sites-cohomology-lemma-direct-image-injective-sheaf}
设 $f : (\Sh(\mathcal{C}), \mathcal{O}_\mathcal{C}) \to
(\Sh(\mathcal{D}), \mathcal{O}_\mathcal{D})$ 为环化拓扑斯的态射。
则对 $\textit{Mod}(\mathcal{O}_\mathcal{C})$ 中的任意内射对象
$\mathcal{I}$，推前 $f_*\mathcal{I}$ 完全无环。
\end{lemma}

\begin{proof}
设 $K$ 为 $\mathcal{D}$ 上的集合层。由《站点上的模》引理
\ref{sites-modules-lemma-morphism-ringed-topoi-comes-from-morphism-ringed-sites}
，可以用``更大''站点替换 $\mathcal{C},\mathcal{D}$，使得 $f$
来自一个由连续函子 $u : \mathcal{D} \to \mathcal{C}$ 诱导的
环化站点态射，并且对 $\mathcal{D}$ 的某个对象 $V$，有
$K = h_V$。

\medskip\noindent
因此，当 $f$ 由环化站点态射给出时，只需证明：对 $q > 0$
及 $\mathcal{D}$ 的所有对象 $V$，$H^q(V, f_*\mathcal{I})$ 为零。
设 $\mathcal{V} = \{V_j \to V\}$ 为 $\mathcal{D}$ 的任意覆盖。
由于 $u$ 连续，$\mathcal{U} = \{u(V_j) \to u(V)\}$
是 $\mathcal{C}$ 的覆盖。于是由 $f_*$ 的定义，有 {\v C}ech 复形的等式
$$
\check{\mathcal{C}}^\bullet(\mathcal{V}, f_*\mathcal{I})
=
\check{\mathcal{C}}^\bullet(\mathcal{U}, \mathcal{I})
$$
。由引理 \ref{sites-cohomology-lemma-injective-module-trivial-cech}，
此复形的正次数上同调为零。结论遂由引理
\ref{sites-cohomology-lemma-cech-vanish-collection} 得出。
\end{proof}

\noindent
对于平坦态射，函子 $f_*$ 保持内射模。特别地，函子
$f_* : \textit{Ab}(\mathcal{C}) \to \textit{Ab}(\mathcal{D})$
总把内射阿贝尔层变为内射阿贝尔层。

\begin{lemma}
\label{sites-cohomology-lemma-pushforward-injective-flat}
设 $f : (\Sh(\mathcal{C}), \mathcal{O}_\mathcal{C}) \to
(\Sh(\mathcal{D}), \mathcal{O}_\mathcal{D})$ 为环化拓扑斯态射。
若 $f$ 平坦，则对任意内射 $\mathcal{O}_\mathcal{C}$-模
$\mathcal{I}$，$f_*\mathcal{I}$ 是内射
$\mathcal{O}_\mathcal{D}$-模。
\end{lemma}

\begin{proof}
此时函子 $f^*$ 正合，见《站点上的模》引理
\ref{sites-modules-lemma-flat-pullback-exact}。
因此结论由《同调代数》引理
\ref{homology-lemma-adjoint-preserve-injectives} 得出。
\end{proof}

\begin{lemma}
\label{sites-cohomology-lemma-limp-acyclic}
设 $(\Sh(\mathcal{C}), \mathcal{O}_\mathcal{C})$ 为环化拓扑斯。
完全无环层对下列函子右无环：
\begin{enumerate}
\item 对 $\mathcal{C}$ 的任意对象 $U$，函子 $H^0(U, -)$；
\item 对任意集合预层 $K$，函子 $\mathcal{F} \mapsto \mathcal{F}(K)$；
\item 整体截面函子 $\Gamma(\mathcal{C}, -)$；
\item 对环化拓扑斯的任意态射
$f : (\Sh(\mathcal{C}), \mathcal{O}_\mathcal{C}) \to
(\Sh(\mathcal{D}), \mathcal{O}_\mathcal{D})$，函子 $f_*$。
\end{enumerate}
\end{lemma}

\begin{proof}
(2) 是完全无环层的定义。正如完全无环层定义后的讨论所指出，
(1) 是 (2) 的推论。(3) 是 (2) 中取 $K = e$ 为
$\Sh(\mathcal{C})$ 的终对象时的特殊情形。

\medskip\noindent
为证明 (4)，由《站点上的模》引理
\ref{sites-modules-lemma-morphism-ringed-topoi-comes-from-morphism-ringed-sites}
，可以假设 $f$ 由站点态射给出。在此情形下，由引理
\ref{sites-cohomology-lemma-higher-direct-images} 中对高阶直接像的描述，
完全无环层的 $R^if_*$ 在 $i > 0$ 时为零。
\end{proof}

\begin{remark}
\label{sites-cohomology-remark-before-Leray}
作为上述结果的推论，《导出范畴》引理
\ref{derived-lemma-compose-derived-functors} 适用于多种情形。
例如，给定环化拓扑斯态射
$f : (\Sh(\mathcal{C}), \mathcal{O}_\mathcal{C}) \to
(\Sh(\mathcal{D}), \mathcal{O}_\mathcal{D})$，有
$$
R\Gamma(\mathcal{D}, Rf_*\mathcal{F}) = R\Gamma(\mathcal{C}, \mathcal{F})
$$
对任意 $\mathcal{O}_\mathcal{C}$-模层 $\mathcal{F}$ 成立。
事实上，对内射 $\mathcal{O}_\mathcal{X}$-模 $\mathcal{I}$，
由引理 \ref{sites-cohomology-lemma-direct-image-injective-sheaf}，
$\mathcal{O}_\mathcal{D}$-模 $f_*\mathcal{I}$ 完全无环；
而由引理 \ref{sites-cohomology-lemma-limp-acyclic}，完全无环层对
$\Gamma(\mathcal{D}, -)$ 无环。
\end{remark}

\begin{lemma}[Leray 谱序列]
\label{sites-cohomology-lemma-Leray}
设 $f : (\Sh(\mathcal{C}), \mathcal{O}_\mathcal{C}) \to
(\Sh(\mathcal{D}), \mathcal{O}_\mathcal{D})$ 为环化拓扑斯态射。
设 $\mathcal{F}^\bullet$ 为 $\mathcal{O}_\mathcal{C}$-模的下有界复形。
存在谱序列
$$
E_2^{p, q} = H^p(\mathcal{D}, R^qf_*(\mathcal{F}^\bullet))
$$
收敛到 $H^{p + q}(\mathcal{C}, \mathcal{F}^\bullet)$。
\end{lemma}

\begin{proof}
这正是《导出范畴》引理
\ref{derived-lemma-grothendieck-spectral-sequence} 中由函子复合
$\Gamma(\mathcal{C}, -) = \Gamma(\mathcal{D}, -) \circ f_*$.
产生的 Grothendieck 谱序列。关于《导出范畴》引理
\ref{derived-lemma-grothendieck-spectral-sequence} 的假设为何满足，
见引理 \ref{sites-cohomology-lemma-direct-image-injective-sheaf} 与
\ref{sites-cohomology-lemma-limp-acyclic}。
\end{proof}

\begin{lemma}
\label{sites-cohomology-lemma-apply-Leray}
设 $f : (\Sh(\mathcal{C}), \mathcal{O}_\mathcal{C}) \to
(\Sh(\mathcal{D}), \mathcal{O}_\mathcal{D})$ 为环化拓扑斯态射，
$\mathcal{F}$ 为 $\mathcal{O}_\mathcal{C}$-模。
\begin{enumerate}
\item 若对 $q > 0$ 有 $R^qf_*\mathcal{F} = 0$，则对所有 $p$，
$H^p(\mathcal{C}, \mathcal{F}) = H^p(\mathcal{D}, f_*\mathcal{F})$。
\item 若对所有 $q$ 及 $p > 0$ 有
$H^p(\mathcal{D}, R^qf_*\mathcal{F}) = 0$，则对所有 $q$，
$H^q(\mathcal{C}, \mathcal{F}) = H^0(\mathcal{D}, R^qf_*\mathcal{F})$。
\end{enumerate}
\end{lemma}

\begin{proof}
这是迫使 Leray 谱序列收敛的两个简单条件。
也可以直接证明这些事实（不使用谱序列），这是层上同调的一个好练习。
\end{proof}

\begin{lemma}[相对 Leray 谱序列]
\label{sites-cohomology-lemma-relative-Leray}
设
$f : (\Sh(\mathcal{C}), \mathcal{O}_\mathcal{C}) \to
(\Sh(\mathcal{D}), \mathcal{O}_\mathcal{D})$
与
$g : (\Sh(\mathcal{D}), \mathcal{O}_\mathcal{D}) \to
(\Sh(\mathcal{E}), \mathcal{O}_\mathcal{E})$
为环化拓扑斯态射。设 $\mathcal{F}$ 为
$\mathcal{O}_\mathcal{C}$-模。存在谱序列，其
$$
E_2^{p, q} = R^pg_*(R^qf_*\mathcal{F})
$$
收敛到 $R^{p + q}(g \circ f)_*\mathcal{F}$。
此谱序列关于 $\mathcal{F}$ 具有函子性，并且对
$\mathcal{O}_\mathcal{C}$-模的下有界复形也有相应版本。
\end{lemma}

\begin{proof}
这是函子复合的 Grothendieck 谱序列；见《导出范畴》引理
\ref{derived-lemma-grothendieck-spectral-sequence} 以及引理
\ref{sites-cohomology-lemma-direct-image-injective-sheaf} 与
\ref{sites-cohomology-lemma-limp-acyclic}。
\end{proof}







\section{基变换映射}
\label{sites-cohomology-section-base-change-map}

\noindent
本节在若干情形下构造基变换映射；一般情形在注记
\ref{sites-cohomology-remark-base-change} 中处理。本节的讨论限制于沿环化站点平坦态射
作基变换的情形，从而避免使用导出拉回。
在陈述结果之前，先讨论导出范畴上的平坦拉回。假设
$g : (\Sh(\mathcal{C}), \mathcal{O}_\mathcal{C})
\to (\Sh(\mathcal{D}), \mathcal{O}_\mathcal{D})$
为环化拓扑斯的平坦态射。由《站点上的模》引理
\ref{sites-modules-lemma-flat-pullback-exact}，函子
$g^* : \textit{Mod}(\mathcal{O}_\mathcal{D}) \to
\textit{Mod}(\mathcal{O}_\mathcal{C})$ 正合。因此它有导出函子
$$
g^* : D(\mathcal{O}_\mathcal{D}) \to D(\mathcal{O}_\mathcal{C})
$$
；它只需拉回 $D(\mathcal{O}_\mathcal{D})$ 中给定对象的一个代表即可计算，
见《导出范畴》引理 \ref{derived-lemma-right-derived-exact-functor}。
它保持上有界与下有界子范畴。因此如上所示，我们以 $g^*$
而非 $Lg^*$ 记此函子。

\begin{lemma}
\label{sites-cohomology-lemma-base-change-map-flat-case}
设
$$
\xymatrix{
(\Sh(\mathcal{C}'), \mathcal{O}_{\mathcal{C}'})
\ar[r]_{g'} \ar[d]_{f'} &
(\Sh(\mathcal{C}), \mathcal{O}_\mathcal{C}) \ar[d]^f \\
(\Sh(\mathcal{D}'), \mathcal{O}_{\mathcal{D}'})
\ar[r]^g &
(\Sh(\mathcal{D}), \mathcal{O}_\mathcal{D})
}
$$
为环化拓扑斯的交换图。设 $\mathcal{F}^\bullet$ 为
$\mathcal{O}_\mathcal{C}$-模的下有界复形。假设 $g$ 与 $g'$ 均平坦。
则 $D^{+}(\mathcal{O}_{\mathcal{D}'})$ 中存在典范基变换映射
$$
g^*Rf_*\mathcal{F}^\bullet
\longrightarrow
R(f')_*(g')^*\mathcal{F}^\bullet
$$
。
\end{lemma}

\begin{proof}
选取内射分解 $\mathcal{F}^\bullet \to \mathcal{I}^\bullet$ 与
$(g')^*\mathcal{F}^\bullet \to \mathcal{J}^\bullet$。
由引理 \ref{sites-cohomology-lemma-pushforward-injective-flat}，
$(g')_*\mathcal{J}^\bullet$ 是表示
$R(g')_*(g')^*\mathcal{F}^\bullet$ 的内射对象复形。因此由
《导出范畴》引理 \ref{derived-lemma-morphisms-lift} 与
\ref{derived-lemma-morphisms-equal-up-to-homotopy}，下图中的箭头
$\beta$ 存在且在同伦意义下唯一：
$$
\xymatrix{
(g')_*(g')^*\mathcal{F}^\bullet \ar[r] &
(g')_*\mathcal{J}^\bullet \\
\mathcal{F}^\bullet \ar[u]^{\text{伴随}} \ar[r] &
\mathcal{I}^\bullet \ar[u]_\beta
}
$$
。推前到 $\mathcal{D}$，得到
$$
f_*\beta :
f_*\mathcal{I}^\bullet
\longrightarrow
f_*(g')_*\mathcal{J}^\bullet
=
g_*(f')_*\mathcal{J}^\bullet
$$
。由 $g^*$ 与 $g_*$ 的伴随性，得到复形态射
$g^*f_*\mathcal{I}^\bullet \to (f')_*\mathcal{J}^\bullet$。
注意，此态射在同伦意义下唯一，因为整个过程中唯一的选取是
态射 $\beta$ 的选取，并且一切都在复形层次上完成。
\end{proof}








\section{上同调与余极限}
\label{sites-cohomology-section-limits}

\noindent
设 $(\mathcal{C}, \mathcal{O})$ 为环化站点。设
$\mathcal{I} \to \textit{Mod}(\mathcal{O})$，
$i \mapsto \mathcal{F}_i$ 为指标范畴 $\mathcal{I}$ 上的图式，
见《范畴》第 \ref{categories-section-limits} 节。
对每个 $i$ 都有典范映射
$\mathcal{F}_i \to \colim_i \mathcal{F}_i$，它诱导上同调之间的映射。
因此，对每个 $p \geq 0$ 及 $\mathcal{C}$ 的每个对象 $U$，
得到典范映射
$$
\colim_i H^p(U, \mathcal{F}_i)
\longrightarrow
H^p(U, \colim_i \mathcal{F}_i)
$$
。一般而言，这些映射不是同构，即使在 $p = 0$ 时也不一定是。

\medskip\noindent
下列引理是《站点》引理
\ref{sites-lemma-directed-colimits-sections} 在上同调中的对应结论。

\begin{lemma}
\label{sites-cohomology-lemma-colim-works-over-collection}
设 $\mathcal{C}$ 为站点，$\text{Cov}_\mathcal{C}$ 为
$\mathcal{C}$ 的覆盖所成的集合（见《站点》定义
\ref{sites-definition-site}）。设
$\mathcal{B} \subset \Ob(\mathcal{C})$，且
$\text{Cov} \subset \text{Cov}_\mathcal{C}$
为子集。假设
\begin{enumerate}
\item 对每个 $\mathcal{U} \in \text{Cov}$，有
$\mathcal{U} = \{U_i \to U\}_{i \in I}$，其中 $I$ 有限，
$U,U_i \in \mathcal{B}$，并且每个
$U_{i_0} \times_U \ldots \times_U U_{i_p} \in \mathcal{B}$；
\item 对每个 $U \in \mathcal{B}$，$\text{Cov}$ 中出现的
$U$ 的覆盖组成 $U$ 的覆盖的一个共尾系。
\end{enumerate}
则映射
$$
\colim_i H^p(U, \mathcal{F}_i)
\longrightarrow
H^p(U, \colim_i \mathcal{F}_i)
$$
对每个 $p \geq 0$、每个 $U \in \mathcal{B}$ 以及每个滤过图式
$\mathcal{I} \to \textit{Ab}(\mathcal{C})$ 均为同构。
\end{lemma}

\begin{proof}
对 $p$ 归纳证明此引理。注意，在 (1) 中我们要求覆盖
$\mathcal{U} \in \text{Cov}$ 有限，故 $\mathcal{B}$ 的所有元素均拟紧。
因此 (2) 与 (1) 推出任意 $U \in \mathcal{B}$ 均满足《站点》引理
\ref{sites-lemma-directed-colimits-sections} (4) 的假设。
于是结论对 $p = 0$ 成立。现在假设引理对 $p$ 成立，并证明
$p + 1$ 的情形。

\medskip\noindent
选取滤过图式
$\mathcal{F} : \mathcal{I} \to \textit{Ab}(\mathcal{C})$,
$i \mapsto \mathcal{F}_i$.
由于 $\textit{Ab}(\mathcal{C})$ 有函子性的内射嵌入，见《内射对象》
定理 \ref{injectives-theorem-sheaves-injectives}，可以找到滤过图式的态射
$\mathcal{F} \to \mathcal{I}$，使得每个
$\mathcal{F}_i \to \mathcal{I}_i$ 都是从阿贝尔层到内射阿贝尔层的
单射。以 $\mathcal{Q}_i$ 记余核，从而有短正合列
$$
0 \to
\mathcal{F}_i \to
\mathcal{I}_i \to
\mathcal{Q}_i \to 0.
$$
层的余极限是预层层次上的余极限的层化，层化正合，并且阿贝尔群的
滤过余极限正合（见《代数》引理
\ref{algebra-lemma-directed-colimit-exact}），故序列
$$
0 \to
\colim_i \mathcal{F}_i \to
\colim_i \mathcal{I}_i \to
\colim_i \mathcal{Q}_i \to 0.
$$
也是短正合列。我们断言，对所有 $U \in \mathcal{B}$ 及所有
$q \geq 1$，有 $H^q(U, \colim_i \mathcal{I}_i) = 0$。
暂且承认此断言，考虑交换图
$$
\xymatrix{
\colim_i H^p(U, \mathcal{I}_i) \ar[d] \ar[r] &
\colim_i H^p(U, \mathcal{Q}_i) \ar[d] \ar[r] &
\colim_i H^{p + 1}(U, \mathcal{F}_i) \ar[d] \ar[r] &
0 \ar[d] \\
H^p(U, \colim_i \mathcal{I}_i) \ar[r] &
H^p(U, \colim_i \mathcal{Q}_i) \ar[r] &
H^{p + 1}(U, \colim_i \mathcal{F}_i) \ar[r] &
0
}
$$
右下角的零来自该断言，右上角的零来自层 $\mathcal{I}_i$ 的内射性。
应用《代数》引理 \ref{algebra-lemma-directed-colimit-exact} 可知上行正合。
因此由蛇引理可得 $p + 1$ 的结论。

\medskip\noindent
还需证明该断言。我们将使用引理
\ref{sites-cohomology-lemma-cech-vanish-collection}。由 $p = 0$ 的结论，
对 $\mathcal{U} \in \text{Cov}$ 有
$$
\check{\mathcal{C}}^\bullet(\mathcal{U}, \colim_i \mathcal{I}_i)
=
\colim_i \check{\mathcal{C}}^\bullet(\mathcal{U}, \mathcal{I}_i)
$$
，因为所有 $U_{j_0} \times_U \ldots \times_U U_{j_p}$
均属于 $\mathcal{B}$。由引理 \ref{sites-cohomology-lemma-injective-trivial-cech}，
{\v C}ech 复形余极限中的每个复形都在次数 $\geq 1$ 无环。
因此由《代数》引理 \ref{algebra-lemma-directed-colimit-exact}，
{\v C}ech 复形
$\check{\mathcal{C}}^\bullet(\mathcal{U}, \colim_i \mathcal{I}_i)$
也在次数 $\geq 1$ 无环。换言之，对所有 $p \geq 1$，有
$\check{H}^p(\mathcal{U}, \colim_i \mathcal{I}_i) = 0$。
故引理 \ref{sites-cohomology-lemma-cech-vanish-collection} 的假设均满足，断言得证。
\end{proof}

\begin{lemma}
\label{sites-cohomology-lemma-colim-global}
设 $\mathcal{C}$ 为站点，$S \subset \Ob(\Sh(\mathcal{C}))$
为子集。以 $*$ 记 $\Sh(\mathcal{C})$ 的终对象。假设
\begin{enumerate}
\item 对某个 $K \in S$，映射 $K \to *$ 为满射；
\item 给定层的满射 $\mathcal{F} \to K$，其中 $K \in S$，
存在 $K' \in S$ 及态射 $K' \to \mathcal{F}$，使得复合
$K' \to K$ 为满射；
\item 给定 $K,K' \in S$，存在满射 $K'' \to K \times K'$，
其中 $K'' \in S$；
\item 给定 $a,b : K \to K'$，其中 $K,K' \in S$，存在满射
$K'' \to \text{等化子}(a, b)$，其中 $K'' \in S$；
\item 每个 $K \in S$ 均拟紧（《站点》定义
\ref{sites-definition-quasi-compact-topos}）。
\end{enumerate}
则对所有 $p \geq 0$，映射
$$
\colim_\lambda H^p(\mathcal{C}, \mathcal{F}_\lambda)
\longrightarrow
H^p(\mathcal{C}, \colim_\lambda \mathcal{F}_\lambda)
$$
对每个滤过图式 $\Lambda \to \textit{Ab}(\mathcal{C})$，
$\lambda \mapsto \mathcal{F}_\lambda$ 均为同构。
\end{lemma}

\begin{proof}
对 $p$ 归纳证明。基本情形 $p = 0$ 由《站点》引理
\ref{sites-lemma-directed-colimits-global-sections} 的 (4) 得出。
下面核对其假设；不过我们建议读者略过这一部分。
假设 $\mathcal{F} \to *$ 为满射。按 (1) 选取 $K \in S$，
使得 $K \to *$ 为满射。于是 $\mathcal{F} \times K \to K$ 为满射。
按 (2) 选取 $K' \to \mathcal{F} \times K$，其中 $K' \in S$，
且 $K' \to K$ 为满射。于是有态射 $K' \to \mathcal{F}$，
并且 $K' \to *$ 为满射。因此《站点》引理
\ref{sites-lemma-directed-colimits-global-sections} 的假设 (4)(a) 满足。
由《站点》引理 \ref{sites-lemma-image-of-quasi-compact} 以及假设
(3)、(5)，对所有 $K \in S$，$K \times K$ 均拟紧。
因此《站点》引理 \ref{sites-lemma-directed-colimits-global-sections}
的假设 (4)(b) 满足。基本情形证毕。

\medskip\noindent
归纳步骤。假设对于 $p \leq p_0$，以及所有可找到满足 (1)--(5)
的集合 $S \subset \Ob(\Sh(\mathcal{C}))$ 的拓扑斯
$\Sh(\mathcal{C})$，结论对 $H^p$ 成立。完全按照引理
\ref{sites-cohomology-lemma-colim-works-over-collection} 的证明论证可知，
只需证明：给定滤过余极限
$\mathcal{I} = \colim \mathcal{I}_\lambda$，其中
$\mathcal{I}_\lambda$ 是内射阿贝尔层，则
$H^{p_0 + 1}(\mathcal{C}, \mathcal{I}) = 0$。
按 (1) 选取满射 $K \to *$，其中 $K \in S$。
以 $K^n$ 记 $K$ 的 $n$ 重自乘积。考虑引理
\ref{sites-cohomology-lemma-cech-to-cohomology-sheaf-sets} 的谱序列
$$
E_1^{p, q} = H^q(K^{p + 1}, \mathcal{I}) \Rightarrow
H^{p + q}(*, \mathcal{I}) = H^{p + q}(\mathcal{C}, \mathcal{I})
$$
。回忆，对 $\mathcal{C}$ 上的任意阿贝尔层 $\mathcal{F}$，有
$H^q(K^{p + 1}, \mathcal{F}) = H^q(\mathcal{C}/K^{p + 1}, j^{-1}\mathcal{F})$,
；见引理 \ref{sites-cohomology-lemma-cohomology-on-sheaf-sets}。由于 $j^{-1}$
与余极限交换，有
$j^{-1}\mathcal{I} = \colim j^{-1}\mathcal{I}_\lambda$
。由引理 \ref{sites-cohomology-lemma-cohomology-of-open}，限制
$j^{-1}\mathcal{I}_\lambda$ 是 $\mathcal{C}/K^{p + 1}$ 上的
内射阿贝尔层。下面将证明归纳假设适用于
$\mathcal{C}/K^{p + 1}$，从而
$H^q(K^{p + 1}, \mathcal{I}) = \colim H^q(K^{p + 1}, \mathcal{I}_\lambda) = 0$
对 $q < p_0 + 1$ 成立（因 $\mathcal{I}_\lambda$ 内射而消失）。由此
$$
H^{p_0 + 1}(\mathcal{C}, \mathcal{I}) =
H^{p_0 + 1}\left(\ldots \to
H^0(K^{p_0}, \mathcal{I}) \to
H^0(K^{p_0 + 1}, \mathcal{I}) \to
H^0(K^{p_0 + 2}, \mathcal{I}) \to \ldots\right)
$$
再次使用归纳假设，右边所示的复形是下列复形关于 $\Lambda$ 的余极限：
$$
\ldots \to
H^0(K^{p_0}, \mathcal{I}_\lambda) \to
H^0(K^{p_0 + 1}, \mathcal{I}_\lambda) \to
H^0(K^{p_0 + 2}, \mathcal{I}_\lambda) \to \ldots
$$
由于 $\mathcal{I}_\lambda$ 是内射阿贝尔层，这些复形正合
（例如可由谱序列得出）。阿贝尔群范畴中的滤过余极限正合，
故得到所需的消失性。

\medskip\noindent
还需证明，对所有 $n \geq 1$，归纳假设均适用于站点
$\mathcal{C}/K^n$。回忆
$\Sh(\mathcal{C}/K^n) = \Sh(\mathcal{C})/K^n$，见《站点》引理
\ref{sites-lemma-localize-topos-site}。因此可以在
$\Sh(\mathcal{C})/K^n$ 中工作。以
$S_n \subset \Ob(\Sh(\mathcal{C}/K^n)$ 记形如
$K' \to K^n$ 的对象所成的集合。逐一核对各项性质：
\begin{enumerate}
\item 由 (3) 及归纳，存在满射 $K' \to K^n$，其中 $K' \in S$。
于是 $(K' \to K^n) \to (K^n \to K^n)$ 是
$\Sh(\mathcal{C})/K^n$ 中的满射，而 $K^n \to K^n$ 是
$\Sh(\mathcal{C})/K^n$ 的终对象。因此 (1) 对 $S_n$ 成立；
\item $S_n$ 的性质 (2) 立即由 $S$ 的性质 (2) 得出；
\item 设 $a : K_1 \to K^n$ 与 $b : K_2 \to K^n$ 属于 $S_n$。
则 $(K_1 \to K^n) \times (K_2 \to K^n)$ 是
$\Sh(\mathcal{C})/K^n$ 的对象
$K_1 \times_{K^n} K_2 \to K^n$。子层
$K_1 \times_{K^n} K_2 \subset K_1 \times K_2$ 是
$a \circ \text{pr}_1$ 与 $b \circ \text{pr}_2$ 的等化子。
写成 $a = (a_1, \ldots, a_n)$ 与 $b = (b_1, \ldots, b_n)$。
选取满射 $K_3 \to K_1 \times K_2$，其中 $K_3 \in S$；
由 $\mathcal{C}$ 的假设 (3)，这是可行的。选取满射
$$
K_4
\longrightarrow
\text{等化子}(K_3 \to K_1 \times K_2 \xrightarrow{a_1, b_1} K)
$$
，其中 $K_4 \in S$。由 $\mathcal{C}$ 的假设 (4)，这是可行的。
选取满射
$$
K_5
\longrightarrow
\text{等化子}(K_4 \to K_1 \times K_2 \xrightarrow{a_2, b_2} K)
$$
，其中 $K_5 \in S$。这仍然可行。以此方式继续，直至得到满射
$$
K_{3 + n}
\longrightarrow
\text{等化子}(K_{2 + n} \to K_1 \times K_2 \xrightarrow{a_n, b_n} K)
$$
，其中 $K_{3 + n} \in S$。由构造，
$K_{3 + n} \to K_1 \times_{K^n} K_2$ 为满射。因此
$(K_{3 + n} \to K^n)$ 属于 $S_n$，并满射到乘积
$(K_1 \to K^n) \times (K_2 \to K^n)$。故 (3) 对 $S_n$ 成立；
\item $S_n$ 的性质 (4) 立即由 $S$ 的性质 (4) 得出；
\item $S_n$ 的性质 (5) 由 $S$ 的性质 (5) 得出。事实上，
$\Sh(\mathcal{C})/K^n$ 的对象 $\mathcal{F} \to K^n$
对应于 $\Sh(\mathcal{C}/K^n)$ 的拟紧对象，当且仅当
$\mathcal{F}$ 是 $\Sh(\mathcal{C})$ 的拟紧对象。
\end{enumerate}
引理证毕。
\end{proof}

\begin{remark}
\label{sites-cohomology-remark-colim-global}
设 $\mathcal{C}$ 为站点，$\mathcal{B} \subset \Ob(\mathcal{C})$
为子集。令 $S \subset \Ob(\Sh(\mathcal{C}))$ 为形如
$$
K = \coprod\nolimits_{i = 1, \ldots, n} h_{U_i}^\#
$$
的层 $K$ 所成的集合，其中 $U_1, \ldots, U_n \in \mathcal{B}$。
可以问：这个集合何时满足引理 \ref{sites-cohomology-lemma-colim-global} 的假设？
一个充分条件是
\begin{enumerate}
\item 对某个 $n \geq 0$ 及 $U_1, \ldots, U_n \in \mathcal{B}$，
映射 $\coprod_{i = 1, \ldots, n} h_{U_i}^\# \to *$ 为满射；
\item $U \in \mathcal{B}$ 的每个覆盖都能由形如
$\{U_i \to U\}_{i = 1, \ldots, n}$ 的覆盖加细，
其中 $U_i \in \mathcal{B}$；
\item 给定 $U,U' \in \mathcal{B}$，存在 $n \geq 0$、
$U_1, \ldots, U_n \in \mathcal{B}$ 以及态射
$U_i \to U$ 与 $U_i \to U'$，使得
$\coprod_{i = 1, \ldots, n} h_{U_i}^\# \to
h_U^\# \times h_{U'}^\#$ 为满射；
\item 给定 $\mathcal{C}$ 中的态射 $a,b : U \to U'$，其中
$U,U' \in \mathcal{B}$，存在 $U_1, \ldots, U_n \in \mathcal{B}$
以及等化 $a,b$ 的态射 $U_i \to U$，使得
$\coprod_{i = 1, \ldots, n} h_{U_i}^\# \to
\text{等化子}(h_a^\#, h_b^\# : h_U^\# \to h_{U'}^\#)$
为满射。
\end{enumerate}
省略详细核对，只指出上述 (2) 保证 $\mathcal{B}$ 的每个元素均拟紧，
从而由《站点》引理 \ref{sites-lemma-coproduct-quasi-compact}，
每个 $K \in S$ 也拟紧。
\end{remark}

\begin{lemma}
\label{sites-cohomology-lemma-higher-direct-image-colimit}
设 $f : \mathcal{C} \to \mathcal{D}$ 为对应于连续函子
$u : \mathcal{D} \to \mathcal{C}$ 的站点态射。
设 $(\mathcal{F}_i, \varphi_{ii'})$ 为 $\mathcal{C}$ 上的阿贝尔层系统，
并令 $\mathcal{F} = \colim \mathcal{F}_i$。设整数 $p \geq 0$。
以 $\mathcal{B}$ 记满足
$H^p(u(V), \mathcal{F}) = \colim H^p(u(V), \mathcal{F}_i)$.
的 $V \in \Ob(\mathcal{D})$ 所成的集合。若 $\mathcal{D}$ 的每个对象
都有一个由 $\mathcal{B}$ 中元素组成的覆盖，则
$R^pf_*\mathcal{F} = \colim R^pf_*\mathcal{F}_i$。
\end{lemma}

\begin{proof}
回忆，$R^pf_*\mathcal{F}$ 是预层 $\mathcal{G}$ 的层化，
其中 $\mathcal{G}$ 把 $V$ 送到 $H^p(u(V), \mathcal{F})$；
见引理 \ref{sites-cohomology-lemma-higher-direct-images}。类似地，
$R^pf_*\mathcal{F}_i$ 是把 $V$ 送到
$H^p(u(V), \mathcal{F}_i)$ 的预层 $\mathcal{G}_i$ 的层化。
回忆，层化是从层到预层的嵌入函子的左伴随，见《站点》第
\ref{sites-section-sheafification} 节。因此层化与余极限交换，
见《范畴》引理 \ref{categories-lemma-adjoint-exact}。
故只需证明预层态射（余极限在预层范畴中取）
$$
\colim \mathcal{G}_i \longrightarrow \mathcal{G}
$$
在层化上诱导同构。这由《站点》引理
\ref{sites-lemma-isomorphism-sheafifications} 及关于
$\mathcal{B}$ 的假设得出。
\end{proof}

\begin{lemma}
\label{sites-cohomology-lemma-colim-sites-injective}
设 $\mathcal{I}$ 为余滤过指标范畴，
$(\mathcal{C}_i, f_a)$ 为 $\mathcal{I}$ 上的站点逆系统，
如《站点》情形 \ref{sites-situation-inverse-limit-sites}。
按照《站点》引理 \ref{sites-lemma-colimit-sites} 与
\ref{sites-lemma-compute-pullback-to-limit}，令
$\mathcal{C} = \colim \mathcal{C}_i$。此外，假设给定
\begin{enumerate}
\item 对所有 $i \in \Ob(\mathcal{I})$，$\mathcal{C}_i$
上的阿贝尔层 $\mathcal{F}_i$；
\item 对 $a : j \to i$，$\mathcal{C}_j$ 上阿贝尔层的态射
$\varphi_a : f_a^{-1}\mathcal{F}_i \to \mathcal{F}_j$
\end{enumerate}
，并且每当 $c = a \circ b$ 时，
$\varphi_c = \varphi_b \circ f_b^{-1}\varphi_a$。
则存在系统的态射
$(\mathcal{F}_i, \varphi_a) \to (\mathcal{G}_i, \psi_a)$，
使得 $\mathcal{F}_i \to \mathcal{G}_i$ 为单射，且
$\mathcal{G}_i$ 是内射阿贝尔层。
\end{lemma}

\begin{proof}
对每个 $i$，选取单射 $\mathcal{F}_i \to \mathcal{A}_i$，
其中 $\mathcal{A}_i$ 是 $\mathcal{C}_i$ 上的内射阿贝尔层。
于是可以考虑映射族
$$
\gamma_i :
\mathcal{F}_i
\longrightarrow
\prod\nolimits_{b : k \to i} f_{b, *}\mathcal{A}_k = \mathcal{G}_i
$$
，其分量映射是态射
$f_b^{-1}\mathcal{F}_i \to \mathcal{F}_k \to \mathcal{A}_k$
的伴随态射。对 $\mathcal{I}$ 中的 $a : j \to i$，存在典范映射
$$
\psi_a : f_a^{-1}\mathcal{G}_i \to \mathcal{G}_j
$$
，其分量为典范映射
$f_b^{-1}f_{a \circ b, *}\mathcal{A}_k \to f_{b, *}\mathcal{A}_k$
，其中 $b : k \to j$。于是得到系统的单射
$(\gamma_i) : (\mathcal{F}_i, \varphi_a) \to (\mathcal{G}_i, \psi_a)$
。注意，$\mathcal{G}_i$ 是 $\mathcal{C}_i$ 上的内射阿贝尔群层；
见引理 \ref{sites-cohomology-lemma-pushforward-injective-flat} 与《同调代数》引理
\ref{homology-lemma-product-injectives}。构造完毕。
\end{proof}

\begin{lemma}
\label{sites-cohomology-lemma-colimit}
在引理 \ref{sites-cohomology-lemma-colim-sites-injective} 的情形下，令
$\mathcal{F} = \colim f_i^{-1}\mathcal{F}_i$。设
$i \in \Ob(\mathcal{I})$，$X_i \in \text{Ob}(\mathcal{C}_i)$。则
$$
\colim_{a : j \to i} H^p(u_a(X_i), \mathcal{F}_j) =
H^p(u_i(X_i), \mathcal{F})
$$
对所有 $p \geq 0$ 成立。
\end{lemma}

\begin{proof}
$p = 0$ 的情形是《站点》引理 \ref{sites-lemma-colimit}。

\medskip\noindent
按照引理 \ref{sites-cohomology-lemma-colim-sites-injective} 选取
$(\mathcal{F}_i, \varphi_a) \to (\mathcal{G}_i, \psi_a)$。
完全依照引理 \ref{sites-cohomology-lemma-colim-works-over-collection} 的证明论证可知，
只需证明对 $p > 0$ 有
$H^p(X, \colim f_i^{-1}\mathcal{G}_i) = 0$。

\medskip\noindent
令 $\mathcal{G} = \colim f_i^{-1}\mathcal{G}_i$。
为证明 $\mathcal{G}$ 在 $\mathcal{C}$ 每个对象上的上同调消失，
只需证明 $\mathcal{G}$ 关于 $\mathcal{C}$ 的任意覆盖
$\mathcal{U}$ 的 {\v C}ech 上同调为零
（引理 \ref{sites-cohomology-lemma-cech-vanish-collection}）。
覆盖 $\mathcal{U}$ 来自某个 $\mathcal{C}_i$ 的覆盖
$\mathcal{U}_i$。由 $p = 0$ 的情形，有
$$
\check{\mathcal{C}}^\bullet(\mathcal{U}, \mathcal{G}) =
\colim_{a : j \to i}
\check{\mathcal{C}}^\bullet(u_a(\mathcal{U}_i), \mathcal{G}_j)
$$
。由引理 \ref{sites-cohomology-lemma-injective-trivial-cech}，右边是无环复形的滤过余极限，
故在正次数无环。见《代数》引理
\ref{algebra-lemma-directed-colimit-exact}。
\end{proof}















\section{平坦分辨}
\label{sites-cohomology-section-flat}

\noindent
本节在环化站点和环化拓扑斯的情形下，重新论证
《上同调》第 \ref{cohomology-section-flat} 节中的论证。

\begin{lemma}
\label{sites-cohomology-lemma-derived-tor-exact}
设 $(\mathcal{C}, \mathcal{O})$ 为环化站点。
设 $\mathcal{G}^\bullet$ 为一个 $\mathcal{O}$-模复形。
函子
$$
K(\textit{Mod}(\mathcal{O}))
\longrightarrow
K(\textit{Mod}(\mathcal{O})),
\quad
\mathcal{F}^\bullet \longmapsto
\text{Tot}(\mathcal{G}^\bullet \otimes_\mathcal{O} \mathcal{F}^\bullet)
$$
以及
$$
K(\textit{Mod}(\mathcal{O}))
\longrightarrow
K(\textit{Mod}(\mathcal{O})),
\quad
\mathcal{F}^\bullet \longmapsto
\text{Tot}(\mathcal{F}^\bullet \otimes_\mathcal{O} \mathcal{G}^\bullet)
$$
都是三角范畴之间的正合函子。
\end{lemma}

\begin{proof}
这由《导出范畴》注记
\ref{derived-remark-double-complex-as-tensor-product-of}.
\end{proof}

\begin{definition}
\label{sites-cohomology-definition-K-flat}
设 $(\mathcal{C}, \mathcal{O})$ 为环化站点。
若对于每个 $\mathcal{O}$-模的无环复形 $\mathcal{F}^\bullet$，复形
$$
\text{Tot}(\mathcal{F}^\bullet \otimes_\mathcal{O} \mathcal{K}^\bullet)
$$
都是无环的，则称 $\mathcal{O}$-模复形 $\mathcal{K}^\bullet$
为 {\it K-平坦}。
\end{definition}

\begin{lemma}
\label{sites-cohomology-lemma-K-flat-quasi-isomorphism}
设 $(\mathcal{C}, \mathcal{O})$ 为环化站点，且 $\mathcal{K}^\bullet$
为 K-平坦复形。则函子
$$
K(\textit{Mod}(\mathcal{O}))
\longrightarrow
K(\textit{Mod}(\mathcal{O})), \quad
\mathcal{F}^\bullet
\longmapsto
\text{Tot}(\mathcal{F}^\bullet \otimes_\mathcal{O} \mathcal{K}^\bullet)
$$
把拟同构映为拟同构。
\end{lemma}

\begin{proof}
这由引理
\ref{sites-cohomology-lemma-derived-tor-exact}
以及拟同构可由其锥为无环复形来刻画这一事实得到。
\end{proof}

\begin{lemma}
\label{sites-cohomology-lemma-restriction-K-flat}
设 $(\mathcal{C}, \mathcal{O})$ 为环化站点，设 $U$ 为 $\mathcal{C}$ 的对象。
若 $\mathcal{K}^\bullet$ 是 K-平坦的 $\mathcal{O}$-模复形，则
$\mathcal{K}^\bullet|_U$ 是 K-平坦的 $\mathcal{O}_U$-模复形。
\end{lemma}

\begin{proof}
设 $\mathcal{G}^\bullet$ 为 $\mathcal{O}_U$-模的正合复形。
由于 $j_{U!}$ 是正合的
(《站点上的模》引理 \ref{sites-modules-lemma-extension-by-zero-exact})
且 $\mathcal{K}^\bullet$ 是 K-平坦的 $\mathcal{O}$-模复形，故复形
$$
j_{U!}(\text{Tot}(\mathcal{G}^\bullet \otimes_{\mathcal{O}_U}
\mathcal{K}^\bullet|_U)) =
\text{Tot}(j_{U!}\mathcal{G}^\bullet \otimes_\mathcal{O} \mathcal{K}^\bullet)
$$
是正合的。这里的等式来自
《站点上的模》引理 \ref{sites-modules-lemma-j-shriek-and-tensor}
以及 $j_{U!}$ 与直和交换（因为它是左伴随）。由
《站点上的模》引理 \ref{sites-modules-lemma-j-shriek-reflects-exactness} 可知 $j_{U!}$ 反映正合性，故结论成立。
\end{proof}

\begin{lemma}
\label{sites-cohomology-lemma-tensor-product-K-flat}
设 $(\mathcal{C}, \mathcal{O})$ 为环化站点。
若 $\mathcal{K}^\bullet$、$\mathcal{L}^\bullet$ 是 $\mathcal{O}$-模的
K-平坦复形，则
$\text{Tot}(\mathcal{K}^\bullet \otimes_\mathcal{O} \mathcal{L}^\bullet)$
是 K-平坦的 $\mathcal{O}$-模复形。
\end{lemma}

\begin{proof}
这由同构
$$
\text{Tot}(\mathcal{M}^\bullet \otimes_\mathcal{O}
\text{Tot}(\mathcal{K}^\bullet \otimes_\mathcal{O} \mathcal{L}^\bullet))
=
\text{Tot}(\text{Tot}(\mathcal{M}^\bullet \otimes_\mathcal{O}
\mathcal{K}^\bullet) \otimes_\mathcal{O} \mathcal{L}^\bullet)
$$
以及定义得到。
\end{proof}

\begin{lemma}
\label{sites-cohomology-lemma-K-flat-two-out-of-three}
设 $(\mathcal{C}, \mathcal{O})$ 为环化站点。
设 $(\mathcal{K}_1^\bullet, \mathcal{K}_2^\bullet, \mathcal{K}_3^\bullet)$
是 $K(\textit{Mod}(\mathcal{O}))$ 中的有区别三角。
若 $\mathcal{K}_i^\bullet$ 中任意两个是 K-平坦的，则第三个也是。
\end{lemma}

\begin{proof}
这由引理 \ref{sites-cohomology-lemma-derived-tor-exact}
以及以下事实得到：在
$K(\textit{Mod}(\mathcal{O}))$
中的有区别三角中，若任意两个复形无环，则第三个也无环。
\end{proof}

\begin{lemma}
\label{sites-cohomology-lemma-K-flat-two-out-of-three-ses}
设 $(\mathcal{C}, \mathcal{O})$ 为环化站点。设
$0 \to \mathcal{K}_1^\bullet \to \mathcal{K}_2^\bullet \to
\mathcal{K}_3^\bullet \to 0$ 是复形的短正合序列，且
$\mathcal{K}_3^\bullet$ 的各项是平坦的 $\mathcal{O}$-模。
若 $\mathcal{K}_i^\bullet$ 中任意两个是 K-平坦的，则第三个也是。
\end{lemma}

\begin{proof}
由《站点上的模》引理 \ref{sites-modules-lemma-flat-tor-zero}，
对每个复形 $\mathcal{L}^\bullet$，我们得到短正合序列
$$
0 \to
\text{Tot}(\mathcal{L}^\bullet \otimes_\mathcal{O} \mathcal{K}_1^\bullet) \to
\text{Tot}(\mathcal{L}^\bullet \otimes_\mathcal{O} \mathcal{K}_2^\bullet) \to
\text{Tot}(\mathcal{L}^\bullet \otimes_\mathcal{O} \mathcal{K}_3^\bullet) \to 0
$$
为复形的短正合序列。因此由上同调层的长正合列和 K-平坦复形的定义，
引理得证。
\end{proof}

\begin{lemma}
\label{sites-cohomology-lemma-bounded-flat-K-flat}
设 $(\mathcal{C}, \mathcal{O})$ 为环化站点。平坦 $\mathcal{O}$-模的
上有界复形是 K-平坦的。
\end{lemma}

\begin{proof}
设 $\mathcal{K}^\bullet$ 为平坦 $\mathcal{O}$-模的上有界复形，
设 $\mathcal{L}^\bullet$ 为 $\mathcal{O}$-模的无环复形。注意
$\mathcal{L}^\bullet = \colim_m \tau_{\leq m}\mathcal{L}^\bullet$
其中余极限逐项取。因此也有
$$
\text{Tot}(\mathcal{K}^\bullet \otimes_\mathcal{O} \mathcal{L}^\bullet)
=
\colim_m \text{Tot}(
\mathcal{K}^\bullet \otimes_\mathcal{O} \tau_{\leq m}\mathcal{L}^\bullet)
$$
逐项成立。因此要证明左边的复形无环，只需证明右边的每个复形无环。
由于 $\tau_{\leq m}\mathcal{L}^\bullet$ 无环，这归结为
$\mathcal{L}^\bullet$ 上有界的情形。
此时《同调》的引理
《同调》引理 \ref{homology-lemma-first-quadrant-ss} 的谱序列为
$$
{}'E_1^{p, q} = H^p(\mathcal{L}^\bullet \otimes_R \mathcal{K}^q)
$$
中的谱序列含有上式，而上式为零，因为 $\mathcal{K}^q$ 平坦且
$\mathcal{L}^\bullet$ 无环。故结论成立。
\end{proof}

\begin{lemma}
\label{sites-cohomology-lemma-colimit-K-flat}
设 $(\mathcal{C}, \mathcal{O})$ 为环化站点。
设 $\mathcal{K}_1^\bullet \to \mathcal{K}_2^\bullet \to \ldots$
为 K-平坦复形组成的系统。
则 $\colim_i \mathcal{K}_i^\bullet$ 是 K-平坦的。
\end{lemma}

\begin{proof}
由于余极限逐项取，显然有
$$
\colim_i \text{Tot}(
\mathcal{F}^\bullet \otimes_\mathcal{O} \mathcal{K}_i^\bullet)
=
\text{Tot}(\mathcal{F}^\bullet \otimes_\mathcal{O}
\colim_i \mathcal{K}_i^\bullet)
$$
因此由滤过余极限正合这一事实，引理得证。
\end{proof}

\begin{lemma}
\label{sites-cohomology-lemma-resolution-by-direct-sums-extensions-by-zero}
设 $(\mathcal{C}, \mathcal{O})$ 为环化站点。对任意 $\mathcal{O}$-模复形
$\mathcal{G}^\bullet$，存在 $\mathcal{O}$-模复形的交换图
$$
\xymatrix{
\mathcal{K}_1^\bullet \ar[d] \ar[r] &
\mathcal{K}_2^\bullet \ar[d] \ar[r] & \ldots \\
\tau_{\leq 1}\mathcal{G}^\bullet \ar[r] &
\tau_{\leq 2}\mathcal{G}^\bullet \ar[r] & \ldots
}
$$
具有如下性质：(1) 垂直箭头是拟同构且逐项满射；
(2) 每个 $\mathcal{K}_n^\bullet$ 是上有界复形，其各项是形如
$j_{U!}\mathcal{O}_U$ 的 $\mathcal{O}$-模的直和；
(3) 映射 $\mathcal{K}_n^\bullet \to \mathcal{K}_{n + 1}^\bullet$ 是逐项可分裂的单射，
余核是形如 $j_{U!}\mathcal{O}_U$ 的 $\mathcal{O}$-模的直和。此外，映射
$\colim \mathcal{K}_n^\bullet \to \mathcal{G}^\bullet$ 是拟同构。
\end{lemma}

\begin{proof}
该图及性质 (1)、(2)、(3) 立即由
《站点上的模》引理 \ref{sites-modules-lemma-module-quotient-flat}
以及《导出范畴》引理 \ref{derived-lemma-special-direct-system} 得到。
诱导映射 $\colim \mathcal{K}_n^\bullet \to \mathcal{G}^\bullet$
是拟同构，因为滤过余极限正合。
\end{proof}

\begin{lemma}
\label{sites-cohomology-lemma-K-flat-resolution}
设 $(\mathcal{C}, \mathcal{O})$ 为环化站点。对任意复形
$\mathcal{G}^\bullet$，存在一个 K-平坦复形 $\mathcal{K}^\bullet$，其各项为平坦的
$\mathcal{O}$-模，并存在逐项满射的拟同构
$\mathcal{K}^\bullet \to \mathcal{G}^\bullet$。
\end{lemma}

\begin{proof}
选取如引理 \ref{sites-cohomology-lemma-resolution-by-direct-sums-extensions-by-zero} 中的图。
每个复形 $\mathcal{K}_n^\bullet$ 都是平坦模的上有界复形，见
《站点上的模》引理 \ref{sites-modules-lemma-j-shriek-flat}。
故由引理 \ref{sites-cohomology-lemma-bounded-flat-K-flat}，$\mathcal{K}_n^\bullet$ 是 K-平坦的。
于是由引理 \ref{sites-cohomology-lemma-colimit-K-flat}，$\colim \mathcal{K}_n^\bullet$ 是 K-平坦的。
由构造，诱导映射 $\colim \mathcal{K}_n^\bullet \to \mathcal{G}^\bullet$
是拟同构且逐项满射。引理
\ref{sites-cohomology-lemma-resolution-by-direct-sums-extensions-by-zero} 的性质 (3) 表明，
$\colim \mathcal{K}_n^m$ 是平坦模的直和，因而平坦，这就证明了最后的断言。
\end{proof}

\begin{lemma}
\label{sites-cohomology-lemma-derived-tor-quasi-isomorphism-other-side}
设 $(\mathcal{C}, \mathcal{O})$ 为环化站点。设
$\alpha : \mathcal{P}^\bullet \to \mathcal{Q}^\bullet$ 为 $\mathcal{O}$-模的
K-平坦复形之间的拟同构。对每个 $\mathcal{O}$-模复形 $\mathcal{F}^\bullet$，
诱导映射
$$
\text{Tot}(\text{id}_{\mathcal{F}^\bullet} \otimes \alpha) :
\text{Tot}(\mathcal{F}^\bullet \otimes_\mathcal{O} \mathcal{P}^\bullet)
\longrightarrow
\text{Tot}(\mathcal{F}^\bullet \otimes_\mathcal{O} \mathcal{Q}^\bullet)
$$
是拟同构。
\end{lemma}

\begin{proof}
选取拟同构 $\mathcal{K}^\bullet \to \mathcal{F}^\bullet$，其中
$\mathcal{K}^\bullet$ 是 K-平坦复形，见
引理 \ref{sites-cohomology-lemma-K-flat-resolution}。
考虑交换图
$$
\xymatrix{
\text{Tot}(\mathcal{K}^\bullet
\otimes_\mathcal{O} \mathcal{P}^\bullet) \ar[r] \ar[d] &
\text{Tot}(\mathcal{K}^\bullet
\otimes_\mathcal{O} \mathcal{Q}^\bullet) \ar[d] \\
\text{Tot}(\mathcal{F}^\bullet
\otimes_\mathcal{O} \mathcal{P}^\bullet) \ar[r] &
\text{Tot}(\mathcal{F}^\bullet
\otimes_\mathcal{O} \mathcal{Q}^\bullet)
}
$$
由引理 \ref{sites-cohomology-lemma-K-flat-quasi-isomorphism}，垂直箭头和上方水平箭头
都是拟同构，故结论成立。
\end{proof}

\noindent
设 $(\mathcal{C}, \mathcal{O})$ 为环化站点。
设 $\mathcal{F}^\bullet$ 为 $D(\mathcal{O})$ 的对象。
选取 K-平坦分辨 $\mathcal{K}^\bullet \to \mathcal{F}^\bullet$，见引理
\ref{sites-cohomology-lemma-K-flat-resolution}。由引理
\ref{sites-cohomology-lemma-derived-tor-exact}，得到三角范畴之间的正合函子
$$
K(\mathcal{O})
\longrightarrow
K(\mathcal{O}),
\quad
\mathcal{G}^\bullet
\longmapsto
\text{Tot}(\mathcal{G}^\bullet \otimes_\mathcal{O} \mathcal{K}^\bullet)
$$
由引理 \ref{sites-cohomology-lemma-K-flat-quasi-isomorphism}，此函子诱导出函子
$D(\mathcal{O}) \to D(\mathcal{O})$，这是因为
$D(\mathcal{O})$ 是 $K(\mathcal{O})$ 在拟同构处的局部化。
由于 $\mathcal{F}^\bullet$ 的 K-平坦分辨范畴是余滤的，且有引理
\ref{sites-cohomology-lemma-derived-tor-quasi-isomorphism-other-side}，所得函子
（在同构意义下）不依赖于 $\mathcal{K}^\bullet$ 的选择。

\begin{definition}
\label{sites-cohomology-definition-derived-tor}
设 $(\mathcal{C}, \mathcal{O})$ 为环化站点。
设 $\mathcal{F}^\bullet$ 为 $D(\mathcal{O})$ 的对象。
上面描述的 {\it 导出张量积}
$$
- \otimes_\mathcal{O}^{\mathbf{L}} \mathcal{F}^\bullet :
D(\mathcal{O})
\longrightarrow
D(\mathcal{O})
$$
是上面描述的三角范畴之间的正合函子。
\end{definition}

\noindent
从我们的显式构造立即可见，对于 $D(\mathcal{O})$ 中的
$\mathcal{G}^\bullet$ 和 $\mathcal{F}^\bullet$，有规范同构
$$
\mathcal{F}^\bullet \otimes_\mathcal{O}^{\mathbf{L}} \mathcal{G}^\bullet
\cong
\mathcal{G}^\bullet \otimes_\mathcal{O}^{\mathbf{L}} \mathcal{F}^\bullet
$$
这里使用《代数进阶》第 \ref{more-algebra-section-sign-rules} 节给出的符号规则。
因此，当我们写
$\mathcal{F}^\bullet \otimes_\mathcal{O}^{\mathbf{L}} \mathcal{G}^\bullet$
时，通常不区分用哪个变量来定义导出张量积。

\begin{definition}
\label{sites-cohomology-definition-tor}
设 $(\mathcal{C}, \mathcal{O})$ 为环化站点。
设 $\mathcal{F}$、$\mathcal{G}$ 为 $\mathcal{O}$-模。
$\mathcal{F}$ 与 $\mathcal{G}$ 的 {\it Tor} 由公式
$$
\text{Tor}_p^\mathcal{O}(\mathcal{F}, \mathcal{G}) =
H^{-p}(\mathcal{F} \otimes_\mathcal{O}^\mathbf{L} \mathcal{G})
$$
定义，其中导出张量积如上所述。
\end{definition}

\noindent
这个定义意味着，对每个 $\mathcal{O}$-模的短正合序列
$0 \to \mathcal{F}_1 \to \mathcal{F}_2 \to \mathcal{F}_3 \to 0$
都有上同调长正合列
$$
\xymatrix{
\mathcal{F}_1 \otimes_\mathcal{O} \mathcal{G} \ar[r] &
\mathcal{F}_2 \otimes_\mathcal{O} \mathcal{G} \ar[r] &
\mathcal{F}_3 \otimes_\mathcal{O} \mathcal{G} \ar[r] & 0 \\
\text{Tor}_1^\mathcal{O}(\mathcal{F}_1, \mathcal{G}) \ar[r] &
\text{Tor}_1^\mathcal{O}(\mathcal{F}_2, \mathcal{G}) \ar[r] &
\text{Tor}_1^\mathcal{O}(\mathcal{F}_3, \mathcal{G}) \ar[ull]
}
$$
对每个 $\mathcal{O}$-模 $\mathcal{G}$ 成立。我们将其称为
该情形所关联的 $\text{Tor}$ 长正合列。

\begin{lemma}
\label{sites-cohomology-lemma-flat-tor-zero}
设 $(\mathcal{C}, \mathcal{O})$ 为环化站点。
设 $\mathcal{F}$ 为 $\mathcal{O}$-模。以下条件等价：
\begin{enumerate}
\item $\mathcal{F}$ 是平坦的 $\mathcal{O}$-模；
\item $\text{Tor}_1^\mathcal{O}(\mathcal{F}, \mathcal{G}) = 0$
对每个 $\mathcal{O}$-模 $\mathcal{G}$ 成立。
\end{enumerate}
\end{lemma}

\begin{proof}
若 $\mathcal{F}$ 平坦，则 $\mathcal{F} \otimes_\mathcal{O} -$
是正合函子且卫星函子消失。反之，设 (2) 成立。若
$\mathcal{G} \to \mathcal{H}$ 为余核为 $\mathcal{Q}$ 的单射，则
$\text{Tor}$ 长正合列表明下列映射的核
$\mathcal{F} \otimes_\mathcal{O} \mathcal{G} \to
\mathcal{F} \otimes_\mathcal{O} \mathcal{H}$
是下列对象的商：
$\text{Tor}_1^\mathcal{O}(\mathcal{F}, \mathcal{Q})$
按假设该对象为零。因此 $\mathcal{F}$ 平坦。
\end{proof}

\begin{lemma}
\label{sites-cohomology-lemma-K-flat-flat-acyclic}
设 $(\mathcal{C}, \mathcal{O})$ 为环化站点。设 $\mathcal{K}^\bullet$
为各项平坦的 K-平坦无环复形。则
$\mathcal{F} = \Ker(\mathcal{K}^n \to \mathcal{K}^{n + 1})$
是平坦的 $\mathcal{O}$-模。
\end{lemma}

\begin{proof}
注意
$$
\ldots \to \mathcal{K}^{n - 2} \to \mathcal{K}^{n - 1} \to
\mathcal{F} \to 0
$$
是我们的模 $\mathcal{F}$ 的平坦分辨。由于平坦模的上有界复形是 K-平坦的
（引理 \ref{sites-cohomology-lemma-bounded-flat-K-flat}），我们可以用此分辨计算
$\text{Tor}_i(\mathcal{F}, \mathcal{G})$，其中 $\mathcal{G}$ 为任意
$\mathcal{O}$-模。一方面
$\mathcal{K}^\bullet \otimes_\mathcal{O}^\mathbf{L} \mathcal{G}$
在 $D(\mathcal{O})$ 中为零，因为 $\mathcal{K}^\bullet$ 无环；另一方面它由
$\mathcal{K}^\bullet \otimes_\mathcal{O} \mathcal{G}$.
因此
$$
\mathcal{K}^{n - 3} \otimes_\mathcal{O} \mathcal{G} \to
\mathcal{K}^{n - 2} \otimes_\mathcal{O} \mathcal{G} \to
\mathcal{K}^{n - 1} \otimes_\mathcal{O} \mathcal{G}
$$
是正合的。所以 $\text{Tor}_1(\mathcal{F}, \mathcal{G}) = 0$，
由引理 \ref{sites-cohomology-lemma-flat-tor-zero} 得出结论。
\end{proof}

\begin{lemma}
\label{sites-cohomology-lemma-factor-through-K-flat}
设 $(\mathcal{C}, \mathcal{O})$ 为环化站点。
设 $a : \mathcal{K}^\bullet \to \mathcal{L}^\bullet$ 为 $\mathcal{O}$-模复形的映射。
若 $\mathcal{K}^\bullet$ K-平坦，则存在复形 $\mathcal{N}^\bullet$ 以及复形映射
$b : \mathcal{K}^\bullet \to \mathcal{N}^\bullet$
和 $c : \mathcal{N}^\bullet \to \mathcal{L}^\bullet$，使得
\begin{enumerate}
\item $\mathcal{N}^\bullet$ 是 K-平坦的；
\item $c$ 是拟同构；
\item $a$ 与 $c \circ b$ 同伦。
\end{enumerate}
若 $\mathcal{K}^\bullet$ 的各项平坦，则可以选取
$\mathcal{N}^\bullet$、$b$ 和 $c$
使 $\mathcal{N}^\bullet$ 也具有同样性质。
\end{lemma}

\begin{proof}
我们将使用同伦范畴 $K(\textit{Mod}(\mathcal{O}))$ 是三角范畴这一事实，见《导出范畴》命题
\ref{derived-proposition-homotopy-category-triangulated}.
选取有区别三角
$\mathcal{K}^\bullet \to \mathcal{L}^\bullet \to
\mathcal{C}^\bullet \to \mathcal{K}^\bullet[1]$.
选取拟同构 $\mathcal{M}^\bullet \to \mathcal{C}^\bullet$，其中
$\mathcal{M}^\bullet$ K-平坦且各项平坦，见引理
\ref{sites-cohomology-lemma-K-flat-resolution}。由三角范畴公理，我们可以将复合映射
$\mathcal{M}^\bullet \to \mathcal{C}^\bullet \to \mathcal{K}^\bullet[1]$
嵌入有区别三角
$\mathcal{K}^\bullet \to \mathcal{N}^\bullet \to
\mathcal{M}^\bullet \to \mathcal{K}^\bullet[1]$.
由引理 \ref{sites-cohomology-lemma-K-flat-two-out-of-three}，$\mathcal{N}^\bullet$ 是 K-平坦的。
再次使用三角范畴公理，可以选取映射 $\mathcal{N}^\bullet \to \mathcal{L}^\bullet$，
使其组成下列有区别三角的态射
$$
\xymatrix{
\mathcal{K}^\bullet \ar[r] \ar[d] &
\mathcal{N}^\bullet \ar[r] \ar[d] &
\mathcal{M}^\bullet \ar[r] \ar[d] &
\mathcal{K}^\bullet[1] \ar[d] \\
\mathcal{K}^\bullet \ar[r] &
\mathcal{L}^\bullet \ar[r] &
\mathcal{C}^\bullet \ar[r] &
\mathcal{K}^\bullet[1]
}
$$
由于三个箭头中有两个是拟同构，故由这些有区别三角所关联的上同调长正合列，
第三个箭头 $\mathcal{N}^\bullet \to \mathcal{L}^\bullet$ 也是拟同构
（也可以将此图映入 $D(\mathcal{O})$，并使用《导出范畴》引理
\ref{derived-lemma-third-isomorphism-triangle}）。这就完成了 (1)、(2)、(3) 的证明。
为证明最后的断言，可以选取 $\mathcal{N}^\bullet$ 使
$\mathcal{N}^n \cong \mathcal{M}^n \oplus \mathcal{K}^n$，见
《导出范畴》引理 \ref{derived-lemma-improve-distinguished-triangle-homotopy}。
若 $\mathcal{K}^\bullet$ 的各项平坦，则由此得到所需的平坦性。
\end{proof}












\section{导出拉回}
\label{sites-cohomology-section-derived-pullback}

\noindent
设
$f : (\Sh(\mathcal{C}), \mathcal{O}) \to
(\Sh(\mathcal{C}'), \mathcal{O}')$
为环化拓扑斯之间的态射。我们可以使用 K-平坦分辨定义导出拉回函子
$$
Lf^* : D(\mathcal{O}') \to D(\mathcal{O})
$$

\begin{lemma}
\label{sites-cohomology-lemma-pullback-K-flat}
设 $f : (\Sh(\mathcal{C}'), \mathcal{O}') \to (\Sh(\mathcal{C}), \mathcal{O})$
为环化拓扑斯之间的态射。设 $\mathcal{K}^\bullet$ 为各项平坦的
$\mathcal{O}$-模 K-平坦复形。则 $f^*\mathcal{K}^\bullet$ 是各项平坦的
$\mathcal{O}'$-模 K-平坦复形。
\end{lemma}

\begin{proof}
由《站点上的模》引理 \ref{sites-modules-lemma-pullback-flat}，
$f^*\mathcal{K}^n$ 的各项是平坦的 $\mathcal{O}'$-模。选取图
$$
\xymatrix{
\mathcal{K}_1^\bullet \ar[d] \ar[r] &
\mathcal{K}_2^\bullet \ar[d] \ar[r] & \ldots \\
\tau_{\leq 1}\mathcal{K}^\bullet \ar[r] &
\tau_{\leq 2}\mathcal{K}^\bullet \ar[r] & \ldots
}
$$
如引理 \ref{sites-cohomology-lemma-resolution-by-direct-sums-extensions-by-zero} 中所示。
下文使用该引理的所有性质而不再另行说明。每个 $\mathcal{K}_n^\bullet$ 是平坦模的
上有界复形，见《站点上的模》引理 \ref{sites-modules-lemma-j-shriek-flat}。
考虑定义 $\mathcal{M}^\bullet$ 的复形短正合序列
$$
0 \to \mathcal{M}^\bullet \to
\colim \mathcal{K}_n^\bullet \to
\mathcal{K}^\bullet \to 0
$$
由引理 \ref{sites-cohomology-lemma-bounded-flat-K-flat} 和 \ref{sites-cohomology-lemma-colimit-K-flat}，
复形 $\colim \mathcal{K}_n^\bullet$ 是 K-平坦的；由《站点上的模》引理
\ref{sites-modules-lemma-colimits-flat}，其各项平坦。由《站点上的模》引理
\ref{sites-modules-lemma-flat-ses}，$\mathcal{M}^\bullet$ 的各项平坦；由引理
\ref{sites-cohomology-lemma-K-flat-two-out-of-three-ses}，$\mathcal{M}^\bullet$ 是 K-平坦的；
再由上同调长正合列，$\mathcal{M}^\bullet$ 无环（因为第二个箭头是拟同构）。拉回
$f^*(\colim \mathcal{K}_n^\bullet) = \colim f^*\mathcal{K}_n^\bullet$
是各项平坦的 $\mathcal{O}'$-模的下有界复形的余极限，因而 K-平坦
（使用上面的同一组引理）。拉回上述短正合序列
$$
0 \to f^*\mathcal{M}^\bullet \to
f^*(\colim \mathcal{K}_n^\bullet) \to
f^*\mathcal{K}^\bullet \to 0
$$
由《站点上的模》引理 \ref{sites-modules-lemma-pullback-ses} 是复形的短正合序列。
因此由引理 \ref{sites-cohomology-lemma-K-flat-two-out-of-three-ses}，只需证明
$f^*\mathcal{M}^\bullet$ K-平坦。这归结为下一段讨论的情形。

\medskip\noindent
设 $\mathcal{K}^\bullet$ 除 K-平坦外还是无环的，且各项平坦。
由引理 \ref{sites-cohomology-lemma-K-flat-flat-acyclic}，$\tau_{\leq n}\mathcal{K}^\bullet$
的所有项都是平坦的 $\mathcal{O}$-模。选取如上的图，并在此图中使用上面证明的
所有性质。记 $\mathcal{M}_n^\bullet$ 为复形映射
$\mathcal{K}_n^\bullet \to \tau_{\leq n}\mathcal{K}^\bullet$，
从而有复形的短正合序列
$$
0 \to \mathcal{M}_n^\bullet \to \mathcal{K}_n^\bullet \to
\tau_{\leq n}\mathcal{K}^\bullet \to 0
$$
由《站点上的模》引理 \ref{sites-modules-lemma-flat-ses}，复形
$\mathcal{M}_n^\bullet$ 的各项平坦。因此在此情形下
$\mathcal{M} = \colim \mathcal{M}_n^\bullet$ 是平坦模的下有界复形的滤过余极限。
于是 $f^*\mathcal{M}^\bullet$ K-平坦（论证同上），结论成立。
\end{proof}

\begin{lemma}
\label{sites-cohomology-lemma-derived-base-change}
设 $f : (\Sh(\mathcal{C}), \mathcal{O}) \to (\Sh(\mathcal{C}'), \mathcal{O}')$
为环化拓扑斯之间的态射。存在三角范畴之间的正合函子
$$
Lf^* : D(\mathcal{O}') \longrightarrow D(\mathcal{O})
$$
使得
$Lf^*\mathcal{K}^\bullet = f^*\mathcal{K}^\bullet$ 对任意各项平坦的
K-平坦复形 $\mathcal{K}^\bullet$ 成立，特别地，对任意平坦
$\mathcal{O}'$-模的上有界复形也成立。
\end{lemma}

\begin{proof}
为此使用《导出范畴》第 \ref{derived-section-derived-functors} 节发展的
一般理论。置 $\mathcal{D} = K(\mathcal{O}')$、$\mathcal{D}' = D(\mathcal{O})$。
记 $F : \mathcal{D} \to \mathcal{D}'$ 为由规则
$F(\mathcal{G}^\bullet) = f^*\mathcal{G}^\bullet$。
令 $S$ 为下列范畴中的拟同构集合：
$\mathcal{D} = K(\mathcal{O}')$.
这给出《导出范畴》情形 \ref{derived-situation-derived-functor} 中的情形，
因而适用《导出范畴》定义 \ref{derived-definition-right-derived-functor-defined}。
我们断言 $LF$ 处处有定义。这由《导出范畴》引理
\ref{derived-lemma-find-existence-computes} 得到，其中
$\mathcal{P} \subset \Ob(\mathcal{D})$ 是各项平坦的 K-平坦复形
$\mathcal{K}^\bullet$ 所成的集合。具体地，(1) 由引理
\ref{sites-cohomology-lemma-K-flat-resolution} 得出；为证明 (2)，需证明对于拟同构
$\mathcal{K}_1^\bullet  \to \mathcal{K}_2^\bullet$ 之间，且两者属于
$\mathcal{P}$，映射 $f^*\mathcal{K}_1^\bullet \to f^*\mathcal{K}_2^\bullet$
是拟同构。为此将其写成
$$
f^{-1}\mathcal{K}_1^\bullet \otimes_{f^{-1}\mathcal{O}'} \mathcal{O}
\longrightarrow
f^{-1}\mathcal{K}_2^\bullet \otimes_{f^{-1}\mathcal{O}'} \mathcal{O}
$$
函子 $f^{-1}$ 正合，故映射
$f^{-1}\mathcal{K}_1^\bullet \to f^{-1}\mathcal{K}_2^\bullet$ 是拟同构。复形
$f^{-1}\mathcal{K}_1^\bullet$ 与 $f^{-1}\mathcal{K}_2^\bullet$
由引理 \ref{sites-cohomology-lemma-pullback-K-flat} 是 $f^{-1}\mathcal{O}'$-模的 K-平坦复形，
因为可以考虑环化拓扑斯态射
$(\Sh(\mathcal{C}), f^{-1}\mathcal{O}') \to
(\Sh(\mathcal{C}'), \mathcal{O}')$。因此引理
\ref{sites-cohomology-lemma-derived-tor-quasi-isomorphism-other-side}
保证显示的映射是拟同构。于是得到导出函子
$$
LF :
D(\mathcal{O}') = S^{-1}\mathcal{D}
\longrightarrow
\mathcal{D}' = D(\mathcal{O})
$$
见《导出范畴》方程 (\ref{derived-equation-everywhere})。
最后，《导出范畴》引理 \ref{derived-lemma-find-existence-computes} 还保证
$LF(\mathcal{K}^\bullet) = F(\mathcal{K}^\bullet) = f^*\mathcal{K}^\bullet$
当 $\mathcal{K}^\bullet$ 属于 $\mathcal{P}$ 时成立。
注意到平坦模的上有界复形由引理 \ref{sites-cohomology-lemma-bounded-flat-K-flat}
属于 $\mathcal{P}$，即可完成证明。
\end{proof}

\begin{lemma}
\label{sites-cohomology-lemma-derived-pullback-composition}
考虑环化拓扑斯态射
$f : (\Sh(\mathcal{C}), \mathcal{O}_\mathcal{C}) \to
(\Sh(\mathcal{D}), \mathcal{O}_\mathcal{D})$
以及
$g : (\Sh(\mathcal{D}), \mathcal{O}_\mathcal{D}) \to
(\Sh(\mathcal{E}), \mathcal{O}_\mathcal{E})$.
则作为函子 $D(\mathcal{O}_\mathcal{E}) \to D(\mathcal{O}_\mathcal{C})$，
$Lf^* \circ Lg^* = L(g \circ f)^*$。
\end{lemma}

\begin{proof}
设 $E$ 为 $D(\mathcal{O}_\mathcal{E})$ 的对象。可用各项平坦的 K-平坦复形
$\mathcal{K}^\bullet$ 表示 $E$，见引理 \ref{sites-cohomology-lemma-K-flat-resolution}。
按构造，$Lg^*E$ 由 $g^*\mathcal{K}^\bullet$ 计算，见引理
\ref{sites-cohomology-lemma-derived-base-change}。由引理 \ref{sites-cohomology-lemma-pullback-K-flat}，复形
$g^*\mathcal{K}^\bullet$ K-平坦且各项平坦。因此 $Lf^*Lg^*E$ 由
$f^*g^*\mathcal{K}^\bullet$ 表示。另一方面 $L(g \circ f)^*E$ 也由
$(g \circ f)^*\mathcal{K}^\bullet = f^*g^*\mathcal{K}^\bullet$ 表示，故结论成立。
\end{proof}

\begin{lemma}
\label{sites-cohomology-lemma-pullback-tensor-product}
设 $f : (\Sh(\mathcal{C}), \mathcal{O}) \to (\Sh(\mathcal{D}), \mathcal{O}')$
为环化拓扑斯之间的态射。存在规范的双函子同构
$$
Lf^*(
\mathcal{F}^\bullet \otimes_{\mathcal{O}'}^{\mathbf{L}} \mathcal{G}^\bullet
) =
Lf^*\mathcal{F}^\bullet 
\otimes_{\mathcal{O}}^{\mathbf{L}}
Lf^*\mathcal{G}^\bullet 
$$
对 $\mathcal{F}^\bullet, \mathcal{G}^\bullet \in \Ob(D(\mathcal{O}'))$ 成立。
\end{lemma}

\begin{proof}
由引理 \ref{sites-cohomology-lemma-derived-base-change} 中导出拉回的构造以及引理
\ref{sites-cohomology-lemma-K-flat-resolution} 中分辨的存在性，可以将 $\mathcal{F}^\bullet$ 和
$\mathcal{G}^\bullet$ 换成各项平坦且 K-平坦的 $\mathcal{O}'$-模复形。在此情形下，
$\mathcal{F}^\bullet \otimes_{\mathcal{O}'}^{\mathbf{L}} \mathcal{G}^\bullet$
就是双复形
$\mathcal{F}^\bullet \otimes_{\mathcal{O}'} \mathcal{G}^\bullet$.
复形
$\text{Tot}(\mathcal{F}^\bullet \otimes_{\mathcal{O}'} \mathcal{G}^\bullet)$
由引理 \ref{sites-cohomology-lemma-tensor-product-K-flat} 和《站点上的模》引理
\ref{sites-modules-lemma-tensor-flats} 可知 K-平坦且各项平坦。
因此本引理的同构来自同构
$$
\text{Tot}(f^*\mathcal{F}^\bullet \otimes_{\mathcal{O}}
f^*\mathcal{G}^\bullet)
\longrightarrow
f^*\text{Tot}(\mathcal{F}^\bullet \otimes_{\mathcal{O}'} \mathcal{G}^\bullet)
$$
其组成部分是《站点上的模》引理
$f^*\mathcal{F}^p \otimes_{\mathcal{O}} f^*\mathcal{G}^q \to
f^*(\mathcal{F}^p \otimes_{\mathcal{O}'} \mathcal{G}^q)$，即
\ref{sites-modules-lemma-tensor-product-pullback} 的同构
\end{proof}

\begin{lemma}
\label{sites-cohomology-lemma-variant-derived-pullback}
设 $f : (\Sh(\mathcal{C}), \mathcal{O}) \to (\Sh(\mathcal{C}'), \mathcal{O}')$
为环化拓扑斯之间的态射。存在规范的双函子同构
$$
\mathcal{F}^\bullet
\otimes_\mathcal{O}^{\mathbf{L}}
Lf^*\mathcal{G}^\bullet
=
\mathcal{F}^\bullet 
\otimes_{f^{-1}\mathcal{O}'}^{\mathbf{L}}
f^{-1}\mathcal{G}^\bullet 
$$
其中 $\mathcal{F}^\bullet$ 属于 $D(\mathcal{O})$，
$\mathcal{G}^\bullet$ 属于 $D(\mathcal{O}')$。
\end{lemma}

\begin{proof}
设 $\mathcal{F}$ 为 $\mathcal{O}$-模，设 $\mathcal{G}$ 为 $\mathcal{O}'$-模。则
$\mathcal{F} \otimes_{\mathcal{O}} f^*\mathcal{G} =
\mathcal{F} \otimes_{f^{-1}\mathcal{O}'} f^{-1}\mathcal{G}$
因为
$f^*\mathcal{G} =
\mathcal{O} \otimes_{f^{-1}\mathcal{O}'} f^{-1}\mathcal{G}$.
由此及定义即可得到引理。
\end{proof}











\begin{lemma}
\label{sites-cohomology-lemma-check-K-flat-stalks}
设 $(\mathcal{C}, \mathcal{O})$ 为环化站点。
设 $\mathcal{K}^\bullet$ 为 $\mathcal{O}$-模复形。
\begin{enumerate}
\item 若 $\mathcal{K}^\bullet$ K-平坦，则对站点 $\mathcal{C}$ 的每个点 $p$，
$\mathcal{O}_p$-模复形 $\mathcal{K}_p^\bullet$ 按《代数进阶》定义
\ref{more-algebra-definition-K-flat} 的意义是 K-平坦的；
\item 若 $\mathcal{C}$ 有足够多的点，则逆命题成立。
\end{enumerate}
\end{lemma}

\begin{proof}
证明 (2)。若 $\mathcal{C}$ 有足够多的点，且对 $\mathcal{C}$ 的所有点 $p$，
$\mathcal{K}_p^\bullet$ 都是 K-平坦的，则 $\mathcal{K}^\bullet$ 是 K-平坦的，
因为 $\otimes$ 和直和与取茎交换，并且可在茎上检验正合性，见
《站点上的模》引理 \ref{sites-modules-lemma-check-exactness-stalks}。

\medskip\noindent
证明 (1)。设 $\mathcal{K}^\bullet$ K-平坦。选取拟同构
$a : \mathcal{L}^\bullet \to \mathcal{K}^\bullet$，使 $\mathcal{L}^\bullet$ K-平坦且各项平坦，见
引理 \ref{sites-cohomology-lemma-K-flat-resolution}。$\mathcal{L}^\bullet$ 的任意拉回都是 K-平坦的，见
引理 \ref{sites-cohomology-lemma-pullback-K-flat}。特别地，茎
$\mathcal{L}_p^\bullet$ 是 $\mathcal{O}_p$-模的 K-平坦复形。
于是 $a$ 的锥 $C(a)$ 是 K-平坦的（引理 \ref{sites-cohomology-lemma-K-flat-two-out-of-three}）
无环 $\mathcal{O}$-模复形，只需证明 $C(a)$ 的茎是 K-平坦的
（由《代数进阶》引理 \ref{more-algebra-lemma-K-flat-two-out-of-three}）。
因此可设 $\mathcal{K}^\bullet$ K-平坦且无环。

\medskip\noindent
设 $\mathcal{K}^\bullet$ 无环且 K-平坦。在继续之前，按《站点》引理
\ref{sites-lemma-topos-good-site} 将站点 $\mathcal{C}$ 换成另一个站点，
以确保 $\mathcal{C}$ 有所有有限极限。这意味着 $p$ 的邻域范畴是滤过的
（《站点》引理 \ref{sites-lemma-neighbourhoods-directed}），并且定义层的茎的余极限
是滤过的。设 $M$ 为有限表示的 $\mathcal{O}_p$-模。只需证明
$\mathcal{K}^\bullet \otimes_{\mathcal{O}_p} M$ 无环，见《代数进阶》引理
\ref{more-algebra-lemma-universally-acyclic-K-flat}。
由于 $\mathcal{O}_p$ 是 $\mathcal{O}(U)$（其中 $U$ 遍历 $p$ 的邻域）的滤过余极限，
可找到 $p$ 的邻域 $(U,x)$ 以及有限表示的 $\mathcal{O}(U)$-模 $M'$，其到
$\mathcal{O}_p$ 的基变换为 $M$，见《代数》引理
\ref{algebra-lemma-colimit-category-fp-modules}。由引理
\ref{sites-cohomology-lemma-restriction-K-flat}，可将 $\mathcal{C},\mathcal{O},\mathcal{K}^\bullet$
替换为 $\mathcal{C}/U,\mathcal{O}_U,\mathcal{K}^\bullet|_U$。
于是可设存在 $\mathcal{O}$-模 $\mathcal{F}$ 使 $M \cong \mathcal{F}_p$。
由于 $\mathcal{K}^\bullet$ K-平坦且无环，复形
$\mathcal{K}^\bullet \otimes_\mathcal{O} \mathcal{F}$ 无环
（按定义它计算导出张量积）。取茎是正合函子，故
$\mathcal{K}^\bullet \otimes_{\mathcal{O}_p} M$ 无环，正是所需。
\end{proof}

\begin{lemma}
\label{sites-cohomology-lemma-pullback-K-flat-points}
设 $f : (\Sh(\mathcal{C}), \mathcal{O}) \to (\Sh(\mathcal{C}'), \mathcal{O}')$
为环化拓扑斯之间的态射。若 $\mathcal{C}$ 有足够多的点，则
$\mathcal{O}'$-模 K-平坦复形的拉回是 $\mathcal{O}$-模的 K-平坦复形。
\end{lemma}

\begin{proof}
这由引理 \ref{sites-cohomology-lemma-check-K-flat-stalks}、《站点上的模》引理
\ref{sites-modules-lemma-pullback-stalk} 以及《代数进阶》引理
\ref{more-algebra-lemma-base-change-K-flat} 得到。
\end{proof}

\begin{lemma}
\label{sites-cohomology-lemma-tensor-pull-compatibility}
设 $f : (\Sh(\mathcal{C}), \mathcal{O}_\mathcal{C}) \to
(\Sh(\mathcal{D}), \mathcal{O}_\mathcal{D})$ 为环化拓扑斯之间的态射。
设 $\mathcal{K}^\bullet$ 和 $\mathcal{M}^\bullet$ 为 $\mathcal{O}_\mathcal{D}$-模复形。
下图
$$
\xymatrix{
Lf^*(\mathcal{K}^\bullet
\otimes_{\mathcal{O}_\mathcal{D}}^\mathbf{L}
\mathcal{M}^\bullet) \ar[r] \ar[d] &
Lf^*\text{Tot}(\mathcal{K}^\bullet
\otimes_{\mathcal{O}_\mathcal{D}}
\mathcal{M}^\bullet) \ar[d] \\
Lf^*\mathcal{K}^\bullet \otimes_{\mathcal{O}_\mathcal{C}}^\mathbf{L}
Lf^*\mathcal{M}^\bullet \ar[d] &
f^*\text{Tot}(\mathcal{K}^\bullet
\otimes_{\mathcal{O}_\mathcal{D}}
\mathcal{M}^\bullet) \ar[d] \\
f^*\mathcal{K}^\bullet \otimes_{\mathcal{O}_\mathcal{C}}^\mathbf{L}
f^*\mathcal{M}^\bullet \ar[r] &
\text{Tot}(f^*\mathcal{K}^\bullet \otimes_{\mathcal{O}_\mathcal{C}}
f^*\mathcal{M}^\bullet)
}
$$
交换。
\end{lemma}

\begin{proof}
我们将使用各项平坦的 K-平坦分辨的存在性
（引理 \ref{sites-cohomology-lemma-K-flat-resolution}），使用导出拉回由此类复形计算这一事实
（引理 \ref{sites-cohomology-lemma-derived-base-change}），以及拉回保持这些性质这一事实
（引理 \ref{sites-cohomology-lemma-pullback-K-flat}）。若选取此类分辨
$\mathcal{P}^\bullet \to \mathcal{K}^\bullet$ 和
$\mathcal{Q}^\bullet \to \mathcal{M}^\bullet$，则有
$$
\xymatrix{
Lf^*\text{Tot}(\mathcal{P}^\bullet
\otimes_{\mathcal{O}_\mathcal{D}}
\mathcal{Q}^\bullet) \ar[r] \ar[d] &
Lf^*\text{Tot}(\mathcal{K}^\bullet
\otimes_{\mathcal{O}_\mathcal{D}}
\mathcal{M}^\bullet) \ar[d] \\
f^*\text{Tot}(\mathcal{P}^\bullet
\otimes_{\mathcal{O}_\mathcal{D}}
\mathcal{Q}^\bullet) \ar[d] \ar[r] &
f^*\text{Tot}(\mathcal{K}^\bullet
\otimes_{\mathcal{O}_\mathcal{D}}
\mathcal{M}^\bullet) \ar[d] \\
\text{Tot}(f^*\mathcal{P}^\bullet \otimes_{\mathcal{O}_\mathcal{C}}
f^*\mathcal{Q}^\bullet) \ar[r] &
\text{Tot}(f^*\mathcal{K}^\bullet \otimes_{\mathcal{O}_\mathcal{C}}
f^*\mathcal{M}^\bullet)
}
$$
交换。不过，由于我们选择了 $\mathcal{P}^\bullet$、$\mathcal{Q}^\bullet$
以及引理 \ref{sites-cohomology-lemma-tensor-product-K-flat}，现在图的左端正是该图的左端。
\end{proof}









\section{无界复形的上同调}
\label{sites-cohomology-section-unbounded}

\noindent
设 $(\mathcal{C}, \mathcal{O})$ 为环化站点。
范畴 $\textit{Mod}(\mathcal{O})$ 是 Grothendieck 阿贝尔范畴：它有所有余极限，
滤过余极限正合，并且有生成元，即
$$
\bigoplus\nolimits_{U \in \Ob(\mathcal{C})} j_{U!}\mathcal{O}_U,
$$
见《站点上的模》第 \ref{sites-modules-section-kernels} 节以及引理
\ref{sites-modules-lemma-j-shriek-flat} 和
\ref{sites-modules-lemma-module-quotient-flat}。
由《内射对象》定理
\ref{injectives-theorem-K-injective-embedding-grothendieck}
对每个 $\mathcal{O}$-模复形 $\mathcal{F}^\bullet$，存在到 K-内射
$\mathcal{O}$-模复形 $\mathcal{I}^\bullet$ 的内射拟同构
$\mathcal{F}^\bullet \to \mathcal{I}^\bullet$，且此嵌入可选为关于
$\mathcal{F}^\bullet$ 的函子性的。由
导出范畴，引理 \ref{derived-lemma-enough-K-injectives-implies}
可知
\begin{enumerate}
\item 任意到三角范畴 $\mathcal{D}$ 的正合函子
$F : K(\textit{Mod}(\mathcal{O})) \to \mathcal{D}$ 都有右导出函子
$RF : D(\mathcal{O}) \to \mathcal{D}$；
\item 对任意到阿贝尔范畴 $\mathcal{A}$ 的加性函子
$F : \textit{Mod}(\mathcal{O}) \to \mathcal{A}$，考虑由 $F$ 诱导的正合函子
$F : K(\textit{Mod}(\mathcal{O})) \to K(\mathcal{A})$，则得到右导出函子
$RF : D(\mathcal{O}) \to D(\mathcal{A})$。
\end{enumerate}
按构造，$RF(\mathcal{F}^\bullet) = F(\mathcal{I}^\bullet)$，
其中 $\mathcal{F}^\bullet \to \mathcal{I}^\bullet$ 如上所述。

\medskip\noindent
下面给出上述构造的一些例子：
\begin{enumerate}
\item 函子 $\Gamma(\mathcal{C}, -) : \textit{Mod}(\mathcal{O}) \to
\text{Mod}_{\Gamma(\mathcal{C}, \mathcal{O})}$ 导出
$$
R\Gamma(\mathcal{C}, -) :
D(\mathcal{O})
\longrightarrow
D(\Gamma(\mathcal{C}, \mathcal{O}))
$$
我们用记号
$H^i(\mathcal{C}, K) = H^i(R\Gamma(\mathcal{C}, K))$ 表示上同调。
\item 对 $\mathcal{C}$ 的对象 $U$，考虑函子
$\Gamma(U, -) : \textit{Mod}(\mathcal{O}) \to
\text{Mod}_{\Gamma(U, \mathcal{O})}$。它导出
$$
R\Gamma(U, -) : D(\mathcal{O}) \to D(\Gamma(U, \mathcal{O}))
$$
我们用记号
$H^i(U, K) = H^i(R\Gamma(U, K))$ 表示上同调。
\item 对环化拓扑斯态射
$f : (\Sh(\mathcal{C}), \mathcal{O}) \to (\Sh(\mathcal{D}), \mathcal{O}')$
考虑函子
$f_* : \textit{Mod}(\mathcal{O}) \to \textit{Mod}(\mathcal{O}')$
它给出全导出直接像
$$
Rf_* : D(\mathcal{O}) \longrightarrow D(\mathcal{O}')
$$
在无界导出范畴上。
\end{enumerate}

\begin{lemma}
\label{sites-cohomology-lemma-adjoint}
设 $f : (\Sh(\mathcal{C}), \mathcal{O}) \to (\Sh(\mathcal{D}), \mathcal{O}')$
为环化拓扑斯之间的态射。上面定义的函子 $Rf_*$ 与引理
\ref{sites-cohomology-lemma-derived-base-change} 中定义的函子 $Lf^*$ 互为伴随：
$$
\Hom_{D(\mathcal{O})}(Lf^*\mathcal{G}^\bullet, \mathcal{F}^\bullet)
=
\Hom_{D(\mathcal{O}')}(\mathcal{G}^\bullet, Rf_*\mathcal{F}^\bullet)
$$
关于 $\mathcal{F}^\bullet \in \Ob(D(\mathcal{O}))$ 和
$\mathcal{G}^\bullet \in \Ob(D(\mathcal{O}'))$ 双函子性地成立。
\end{lemma}

\begin{proof}
这形式地由 $Rf_*$ 和 $Lf^*$ 存在这一事实得到，见《导出范畴》引理
\ref{derived-lemma-derived-adjoint-functors}。
\end{proof}

\begin{lemma}
\label{sites-cohomology-lemma-derived-pushforward-composition}
设
$f : (\Sh(\mathcal{C}), \mathcal{O}_\mathcal{C}) \to
(\Sh(\mathcal{D}), \mathcal{O}_\mathcal{D})$
以及
$g : (\Sh(\mathcal{D}), \mathcal{O}_\mathcal{D}) \to
(\Sh(\mathcal{E}), \mathcal{O}_\mathcal{E})$
为环化拓扑斯态射。则作为函子
$D(\mathcal{O}_\mathcal{C}) \to D(\mathcal{O}_\mathcal{E})$，有
$Rg_* \circ Rf_* = R(g \circ f)_*$。
\end{lemma}

\begin{proof}
由引理 \ref{sites-cohomology-lemma-adjoint}，$Rg_* \circ Rf_*$ 是 $Lf^* \circ Lg^*$ 的伴随。
由引理 \ref{sites-cohomology-lemma-derived-pullback-composition}，有
$Lf^* \circ Lg^* = L(g \circ f)^*$，因此由伴随函子的唯一性，
$Rg_* \circ Rf_* = R(g \circ f)_*$。
\end{proof}

\begin{remark}
\label{sites-cohomology-remark-base-change}
无界导出函子 $Lf^*$ 和 $Rf_*$ 的构造使我们能够在完全一般性下构造基变换映射。
具体地，设
$$
\xymatrix{
(\Sh(\mathcal{C}'), \mathcal{O}_{\mathcal{C}'})
\ar[r]_{g'} \ar[d]_{f'} &
(\Sh(\mathcal{C}), \mathcal{O}_\mathcal{C}) \ar[d]^f \\
(\Sh(\mathcal{D}'), \mathcal{O}_{\mathcal{D}'})
\ar[r]^g &
(\Sh(\mathcal{D}), \mathcal{O}_\mathcal{D})
}
$$
是环化拓扑斯的交换图。设 $K$ 为 $D(\mathcal{O}_\mathcal{C})$ 的对象。
则存在规范的基变换映射
$$
Lg^*Rf_*K \longrightarrow R(f')_*L(g')^*K
$$
它位于 $D(\mathcal{O}_{\mathcal{D}'})$ 中。具体地，该映射是下列映射的伴随：
$L(f')^*Lg^*Rf_*K \to L(g')^*K$.
由于 $L(f')^* \circ Lg^* = L(g')^* \circ Lf^*$，可见这与映射
$L(g')^*Lf^*Rf_*K \to L(g')^*K$ 相同；我们可以将其取为伴随映射
$Lf^*Rf_*K \to K$ 在 $L(g')^*$ 下的像。
\end{remark}

\begin{remark}
\label{sites-cohomology-remark-compose-base-change}
考虑交换图
$$
\xymatrix{
(\Sh(\mathcal{B}'), \mathcal{O}_{\mathcal{B}'})
\ar[r]_k \ar[d]_{f'} &
(\Sh(\mathcal{B}), \mathcal{O}_\mathcal{B}) \ar[d]^f \\
(\Sh(\mathcal{C}'), \mathcal{O}_{\mathcal{C}'})
\ar[r]^l \ar[d]_{g'} &
(\Sh(\mathcal{C}), \mathcal{O}_\mathcal{C}) \ar[d]^g \\
(\Sh(\mathcal{D}'), \mathcal{O}_{\mathcal{D}'})
\ar[r]^m &
(\Sh(\mathcal{D}), \mathcal{O}_\mathcal{D}) \\
}
$$
它是环化拓扑斯的图。则《基变换》注记
\ref{sites-cohomology-remark-base-change} 中两个方形的基变换映射复合，给出外部矩形的基变换映射。
更精确地，复合
\begin{align*}
Lm^* \circ R(g \circ f)_*
& =
Lm^* \circ Rg_* \circ Rf_* \\
& \to Rg'_* \circ Ll^* \circ Rf_* \\
& \to Rg'_* \circ Rf'_* \circ Lk^* \\
& = R(g' \circ f')_* \circ Lk^*
\end{align*}
是该矩形的基变换映射。验证从略。
\end{remark}

\begin{remark}
\label{sites-cohomology-remark-compose-base-change-horizontal}
考虑交换图
$$
\xymatrix{
(\Sh(\mathcal{C}''), \mathcal{O}_{\mathcal{C}''})
\ar[r]_{g'} \ar[d]_{f''} &
(\Sh(\mathcal{C}'), \mathcal{O}_{\mathcal{C}'})
\ar[r]_g \ar[d]_{f'} &
(\Sh(\mathcal{C}), \mathcal{O}_\mathcal{C}) \ar[d]^f \\
(\Sh(\mathcal{D}''), \mathcal{O}_{\mathcal{D}''})
\ar[r]^{h'} &
(\Sh(\mathcal{D}'), \mathcal{O}_{\mathcal{D}'})
\ar[r]^h &
(\Sh(\mathcal{D}), \mathcal{O}_\mathcal{D})
}
$$
它是环化拓扑斯的图。则《基变换》注记
\ref{sites-cohomology-remark-base-change} 中两个方形的基变换映射复合，给出外部矩形的基变换映射。
更精确地，复合
\begin{align*}
L(h \circ h')^* \circ Rf_*
& =
L(h')^* \circ Lh^* \circ Rf_* \\
& \to L(h')^* \circ Rf'_* \circ Lg^* \\
& \to Rf''_* \circ L(g')^* \circ Lg^* \\
& = Rf''_* \circ L(g \circ g')^*
\end{align*}
是该矩形的基变换映射。验证从略。
\end{remark}



\begin{lemma}
\label{sites-cohomology-lemma-adjoints-push-pull-compatibility}
设 $f : (\Sh(\mathcal{C}), \mathcal{O}_\mathcal{C}) \to
(\Sh(\mathcal{D}), \mathcal{O}_\mathcal{D})$ 为环化拓扑斯之间的态射。
设 $\mathcal{K}^\bullet$ 为 $\mathcal{O}_\mathcal{C}$-模复形。下图
$$
\xymatrix{
Lf^*f_*\mathcal{K}^\bullet \ar[r] \ar[d] &
f^*f_*\mathcal{K}^\bullet \ar[d] \\
Lf^*Rf_*\mathcal{K}^\bullet \ar[r] &
\mathcal{K}^\bullet
}
$$
来自复形上的 $Lf^* \to f^*$、复形上的 $f_* \to Rf_*$ 以及伴随
$Lf^* \circ Rf_* \to \text{id}$，并且在 $D(\mathcal{O}_\mathcal{C})$ 中交换。
\end{lemma}

\begin{proof}
我们将使用 K-平坦分辨和 K-内射分辨的存在性，见引理
\ref{sites-cohomology-lemma-K-flat-resolution}、\ref{sites-cohomology-lemma-derived-base-change}、
\ref{sites-cohomology-lemma-pullback-K-flat} 以及上面的讨论。
选取拟同构
$\mathcal{K}^\bullet \to \mathcal{I}^\bullet$，其中 $\mathcal{I}^\bullet$
是 $\mathcal{O}_\mathcal{C}$-模的 K-内射复形。选取拟同构
$\mathcal{Q}^\bullet \to f_*\mathcal{I}^\bullet$
其中 $\mathcal{Q}^\bullet$ 是各项平坦的 $\mathcal{O}_\mathcal{D}$-模 K-平坦复形。
可选取各项平坦的 $\mathcal{O}_\mathcal{D}$-模 K-平坦复形
$\mathcal{P}^\bullet$ 以及复形态射图
$$
\xymatrix{
\mathcal{P}^\bullet \ar[r] \ar[d] &
f_*\mathcal{K}^\bullet \ar[d] \\
\mathcal{Q}^\bullet \ar[r] & f_*\mathcal{I}^\bullet
}
$$
在同伦意义下交换，且上方水平箭头是拟同构。具体地，首先可对某个复形
$\mathcal{P}^\bullet$ 选取这样的图，因为拟同构在复形的同伦范畴中形成乘法系统；
然后可用各项平坦的 K-平坦复形对 $\mathcal{P}^\bullet$ 作分辨。
取拉回得到复形态射图
$$
\xymatrix{
f^*\mathcal{P}^\bullet \ar[r] \ar[d] &
f^*f_*\mathcal{K}^\bullet \ar[d] \ar[r] &
\mathcal{K}^\bullet \ar[d] \\
f^*\mathcal{Q}^\bullet \ar[r] &
f^*f_*\mathcal{I}^\bullet \ar[r] &
\mathcal{I}^\bullet
}
$$
在同伦意义下交换。外部矩形证明了引理中的断言。
\end{proof}



\begin{remark}
\label{sites-cohomology-remark-cup-product}
设 $f : (\Sh(\mathcal{C}), \mathcal{O}_\mathcal{C}) \to
(\Sh(\mathcal{D}), \mathcal{O}_\mathcal{D})$ 为环化拓扑斯之间的态射。
$Lf^*$ 与 $Rf_*$ 的伴随性允许我们构造相对杯积
$$
Rf_*K \otimes_{\mathcal{O}_\mathcal{D}}^\mathbf{L} Rf_*L
\longrightarrow
Rf_*(K \otimes_{\mathcal{O}_\mathcal{C}}^\mathbf{L} L)
$$
它位于 $D(\mathcal{O}_\mathcal{D})$ 中，对所有
$K,L \in D(\mathcal{O}_\mathcal{C})$ 成立。具体地，该映射是下列映射的伴随：
$Lf^*(Rf_*K \otimes_{\mathcal{O}_\mathcal{D}}^\mathbf{L} Rf_*L) \to
K \otimes_{\mathcal{O}_\mathcal{C}}^\mathbf{L} L$，其中可取同构
$Lf^*(Rf_*K \otimes_{\mathcal{O}_\mathcal{D}}^\mathbf{L} Rf_*L) =
Lf^*Rf_*K \otimes_{\mathcal{O}_\mathcal{C}}^\mathbf{L} Lf^*Rf_*L$
(引理 \ref{sites-cohomology-lemma-pullback-tensor-product})
与映射
$Lf^*Rf_*K \otimes_{\mathcal{O}_\mathcal{C}}^\mathbf{L} Lf^*Rf_*L
\to K \otimes_{\mathcal{O}_\mathcal{C}}^\mathbf{L} L$
后者来自余单位 $Lf^* \circ Rf_* \to \text{id}$。
\end{remark}

\begin{lemma}
\label{sites-cohomology-lemma-torsion}
设 $\mathcal{C}$ 为站点。设 $\mathcal{A} \subset \textit{Ab}(\mathcal{C})$
为由挠阿贝尔层组成的 Serre 子范畴。
则函子 $D(\mathcal{A}) \to D_\mathcal{A}(\mathcal{C})$ 是等价。
\end{lemma}

\begin{proof}
关键观察是，内射阿贝尔层 $\mathcal{I}$ 是可除的。具体地，若
$s \in \mathcal{I}(U)$ 是局部截面，则将 $s$ 解释为映射
$s : j_{U!}\mathbf{Z} \to \mathcal{I}$，并将内射对象的定义性质应用于
层的单射 $n : j_{U!}\mathbf{Z} \to j_{U!}\mathbf{Z}$，
从而得到存在 $s' \in \mathcal{I}(U)$ 使 $ns' = s$。

\medskip\noindent
对层 $\mathcal{F}$，记其挠子层为 $\mathcal{F}_{tor}$。
我们断言，若 $\mathcal{I}^\bullet$ 是内射阿贝尔层复形且其上同调层为挠的，则
$$
\mathcal{I}^\bullet_{tor} \to \mathcal{I}^\bullet
$$
是拟同构。具体地，由于 $\mathbf{Q}$ 在 $\mathbf{Z}$ 上平坦，有
$$
H^p(\mathcal{I}^\bullet) \otimes_\mathbf{Z} \mathbf{Q} =
H^p(\mathcal{I}^\bullet \otimes_\mathbf{Z} \mathbf{Q})
$$
它为零，因为上同调层是挠的。由上面证明的可除性，
$\mathcal{I}^\bullet \to \mathcal{I}^\bullet \otimes_\mathbf{Z} \mathbf{Q}$
是满射且核为 $\mathcal{I}^\bullet_{tor}$。由所得短正合序列所关联的
上同调层长正合列，断言成立。

\medskip\noindent
为证明引理，构造右伴随
$T : D(\mathcal{C}) \to D(\mathcal{A})$。具体地，对 $D(\mathcal{C})$ 中的 $K$，
可用上同调层为内射的 K-内射复形 $\mathcal{I}^\bullet$ 表示，见《内射对象》定理
\ref{injectives-theorem-K-injective-embedding-grothendieck}.
令 $T(K) = \mathcal{I}^\bullet_{tor}$，换言之，$T$ 是取挠子的右导出函子。
函子 $T$ 是下列函子的右伴随：
$i : D(\mathcal{A}) \to D_\mathcal{A}(\mathcal{C})$.
这立即由如下观察得到：若 $\mathcal{F}^\bullet$ 是挠层复形，则
$$
\Hom_{K(\mathcal{A})}(\mathcal{F}^\bullet, I^\bullet_{tor}) =
\Hom_{K(\textit{Ab}(\mathcal{C}))}(\mathcal{F}^\bullet, I^\bullet)
$$
特别地，$\mathcal{I}^\bullet_{tor}$ 是 $\mathcal{A}$ 的 K-内射复形。
细节从略；如有疑问，也可由更一般的《导出范畴》引理
\ref{derived-lemma-derived-adjoint-functors} 得到。
上面的断言给出，对 $D(\mathcal{A})$ 中的 $L$ 有 $L = T(i(L))$，
且当 $K \in D_\mathcal{A}(\mathcal{C})$ 时有 $i(T(K)) = K$。
由《范畴》引理 \ref{categories-lemma-adjoint-fully-faithful} 得到结论。
\end{proof}




\section{K-内射复形的一些性质}
\label{sites-cohomology-section-properties-K-injective}

\noindent
设 $(\mathcal{C}, \mathcal{O})$ 为环化站点。设 $U$ 为 $\mathcal{C}$ 的对象。记
$j : (\Sh(\mathcal{C}/U), \mathcal{O}_U) \to (\Sh(\mathcal{C}), \mathcal{O})$
为相应的局部化态射。拉回函子 $j^*$ 是正合的，因为它就是限制函子。
因此，导出拉回 $Lj^*$ 对任意复形都只需通过限制该复形来计算。我们常将相应函子
$$
D(\mathcal{O}) \to D(\mathcal{O}_U),
\quad
E \mapsto j^*E = E|_U
$$
类似地，延零函子
$j_! : \textit{Mod}(\mathcal{O}_U) \to \textit{Mod}(\mathcal{O})$（见《站点上的模》定义
\ref{sites-modules-definition-localize-ringed-site}）是正合函子
（《站点上的模》引理 \ref{sites-modules-lemma-extension-by-zero-exact}）。
因此它诱导函子
$$
j_! : D(\mathcal{O}_U) \to D(\mathcal{O}), \quad
F \mapsto j_!F
$$
只需将 $j_!$ 作用于表示对象 $F$ 的任意复形即可。

\begin{lemma}
\label{sites-cohomology-lemma-restrict-K-injective-to-open}
设 $(\mathcal{C}, \mathcal{O})$ 为环化站点，设 $U$ 为 $\mathcal{C}$ 的对象。
$\mathcal{O}$-模 K-内射复形限制到 $\mathcal{C}/U$ 后，是
$\mathcal{O}_U$-模的 K-内射复形。
\end{lemma}

\begin{proof}
立即由《导出范畴》引理
\ref{derived-lemma-adjoint-preserve-K-injectives} 以及限制函子有正合左伴随
$j_!$ 这一事实得到。见上面的讨论。
\end{proof}

\begin{lemma}
\label{sites-cohomology-lemma-unbounded-cohomology-of-open}
设 $(\mathcal{C}, \mathcal{O})$ 为环化站点。设 $U \in \Ob(\mathcal{C})$。
对 $D(\mathcal{O})$ 中的 $K$，有
$H^p(U, K) = H^p(\mathcal{C}/U, K|_{\mathcal{C}/U})$.
\end{lemma}

\begin{proof}
设 $\mathcal{I}^\bullet$ 为表示 $K$ 的 $\mathcal{O}$-模 K-内射复形。则
$$
H^q(U, K) = H^q(\Gamma(U, \mathcal{I}^\bullet)) =
H^q(\Gamma(\mathcal{C}/U, \mathcal{I}^\bullet|_{\mathcal{C}/U}))
$$
由上同调的构造成立。由引理 \ref{sites-cohomology-lemma-restrict-K-injective-to-open}，复形
$\mathcal{I}^\bullet|_{\mathcal{C}/U}$ 是表示 $K|_{\mathcal{C}/U}$ 的 K-内射复形，
引理得证。
\end{proof}

\begin{lemma}
\label{sites-cohomology-lemma-sheafification-cohomology}
设 $(\mathcal{C}, \mathcal{O})$ 为环化站点。设 $K$ 为 $D(\mathcal{O})$ 的对象。
下列预层的层化
$$
U \mapsto H^q(U, K) = H^q(\mathcal{C}/U, K|_{\mathcal{C}/U})
$$
就是 $K$ 的第 $q$ 个上同调层 $H^q(K)$。
\end{lemma}

\begin{proof}
等式 $H^q(U, K) = H^q(\mathcal{C}/U, K|_{\mathcal{C}/U})$
由引理 \ref{sites-cohomology-lemma-unbounded-cohomology-of-open} 成立。选取表示 $K$ 的 K-内射复形
$\mathcal{I}^\bullet$。则
$$
H^q(U, K) =
\frac{\Ker(\mathcal{I}^q(U) \to \mathcal{I}^{q + 1}(U))}
{\Im(\mathcal{I}^{q - 1}(U) \to \mathcal{I}^q(U))}.
$$
由上同调的构造成立。由于
$H^q(K) = \Ker(\mathcal{I}^q \to \mathcal{I}^{q + 1})/
\Im(\mathcal{I}^{q - 1} \to \mathcal{I}^q)$，结论显然。
\end{proof}

\begin{lemma}
\label{sites-cohomology-lemma-restrict-direct-image-open}
设 $f : (\mathcal{C}, \mathcal{O}_\mathcal{C}) \to
(\mathcal{D}, \mathcal{O}_\mathcal{D})$ 为对应于连续函子 $u : \mathcal{D} \to \mathcal{C}$
的环化站点态射。给定 $V \in \mathcal{D}$，置 $U = u(V)$，并记
$g : (\mathcal{C}/U, \mathcal{O}_U) \to (\mathcal{D}/V, \mathcal{O}_V)$
为诱导的环化站点态射（《站点上的模》引理
\ref{sites-modules-lemma-localize-morphism-ringed-sites}）。
则对 $D(\mathcal{O}_\mathcal{C})$ 中的 $E$，有
$(Rf_*E)|_{\mathcal{D}/V} = Rg_*(E|_{\mathcal{C}/U})$。
\end{lemma}

\begin{proof}
用 $\mathcal{O}_\mathcal{C}$-模 K-内射复形 $\mathcal{I}^\bullet$ 表示 $E$。则
$Rf_*(E) = f_*\mathcal{I}^\bullet$
且由引理 \ref{sites-cohomology-lemma-restrict-K-injective-to-open}，
$Rg_*(E|_{\mathcal{C}/U}) = g_*(\mathcal{I}^\bullet|_{\mathcal{C}/U})$。
由于显然有
$(f_*\mathcal{F})|_{\mathcal{D}/V} = g_*(\mathcal{F}|_{\mathcal{C}/U})$
对 $\mathcal{C}$ 上任意层 $\mathcal{F}$ 成立
（见《站点上的模》引理 \ref{sites-modules-lemma-localize-morphism-ringed-sites}，
或更基本的《站点》引理 \ref{sites-lemma-localize-morphism}），故结论成立。
\end{proof}

\begin{lemma}
\label{sites-cohomology-lemma-Leray-unbounded}
设 $f : (\mathcal{C}, \mathcal{O}_\mathcal{C}) \to
(\mathcal{D}, \mathcal{O}_\mathcal{D})$ 为对应于连续函子 $u : \mathcal{D} \to \mathcal{C}$ 的环化站点态射
。
作为函子 $D(\mathcal{O}_\mathcal{C}) \to D(\Gamma(\mathcal{O}_\mathcal{D}))$，有
$R\Gamma(\mathcal{D}, -) \circ Rf_* = R\Gamma(\mathcal{C}, -)$。
更一般地，对 $V \in \mathcal{D}$ 且 $U = u(V)$，有
$R\Gamma(U, -) = R\Gamma(V, -) \circ Rf_*$。
\end{lemma}

\begin{proof}
考虑点拓扑斯 $pt$，其 $\mathcal{O}_{pt}$ 由环
$\Gamma(\mathcal{O}_\mathcal{D})$ 给出。存在规范态射
$(\mathcal{D}, \mathcal{O}_\mathcal{D}) \to (pt, \mathcal{O}_{pt})$
它是环化拓扑斯态射，并在结构层的全局截面上诱导出该识别。于是
$D(\mathcal{O}_{pt}) = D(\Gamma(\mathcal{O}_\mathcal{D}))$.
断言
$R\Gamma(\mathcal{D}, -) \circ Rf_* = R\Gamma(\mathcal{C}, -)$
由引理 \ref{sites-cohomology-lemma-derived-pushforward-composition} 应用于
$$
(\mathcal{C}, \mathcal{O}_\mathcal{C}) \to
(\mathcal{D}, \mathcal{O}_\mathcal{D}) \to (pt, \mathcal{O}_{pt})
$$
第二个（更一般的）断言在应用引理 \ref{sites-cohomology-lemma-restrict-direct-image-open} 后由第一个断言得到。
\end{proof}

\begin{lemma}
\label{sites-cohomology-lemma-unbounded-describe-higher-direct-images}
设 $f : (\mathcal{C}, \mathcal{O}_\mathcal{C}) \to
(\mathcal{D}, \mathcal{O}_\mathcal{D})$ 为对应于连续函子 $u : \mathcal{D} \to \mathcal{C}$
的环化站点态射。设 $K \in D(\mathcal{O}_\mathcal{C})$。则 $H^i(Rf_*K)$ 是与下列预层
关联的层：
$$
V \mapsto H^i(u(V), K) = H^i(V, Rf_*K)
$$
\end{lemma}

\begin{proof}
等式 $H^i(u(V), K) = H^i(V, Rf_*K)$ 由引理
\ref{sites-cohomology-lemma-Leray-unbounded} 的第二个断言取上同调得到。层化的断言
由引理 \ref{sites-cohomology-lemma-sheafification-cohomology} 得到。
\end{proof}

\begin{lemma}
\label{sites-cohomology-lemma-modules-abelian-unbounded}
设 $(\mathcal{C}, \mathcal{O}_\mathcal{C})$ 为环化站点。
设 $K$ 为 $D(\mathcal{O}_\mathcal{C})$ 中的对象，并记 $K_{ab}$ 为其在
$D(\underline{\mathbf{Z}}_\mathcal{C})$ 中的像。
\begin{enumerate}
\item 存在典范映射
$R\Gamma(\mathcal{C}, K) \to R\Gamma(\mathcal{C}, K_{ab})$
，它是 $D(\textit{Ab})$ 中的同构。
\item 对任意 $U \in \mathcal{C}$，存在典范映射
$R\Gamma(U, K) \to R\Gamma(U, K_{ab})$
，它是 $D(\textit{Ab})$ 中的同构。
\item 设 $f : (\mathcal{C}, \mathcal{O}_\mathcal{C}) \to
(\mathcal{D}, \mathcal{O}_\mathcal{D})$ 为环化站点态射。存在典范映射
$Rf_*K \to Rf_*(K_{ab})$，它是
$D(\underline{\mathbf{Z}}_\mathcal{D})$ 中的同构。
\end{enumerate}
\end{lemma}

\begin{proof}
映射构造如下。选取表示 $K$ 的 K-内射复形
$\mathcal{I}^\bullet$。选取拟同构
$\mathcal{I}^\bullet \to \mathcal{J}^\bullet$，其中 $\mathcal{J}^\bullet$
是阿贝尔群的 K-内射复形。于是 (1) 中的映射为
$\Gamma(\mathcal{C}, \mathcal{I}^\bullet) \to
\Gamma(\mathcal{C}, \mathcal{J}^\bullet)$
(2) 中的映射为
$\Gamma(U, \mathcal{I}^\bullet) \to \Gamma(U, \mathcal{J}^\bullet)$
而 (3) 中的映射为
$f_*\mathcal{I}^\bullet \to f_*\mathcal{J}^\bullet$.
要证明这些映射是同构，只需证明它们在上同调群和上同调层上诱导同构。
由引理 \ref{sites-cohomology-lemma-unbounded-cohomology-of-open} 和
\ref{sites-cohomology-lemma-unbounded-describe-higher-direct-images}，只需证明映射
$$
H^0(\mathcal{C}, K) \longrightarrow H^0(\mathcal{C}, K_{ab})
$$
是同构。注意
$$
H^0(\mathcal{C}, K) =
\Hom_{D(\mathcal{O}_\mathcal{C})}(\mathcal{O}_\mathcal{C}, K)
$$
，另一个群同样如此。任取表示 $K$ 的 $\mathcal{O}_\mathcal{C}$-模复形
$\mathcal{K}^\bullet$。由将导出范畴构造为局部化可知
$$
\Hom_{D(\mathcal{O}_\mathcal{C})}(\mathcal{O}_\mathcal{C}, K) =
\colim_{s : \mathcal{F}^\bullet \to \mathcal{O}_\mathcal{C}}
\Hom_{K(\mathcal{O}_\mathcal{C})}(\mathcal{F}^\bullet, \mathcal{K}^\bullet)
$$
其中余极限取遍 $\mathcal{O}_\mathcal{C}$-模复形的拟同构 $s$。类似地，有
$$
\Hom_{D(\underline{\mathbf{Z}}_\mathcal{C})}
(\underline{\mathbf{Z}}_\mathcal{C}, K) =
\colim_{s : \mathcal{G}^\bullet \to \underline{\mathbf{Z}}_\mathcal{C}}
\Hom_{K(\underline{\mathbf{Z}}_\mathcal{C})}
(\mathcal{G}^\bullet, \mathcal{K}^\bullet)
$$
接着注意到，拟同构
$s : \mathcal{G}^\bullet \to \underline{\mathbf{Z}}_\mathcal{C}$
其中 $\mathcal{G}^\bullet$ 为平坦的 $\underline{\mathbf{Z}}_\mathcal{C}$-模的上有界复形，
在该系统中共尾。（这由《站点上的模》引理
\ref{sites-modules-lemma-module-quotient-flat} 和《导出范畴》引理
\ref{derived-lemma-subcategory-left-resolution} 得到；见第 \ref{sites-cohomology-section-flat} 节的讨论。）
因此可以构造映射
$H^0(\mathcal{C}, K) \longrightarrow H^0(\mathcal{C}, K_{ab})$
的逆映射：将 $\xi \in H^0(\mathcal{C}, K_{ab})$ 表示为一对
$$
(s : \mathcal{G}^\bullet \to \underline{\mathbf{Z}}_\mathcal{C},
a : \mathcal{G}^\bullet \to \mathcal{K}^\bullet)
$$
其中 $\mathcal{G}^\bullet$ 是平坦的 $\underline{\mathbf{Z}}_\mathcal{C}$-模的上有界复形，
并将其送到
$$
(\mathcal{G}^\bullet \otimes_{\underline{\mathbf{Z}}_\mathcal{C}}
\mathcal{O}_\mathcal{C}
\to \mathcal{O}_\mathcal{C},
\mathcal{G}^\bullet  \otimes_{\underline{\mathbf{Z}}_\mathcal{C}}
\mathcal{O}_\mathcal{C}
\to \mathcal{K}^\bullet)
$$
这里只需注意，第一个箭头由引理
\ref{sites-cohomology-lemma-derived-tor-quasi-isomorphism-other-side} 和
\ref{sites-cohomology-lemma-bounded-flat-K-flat} 是拟同构。具体验证该构造确为逆映射的细节从略。
\end{proof}

\begin{lemma}
\label{sites-cohomology-lemma-adjoint-lower-shriek-restrict}
设 $(\mathcal{C}, \mathcal{O})$ 为环化站点。设 $U$ 为 $\mathcal{C}$ 的对象。记
$j : (\Sh(\mathcal{C}/U), \mathcal{O}_U) \to (\Sh(\mathcal{C}), \mathcal{O})$
为相应的局部化态射。限制函子 $D(\mathcal{O}) \to D(\mathcal{O}_U)$ 是
延零函子 $j_! : D(\mathcal{O}_U) \to D(\mathcal{O})$ 的右伴随。
\end{lemma}

\begin{proof}
需证明
$$
\Hom_{D(\mathcal{O})}(j_!E, F) = \Hom_{D(\mathcal{O}_U)}(E, F|_U)
$$
选取表示 $E$ 的 $\mathcal{O}_U$-模复形 $\mathcal{E}^\bullet$，并选取表示 $F$ 的
K-内射复形 $\mathcal{I}^\bullet$。由引理
\ref{sites-cohomology-lemma-restrict-K-injective-to-open}，$\mathcal{I}^\bullet|_U$ 也是 K-内射的。
因此上式变为
$$
\Hom_{D(\mathcal{O})}(j_!\mathcal{E}^\bullet, \mathcal{I}^\bullet) =
\Hom_{D(\mathcal{O}_U)}(\mathcal{E}^\bullet, \mathcal{I}^\bullet|_U)
$$
这成立是因为 $|_U$ 与 $j_!$ 互为伴随函子
（《站点上的模》引理 \ref{sites-modules-lemma-extension-by-zero}），以及
《导出范畴》引理 \ref{derived-lemma-K-injective}。
\end{proof}

\begin{lemma}
\label{sites-cohomology-lemma-j-shriek-and-tensor}
设 $(\mathcal{C}, \mathcal{O})$ 为环化站点。设 $U \in \Ob(\mathcal{C})$。
对 $L \in D(\mathcal{O}_U)$ 和 $K \in D(\mathcal{O})$，有
$j_!L \otimes_\mathcal{O}^\mathbf{L} K =
j_!(L \otimes_{\mathcal{O}_U}^\mathbf{L} K|_U)$.
\end{lemma}

\begin{proof}
用 $\mathcal{O}_U$-模复形表示 $L$，用 $\mathcal{O}$-模 K-平坦复形表示 $K$，
并应用《站点上的模》引理 \ref{sites-modules-lemma-j-shriek-and-tensor}。
细节从略。
\end{proof}

\begin{lemma}
\label{sites-cohomology-lemma-K-injective-flat}
设 $f : (\Sh(\mathcal{C}), \mathcal{O}_\mathcal{C}) \to
(\Sh(\mathcal{D}), \mathcal{O}_\mathcal{D})$ 为环化拓扑斯之间的平坦态射。
若 $\mathcal{I}^\bullet$ 是 $\mathcal{O}_\mathcal{C}$-模的 K-内射复形，则
$f_*\mathcal{I}^\bullet$ 是 $\mathcal{O}_\mathcal{D}$-模的 K-内射复形。
\end{lemma}

\begin{proof}
这是因为
$$
\Hom_{K(\mathcal{O}_\mathcal{D})}(\mathcal{F}^\bullet, f_*\mathcal{I}^\bullet)
=
\Hom_{K(\mathcal{O}_\mathcal{C})}(f^*\mathcal{F}^\bullet, \mathcal{I}^\bullet)
$$
由《站点上的模》引理
\ref{sites-modules-lemma-adjoint-pullback-pushforward-modules} 以及 $f$ 平坦从而
$f^*$ 正合这一事实可知。
\end{proof}

\begin{lemma}
\label{sites-cohomology-lemma-hom-K-injective}
设 $\mathcal{C}$ 为站点。设 $\mathcal{O} \to \mathcal{O}'$ 为层环映射。
若 $\mathcal{I}^\bullet$ 是 $\mathcal{O}$-模的 K-内射复形，则
$\SheafHom_\mathcal{O}(\mathcal{O}', \mathcal{I}^\bullet)$
是 $\mathcal{O}'$-模的 K-内射复形。
\end{lemma}

\begin{proof}
这是因为
$\Hom_{K(\mathcal{O}')}(\mathcal{G}^\bullet,
\Hom_\mathcal{O}(\mathcal{O}', \mathcal{I}^\bullet)) =
\Hom_{K(\mathcal{O})}(\mathcal{G}^\bullet, \mathcal{I}^\bullet)$
由《站点上的模》引理 \ref{sites-modules-lemma-adjoint-hom-restrict} 成立。
\end{proof}






\section{局部化与上同调}
\label{sites-cohomology-section-localization}

\noindent
设 $\mathcal{C}$ 为站点，设 $f : X \to Y$ 为 $\mathcal{C}$ 的态射。
则得到拓扑斯态射
$$
j_{X/Y} : \Sh(\mathcal{C}/X) \longrightarrow \Sh(\mathcal{C}/Y)
$$
见《站点》第 \ref{sites-section-localize} 节和
\ref{sites-section-more-localize} 节。对于这类拓扑斯态射，某些上同调问题更容易处理。
下面给出一个得到平凡类型基变换定理的例子。

\begin{lemma}
\label{sites-cohomology-lemma-localize-cartesian-square}
设 $\mathcal{C}$ 为站点。设
$$
\xymatrix{
X' \ar[d] \ar[r] & X \ar[d] \\
Y' \ar[r] & Y
}
$$
为 $\mathcal{C}$ 的笛卡尔图。则有
$j_{Y'/Y}^{-1} \circ Rj_{X/Y, *} = Rj_{X'/Y', *} \circ j_{X'/X}^{-1}$
作为函子 $D(\mathcal{C}/X) \to D(\mathcal{C}/Y')$ 成立。
\end{lemma}

\begin{proof}
设 $E \in D(\mathcal{C}/X)$。选取表示 $E$ 的
$\mathcal{C}/X$ 上阿贝尔层的 K-内射复形 $\mathcal{I}^\bullet$。
由引理 \ref{sites-cohomology-lemma-restrict-K-injective-to-open}，$j_{X'/X}^{-1}\mathcal{I}^\bullet$
也是 K-内射的。因此可用 $j_{X'/Y', *}j_{X'/X}^{-1}\mathcal{I}^\bullet$
计算 $Rj_{X'/Y'}(j_{X'/X}^{-1}E)$。于是由《站点》引理
\ref{sites-lemma-localize-cartesian-square} 得到等式。
\end{proof}

\noindent
若 $(\mathcal{C}, \mathcal{O})$ 为环化站点且 $f : X \to Y$ 为 $\mathcal{C}$ 的态射，
则 $j_{X/Y}$ 变为环化拓扑斯态射
$$
j_{X/Y} :
(\Sh(\mathcal{C}/X), \mathcal{O}_X)
\longrightarrow
(\Sh(\mathcal{C}/Y), \mathcal{O}_Y)
$$
见《站点上的模》引理 \ref{sites-modules-lemma-relocalize}。

\begin{lemma}
\label{sites-cohomology-lemma-localize-cartesian-square-modules}
设 $(\mathcal{C}, \mathcal{O})$ 为环化站点。设
$$
\xymatrix{
X' \ar[d] \ar[r] & X \ar[d] \\
Y' \ar[r] & Y
}
$$
为 $\mathcal{C}$ 的笛卡尔图。则有
$j_{Y'/Y}^* \circ Rj_{X/Y, *} = Rj_{X'/Y', *} \circ j_{X'/X}^*$
作为函子 $D(\mathcal{O}_X) \to D(\mathcal{O}_{Y'})$ 成立。
\end{lemma}

\begin{proof}
由于 $j_{Y'/Y}^{-1}\mathcal{O}_Y = \mathcal{O}_{Y'}$，有
$j_{Y'/Y}^* = Lj_{Y'/Y}^* = j_{Y'/Y}^{-1}$。类似地，
$j_{X'/X}^* = Lj_{X'/X}^* = j_{X'/X}^{-1}$。因此由引理
\ref{sites-cohomology-lemma-modules-abelian-unbounded}，只需在阿贝尔层的导出范畴上证明结果，
而这已由引理 \ref{sites-cohomology-lemma-localize-cartesian-square} 完成。
\end{proof}












\section{逆系统与上同调}
\label{sites-cohomology-section-inverse-systems}

\noindent
本节证明关于模层逆系统的一些结果。

\begin{lemma}
\label{sites-cohomology-lemma-ML-general}
设 $I$ 为环 $A$ 的理想，设 $\mathcal{C}$ 为站点。设
$$
\ldots \to \mathcal{F}_3 \to \mathcal{F}_2 \to \mathcal{F}_1
$$
为 $\mathcal{C}$ 上 $A$-模层的逆系统，满足
$\mathcal{F}_n = \mathcal{F}_{n + 1}/I^n\mathcal{F}_{n + 1}$。
设 $p \geq 0$。假设
$$
\bigoplus\nolimits_{n \geq 0} H^{p + 1}(\mathcal{C}, I^n\mathcal{F}_{n + 1})
$$
作为分次 $\bigoplus_{n \geq 0} I^n/I^{n + 1}$-模满足升链条件。
则逆系统 $M_n = H^p(\mathcal{C}, \mathcal{F}_n)$ 满足 Mittag-Leffler 条件
\footnote{事实上，存在 $c \geq 0$，使得对所有 $n \geq c$，
$\Im(M_n \to M_{n - c})$ 都是稳定像。}。
\end{lemma}

\begin{proof}
置 $N_n = H^{p + 1}(\mathcal{C}, I^n\mathcal{F}_{n + 1})$，并令
$\delta_n : M_n \to N_n$ 为由短正合序列
$0 \to I^n\mathcal{F}_{n + 1} \to \mathcal{F}_{n + 1} \to \mathcal{F}_n \to 0$.
诱导的上同调边界映射。则 $\bigoplus \Im(\delta_n) \subset \bigoplus N_n$
是分次子模。具体地，若 $s \in M_n$ 且 $f \in I^m$，则有交换图
$$
\xymatrix{
0 \ar[r] &
I^n\mathcal{F}_{n + 1} \ar[d]_f \ar[r] &
\mathcal{F}_{n + 1} \ar[d]_f \ar[r] &
\mathcal{F}_n \ar[d]_f \ar[r] & 0 \\
0 \ar[r] &
I^{n + m}\mathcal{F}_{n + m + 1} \ar[r] &
\mathcal{F}_{n + m + 1} \ar[r] &
\mathcal{F}_{n + m} \ar[r] & 0
}
$$
中间的垂直映射通过将 $\mathcal{F}_{n + 1}$ 的局部截面提升为
$\mathcal{F}_{n + m + 1}$ 的截面，再乘以 $f$ 给出；其他垂直箭头类似。
因此 $\delta_{n + m}(fs) = f \delta_n(s)$。由假设，可找到
$s_j \in M_{n_j}$（$j = 1, \ldots, N$），使 $\delta_{n_j}(s_j)$
作为分次模生成 $\bigoplus \Im(\delta_n)$。令 $n > c = \max(n_j)$。
对 $s \in M_n$，可找到 $f_j \in I^{n - n_j}$ 使
$\delta_n(s) = \sum f_j \delta_{n_j}(s_j)$。于是
$\delta(s - \sum f_j s_j) = 0$，即存在 $s' \in M_{n + 1}$
映到 $M_n$ 中的 $s - \sum f_js_j$。从而
$$
\Im(M_{n + 1} \to M_{n - c}) = \Im(M_n \to M_{n - c})
$$
这是因为元素 $f_js_j$ 映到 $M_{n - c}$ 中为零。引理得证。
\end{proof}

\begin{lemma}
\label{sites-cohomology-lemma-ML-general-better}
设 $I$ 为环 $A$ 的理想，设 $\mathcal{C}$ 为站点。设
$$
\ldots \to \mathcal{F}_3 \to \mathcal{F}_2 \to \mathcal{F}_1
$$
为 $\mathcal{C}$ 上 $A$-模的逆系统，满足
$\mathcal{F}_n = \mathcal{F}_{n + 1}/I^n\mathcal{F}_{n + 1}$。
设 $p \geq 0$。对给定的 $n$，定义
$$
N_n =
\bigcap\nolimits_{m \geq n}
\Im\left(
H^{p + 1}(\mathcal{C}, I^n\mathcal{F}_{m + 1}) \to
H^{p + 1}(\mathcal{C}, I^n\mathcal{F}_{n + 1})
\right)
$$
若 $\bigoplus N_n$ 作为分次
$\bigoplus_{n \geq 0} I^n/I^{n + 1}$-模满足升链条件，则逆系统
$M_n = H^p(\mathcal{C}, \mathcal{F}_n)$ 满足 Mittag-Leffler 条件
\footnote{事实上，存在 $c \geq 0$，使得对所有 $n \geq c$，
$\Im(M_n \to M_{n - c})$ 都是稳定像。}。
\end{lemma}

\begin{proof}
证明与引理 \ref{sites-cohomology-lemma-ML-general} 的证明完全相同。事实上，只要证明
$\bigoplus N_n$ 是
$\bigoplus H^{p + 1}(\mathcal{C}, I^n\mathcal{F}_{n + 1})$ 的分次
$\bigoplus_{n \geq 0} I^n/I^{n + 1}$-子模，且边界映射
$\delta_n : M_n \to H^{p + 1}(\mathcal{C}, I^n\mathcal{F}_{n + 1})$
的像包含于 $N_n$，上面的论证即给出结果。

\medskip\noindent
设 $\xi \in N_n$ 且 $f \in I^k$。选取 $m \gg n + k$。选取
$\xi' \in H^{p + 1}(\mathcal{C}, I^n\mathcal{F}_{m + 1})$ 提升 $\xi$。
考虑图
$$
\xymatrix{
0 \ar[r] &
I^n\mathcal{F}_{m + 1} \ar[d]_f \ar[r] &
\mathcal{F}_{m + 1} \ar[d]_f \ar[r] &
\mathcal{F}_n \ar[d]_f \ar[r] & 0 \\
0 \ar[r] &
I^{n + k}\mathcal{F}_{m + 1} \ar[r] &
\mathcal{F}_{m + 1} \ar[r] &
\mathcal{F}_{n + k} \ar[r] & 0
}
$$
它按引理 \ref{sites-cohomology-lemma-ML-general} 的证明构造。得到上同调上的诱导映射，并注意到
$f \xi' \in H^{p + 1}(\mathcal{C}, I^{n + k}\mathcal{F}_{m + 1})$
映到 $f \xi$。由于对所有 $m \gg n + k$ 均如此，故 $f\xi \in N_{n + k}$，如所需。

\medskip\noindent
为证明边界映射 $\delta_n$ 的像包含于 $N_n$，考虑图
$$
\xymatrix{
0 \ar[r] &
I^n\mathcal{F}_{m + 1} \ar[d] \ar[r] &
\mathcal{F}_{m + 1} \ar[d] \ar[r] &
\mathcal{F}_n \ar[d] \ar[r] & 0 \\
0 \ar[r] &
I^n\mathcal{F}_{n + 1} \ar[r] &
\mathcal{F}_{n + 1} \ar[r] &
\mathcal{F}_n \ar[r] & 0
}
$$
其中 $m \geq n$。考察上同调上的诱导映射即可得出结论。
\end{proof}

\begin{lemma}
\label{sites-cohomology-lemma-topology-I-adic-general}
设 $I$ 为环 $A$ 的理想，设 $\mathcal{C}$ 为站点。设
$$
\ldots \to \mathcal{F}_3 \to \mathcal{F}_2 \to \mathcal{F}_1
$$
为 $\mathcal{C}$ 上 $A$-模层的逆系统，满足
$\mathcal{F}_n = \mathcal{F}_{n + 1}/I^n\mathcal{F}_{n + 1}$.
设 $p \geq 0$。假设
$$
\bigoplus\nolimits_{n \geq 0} H^p(\mathcal{C}, I^n\mathcal{F}_{n + 1})
$$
作为分次
$\bigoplus_{n \geq 0} I^n/I^{n + 1}$-模满足升链条件。则 $M = \lim H^p(\mathcal{C}, \mathcal{F}_n)$
上的极限拓扑就是 $I$-进拓扑。
\end{lemma}

\begin{proof}
对 $n \geq 1$ 置 $F^n = \Ker(M \to H^p(\mathcal{C}, \mathcal{F}_n))$，并置
$F^0 = M$。
注意到 $I F^n \subset F^{n + 1}$。特别地，$I^n M \subset F^n$。
因此 $I$-进拓扑比极限拓扑精细。反过来，我们将证明：对给定的 $n$，
存在 $m \geq n$ 使得 $F^m \subset I^nM$\footnote{事实上，
存在 $c \geq 0$，使得对所有 $n$ 都有 $F^{n + c} \subset I^nM$。}。
我们有单射
$$
F^n/F^{n + 1} \longrightarrow H^p(\mathcal{C}, \mathcal{F}_{n + 1})
$$
其像包含于
$H^p(\mathcal{C}, I^n\mathcal{F}_{n + 1}) \to
H^p(\mathcal{C}, \mathcal{F}_{n + 1})$.
记
$$
E_n \subset H^p(\mathcal{C}, I^n\mathcal{F}_{n + 1})
$$
为 $F^n/F^{n + 1}$ 的逆像。于是 $\bigoplus E_n$ 是
$\bigoplus H^p(\mathcal{C}, I^n\mathcal{F}_{n + 1})$ 的分次
$\bigoplus I^n/I^{n + 1}$-子模，并且
$\bigoplus E_n \to \bigoplus F^n/F^{n + 1}$ 是分次模同态；细节略去。
由假设，$\bigoplus E_n$ 在 $\bigoplus I^n/I^{n + 1}$ 上由有限个齐次元素生成。
由于 $E_n \to F^n/F^{n + 1}$ 是满射，可知
$\bigoplus F^n/F^{n + 1}$ 也具有同样性质。
因此可以找到 $r$、$c_1, \ldots, c_r \geq 0$ 以及
$a_i \in F^{c_i}$，使它们在 $\bigoplus F^n/F^{n + 1}$ 中的像生成该模。
置 $c = \max(c_i)$。

\medskip\noindent
对 $n \geq c$，我们断言 $I F^n = F^{n + 1}$。该断言表明
$F^{n + c} = I^nF^c \subset I^nM$，正如所需。为证明该断言，
设 $a \in F^{n + 1}$。$a$ 在 $F^{n + 1}/F^{n + 2}$ 中的像是
我们的 $a_i$ 的线性组合。因此，对某些 $f_i \in I^{n + 1 - c_i}$，有
$a  - \sum f_i a_i \in F^{n + 2}$。由于 $n \geq c_i$，
$I^{n + 1 - c_i} = I \cdot I^{n - c_i}$，所以可以写成
$f_i = \sum g_{i, j} h_{i, j}$，其中 $g_{i, j} \in I$
且 $h_{i, j}a_i \in F^n$。于是
$F^{n + 1} = F^{n + 2} + IF^n$。
简单的归纳论证给出，对所有 $e > 0$ 都有 $F^{n + 1} = F^{n + e} + IF^n$。
由此可知 $IF^n$ 在 $F^{n + 1}$ 中稠密。
取 $I$ 的生成元 $k_1, \ldots, k_r$，并考虑连续映射
$$
u : (F^n)^{\oplus r} \longrightarrow F^{n + 1},\quad
(x_1, \ldots, x_r) \mapsto \sum k_i x_i
$$
（在极限拓扑下）。
由上述论证，对所有 $m \geq n$，$u$ 在 $(F^m)^{\oplus r}$ 上的像在
$F^{m + 1}$ 中稠密。由《更多代数》引理
\ref{more-algebra-lemma-open-mapping}（开映射引理），可知 $u$ 是开映射。
因此 $u$ 满射，从而当 $n \geq c$ 时 $IF^n = F^{n + 1}$。
证明完毕。
\end{proof}

\begin{lemma}
\label{sites-cohomology-lemma-topology-I-adic-general-better}
设 $I$ 为环 $A$ 的理想，设 $\mathcal{C}$ 为站点。设
$$
\ldots \to \mathcal{F}_3 \to \mathcal{F}_2 \to \mathcal{F}_1
$$
为 $\mathcal{C}$ 上 $A$-模层的逆系统，满足
$\mathcal{F}_n = \mathcal{F}_{n + 1}/I^n\mathcal{F}_{n + 1}$.
设 $p \geq 0$。对给定的 $n$ 定义
$$
N_n =
\bigcap\nolimits_{m \geq n}
\Im\left(
H^p(\mathcal{C}, I^n\mathcal{F}_{m + 1}) \to
H^p(\mathcal{C}, I^n\mathcal{F}_{n + 1})
\right)
$$
若 $\bigoplus N_n$ 作为分次
$\bigoplus_{n \geq 0} I^n/I^{n + 1}$-模满足升链条件，则
$M = \lim H^p(\mathcal{C}, \mathcal{F}_n)$ 上的极限拓扑就是 $I$-进拓扑。
\end{lemma}

\begin{proof}
证明与引理 \ref{sites-cohomology-lemma-topology-I-adic-general} 的证明完全相同。
事实上，只要证明
$\bigoplus N_n$ 是 $\bigoplus H^{p + 1}(\mathcal{C}, I^n\mathcal{F}_{n + 1})$
的分次 $\bigoplus_{n \geq 0} I^n/I^{n + 1}$-子模，并且
$F^n/F^{n + 1} \subset H^p(\mathcal{C}, \mathcal{F}_{n + 1})$
包含于 $N_n \to H^p(\mathcal{C}, \mathcal{F}_{n + 1})$ 的像中，
就可由那里给出的论证得到结果。在引理
\ref{sites-cohomology-lemma-ML-general-better} 的证明中，我们已经看到了关于模结构的断言。

\medskip\noindent
设 $t \in F^n$。选取元素
$s \in H^p(\mathcal{C}, I^n\mathcal{F}_{n + 1})$
使其映到 $t$ 在 $H^p(\mathcal{C}, \mathcal{F}_{n + 1})$ 中的像。
我们需要证明 $s \in N_n$。现在 $F^n$ 是映射
$M \to H^p(\mathcal{C}, \mathcal{F}_n)$ 的核，因此对所有 $m \geq n$，可将 $t$
映到某个 $t_m \in H^p(\mathcal{C}, \mathcal{F}_{m + 1})$，且该元素在
$H^p(\mathcal{C}, \mathcal{F}_n)$ 中映到零。考虑上同调序列
$$
H^{p - 1}(\mathcal{C}, \mathcal{F}_n) \to
H^p(\mathcal{C}, I^n\mathcal{F}_{m + 1}) \to
H^p(\mathcal{C}, \mathcal{F}_{m + 1}) \to
H^p(\mathcal{C}, \mathcal{F}_n)
$$
它来自短正合序列
$0 \to I^n\mathcal{F}_{m + 1} \to \mathcal{F}_{m + 1} \to \mathcal{F}_n \to 0$.
可以选取 $s_m \in H^p(\mathcal{C}, I^n\mathcal{F}_{m + 1})$
映到 $t_m$。将上述序列与 $m = n$ 时的序列比较，可知
$s_m$ 映到 $s$，相差一个如下映射的像中的元素：
$H^{p - 1}(\mathcal{C}, \mathcal{F}_n) \to
H^p(\mathcal{C}, I^n\mathcal{F}_{n + 1})$.
然而，该映射分解经过映射
$H^p(\mathcal{C}, I^n\mathcal{F}_{m + 1}) \to
H^p(\mathcal{C}, I^n\mathcal{F}_{n + 1})$
，因此可知 $s$ 属于所需的像。
\end{proof}












\section{导出极限与同伦极限}
\label{sites-cohomology-section-derived-limits}

\noindent
设 $\mathcal{C}$ 为站点。考虑范畴 $\mathcal{C} \times \mathbf{N}$，其中
$\Mor((U, n), (V, m)) = \emptyset$（若 $n > m$），而
$\Mor((U, n), (V, m)) = \Mor(U, V)$（其他情形）。通过规定覆盖为满足
$\{U_i \to U\}$ 是 $\mathcal{C}$ 的覆盖的族
$\{(U_i, n) \to (U, n)\}$，我们给它赋予站点结构。
读者立即验证，$\mathcal{C} \times \mathbf{N}$ 上的层
正好等同于 $\mathcal{C}$ 上层的逆系统。
特别地，$\textit{Ab}(\mathcal{C} \times \mathbf{N})$
就是 $\mathcal{C}$ 上阿贝尔层的逆系统范畴。
现在考虑函子
$$
\lim : \textit{Ab}(\mathcal{C} \times \mathbf{N}) \to \textit{Ab}(\mathcal{C})
$$
它将逆系统取为其极限。这不过就是 $g_*$，其中
$g : \Sh(\mathcal{C} \times \mathbf{N}) \to \Sh(\mathcal{C})$
是与连续且余连续函子 $\mathcal{C} \times \mathbf{N} \to \mathcal{C}$
相伴的拓扑斯态射。（注意，$g^{-1}$ 将 $\mathcal{C}$ 上的层赋给
相应的常值逆系统。）

\medskip\noindent
由上面说明的一般构造，我们得到导出函子
$$
R\lim = Rg_* : D(\mathcal{C} \times \mathbf{N}) \to D(\mathcal{C}).
$$
如所示，该函子通常记作 $R\lim$。

\medskip\noindent
另一方面，连续且余连续函子
$\mathcal{C} \to \mathcal{C} \times \mathbf{N}$、$U \mapsto (U, n)$
定义了拓扑斯态射
$i_n : \Sh(\mathcal{C}) \to \Sh(\mathcal{C} \times \mathbf{N})$。
当然，$i_n^{-1}$ 是取逆系统第 $n$ 项的函子。因此有函子变换
$i_{n + 1}^{-1} \to i_n^{-1}$。于是给定
$K \in D(\mathcal{C} \times \mathbf{N})$，我们得到
$K_n = i_n^{-1}K \in D(\mathcal{C})$ 以及映射 $K_{n + 1} \to K_n$。
在《导出范畴》定义 \ref{derived-definition-derived-limit} 中，
我们定义了同伦极限的概念
$$
R\lim K_n \in D(\mathcal{C})
$$
我们断言这两个概念一致（在有意义的范围内）。

\begin{lemma}
\label{sites-cohomology-lemma-derived-limit-is-ok}
设 $\mathcal{C}$ 为站点，设 $K$ 为
$D(\mathcal{C} \times \mathbf{N})$ 的对象。按上面的方式置
$K_n = i_n^{-1}K$。则
$$
R\lim K \cong R\lim K_n
$$
在 $D(\mathcal{C})$ 中成立。
\end{lemma}

\begin{proof}
为在 $D(\mathcal{C} \times \mathbf{N})$ 的对象 $K$ 上计算 $R\lim$，
我们选取一个 K-内射表示 $\mathcal{I}^\bullet$，其各项是
$\textit{Ab}(\mathcal{C} \times \mathbf{N})$ 中的内射对象；见《内射对象》定理
\ref{injectives-theorem-K-injective-embedding-grothendieck}.
我们可以并且确实把 $\mathcal{I}^\bullet$ 看作复形
$(\mathcal{I}_n^\bullet)$ 的逆系统，于是得到
$$
R\lim K = \lim \mathcal{I}_n^\bullet
$$
其中右端是逐项逆极限。

\medskip\noindent
设 $\mathcal{J} = (\mathcal{J}_n)$ 是
$\textit{Ab}(\mathcal{C} \times \mathbf{N})$ 中的内射对象。态射
$(U, n) \to (U, n + 1)$ 是 $\mathcal{C} \times \mathbf{N}$ 中的单态射，故
$\mathcal{J}(U, n + 1) \to \mathcal{J}(U, n)$ 是满射
（引理 \ref{sites-cohomology-lemma-restriction-along-monomorphism-surjective}）。
因此 $\mathcal{J}_{n + 1} \to \mathcal{J}_n$ 作为预层映射是满射。

\medskip\noindent
注意函子 $i_n^{-1}$ 有一个正合左伴随 $i_{n, !}$。
具体地，$i_{n, !}\mathcal{F}$ 是逆系统
$\ldots 0 \to 0 \to \mathcal{F} \to \ldots \to \mathcal{F}$.
因此，由《导出范畴》引理
\ref{derived-lemma-adjoint-preserve-K-injectives}，复形
$i_n^{-1}\mathcal{I}^\bullet = \mathcal{I}_n^\bullet$ 是 K-内射的。

\medskip\noindent
由于我们选取的 K-内射复形各项都是内射的，我们得到
$$
0 \to  \lim \mathcal{I}_n^\bullet \to \prod \mathcal{I}_n^\bullet
\to \prod \mathcal{I}_n^\bullet \to 0
$$
是阿贝尔层复形的短正合序列，因为它也是阿贝尔预层复形的短正合序列。
此外，中间和右端的乘积表示 $D(\mathcal{C})$ 中的乘积；见
《内射对象》引理 \ref{injectives-lemma-derived-products} 及其证明
（这里正是使用 $\mathcal{I}_n^\bullet$ 为 K-内射的地方）。
因此，按三角范畴中同伦极限的定义，$R\lim K$ 是逆系统 $(K_n)$ 的同伦极限。
\end{proof}

\begin{lemma}
\label{sites-cohomology-lemma-RGamma-commutes-with-Rlim}
设 $(\mathcal{C}, \mathcal{O})$ 为环化站点。函子
$R\Gamma(\mathcal{C}, -)$ 以及对 $U \in \Ob(\mathcal{C})$ 的
$R\Gamma(U, -)$ 都与 $R\lim$ 交换。此外，对
$$
0 \to
R^1\lim H^{m - 1}(U, K_n) \to H^m(U, R\lim K_n) \to
\lim H^m(U, K_n) \to 0
$$
对 $D(\mathcal{O})$ 中任意逆系统 $(K_n)$ 及 $m \in \mathbf{Z}$，都有短正合序列。
对 $H^m(\mathcal{C}, R\lim K_n)$ 也有类似结论。
\end{lemma}

\begin{proof}
第一个断言由《内射对象》引理
\ref{injectives-lemma-RF-commutes-with-Rlim} 得到。然后可将《更多代数》备注
\ref{more-algebra-remark-compare-derived-limit} 应用于
$R\lim R\Gamma(U, K_n) = R\Gamma(U, R\lim K_n)$，从而得到这些短正合序列。
\end{proof}

\begin{lemma}
\label{sites-cohomology-lemma-Rf-commutes-with-Rlim}
设 $f : (\Sh(\mathcal{C}), \mathcal{O}) \to (\Sh(\mathcal{C}'), \mathcal{O}')$
为环化拓扑斯态射。则 $Rf_*$ 与 $R\lim$ 交换，即
$Rf_*$ 与导出极限交换。
\end{lemma}

\begin{proof}
设 $(K_n)$ 为 $D(\mathcal{O})$ 中对象的逆系统。对 $n$ 归纳，
可以选取 $\mathcal{O}$-模的实际复形 $\mathcal{K}_n^\bullet$ 以及复形映射
$\mathcal{K}_{n + 1}^\bullet \to \mathcal{K}_n^\bullet$，它们在
$D(\mathcal{O})$ 中表示映射 $K_{n + 1} \to K_n$。换言之，存在
$D(\mathcal{C} \times \mathbf{N})$ 中的对象 $K$，其对应的逆系统就是给定的逆系统。
接下来考虑交换图
$$
\xymatrix{
\Sh(\mathcal{C} \times \mathbf{N}) \ar[r]_g \ar[d]_{f \times 1} &
\Sh(\mathcal{C}) \ar[d]_f \\
\Sh(\mathcal{C}' \times \mathbf{N}) \ar[r]^{g'} &
\Sh(\mathcal{C}')
}
$$
其中各箭头为拓扑斯态射。于是
$R\lim R(f \times 1)_*K = Rf_* R\lim K$。逐一展开定义并使用引理
\ref{sites-cohomology-lemma-derived-limit-is-ok}，得到
$R\lim (Rf_*K_n) = Rf_*(R\lim K_n)$。

\medskip\noindent
当 $\mathcal{C}$ 有足够多的点时的另一证明。考虑定义性的有区别三角
$$
R\lim K_n \to \prod K_n \to \prod K_n
$$
它位于 $D(\mathcal{O})$ 中。应用正合函子 $Rf_*$，得到有区别三角
$$
Rf_*(R\lim K_n) \to Rf_*\left(\prod K_n\right) \to Rf_*\left(\prod K_n\right)
$$
它位于 $D(\mathcal{O}')$ 中。因此只需证明 $Rf_*$ 与导出范畴中的乘积交换
即可（这些乘积不只是复形的乘积，见《内射对象》引理
\ref{injectives-lemma-derived-products}）。然而，由于按引理
\ref{sites-cohomology-lemma-adjoint}，$Rf_*$ 是右伴随，这一点形式上成立（见《范畴》引理
\ref{categories-lemma-adjoint-exact}）。
注意：不能直接应用《范畴》引理
\ref{categories-lemma-adjoint-exact}，因为 $R\lim K_n$ 不是
$D(\mathcal{O})$ 中的极限。
\end{proof}

\begin{remark}
\label{sites-cohomology-remark-discuss-derived-limit}
设 $(\mathcal{C}, \mathcal{O})$ 为环化站点，设 $(K_n)$ 为
$D(\mathcal{O})$ 中的逆系统。置 $K = R\lim K_n$。对每个 $n$ 和 $m$，
令 $\mathcal{H}^m_n = H^m(K_n)$ 为 $K_n$ 的第 $m$ 个上同调层，并类似地置
$\mathcal{H}^m = H^m(K)$。记 $\underline{\mathcal{H}}^m_n$ 为预层
$$
U \longmapsto \underline{\mathcal{H}}^m_n(U) = H^m(U, K_n)
$$
类似地置 $\underline{\mathcal{H}}^m(U) = H^m(U, K)$。
由引理 \ref{sites-cohomology-lemma-sheafification-cohomology} 可知，
$\mathcal{H}^m_n$ 是 $\underline{\mathcal{H}}^m_n$ 的层化，
$\mathcal{H}^m$ 是 $\underline{\mathcal{H}}^m$ 的层化。
有如下交换图
$$
\xymatrix{
K \ar@{=}[d] &
\underline{\mathcal{H}}^m \ar[d] \ar[r] & 
\mathcal{H}^m \ar[d] \\
R\lim K_n &
\lim \underline{\mathcal{H}}^m_n \ar[r] & 
\lim \mathcal{H}^m_n
}
$$
一般而言，$\lim \mathcal{H}^m_n$ 不一定是
$\lim \underline{\mathcal{H}}^m_n$ 的层化。
若 $U \in \mathcal{C}$，则有短正合序列
\begin{equation}
\label{sites-cohomology-equation-ses-Rlim-over-U}
0 \to
R^1\lim \underline{\mathcal{H}}^{m - 1}_n(U) \to
\underline{\mathcal{H}}^m(U) \to
\lim \underline{\mathcal{H}}^m_n(U) \to 0
\end{equation}
这是由引理 \ref{sites-cohomology-lemma-RGamma-commutes-with-Rlim} 得到的。
\end{remark}

\noindent
下面的引理适用于代数空间或代数叠上的拟凝聚模逆系统，且其过渡映射为满射。

\begin{lemma}
\label{sites-cohomology-lemma-inverse-limit-is-derived-limit}
设 $(\mathcal{C}, \mathcal{O})$ 为环化站点，设 $(\mathcal{F}_n)$ 为
$\mathcal{O}$-模的逆系统。设 $\mathcal{B} \subset \Ob(\mathcal{C})$ 为子集。假设
\begin{enumerate}
\item $\mathcal{C}$ 的每个对象都有一个覆盖，其成员属于 $\mathcal{B}$；
\item 对 $p > 0$ 及 $U \in \mathcal{B}$，有 $H^p(U, \mathcal{F}_n) = 0$；
\item 对 $U \in \mathcal{B}$，逆系统 $\mathcal{F}_n(U)$ 的 $R^1\lim$ 消失。
\end{enumerate}
则 $R\lim \mathcal{F}_n = \lim \mathcal{F}_n$，且对 $p > 0$ 及
$U \in \mathcal{B}$，有 $H^p(U, \lim \mathcal{F}_n) = 0$。
\end{lemma}

\begin{proof}
置 $K_n = \mathcal{F}_n$，并置 $K = R\lim \mathcal{F}_n$。使用备注
\ref{sites-cohomology-remark-discuss-derived-limit} 的记号以及假设 (2)，可知对 $U \in \mathcal{B}$，
当 $m \not = 0$ 时 $\underline{\mathcal{H}}_n^m(U) = 0$，而
$\underline{\mathcal{H}}_n^0(U) = \mathcal{F}_n(U)$。
由等式（\ref{sites-cohomology-equation-ses-Rlim-over-U}）及假设 (3)，可知当 $m \not = 0$ 时
$\underline{\mathcal{H}}^m(U) = 0$，而当 $m = 0$ 时它等于
$\lim \mathcal{F}_n(U)$。利用 (1) 进行层化，得到当 $m \not = 0$ 时
$\mathcal{H}^m = 0$，而当 $m = 0$ 时它等于 $\lim \mathcal{F}_n$。
因此 $K = \lim \mathcal{F}_n$。由于对 $m > 0$ 有
$H^m(U, K) = \underline{\mathcal{H}}^m(U) = 0$（见上文），
故第二个断言成立。
\end{proof}

\begin{lemma}
\label{sites-cohomology-lemma-cohomology-derived-limit-injective}
设 $(\mathcal{C}, \mathcal{O})$ 为环化站点，设 $(K_n)$ 为
$D(\mathcal{O})$ 中的逆系统。设 $V \in \Ob(\mathcal{C})$ 且
$m \in \mathbf{Z}$。假设存在整数 $n(V)$ 以及 $V$ 的覆盖的共终系统
$\text{Cov}_V$，使得对 $\{V_i \to V\} \in \text{Cov}_V$，有
\begin{enumerate}
\item $R^1\lim H^{m - 1}(V_i, K_n) = 0$；
\item 当 $n \geq n(V)$ 时，$H^m(V_i, K_n) \to H^m(V_i, K_{n(V)})$
是单射。
\end{enumerate}
则截面上的映射 $H^m(R\lim K_n)(V) \to H^m(K_{n(V)})(V)$ 是单射。
\end{lemma}

\begin{proof}
设 $\gamma \in H^m(R\lim K_n)(V)$ 映到
$H^m(K_{n(V)})(V)$ 中的零元。由于 $H^m(R\lim K_n)$ 是预层
$U \mapsto H^m(U, R\lim K_n)$ 的层化（由引理
\ref{sites-cohomology-lemma-sheafification-cohomology}），可以选取 $\{V_i \to V\} \in \text{Cov}_V$
以及元素 $\tilde\gamma_i \in H^m(V_i, R\lim K_n)$，使其映到
$\gamma|_{V_i}$。
于是 $\tilde\gamma_i$ 映到 $\tilde\gamma_{i, n(V)} \in H^m(V_i, K_{n(V)})$。
由于 $H^m(K_{n(V)})$ 是预层 $U \mapsto H^m(U, K_{n(V)})$ 的层化
（再次使用引理 \ref{sites-cohomology-lemma-sheafification-cohomology}），可知将
$\{V_i \to V\}$ 替换为其细化后，可以假设对所有 $i$ 都有
$\tilde\gamma_{i, n(V)} = 0$。
对该覆盖，考虑引理 \ref{sites-cohomology-lemma-RGamma-commutes-with-Rlim} 给出的短正合序列
$$
0 \to
R^1\lim H^{m - 1}(V_i, K_n) \to H^m(V_i, R\lim K_n) \to
\lim H^m(V_i, K_n) \to 0
$$
。由假设 (1)，左端的群为零；由假设 (2)，右端的群到
$H^m(V_i, K_{n(V)})$ 的映射为单射。因此 $\tilde\gamma_i = 0$，
从而 $\gamma = 0$，正如所需。
\end{proof}

\begin{lemma}
\label{sites-cohomology-lemma-is-limit-per-object}
设 $(\mathcal{C}, \mathcal{O})$ 为环化站点，设 $E \in D(\mathcal{O})$。
设 $\mathcal{B} \subset \Ob(\mathcal{C})$ 为子集。假设
\begin{enumerate}
\item $\mathcal{C}$ 的每个对象都有一个覆盖，其成员属于 $\mathcal{B}$；
\item 对每个 $V \in \mathcal{B}$，存在函数
$p(V, -) : \mathbf{Z} \to \mathbf{Z}$ 以及 $V$ 的覆盖的共终系统 $\text{Cov}_V$，使得
$$
H^p(V_i, H^{m - p}(E)) = 0
$$
对所有 $\{V_i \to V\} \in \text{Cov}_V$ 及满足 $p > p(V, m)$ 的整数 $p,m$ 都成立。
\end{enumerate}
则《导出范畴》备注
\ref{derived-remark-map-into-derived-limit-truncations} 中的映射
$E \to R\lim \tau_{\geq -n} E$ 在 $D(\mathcal{O})$ 中是同构。
\end{lemma}

\begin{proof}
置 $K_n = \tau_{\geq -n}E$，并置 $K = R\lim K_n$。
典范映射 $E \to K$ 来自典范映射 $E \to K_n = \tau_{\geq -n}E$。
我们需要证明 $E \to K$ 在上同调层上诱导同构
$H^m(E) \to H^m(K)$。在证明的其余部分固定 $m$。若 $n \geq -m$，则
$E \to \tau_{\geq -n}E = K_n$ 诱导同构 $H^m(E) \to H^m(K_n)$。
要完成证明，只需证明对每个 $V \in \mathcal{B}$，存在整数 $n(V) \geq -m$，使映射
 $H^m(K)(V) \to H^m(K_{n(V)})(V)$ 为单射。事实上，此时复合映射
$$
H^m(E)(V) \to H^m(K)(V) \to H^m(K_{n(V)})(V)
$$
是双射，而第二个箭头是单射，故第一个箭头是双射。由性质 (1) 可知
$H^m(E) \to H^m(K)$ 是同构。置
$$
n(V) = 1 + \max\{-m, p(V, m - 1) - m, -1 + p(V, m) - m, -2 + p(V, m + 1) - m\}.
$$
从而无论如何都有 $n(V) \geq -m$。断言：当
$$
H^{m - 1}(V_i, K_{n + 1}) \to H^{m - 1}(V_i, K_n)
\quad\text{且}\quad
H^m(V_i, K_{n + 1}) \to H^m(V_i, K_n)
$$
对 $n \geq n(V)$ 且 $\{V_i \to V\} \in \text{Cov}_V$ 时，上述映射都是同构。
该断言意味着引理 \ref{sites-cohomology-lemma-cohomology-derived-limit-injective}
中的条件 (1) 和 (2) 得到满足，因而推出所需的单射性。
回忆《导出范畴》备注
\ref{derived-remark-truncation-distinguished-triangle}，我们有有区别三角
$$
H^{-n - 1}(E)[n + 1] \to
K_{n + 1} \to K_n \to H^{-n - 1}(E)[n + 2]
$$
考察相应的上同调长正合序列可知，只要
$$
H^{m + n}(V_i, H^{-n - 1}(E)),\quad
H^{m + n + 1}(V_i, H^{-n - 1}(E)),\quad
H^{m + n + 2}(V_i, H^{-n - 1}(E))
$$
在 $n \geq n(V)$ 且 $\{V_i \to V\} \in \text{Cov}_V$ 时为零，断言就成立。
这由我们对 $n(V)$ 的选择以及引理中的假设推出。
\end{proof}

\begin{lemma}
\label{sites-cohomology-lemma-is-limit-spaltenstein}
设 $(\mathcal{C}, \mathcal{O})$ 为环化站点，设 $E \in D(\mathcal{O})$。
设 $\mathcal{B} \subset \Ob(\mathcal{C})$ 为子集。假设
\begin{enumerate}
\item $\mathcal{C}$ 的每个对象都有一个覆盖，其成员属于 $\mathcal{B}$；
\item 对每个 $V \in \mathcal{B}$，存在整数 $d_V \geq 0$ 以及 $V$ 的覆盖的共终系统
$\text{Cov}_V$，使得
$$
H^p(V_i, H^q(E)) = 0 \text{（对 }
\{V_i \to V\} \in \text{Cov}_V,\ p > d_V, \text{ 且 }q < 0\text{）}
$$
\end{enumerate}
则《导出范畴》备注
\ref{derived-remark-map-into-derived-limit-truncations} 中的映射
$E \to R\lim \tau_{\geq -n} E$ 在 $D(\mathcal{O})$ 中是同构。
\end{lemma}

\begin{proof}
这由引理 \ref{sites-cohomology-lemma-is-limit-per-object}（取
$p(V, m) = d_V + \max(0, m)$）推出。
\end{proof}

\begin{lemma}
\label{sites-cohomology-lemma-is-limit}
设 $(\mathcal{C}, \mathcal{O})$ 为环化站点，设 $E \in D(\mathcal{O})$。
假设存在函数 $p(-) : \mathbf{Z} \to \mathbf{Z}$ 及子集
$\mathcal{B} \subset \Ob(\mathcal{C})$，使得
\begin{enumerate}
\item $\mathcal{C}$ 的每个对象都有一个覆盖，其成员属于 $\mathcal{B}$；
\item 对 $p > p(m)$ 及 $V \in \mathcal{B}$，有
$H^p(V, H^{m - p}(E)) = 0$。
\end{enumerate}
则《导出范畴》备注
\ref{derived-remark-map-into-derived-limit-truncations} 中的映射
$E \to R\lim \tau_{\geq -n} E$ 在 $D(\mathcal{O})$ 中是同构。
\end{lemma}

\begin{proof}
应用引理 \ref{sites-cohomology-lemma-is-limit-per-object}，取 $p(V, m) = p(m)$，并令
$\text{Cov}_V$ 等于所有满足对每个 $i$ 都有 $V_i \in \mathcal{B}$ 的
覆盖 $\{V_i \to V\}$ 所成的集合。
\end{proof}

\begin{lemma}
\label{sites-cohomology-lemma-is-limit-dimension}
设 $(\mathcal{C}, \mathcal{O})$ 为环化站点，设 $E \in D(\mathcal{O})$。
假设存在整数 $d \geq 0$ 及子集 $\mathcal{B} \subset \Ob(\mathcal{C})$，使得
\begin{enumerate}
\item $\mathcal{C}$ 的每个对象都有一个覆盖，其成员属于 $\mathcal{B}$；
\item 对 $p > d$、$q < 0$ 及 $V \in \mathcal{B}$，有
$H^p(V, H^q(E)) = 0$。
\end{enumerate}
则《导出范畴》备注
\ref{derived-remark-map-into-derived-limit-truncations} 中的映射
$E \to R\lim \tau_{\geq -n} E$ 在 $D(\mathcal{O})$ 中是同构。
\end{lemma}

\begin{proof}
应用引理 \ref{sites-cohomology-lemma-is-limit-spaltenstein}，取 $d_V = d$，并令
$\text{Cov}_V$ 等于所有满足对每个 $i$ 都有 $V_i \in \mathcal{B}$ 的
覆盖 $\{V_i \to V\}$ 所成的集合。
\end{proof}

\noindent
上述引理可用于在某些情形计算上同调。

\begin{lemma}
\label{sites-cohomology-lemma-cohomology-over-U-trivial}
设 $(\mathcal{C}, \mathcal{O})$ 为环化站点，设 $K$
为 $D(\mathcal{O})$ 的对象。设 $\mathcal{B} \subset \Ob(\mathcal{C})$ 为子集。假设
\begin{enumerate}
\item $\mathcal{C}$ 的每个对象都有一个覆盖，其成员属于 $\mathcal{B}$；
\item 对所有 $p > 0$、$q \in \mathbf{Z}$ 及 $U \in \mathcal{B}$，有
$H^p(U, H^q(K)) = 0$。
\end{enumerate}
则对 $q \in \mathbf{Z}$ 及 $U \in \mathcal{B}$，有
$H^q(U, K) = H^0(U, H^q(K))$。
\end{lemma}

\begin{proof}
注意，由引理 \ref{sites-cohomology-lemma-is-limit-dimension}（取 $d = 0$），有
$K = R\lim \tau_{\geq -n} K$。令 $U \in \mathcal{B}$。由等式
（\ref{sites-cohomology-equation-ses-Rlim-over-U}）得到短正合序列
$$
0 \to R^1\lim H^{q - 1}(U, \tau_{\geq -n}K) \to
H^q(U, K) \to \lim H^q(U, \tau_{\geq -n}K) \to 0
$$
由《导出范畴》引理
\ref{derived-lemma-two-ss-complex-functor} 的谱序列，条件 (2) 蕴含对所有 $q$ 有
$H^q(U, \tau_{\geq -n} K) = H^0(U, H^q(\tau_{\geq -n} K))$。
该谱序列收敛，因为 $\tau_{\geq -n}K$ 下有界。若 $n > -q$，则
$H^q(\tau_{\geq -n}K) = H^q(K)$。因此，所显示短正合序列左端和右端的系统
最终分别恒定为 $H^0(U, H^{q - 1}(K))$ 和 $H^0(U, H^q(K))$，引理得证。
\end{proof}

\noindent
下面是另一种可以描述导出极限的情形。

\begin{lemma}
\label{sites-cohomology-lemma-derived-limit-suitable-system}
设 $(\mathcal{C}, \mathcal{O})$ 为环化站点。设 $(K_n)$
为 $D(\mathcal{O})$ 中对象的逆系统。设 $\mathcal{B} \subset \Ob(\mathcal{C})$ 为子集。假设
\begin{enumerate}
\item $\mathcal{C}$ 的每个对象都有一个覆盖，其成员属于 $\mathcal{B}$；
\item 对所有 $U \in \mathcal{B}$ 及 $q \in \mathbf{Z}$，都有
\begin{enumerate}
\item 对 $p > 0$，有 $H^p(U, H^q(K_n)) = 0$；
\item 逆系统 $H^0(U, H^q(K_n))$ 的 $R^1\lim$ 消失。
\end{enumerate}
\end{enumerate}
则对 $q \in \mathbf{Z}$，有 $H^q(R\lim K_n) = \lim H^q(K_n)$。
\end{lemma}

\begin{proof}
置 $K = R\lim K_n$。我们使用备注
\ref{sites-cohomology-remark-discuss-derived-limit} 中的记号。令 $U \in \mathcal{B}$。
由引理 \ref{sites-cohomology-lemma-cohomology-over-U-trivial} 及 (2)(a)，有
$H^q(U, K_n) = H^0(U, H^q(K_n))$。
由于函子 $R\Gamma(U, -)$ 与导出极限交换，有
$$
H^q(U, K) = H^q(R\lim R\Gamma(U, K_n)) = \lim H^0(U, H^q(K_n))
$$
其中最后一个等号由《更多代数》备注
\ref{more-algebra-remark-compare-derived-limit} 及假设 (2)(b) 得到。因此 $H^q(U, K)$
是层 $H^q(K_n)$ 在 $U$ 上截面的逆极限。由于 $\lim H^q(K_n)$ 是一个层，
利用假设 (1) 可知，$H^q(K)$（它是预层 $U \mapsto H^q(U, K)$ 的层化）
等于 $\lim H^q(K_n)$。引理得证。
\end{proof}







\section{构造 K-内射分辨}
\label{sites-cohomology-section-K-injective}

\noindent
设 $(\mathcal{C}, \mathcal{O})$ 为环化站点。设 $\mathcal{F}^\bullet$ 为
$\mathcal{O}$-模复形。范畴 $\textit{Mod}(\mathcal{O})$ 有足够多的内射对象，
因此可使用《导出范畴》引理 \ref{derived-lemma-special-inverse-system}
构造如下图
$$
\xymatrix{
\ldots \ar[r] &
\tau_{\geq -2}\mathcal{F}^\bullet \ar[r] \ar[d] &
\tau_{\geq -1}\mathcal{F}^\bullet \ar[d] \\
\ldots \ar[r] & \mathcal{I}_2^\bullet \ar[r] & \mathcal{I}_1^\bullet
}
$$
它位于 $\mathcal{O}$-模复形范畴中，并满足
\begin{enumerate}
\item 竖直箭头是拟同构；
\item $\mathcal{I}_n^\bullet$ 是下有界的内射对象复形；
\item 箭头 $\mathcal{I}_{n + 1}^\bullet \to \mathcal{I}_n^\bullet$
是逐项可裂满射。
\end{enumerate}
 $\mathcal{O}$-模范畴有极限（在预层层面计算），因此可以形成逐项极限
$\mathcal{I}^\bullet = \lim_n \mathcal{I}_n^\bullet$。由《导出范畴》引理
\ref{derived-lemma-bounded-below-injectives-K-injective} 以及
\ref{derived-lemma-limit-K-injectives}
这是 K-内射复形。一般而言，典范映射
\begin{equation}
\label{sites-cohomology-equation-into-candidate-K-injective}
\mathcal{F}^\bullet \to \mathcal{I}^\bullet
\end{equation}
不一定是拟同构。下面的引理给出使其成为拟同构的一些条件。

\begin{lemma}
\label{sites-cohomology-lemma-K-injective}
在上述情形下，记
$\mathcal{H}^m = H^m(\mathcal{F}^\bullet)$ 为第 $m$ 个上同调层。
设 $\mathcal{B} \subset \Ob(\mathcal{C})$ 为子集，设 $d \in \mathbf{N}$。假设
\begin{enumerate}
\item $\mathcal{C}$ 的每个对象都有一个覆盖，其成员属于 $\mathcal{B}$；
\item 对每个 $U \in \mathcal{B}$，当 $p > d$ 且 $q < 0$ 时有
$H^p(U, \mathcal{H}^q) = 0$\footnote{只要
$\forall m$、$\exists p(m)$，且当 $p > p(m)$ 时有
$H^p(U. \mathcal{H}^{m - p}) = 0$ 即可；见引理
\ref{sites-cohomology-lemma-is-limit}。}。
\end{enumerate}
则（\ref{sites-cohomology-equation-into-candidate-K-injective}）是拟同构。
\end{lemma}

\begin{proof}
由《导出范畴》引理 \ref{derived-lemma-difficulty-K-injectives}，只需证明映射
$\mathcal{F}^\bullet \to R\lim \tau_{\geq -n} \mathcal{F}^\bullet$
是同构。这由引理 \ref{sites-cohomology-lemma-is-limit-dimension} 得到。
\end{proof}

\noindent
下面是关于模复形逆系统逐项极限的上同调层的技术引理。我们应尽可能避免使用此引理，
而应改用导出逆极限的论证。

\begin{lemma}
\label{sites-cohomology-lemma-inverse-limit-complexes}
设 $(\mathcal{C}, \mathcal{O})$ 为环化站点。设
$(\mathcal{F}_n^\bullet)$ 为 $\mathcal{O}$-模复形的逆系统。设 $m \in \mathbf{Z}$。
给定 $\mathcal{B} \subset \Ob(\mathcal{C})$ 及整数 $n_0$，使得
\begin{enumerate}
\item $\mathcal{C}$ 的每个对象都有一个覆盖，其成员属于 $\mathcal{B}$；
\item 对每个 $U \in \mathcal{B}$，
\begin{enumerate}
\item 阿贝尔群系统
$\mathcal{F}_n^{m - 2}(U)$ 和 $\mathcal{F}_n^{m - 1}(U)$ 的 $R^1\lim$ 消失
（例如，它们具有 Mittag-Leffler 性质）；
\item 阿贝尔群系统 $H^{m - 1}(\mathcal{F}_n^\bullet(U))$ 的 $R^1\lim$ 消失
（例如，它具有 Mittag-Leffler 性质）；
\item 有
$H^m(\mathcal{F}_n^\bullet(U)) = H^m(\mathcal{F}_{n_0}^\bullet(U))$
对所有 $n \geq n_0$ 都成立。
\end{enumerate}
\end{enumerate}
则映射 $H^m(\mathcal{F}^\bullet) \to \lim H^m(\mathcal{F}_n^\bullet)
\to H^m(\mathcal{F}_{n_0}^\bullet)$ 是层的同构，其中
$\mathcal{F}^\bullet = \lim \mathcal{F}_n^\bullet$ 是逐项逆极限。
\end{lemma}

\begin{proof}
令 $U \in \mathcal{B}$。注意 $H^m(\mathcal{F}^\bullet(U))$ 是下列复形
的上同调：
$$
\lim_n \mathcal{F}_n^{m - 2}(U) \to
\lim_n \mathcal{F}_n^{m - 1}(U) \to
\lim_n \mathcal{F}_n^m(U) \to
\lim_n \mathcal{F}_n^{m + 1}(U)
$$
中从左起第三个位置。由假设 (2)(a) 和 (2)(b)，可应用《更多代数》引理
\ref{more-algebra-lemma-apply-Mittag-Leffler-again}，从而得到
$$
H^m(\mathcal{F}^\bullet(U)) = \lim H^m(\mathcal{F}_n^\bullet(U))
$$
由假设 (2)(c)，得到
$$
H^m(\mathcal{F}^\bullet(U)) = H^m(\mathcal{F}_n^\bullet(U))
$$
对所有 $n \geq n_0$ 成立。由假设 (1)，预层
$U \mapsto H^m(\mathcal{F}^\bullet(U))$ 的层化等于对所有 $n \geq n_0$ 的预层
$U \mapsto H^m(\mathcal{F}_n^\bullet(U))$ 的层化。因此层的逆系统
$H^m(\mathcal{F}_n^\bullet)$ 在 $n \geq n_0$ 时恒定，值为
$H^m(\mathcal{F}^\bullet)$，引理得证。
\end{proof}







\section{有界上同调维数}
\label{sites-cohomology-section-bounded}

\noindent
本节讨论函子 $Rf_*$ 何时具有有界上同调维数。
当考虑无界复形时，这是一个相当微妙的问题。

\begin{situation}
\label{sites-cohomology-situation-olsson-laszlo}
设 $\mathcal{C}$ 为站点，设 $\mathcal{O}$ 为 $\mathcal{C}$ 上的层环。
设 $\mathcal{A} \subset \textit{Mod}(\mathcal{O})$ 为弱 Serre 子范畴。
我们假设以下条件成立：存在子集 $\mathcal{B} \subset \Ob(\mathcal{C})$，使得
\begin{enumerate}
\item $\mathcal{C}$ 的每个对象都有一个覆盖，其成员属于 $\mathcal{B}$；
\item 对每个 $V \in \mathcal{B}$，存在整数 $d_V$ 及 $V$ 的覆盖的共终系统
$\text{Cov}_V$，使得
$$
H^p(V_i, \mathcal{F}) = 0 \text{（对 }
\{V_i \to V\} \in \text{Cov}_V,\ p > d_V, \text{ 且 }
\mathcal{F} \in \Ob(\mathcal{A})
\text{）}
$$
\end{enumerate}
\end{situation}

\begin{lemma}
\label{sites-cohomology-lemma-olsson-laszlo}
\begin{reference}
这是带有略作修改假设的 \cite[Proposition 2.1.4]{six-I}；它是站点情形中
\cite[Proposition 3.13]{Spaltenstein} 的类似结果。
\end{reference}
在情形 \ref{sites-cohomology-situation-olsson-laszlo} 中，对任意
$E \in D_\mathcal{A}(\mathcal{O})$，《导出范畴》备注
\ref{derived-remark-map-into-derived-limit-truncations} 中的映射
$E \to R\lim \tau_{\geq -n} E$ 在 $D(\mathcal{O})$ 中是同构。
\end{lemma}

\begin{proof}
立即由引理 \ref{sites-cohomology-lemma-is-limit-spaltenstein} 得到。
\end{proof}

\begin{lemma}
\label{sites-cohomology-lemma-olsson-laszlo-modified}
在情形 \ref{sites-cohomology-situation-olsson-laszlo} 中，设
$(K_n)$ 为 $D_\mathcal{A}^+(\mathcal{O})$ 中的逆系统。
假设对每个 $j$，$\mathcal{A}$ 中的逆系统 $(H^j(K_n))$
最终恒定为 $\mathcal{H}^j$。则对所有 $j$，有
$H^j(R\lim K_n) = \mathcal{H}^j$。
\end{lemma}

\begin{proof}
设 $V \in \mathcal{B}$。设 $\{V_i \to V\}$ 属于情形
\ref{sites-cohomology-situation-olsson-laszlo} 中的集合 $\text{Cov}_V$。
由于 $K_n$ 下有界，有谱序列
$$
E_2^{p, q} = H^p(V_i, H^q(K_n))
$$
收敛到 $H^{p + q}(V_i, K_n)$。见《导出范畴》引理
\ref{derived-lemma-two-ss-complex-functor}。由假设，当 $p > d_V$ 时
$E_2^{p, q} = 0$。取 $n_0$ 使得
$$
\begin{matrix}
\mathcal{H}^{j + 1} & = & H^{j + 1}(K_n), \\
\mathcal{H}^j & = & H^j(K_n), \\
\ldots, \\
\mathcal{H}^{j - d_V - 2} & = & H^{j - d_V - 2}(K_n)
\end{matrix}
$$
对所有 $n \geq n_0$ 成立。比较上述关于 $K_n$ 和 $K_{n_0}$ 的谱序列，
可知当 $n \geq n_0$ 时，上同调群 $H^{j - 1}(V_i, K_n)$ 和 $H^j(V_i, K_n)$
与 $n$ 无关。因此对充分大的 $n$（取决于 $V$），截面上的映射
$H^j(R\lim K_n)(V) \to H^j(K_n)(V)$ 是单射，见引理
\ref{sites-cohomology-lemma-cohomology-derived-limit-injective}。
由于 $\mathcal{C}$ 的每个对象都可由 $\mathcal{B}$ 中的元素覆盖，故映射
$H^j(R\lim K_n) \to \mathcal{H}^j$ 是单射。

\medskip\noindent
满射性以类似方式证明。具体地，取 $U \in \Ob(\mathcal{C})$ 及
$\gamma \in \mathcal{H}^j(U)$。我们希望在用覆盖的成员替换 $U$ 后，
将 $\gamma$ 提升为 $H^j(R\lim K_n)$ 的截面。因此由情形
\ref{sites-cohomology-situation-olsson-laszlo} 的性质 (1)，可设 $U = V \in \mathcal{B}$。
取 $n_0$ 使得
$$
\begin{matrix}
\mathcal{H}^{j + 1} & = & H^{j + 1}(K_n), \\
\mathcal{H}^j & = & H^j(K_n), \\
\ldots, \\
\mathcal{H}^{j - d_V - 2} & = & H^{j - d_V - 2}(K_n)
\end{matrix}
$$
对所有 $n \geq n_0$ 成立。取 $\text{Cov}_V$ 中的覆盖
$\{V_i \to V\}$，使得
$\gamma|_{V_i} \in \mathcal{H}^j(V_i) = H^j(K_{n_0})(V_i)$
可提升为 $\gamma_{n_0, i} \in H^j(V_i, K_{n_0})$。
这是可能的，因为由引理 \ref{sites-cohomology-lemma-sheafification-cohomology}，
$H^j(K_{n_0})$ 是预层 $U \mapsto H^j(U, K_{n_0})$ 的层化。
由证明第一段的讨论，$H^{j - 1}(V_i, K_n)$ 和 $H^j(V_i, K_n)$
在 $n \geq n_0$ 时与 $n$ 无关。因此由引理
\ref{sites-cohomology-lemma-RGamma-commutes-with-Rlim}，$\gamma_{n_0, i}$ 可提升为
$\gamma_i \in H^j(V_i, R\lim K_n)$。证明完毕。
\end{proof}

\begin{lemma}
\label{sites-cohomology-lemma-olsson-laszlo-map-version-one}
\begin{reference}
这是带有略作修改假设的 \cite[Lemma 2.1.10]{six-I} 的一个版本。
\end{reference}
设 $f : (\Sh(\mathcal{C}), \mathcal{O}) \to (\Sh(\mathcal{C}'), \mathcal{O}')$
为环化拓扑斯态射。设 $\mathcal{A} \subset \textit{Mod}(\mathcal{O})$
和 $\mathcal{A}' \subset \textit{Mod}(\mathcal{O}')$ 为弱 Serre 子范畴。
假设存在整数 $N$ 使得
\begin{enumerate}
\item $\mathcal{C}, \mathcal{O}, \mathcal{A}$ 满足情形
\ref{sites-cohomology-situation-olsson-laszlo} 的假设；
\item $\mathcal{C}', \mathcal{O}', \mathcal{A}'$ 满足情形
\ref{sites-cohomology-situation-olsson-laszlo} 的假设；
\item 对 $p \geq 0$ 及 $\mathcal{F} \in \Ob(\mathcal{A})$，有
$R^pf_*\mathcal{F} \in \Ob(\mathcal{A}')$；
\item 对 $p > N$ 及 $\mathcal{F} \in \Ob(\mathcal{A})$，有
$R^pf_*\mathcal{F} = 0$；
\end{enumerate}
则对 $D_\mathcal{A}(\mathcal{O})$ 中的 $K$，有
\begin{enumerate}
\item[(a)] $Rf_*K$ 属于 $D_{\mathcal{A}'}(\mathcal{O}')$；
\item[(b)] 当 $j \geq N - n$ 时，映射
$H^j(Rf_*K) \to H^j(Rf_*(\tau_{\geq -n}K))$ 是同构。
\end{enumerate}
\end{lemma}

\begin{proof}
由引理 \ref{sites-cohomology-lemma-olsson-laszlo}，有 $K = R\lim \tau_{\geq -n}K$。
由引理 \ref{sites-cohomology-lemma-Rf-commutes-with-Rlim}，有
$Rf_*K = R\lim Rf_*\tau_{\geq -n}K$。
复形 $Rf_*\tau_{\geq -n}K$ 下有界。谱序列
$$
E_2^{p, q} = R^pf_*H^q(\tau_{\geq -n}K)
$$
收敛到 $H^{p + q}(Rf_*\tau_{\geq -n}K)$
（《导出范畴》引理 \ref{derived-lemma-two-ss-complex-functor}）。由假设 (3)，
$Rf_*\tau_{\geq -n}K$ 属于 $D^+_{\mathcal{A}'}(\mathcal{O}')$，
见《同调代数》引理 \ref{homology-lemma-biregular-ss-converges}。
注意，当 $m \geq n$ 时映射
$$
Rf_*(\tau_{\geq -m}K) \longrightarrow Rf_*(\tau_{\geq -n}K)
$$
由上述谱序列在 $j \geq -n + N$ 的次数上诱导上同调层的同构。因此可应用引理
\ref{sites-cohomology-lemma-olsson-laszlo-modified} 得到结论。
\end{proof}

\noindent
有时我们需要上述引理的一个变体，其中的假设略有不同。

\begin{situation}
\label{sites-cohomology-situation-olsson-laszlo-prime}
设
$f : (\mathcal{C}, \mathcal{O}) \to (\mathcal{C}', \mathcal{O}')$
为环化站点态射。设 $u : \mathcal{C}' \to \mathcal{C}$
为相应的站点连续函子。设 $\mathcal{A} \subset \textit{Mod}(\mathcal{O})$
为弱 Serre 子范畴。我们假设以下条件成立：存在子集
$\mathcal{B}' \subset \Ob(\mathcal{C}')$，使得
\begin{enumerate}
\item $\mathcal{C}'$ 的每个对象都有一个覆盖，其成员属于 $\mathcal{B}'$；
\item 对每个 $V' \in \mathcal{B}'$，存在整数 $d_{V'}$ 及 $V'$ 的覆盖的共终系统
$\text{Cov}_{V'}$，使得
$$
H^p(u(V'_i), \mathcal{F}) = 0 \text{（对 }
\{V'_i \to V'\} \in \text{Cov}_{V'},\ p > d_{V'}, \text{ 且 }
\mathcal{F} \in \Ob(\mathcal{A})\text{）}
$$
\end{enumerate}
\end{situation}

\begin{lemma}
\label{sites-cohomology-lemma-olsson-laszlo-map-version-two}
\begin{reference}
这是带有略作修改假设的 \cite[Lemma 2.1.10]{six-I} 的一个版本。
\end{reference}
设 $f : (\mathcal{C}, \mathcal{O}) \to (\mathcal{C}', \mathcal{O}')$
为环化站点态射。此外假设存在整数 $N$，使得
\begin{enumerate}
\item $\mathcal{C}, \mathcal{O}, \mathcal{A}$ 满足情形
\ref{sites-cohomology-situation-olsson-laszlo} 的假设；
\item $f : (\mathcal{C}, \mathcal{O}) \to (\mathcal{C}', \mathcal{O}')$
和 $\mathcal{A}$ 满足情形 \ref{sites-cohomology-situation-olsson-laszlo-prime} 的假设；
\item 对 $p > N$ 及 $\mathcal{F} \in \Ob(\mathcal{A})$，有
$R^pf_*\mathcal{F} = 0$；
\end{enumerate}
则对 $D_\mathcal{A}(\mathcal{O})$ 中的 $K$，当 $j \geq N - n$ 时，映射
$H^j(Rf_*K) \to H^j(Rf_*(\tau_{\geq -n}K))$ 是同构。
\end{lemma}

\begin{proof}
设 $K$ 属于 $D_\mathcal{A}(\mathcal{O})$。
根据引理 \ref{sites-cohomology-lemma-olsson-laszlo}，有 $K = R\lim \tau_{\geq -n}K$。
根据引理 \ref{sites-cohomology-lemma-Rf-commutes-with-Rlim}，有
$Rf_*K = R\lim Rf_*(\tau_{\geq -n}K)$。
取 $V' \in \mathcal{B}'$，并取 $\{V'_i \to V'\}$ 为
$\text{Cov}_{V'}$ 的一个元素。则有
$$
H^j(V'_i, Rf_*K) = H^j(u(V'_i), K)
\quad\text{且}\quad
H^j(V'_i, Rf_*(\tau_{\geq -n}K)) = H^j(u(V'_i), \tau_{\geq -n}K)
$$
情形 \ref{sites-cohomology-situation-olsson-laszlo-prime} 中的假设意味着，当 $n$ 足够大时，
最后一个群与 $n$ 无关；所需的 $n$ 只依赖于 $j$ 和 $d_{V'}$。细节略去。
对 $j$ 和 $j - 1$ 应用此结论，并通过引理
\ref{sites-cohomology-lemma-RGamma-commutes-with-Rlim} 得到
$$
H^j(V'_i, Rf_*K) = \lim H^j(V'_i, Rf_*(\tau_{\geq -n} K))
$$
且当 $n$ 大于某个仅依赖于 $d_{V'}$ 和 $j$ 的常数时，右侧系统为恒定系统。
因此引理 \ref{sites-cohomology-lemma-cohomology-derived-limit-injective} 蕴含
$$
H^j(Rf_*K)(V') \longrightarrow
\left(\lim H^j(Rf_*(\tau_{\geq -n}K))\right)(V')
$$
为单射。由于 $V' \in \mathcal{B}'$ 中的对象覆盖 $\mathcal{C}'$ 的每个对象，
我们得到映射
$H^j(Rf_*K) \to \lim H^j(Rf_*(\tau_{\geq -n}K))$
为单射。谱序列
$$
E_2^{p, q} = R^pf_*H^q(\tau_{\geq -n}K)
$$
收敛到 $H^{p + q}(Rf_*(\tau_{\geq -n}K))$
（导出范畴，引理 \ref{derived-lemma-two-ss-complex-functor}），
结合假设 (3) 可知
$H^j(Rf_*(\tau_{\geq -n}K))$ 在 $n \geq N - j$ 时为恒定群。
因此，当 $j \geq N - n$ 时，
$H^j(Rf_*K) \to H^j(Rf_*(\tau_{\geq -n}K))$ 为单射。

\medskip\noindent
于是我们已经证明，将引理最后一行的 ``isomorphism'' 替换为 ``injective''
后命题成立。现在取满足 $j \geq N - n$ 的 $j$ 和 $n$，考虑有区别三角
$$
\tau_{\leq -n - 1}K \to K \to \tau_{\geq -n}K \to (\tau_{\leq -n - 1}K)[1]
$$
参见导出范畴，注记 \ref{derived-remark-truncation-distinguished-triangle}。
由于 $\tau_{\geq -n}\tau_{\leq -n -1}K = 0$，对
$\tau_{-n - 1}K$ 已证明的单射性给出
$$
0 = H^j(Rf_*(\tau_{\leq -n - 1}K)) =
H^{j + 1}(Rf_*(\tau_{\leq -n - 1}K)) =
H^{j + 2}(Rf_*(\tau_{\leq -n - 1}K)) = \ldots
$$
由与该有区别三角相伴的上同调长正合序列
$$
Rf_*(\tau_{\leq -n - 1}K) \to Rf_*K \to Rf_*(\tau_{\geq -n}K) \to
Rf_*(\tau_{\leq -n - 1}K)[1]
$$
推出 $H^j(Rf_*K) \to H^j(Rf_*(\tau_{\geq -n}K))$
为同构。
\end{proof}








\section{Mayer-Vietoris}
\label{sites-cohomology-section-mayer-vietoris}

\noindent
通常的 Mayer--Vietoris 陈述及证明请参见上同调，第
\ref{cohomology-section-mayer-vietoris} 节。

\medskip\noindent
设 $(\mathcal{C}, \mathcal{O})$ 为环化站点。考虑范畴
$\mathcal{C}$ 中的交换图
$$
\xymatrix{
E \ar[d] \ar[r] & Y \ar[d] \\
Z \ar[r] & X
}
$$
在此情形，给定 $D(\mathcal{O})$ 中的对象 $K$，我们得到一个看似有区别三角
的开头
$$
R\Gamma(X, K) \to
R\Gamma(Z, K) \oplus
R\Gamma(Y, K) \to
R\Gamma(E, K)
$$
下面的引理将这一点说得更精确。

\begin{lemma}
\label{sites-cohomology-lemma-c-square}
在上述情形，选取表示 $K$ 的 $\mathcal{O}$-模 K-内射复形
$\mathcal{I}^\bullet$。将四个箭头之一的典范映射乘以 $-1$，得到复形映射
$$
\mathcal{I}^\bullet(X) \xrightarrow{\alpha}
\mathcal{I}^\bullet(Z) \oplus
\mathcal{I}^\bullet(Y) \xrightarrow{\beta}
\mathcal{I}^\bullet(E)
$$
且 $\beta \circ \alpha = 0$。于是有典范映射
$$
c^K_{X, Z, Y, E} :
\mathcal{I}^\bullet(X)
\longrightarrow
C(\beta)^\bullet[-1]
$$
所谓典范，是指表示 $K$ 的 K-内射复形取不同选择时，在阿贝尔群的导出范畴中
得到同构的箭头。若 $c^K_{X, Z, Y, E}$ 为同构，则利用其逆映射得到典范的
有区别三角
$$
R\Gamma(X, K) \to
R\Gamma(Z, K) \oplus
R\Gamma(Y, K) \to
R\Gamma(E, K) \to
R\Gamma(X, K)[1]
$$
所有这些构造关于 $K$ 都是函子性的。
\end{lemma}

\begin{proof}
此引理是自明的。例如，若 $\mathcal{J}^\bullet$ 是表示 $K$ 的另一个
K-内射复形，则可选取拟同构
$\mathcal{I}^\bullet \to \mathcal{J}^\bullet$
它在所涉及的全部复形之间确定拟同构。细节略去。关于锥的构造以及它与
有区别三角的关系，参见导出范畴，第
\ref{derived-section-cones} 节和
\ref{derived-section-homotopy-triangulated} 节。
\end{proof}

\begin{lemma}
\label{sites-cohomology-lemma-two-out-of-three-blow-up-square}
在上述情形，设
$K_1 \to K_2 \to K_3 \to K_1[1]$ 为 $D(\mathcal{O})$ 中的有区别三角。
若对 $\{1, 2, 3\}$ 中的两个 $i$，$c^{K_i}_{X, Z, Y, E}$ 是拟同构，
则对第三个 $i$ 它也是拟同构。
\end{lemma}

\begin{proof}
旋转该三角后，不妨设 $c^{K_1}_{X, Z, Y, E}$ 和
$c^{K_2}_{X, Z, Y, E}$ 为拟同构。选取表示 $K_1 \to K_2$ 的
$\mathcal{O}$-模 K-内射复形之间的映射
$f : \mathcal{I}^\bullet_1 \to \mathcal{I}^\bullet_2$。
则 $K_3$ 由 K-内射复形 $C(f)^\bullet$ 表示，参见
导出范畴，引理 \ref{derived-lemma-triangle-K-injective}。
于是态射 $c^{K_3}_{X, Z, Y, E}$ 是同构，因为它是
$K(\textit{Ab})$ 中有区别三角映射的第三条边，而另外两条边是拟同构。
细节略去；使用导出范畴，引理
\ref{derived-lemma-third-isomorphism-triangle}。
\end{proof}

\noindent
下面给出何时上述构造确实产生有区别三角的判据。

\begin{lemma}
\label{sites-cohomology-lemma-square-triangle}
在上述情形，假设
\begin{enumerate}
\item $h_X^\# = h_Y^\# \amalg_{h_E^\#} h_Z^\#$，且
\item $h_E^\# \to h_Y^\#$ 是单射。
\end{enumerate}
则引理 \ref{sites-cohomology-lemma-c-square} 的构造产生有区别三角
$$
R\Gamma(X, K) \to
R\Gamma(Z, K) \oplus
R\Gamma(Y, K) \to
R\Gamma(E, K) \to R\Gamma(X, K)[1]
$$
且对 $D(\mathcal{C})$ 中的 $K$ 是函子性的。
\end{lemma}

\begin{proof}
我们可以用各项均为内射阿贝尔层的 K-内射复形表示 $K$，参见
第 \ref{sites-cohomology-section-unbounded} 节。因此只需证明：若 $\mathcal{I}$ 是内射阿贝尔
层，则
$$
0 \to \mathcal{I}(X) \to
\mathcal{I}(Z) \oplus \mathcal{I}(Y) \to
\mathcal{I}(E) \to 0
$$
是短正合序列。第一个箭头是单射，因为由条件 (1)，映射
$h_Y \amalg h_Z \to h_X$ 层化后成为满射；这意味着
$\{Y \to X, Z \to X\}$ 可以由 $X$ 的一个覆盖细化。
最后一个箭头是满射，因为 $\mathcal{I}(Y) \to \mathcal{I}(E)$ 是满射。
具体地，有
$\mathcal{I}(E) = \Hom(\mathbf{Z}_E^\#, \mathcal{I})$,
$\mathcal{I}(Y) = \Hom(\mathbf{Z}_Y^\#, \mathcal{I})$,
由 (2)，映射 $\mathbf{Z}_E^\# \to \mathbf{Z}_Y^\#$ 是单射，且
$\mathcal{I}$ 是内射阿贝尔层。请与《站点上的模》，第
\ref{sites-modules-section-free-abelian-sheaf} 节相比较。
最后，设 $s \in \mathcal{I}(Y)$ 且
$t \in \mathcal{F}(Z)$ 映射到 $\mathcal{I}(E)$ 中的同一元素。
则 $s$ 和 $t$ 定义映射
$$
s \amalg t : h_Y^\# \amalg h_Z^\# \longrightarrow \mathcal{I}
$$
该映射按假设可分解经过 $h_Y^\# \amalg_{h_E^\#} h_Z^\#$。
于是由假设 (1) 得到唯一映射
$h_X^\# \to \mathcal{I}$，它对应于 $\mathcal{I}(X)$ 中的一个元素，
其在 $Y$ 上限制为 $s$，在 $Z$ 上限制为 $t$。
\end{proof}

\begin{lemma}
\label{sites-cohomology-lemma-square-triangle-general}
设 $\mathcal{C}$ 为站点。考虑交换图
$$
\xymatrix{
\mathcal{D} \ar[r] \ar[d] & \mathcal{F} \ar[d] \\
\mathcal{E} \ar[r] & \mathcal{G}
}
$$
其由 $\mathcal{C}$ 上的集合预层组成，并假设
\begin{enumerate}
\item $\mathcal{G}^\# =
\mathcal{E}^\# \amalg_{\mathcal{D}^\#} \mathcal{F}^\#$，且
\item $\mathcal{D}^\# \to \mathcal{F}^\#$ 是单射。
\end{enumerate}
则存在典范有区别三角
$$
R\Gamma(\mathcal{G}, K) \to
R\Gamma(\mathcal{E}, K) \oplus
R\Gamma(\mathcal{F}, K) \to
R\Gamma(\mathcal{D}, K) \to
R\Gamma(\mathcal{G}, K)[1]
$$
它关于 $K \in D(\mathcal{C})$ 是函子性的，其中 $R\Gamma(\mathcal{G}, -)$
是第 \ref{sites-cohomology-section-limp} 节所讨论的上同调。
\end{lemma}

\begin{proof}
由于层化是正合的，且
$R\Gamma(\mathcal{G}, -) = R\Gamma(\mathcal{G}^\#, -)$，
不妨假设 $\mathcal{D}, \mathcal{E}, \mathcal{F}, \mathcal{G}$ 是集合层。
此外，上同调 $R\Gamma(\mathcal{G}, -)$ 只依赖于拓扑斯，而不依赖于底层站点。
因此由
站点，引理 \ref{sites-lemma-topos-good-site}
可以将 $\mathcal{C}$ 替换为具有次典范拓扑的一个较大站点，
使得 $\mathcal{G} = h_X$、$\mathcal{F} = h_Y$、$\mathcal{E} = h_Z$ 且
$\mathcal{D} = h_E$，其中 $X, Y, Z, E$ 是 $\mathcal{C}$ 的对象。
此时结论由引理 \ref{sites-cohomology-lemma-square-triangle} 得出。
\end{proof}







\section{比较两个拓扑}
\label{sites-cohomology-section-compare}

\noindent
设 $\mathcal{C}$ 为范畴。令
$\text{Cov}(\mathcal{C}) \supset \text{Cov}'(\mathcal{C})$
为两种赋予 $\mathcal{C}$ 站点结构的方式。记 $\tau$ 为对应于
$\text{Cov}(\mathcal{C})$ 的拓扑，记 $\tau'$ 为对应于
$\text{Cov}'(\mathcal{C})$ 的拓扑。于是 $\mathcal{C}$ 上的恒等函子定义站点态射
$$
\epsilon : \mathcal{C}_\tau \longrightarrow \mathcal{C}_{\tau'}
$$
其中 $\epsilon_*$ 是底层预层上的恒等函子，而 $\epsilon^{-1}$ 是
$\tau'$-层的 $\tau$-层化。参见站点，例
\ref{sites-example-finer-topology} 和
\ref{sites-example-finer-topology-bis}。
在上述情形中有以下结论：
\begin{enumerate}
\item $\epsilon_* : \Sh(\mathcal{C}_\tau) \to \Sh(\mathcal{C}_{\tau'})$
是完全忠实的，且 $\epsilon^{-1} \circ \epsilon_* = \text{id}$；
\item $\epsilon_* : \textit{Ab}(\mathcal{C}_\tau) \to
\textit{Ab}(\mathcal{C}_{\tau'})$
是完全忠实的，且 $\epsilon^{-1} \circ \epsilon_* = \text{id}$；
\item $R\epsilon_* : D(\mathcal{C}_\tau) \to D(\mathcal{C}_{\tau'})$
是完全忠实的，且 $\epsilon^{-1} \circ R\epsilon_* = \text{id}$；
\item 若 $\mathcal{O}$ 是 $\tau$-拓扑下的环层，则 $\mathcal{O}$ 也是
$\tau'$-拓扑下的层，且 $\epsilon$ 成为环化站点的平坦态射
$$
\epsilon :
(\mathcal{C}_\tau, \mathcal{O}_\tau)
\longrightarrow
(\mathcal{C}_{\tau'}, \mathcal{O}_{\tau'})
$$
\item $\epsilon_* : \textit{Mod}(\mathcal{O}_\tau) \to
\textit{Mod}(\mathcal{O}_{\tau'})$
是完全忠实的，且 $\epsilon^* \circ \epsilon_* = \text{id}$；
\item $R\epsilon_* : D(\mathcal{O}_\tau) \to D(\mathcal{O}_{\tau'})$
是完全忠实的，且 $\epsilon^* \circ R\epsilon_* = \text{id}$。
\end{enumerate}
说明如下。

\medskip\noindent
(1)。设 $\mathcal{F}$ 是 $\tau$-拓扑下的集合层。
那么 $\epsilon_*\mathcal{F}$ 只是将 $\mathcal{F}$ 看作 $\tau'$-拓扑下的层。
应用 $\epsilon^{-1}$ 意味着对 $\mathcal{F}$ 作 $\tau$-层化；由于
$\mathcal{F}$ 已经是 $\tau$-层，这一操作不产生任何改变。因此
$\epsilon^{-1}(\epsilon_*\mathcal{F})) = \mathcal{F}$.
完全忠实性由范畴，引理
\ref{categories-lemma-adjoint-fully-faithful}.

\medskip\noindent
(2)。这是 (1) 的推论，因为阿贝尔层的拉回和推前与对底层集合层
进行相同操作是一回事。

\medskip\noindent
(3)。设 $K$ 是 $D(\mathcal{C}_\tau)$ 中的对象。为计算
$R\epsilon_*K$，选取表示 $K$ 的 K-内射复形
$\mathcal{I}^\bullet$，并令 $R\epsilon_*K = \epsilon_*\mathcal{I}^\bullet$。
由于 $\epsilon^{-1} : D(\mathcal{C}_{\tau'}) \to D(\mathcal{C}_\tau)$
对对象 $L$ 的计算方式是对表示 $L$ 的任意阿贝尔层复形应用正合函子
$\epsilon^{-1}$，故 $\epsilon^{-1}R\epsilon_*K$ 由
$\epsilon^{-1}\epsilon_*\mathcal{I}^\bullet$ 表示。由
(1) 有
$\mathcal{I}^\bullet = \epsilon^{-1}\epsilon_*\mathcal{I}^\bullet$.
换言之，$\epsilon^{-1} \circ R\epsilon_* = \text{id}$，结论如前。

\medskip\noindent
(4)。注意 $\epsilon^{-1}\mathcal{O}_{\tau'} = \mathcal{O}_\tau$，参见
(1) 中的讨论。因此 $\epsilon$ 是环化站点的平坦态射，参见
《站点上的模》，定义 \ref{sites-modules-definition-flat-morphism}。
不仅如此，显然在 $\mathcal{O}_{\tau'}$-模上有
$\epsilon^* = \epsilon^{-1}$（作为模的拉回与底层阿贝尔层的拉回具有相同
的底层阿贝尔层）。

\medskip\noindent
(5)。这由 (2) 及 (4) 中的论述立即得到。

\medskip\noindent
(6)。这与 (3) 类似，细节略去。












\section{上同调下降的形式性}
\label{sites-cohomology-section-formal-cohomological-descent}

\noindent
本节只讨论环化拓扑斯的态射在多大程度上确定从下方导出范畴
到上方导出范畴的嵌入。下面是一个典型结果。

\begin{lemma}
\label{sites-cohomology-lemma-downstairs}
设 $f : (\Sh(\mathcal{C}), \mathcal{O}_\mathcal{C}) \to
(\Sh(\mathcal{D}), \mathcal{O}_\mathcal{D})$ 为环化拓扑斯的态射。
考虑 $D(\mathcal{O}_\mathcal{D})$ 的满子范畴
$D' \subset D(\mathcal{O}_\mathcal{D})$，其对象 $K$ 满足
$$
K \longrightarrow Rf_*Lf^*K
$$
是同构。则 $D'$ 是 $D(\mathcal{O}_\mathcal{D})$ 的饱和三角严格满子范畴，
且函子 $Lf^* : D' \to D(\mathcal{O}_\mathcal{C})$ 完全忠实。
\end{lemma}

\begin{proof}
关于本情形中“饱和”的定义，参见导出范畴，定义
\ref{derived-definition-saturated}。关于三角子范畴，参见
导出范畴，引理 \ref{derived-lemma-triangulated-subcategory}。
引理中的典范映射是伴随函子对 $(Lf^*, Rf_*)$ 的单位，参见引理
\ref{sites-cohomology-lemma-adjoint}。有了这一点，$D'$ 为饱和三角子范畴的证明略去；
它形式上来自 $Lf^*$ 和 $Rf_*$ 是三角范畴之间的正合函子这一事实。
最后一部分形式上来自 $Lf^*$ 与 $Rf_*$ 互为伴随；可与范畴，引理
\ref{categories-lemma-adjoint-fully-faithful} 比较。
\end{proof}

\begin{lemma}
\label{sites-cohomology-lemma-upstairs}
设 $f : (\Sh(\mathcal{C}), \mathcal{O}_\mathcal{C}) \to
(\Sh(\mathcal{D}), \mathcal{O}_\mathcal{D})$ 为环化拓扑斯的态射。
考虑 $D(\mathcal{O}_\mathcal{C})$ 的满子范畴
$D' \subset D(\mathcal{O}_\mathcal{C})$，其对象 $K$ 满足
$$
Lf^*Rf_*K \longrightarrow K
$$
是同构。则 $D'$ 是 $D(\mathcal{O}_\mathcal{C})$ 的饱和三角严格满子范畴，
且函子 $Rf_* : D' \to D(\mathcal{O}_\mathcal{D})$ 完全忠实。
\end{lemma}

\begin{proof}
关于本情形中“饱和”的定义，参见导出范畴，定义
\ref{derived-definition-saturated}。关于三角子范畴，参见
导出范畴，引理 \ref{derived-lemma-triangulated-subcategory}。
引理中的典范映射是伴随函子对 $(Lf^*, Rf_*)$ 的余单位，参见引理
\ref{sites-cohomology-lemma-adjoint}。有了这一点，$D'$ 为饱和三角子范畴的证明略去；
它形式上来自 $Lf^*$ 和 $Rf_*$ 是三角范畴之间的正合函子这一事实。
最后一部分形式上来自 $Lf^*$ 与 $Rf_*$ 互为伴随；可与范畴，引理
\ref{categories-lemma-adjoint-fully-faithful} 比较。
\end{proof}

\begin{lemma}
\label{sites-cohomology-lemma-bounded-in-image-upstairs}
设 $f : (\Sh(\mathcal{C}), \mathcal{O}_\mathcal{C}) \to
(\Sh(\mathcal{D}), \mathcal{O}_\mathcal{D})$ 为环化拓扑斯的态射。
设 $K$ 是 $D(\mathcal{O}_\mathcal{C})$ 中的对象。假设
\begin{enumerate}
\item $f$ 是平坦的；
\item $K$ 下有界；
\item $f^*Rf_*H^q(K) \to H^q(K)$ 是同构。
\end{enumerate}
则 $f^*Rf_*K \to K$ 是同构。
\end{lemma}

\begin{proof}
注意，$f^*Rf_*K \to K$ 是同构，当且仅当它在上同调层 $H^j$ 上是同构。
注意
$H^j(f^*Rf_*K) = f^*H^j(Rf_*K) = f^*H^j(Rf_*\tau_{\leq j}K) =
H^j(f^*Rf_*\tau_{\leq j}K)$.
因此不妨假设 $K$ 有界。于是性质 (3) 告诉我们，上同调层属于
引理 \ref{sites-cohomology-lemma-upstairs} 中的三角子范畴
$D' \subset D(\mathcal{O}_\mathcal{C})$。因此 $K$ 也属于该子范畴。
\end{proof}

\begin{lemma}
\label{sites-cohomology-lemma-bounded-in-image-downstairs}
设 $f : (\Sh(\mathcal{C}), \mathcal{O}_\mathcal{C}) \to
(\Sh(\mathcal{D}), \mathcal{O}_\mathcal{D})$ 为环化拓扑斯的态射。
设 $K$ 是 $D(\mathcal{O}_\mathcal{D})$ 中的对象。假设
\begin{enumerate}
\item $f$ 是平坦的；
\item $K$ 下有界；
\item $H^q(K) \to Rf_*f^*H^q(K)$ 是同构。
\end{enumerate}
则 $K \to Rf_*f^*K$ 是同构。
\end{lemma}

\begin{proof}
注意，$K \to Rf_*f^*K$ 是同构，当且仅当它在上同调层 $H^j$ 上是同构。
注意
$H^j(Rf_*f^*K) = H^j(Rf_*\tau_{\leq j}f^*K) = H^j(Rf_*f^*\tau_{\leq j}K)$.
因此不妨假设 $K$ 有界。于是性质 (3) 告诉我们，上同调层属于
引理 \ref{sites-cohomology-lemma-downstairs} 中的三角子范畴
$D' \subset D(\mathcal{O}_\mathcal{D})$。因此 $K$ 也属于该子范畴。
\end{proof}

\begin{lemma}
\label{sites-cohomology-lemma-equivalence-bounded}
设 $f : (\Sh(\mathcal{C}), \mathcal{O}) \to (\Sh(\mathcal{C}'), \mathcal{O}')$
为环化拓扑斯的态射。设 $\mathcal{A} \subset \textit{Mod}(\mathcal{O})$
及 $\mathcal{A}' \subset \textit{Mod}(\mathcal{O}')$ 为弱 Serre 子范畴。假设
\begin{enumerate}
\item $f$ 是平坦的；
\item $f^*$ 诱导范畴等价
$\mathcal{A}' \to \mathcal{A}$；
\item 对 $\mathcal{F}' \in \Ob(\mathcal{A}')$，映射
$\mathcal{F}' \to Rf_*f^*\mathcal{F}'$ 是同构。
\end{enumerate}
则
$f^* : D_{\mathcal{A}'}^+(\mathcal{O}') \to D_\mathcal{A}^+(\mathcal{O})$
是范畴等价，其拟逆为
$Rf_* : D_\mathcal{A}^+(\mathcal{O}) \to D_{\mathcal{A}'}^+(\mathcal{O}')$.
\end{lemma}

\begin{proof}
由假设 (2)、(3) 以及引理
\ref{sites-cohomology-lemma-bounded-in-image-downstairs} 和 \ref{sites-cohomology-lemma-downstairs}，可知
$f^* : D_{\mathcal{A}'}^+(\mathcal{O}') \to D_\mathcal{A}^+(\mathcal{O})$
是完全忠实的。对 $\mathcal{F} \in \Ob(\mathcal{A})$，可写成
$\mathcal{F} = f^*\mathcal{F}'$。于是
$Rf_*\mathcal{F} = Rf_* f^*\mathcal{F}' = \mathcal{F}'$.
特别地，当 $p > 0$ 时 $R^pf_*\mathcal{F} = 0$，且
$f_*\mathcal{F} \in \Ob(\mathcal{A}')$。因此对任意
$K \in D^+_\mathcal{A}(\mathcal{O})$，利用收敛到 $R^{p + q}f_*K$ 的谱序列
$E_2^{p, q} = R^pf_*H^q(K)$，可知
$Rf_*K \in D^+_{\mathcal{A}'}(\mathcal{O}')$。
当然，由引理
\ref{sites-cohomology-lemma-bounded-in-image-upstairs} 和 \ref{sites-cohomology-lemma-upstairs} 可知
$Rf_* : D_\mathcal{A}^+(\mathcal{O}) \to D_{\mathcal{A}'}^+(\mathcal{O}')$
是完全忠实的。由于 $f^*$ 与 $Rf_*$ 互为伴随，遂得到引理结论；
例如可用范畴，引理 \ref{categories-lemma-adjoint-fully-faithful}。
\end{proof}

\begin{lemma}
\label{sites-cohomology-lemma-equivalence-unbounded-one}
\begin{reference}
这与 \cite[Theorem 2.2.3]{six-I} 类似。
\end{reference}
设 $f : (\Sh(\mathcal{C}), \mathcal{O}) \to (\Sh(\mathcal{C}'), \mathcal{O}')$
为环化拓扑斯的态射。设 $\mathcal{A} \subset \textit{Mod}(\mathcal{O})$
及 $\mathcal{A}' \subset \textit{Mod}(\mathcal{O}')$ 为弱 Serre 子范畴。假设
\begin{enumerate}
\item $f$ 是平坦的；
\item $f^*$ 诱导范畴等价 $\mathcal{A}' \to \mathcal{A}$；
\item 对 $\mathcal{F}' \in \Ob(\mathcal{A}')$，映射
$\mathcal{F}' \to Rf_*f^*\mathcal{F}'$ 是同构；
\item $\mathcal{C}, \mathcal{O}, \mathcal{A}$ 满足情形
\ref{sites-cohomology-situation-olsson-laszlo} 的假设；
\item $\mathcal{C}', \mathcal{O}', \mathcal{A}'$ 满足情形
\ref{sites-cohomology-situation-olsson-laszlo} 的假设。
\end{enumerate}
则 $f^* : D_{\mathcal{A}'}(\mathcal{O}') \to D_\mathcal{A}(\mathcal{O})$
是范畴等价，其拟逆为
$Rf_* : D_\mathcal{A}(\mathcal{O}) \to D_{\mathcal{A}'}(\mathcal{O}')$.
\end{lemma}

\begin{proof}
由于 $f^*$ 是正合的，显然 $f^*$ 定义函子
$f^* : D_{\mathcal{A}'}(\mathcal{O}') \to D_\mathcal{A}(\mathcal{O})$
如引理所述，且该函子还与截断函子 $\tau_{\geq -n}$ 交换。
我们已经知道，在相应的下有界范畴上，$f^*$ 与 $Rf_*$ 是互为拟逆的
等价，参见引理 \ref{sites-cohomology-lemma-equivalence-bounded}。
由引理 \ref{sites-cohomology-lemma-olsson-laszlo-map-version-one}（取 $N = 0$），可知
$Rf_*$ 确实定义函子
$Rf_* : D_\mathcal{A}(\mathcal{O}) \to D_{\mathcal{A}'}(\mathcal{O}')$
且该函子还与截断函子 $\tau_{\geq -n}$ 交换。
因此对 $D_\mathcal{A}(\mathcal{O})$ 中的 $K$，映射
$f^*Rf_*K \to K$ 是同构，因为它在各截断上是同构。
类似地，对 $D_{\mathcal{A}'}(\mathcal{O}')$ 中的 $K'$，映射
$K' \to Rf_*f^*K'$ 是同构，因为它在各截断上是同构。
证明完毕。
\end{proof}

\begin{lemma}
\label{sites-cohomology-lemma-equivalence-unbounded-two}
\begin{reference}
这与 \cite[Theorem 2.2.3]{six-I} 类似。
\end{reference}
设 $f : (\mathcal{C}, \mathcal{O}) \to (\mathcal{C}', \mathcal{O}')$
为环化站点态射。设 $\mathcal{A} \subset \textit{Mod}(\mathcal{O})$
及 $\mathcal{A}' \subset \textit{Mod}(\mathcal{O}')$ 为弱 Serre 子范畴。假设
\begin{enumerate}
\item $f$ 是平坦的；
\item $f^*$ 诱导范畴等价 $\mathcal{A}' \to \mathcal{A}$；
\item 对 $\mathcal{F}' \in \Ob(\mathcal{A}')$，映射
$\mathcal{F}' \to Rf_*f^*\mathcal{F}'$ 是同构；
\item $\mathcal{C}, \mathcal{O}, \mathcal{A}$ 满足情形
\ref{sites-cohomology-situation-olsson-laszlo} 的假设；
\item $f : (\mathcal{C}, \mathcal{O}) \to (\mathcal{C}', \mathcal{O}')$
和 $\mathcal{A}$ 满足情形
\ref{sites-cohomology-situation-olsson-laszlo-prime} 的假设。
\end{enumerate}
则 $f^* : D_{\mathcal{A}'}(\mathcal{O}') \to D_\mathcal{A}(\mathcal{O})$
是范畴等价，其拟逆为
$Rf_* : D_\mathcal{A}(\mathcal{O}) \to D_{\mathcal{A}'}(\mathcal{O}')$.
\end{lemma}

\begin{proof}
本引理的证明与引理 \ref{sites-cohomology-lemma-equivalence-unbounded-one} 的证明完全相同，
区别仅在于将对引理 \ref{sites-cohomology-lemma-olsson-laszlo-map-version-one} 的引用
替换为对
Lemma \ref{sites-cohomology-lemma-olsson-laszlo-map-version-two}.
\end{proof}








\section{比较两个拓扑，II}
\label{sites-cohomology-section-compare-II}

\noindent
设 $\mathcal{C}$ 为范畴。令
$\text{Cov}(\mathcal{C}) \supset \text{Cov}'(\mathcal{C})$
为两种赋予 $\mathcal{C}$ 站点结构的方式。记 $\tau$ 为对应于
$\text{Cov}(\mathcal{C})$ 的拓扑，记 $\tau'$ 为对应于
$\text{Cov}'(\mathcal{C})$ 的拓扑。于是 $\mathcal{C}$ 上的恒等函子定义站点态射
$$
\epsilon : \mathcal{C}_\tau \longrightarrow \mathcal{C}_{\tau'}
$$
其中 $\epsilon_*$ 是底层预层上的恒等函子，而 $\epsilon^{-1}$ 是
$\tau'$-层的 $\tau$-层化（因而显然正合）。设 $\mathcal{O}$ 是
$\tau$-拓扑下的环层。则 $\mathcal{O}$ 也是 $\tau'$-拓扑下的层，
且 $\epsilon$ 成为环化站点态射
$$
\epsilon :
(\mathcal{C}_\tau, \mathcal{O}_\tau)
\longrightarrow
(\mathcal{C}_{\tau'}, \mathcal{O}_{\tau'})
$$
更多讨论参见第 \ref{sites-cohomology-section-compare} 节。

\begin{lemma}
\label{sites-cohomology-lemma-compare-topologies-derived-adequate-modules}
设 $\epsilon : (\mathcal{C}_\tau, \mathcal{O}_\tau) \to
(\mathcal{C}_{\tau'}, \mathcal{O}_{\tau'})$ 如上。令
$\mathcal{B} \subset \Ob(\mathcal{C})$ 为子集。令
$\mathcal{A} \subset \textit{PMod}(\mathcal{O})$ 为满子范畴。假设
\begin{enumerate}
\item $\mathcal{A}$ 的每个对象都是 $\tau$-拓扑下的层；
\item $\mathcal{A}$ 是 $\textit{Mod}(\mathcal{O}_\tau)$ 的弱 Serre 子范畴；
\item $\mathcal{C}$ 的每个对象都有一个 $\tau'$-覆盖，其成员属于
$\mathcal{B}$；
\item 对每个 $U \in \mathcal{B}$，有
$H^p_\tau(U, \mathcal{F}) = 0$（对所有 $\mathcal{F} \in \mathcal{A}$ 及 $p > 0$）。
\end{enumerate}
则 $\mathcal{A}$ 是
$\textit{Mod}(\mathcal{O}_{\tau'})$ 的弱 Serre 子范畴，且存在三角范畴等价
$D_\mathcal{A}(\mathcal{O}_\tau) = D_\mathcal{A}(\mathcal{O}_{\tau'})$
由 $\epsilon^*$ 和 $R\epsilon_*$ 给出。
\end{lemma}

\begin{proof}
由于 $\epsilon^{-1}\mathcal{O}_{\tau'} = \mathcal{O}_\tau$，
可知 $\epsilon$ 是环化站点的平坦态射；事实上，在模层上
$\epsilon^{-1} = \epsilon^*$。由性质 (1)，可将 $\mathcal{A}$ 的每个对象
视为 $\mathcal{O}_\tau$-模层，也视为 $\mathcal{O}_{\tau'}$-模层。
换言之，有完全忠实的包含函子
$$
\mathcal{A} \to \textit{Mod}(\mathcal{O}_\tau) \to
\textit{Mod}(\mathcal{O}_{\tau'})
$$
为避免混淆，记 $\mathcal{A}' \subset \textit{Mod}(\mathcal{O}_{\tau'})$
为 $\mathcal{A}$ 的像。显然，$\epsilon_* : \mathcal{A} \to \mathcal{A}'$
和 $\epsilon^* : \mathcal{A}' \to \mathcal{A}$ 互为拟逆等价（参见引理前的
讨论，并注意 $\mathcal{A}'$ 的对象是 $\tau$-拓扑下的层）。

\medskip\noindent
条件 (3) 和 (4) 蕴含，当 $p > 0$ 且 $\mathcal{F} \in \Ob(\mathcal{A})$ 时，
$R^p\epsilon_*\mathcal{F} = 0$。这是因为 $R^p\epsilon_*$ 是预层
$U \mapsto H^p_\tau(U, \mathcal{F})$ 所关联的层，参见引理
\ref{sites-cohomology-lemma-higher-direct-images}。因此，对 $\mathcal{A}$ 中的任意正合复形
（这等价于 $\textit{Mod}(\mathcal{O}_\tau)$ 中各项属于 $\mathcal{A}$ 的正合复形，
参见同调，引理 \ref{homology-lemma-characterize-weak-serre-subcategory}），
应用函子 $\epsilon_*$ 后仍为正合。

\medskip\noindent
考虑正合序列
$$
\mathcal{F}'_0 \to \mathcal{F}'_1 \to
\mathcal{F}'_2 \to \mathcal{F}'_3 \to
\mathcal{F}'_4
$$
它位于 $\textit{Mod}(\mathcal{O}_{\tau'})$ 中，且
$\mathcal{F}'_0, \mathcal{F}'_1, \mathcal{F}'_3, \mathcal{F}'_4$ 属于
$\mathcal{A}'$。应用正合函子 $\epsilon^*$ 得到正合序列
$$
\epsilon^*\mathcal{F}'_0 \to \epsilon^*\mathcal{F}'_1 \to
\epsilon^*\mathcal{F}'_2 \to \epsilon^*\mathcal{F}'_3 \to
\epsilon^*\mathcal{F}'_4
$$
它位于 $\textit{Mod}(\mathcal{O}_\tau)$ 中。由于 $\mathcal{A}$ 是
弱 Serre 子范畴，且
$\epsilon^*\mathcal{F}'_0, \epsilon^*\mathcal{F}'_1,
\epsilon^*\mathcal{F}'_3, \epsilon^*\mathcal{F}'_4$ 属于
$\mathcal{A}$。由同调（参见定义
\ref{homology-definition-serre-subcategory}），得到
$\epsilon^*\mathcal{F}_2 \in \mathcal{A}$。考虑序列映射
$$
\xymatrix{
\mathcal{F}'_0 \ar[r] \ar[d] &
\mathcal{F}'_1 \ar[r] \ar[d] &
\mathcal{F}'_2 \ar[r] \ar[d] &
\mathcal{F}'_3 \ar[r] \ar[d] &
\mathcal{F}'_4 \ar[d] \\
\epsilon_*\epsilon^*\mathcal{F}'_0 \ar[r] &
\epsilon_*\epsilon^*\mathcal{F}'_1 \ar[r] &
\epsilon_*\epsilon^*\mathcal{F}'_2 \ar[r] &
\epsilon_*\epsilon^*\mathcal{F}'_3 \ar[r] &
\epsilon_*\epsilon^*\mathcal{F}'_4
}
$$
由前一段讨论，下行正合。指标为 $0$、$1$、$3$、$4$ 的竖直箭头
由第一段讨论是同构。由 $5$ 引理（同调，引理
\ref{homology-lemma-five-lemma}）可得
$\mathcal{F}'_2 \cong \epsilon_*\epsilon^*\mathcal{F}'_2$，
从而 $\mathcal{F}'_2 \in \mathcal{A}'$。这样便知
$\mathcal{A}'$ 是 $\textit{Mod}(\mathcal{O}_{\tau'})$ 的弱 Serre 子范畴，
参见同调，定义 \ref{homology-definition-serre-subcategory}。

\medskip\noindent
此时讨论导出范畴
$D_\mathcal{A}(\mathcal{O}_\tau)$ 和
$D_{\mathcal{A}'}(\mathcal{O}_{\tau'})$ 才有意义，参见导出范畴，
第 \ref{derived-section-triangulated-sub} 节。为完成证明，我们说明引理
\ref{sites-cohomology-lemma-equivalence-unbounded-two} 的条件 (1)--(5) 得到满足。
上面已经验证了 (1)、(2)、(3)。注意，由于每个对象都有由
$\mathcal{B}$ 中对象组成的 $\tau'$-覆盖，当然每个对象也有由
$\mathcal{B}$ 中对象组成的 $\tau$-覆盖。因此引理
\ref{sites-cohomology-lemma-equivalence-unbounded-two} 的条件 (4) 得到满足。
同理，条件 (5) 也得到满足。
\end{proof}

\begin{lemma}
\label{sites-cohomology-lemma-descent-squares}
设 $\epsilon : (\mathcal{C}_\tau, \mathcal{O}_\tau) \to
(\mathcal{C}_{\tau'}, \mathcal{O}_{\tau'})$ 如上。令 $A$ 为集合，并对
$\alpha \in A$ 令
$$
\xymatrix{
E_\alpha \ar[d] \ar[r] & Y_\alpha \ar[d] \\
Z_\alpha \ar[r] & X_\alpha
}
$$
为范畴 $\mathcal{C}$ 中的交换图。假设
\begin{enumerate}
\item $\tau'$-层 $\mathcal{F}'$ 是 $\tau$-层，当且仅当对所有 $\alpha$ 有
$\mathcal{F}'(X_\alpha) =
\mathcal{F}'(Z_\alpha) \times_{\mathcal{F}'(E_\alpha)} 
\mathcal{F}'(Y_\alpha)$；
\item 对于 $D(\mathcal{O}_{\tau'})$ 中属于 $R\epsilon_*$ 本质像的 $K'$，
引理 \ref{sites-cohomology-lemma-c-square} 的映射
$c^{K'}_{X_\alpha, Z_\alpha, Y_\alpha, E_\alpha}$
对所有 $\alpha$ 都是同构。
\end{enumerate}
则对 $K' \in D^+(\mathcal{O}_{\tau'})$，它属于 $R\epsilon_*$ 的本质像，
当且仅当映射 $c^{K'}_{X_\alpha, Z_\alpha, Y_\alpha, E_\alpha}$
对所有 $\alpha$ 都是同构。
\end{lemma}

\begin{proof}
“仅当”方向由假设 (2) 蕴含。
另一方面，若 $K'$ 只有一个非零上同调层，
则“若”方向由假设 (1) 得出。
一般地，我们将用归纳法证明“若”方向。称
$D^+(\mathcal{O}_{\tau'})$ 中的对象 $K'$ 满足 (P)，如果映射
满足 (P)，如果映射 $c^{K'}_{X_\alpha, Z_\alpha, Y_\alpha, E_\alpha}$
对所有 $\alpha \in A$ 都是同构。

\medskip\noindent
具体地，设 $K'$ 是满足 (P) 的 $D^+(\mathcal{O}_{\tau'})$ 中对象。
选取有区别三角
$$
K' \to R\epsilon_*\epsilon^{-1}K' \to M' \to K'[1]
$$
它位于 $D^+(\mathcal{O}_{\tau'})$ 中，其中第一个箭头是伴随映射。
由 (2) 和引理 \ref{sites-cohomology-lemma-two-out-of-three-blow-up-square} 可知 $M'$ 满足 (P)。另一方面，应用

$\epsilon^{-1}$，并利用第 \ref{sites-cohomology-section-compare} 节中
$\epsilon^{-1}R\epsilon_* = \text{id}$，得到 $\epsilon^{-1}M' = 0$。
下一段将证明 $M' = 0$，从而完成证明。

\medskip\noindent
设 $K'$ 是满足 (P) 且 $\epsilon^{-1}K' = 0$ 的 $D^+(\mathcal{O}_{\tau'})$ 中对象。我们将证明 $K' = 0$。

具体地，给定 $n \in \mathbf{Z}$，若 $i < n$ 时 $H^i(K') = 0$，
我们将证明 $H^n(K') = 0$。对 $\alpha \in A$，由引理 \ref{sites-cohomology-lemma-c-square} 有有区别三角

$$
R\Gamma_{\tau'}(X_\alpha, K') \to
R\Gamma_{\tau'}(Z_\alpha, K') \oplus
R\Gamma_{\tau'}(Y_\alpha, K') \to
R\Gamma_{\tau'}(E_\alpha, K') \to
R\Gamma_{\tau'}(X_\alpha, K')[1]
$$
取次数 $n$ 的上同调，并利用已假设的 $K'$ 上同调层消失，得到正合序列


$$
0 \to
H^n_{\tau'}(X_\alpha, K') \to
H^n_{\tau'}(Z_\alpha, K') \oplus
H^n_{\tau'}(Y_\alpha, K') \to
H^n_{\tau'}(E_\alpha, K')
$$
这与正合序列
$$
0 \to
\Gamma(X_\alpha, H^n(K')) \to
\Gamma(Z_\alpha, H^n(K')) \oplus
\Gamma(Y_\alpha, H^n(K')) \to
\Gamma(E_\alpha, H^n(K'))
$$
由假设 (1)，可得 $H^n(K')$ 是 $\tau$-层。然而，由于
$\tau$-层化 $\epsilon^{-1}H^n(K') = H^n(\epsilon^{-1}K')$，且
$\epsilon^{-1}H^n(K') = H^n(\epsilon^{-1}K')$ 为 $0$，因为
$\epsilon^{-1}K' = 0$，故其为 $0$；于是按要求
$H^n(K') = 0$。
\end{proof}

\begin{lemma}
\label{sites-cohomology-lemma-descent-squares-helper}
设 $\epsilon : (\mathcal{C}_\tau, \mathcal{O}_\tau) \to
(\mathcal{C}_{\tau'}, \mathcal{O}_{\tau'})$ 如上。令
$$
\xymatrix{
E \ar[d] \ar[r] & Y \ar[d] \\
Z \ar[r] & X
}
$$
为范畴 $\mathcal{C}$ 中满足下列条件的交换图：
\begin{enumerate}
\item $h_X^\# = h_Y^\# \amalg_{h_E^\#} h_Z^\#$；
\item $h_E^\# \to h_Y^\#$ 是单射。
\end{enumerate}
其中上标井号表示 $\tau$-层化。则对属于 $R\epsilon_*$ 本质像的
$K' \in D(\mathcal{O}_{\tau'})$，引理 \ref{sites-cohomology-lemma-c-square} 的映射
$c^{K'}_{X, Z, Y, E}$（使用 $\tau'$-拓扑）是同构。
\end{lemma}


\begin{proof}
此辅助引理几乎立即由引理 \ref{sites-cohomology-lemma-square-triangle} 得出，
我们强烈建议读者跳过证明。设 $K' = R\epsilon_*K$。
选取表示 $K$ 的 $\mathcal{O}_\tau$-模 K-内射复形
$\mathcal{J}^\bullet$。则 $\epsilon_*\mathcal{J}^\bullet$ 是表示 $K'$ 的
$\mathcal{O}_{\tau'}$-模 K-内射复形，参见引理
\ref{sites-cohomology-lemma-K-injective-flat}。接着，
$$
0 \to
\mathcal{J}^\bullet(X) \xrightarrow{\alpha}
\mathcal{J}^\bullet(Z) \oplus
\mathcal{J}^\bullet(Y) \xrightarrow{\beta}
\mathcal{J}^\bullet(E)
\to 0
$$
是阿贝尔群复形的短正合序列，参见引理
\ref{sites-cohomology-lemma-square-triangle} 及其证明。由于这正是定义
$c^{K'}_{X, Z, Y, E}$ 时所用的阿贝尔群复形序列，结论成立。
\end{proof}







\section{比较上同调}
\label{sites-cohomology-section-compare-general}

\noindent
我们发展一些一般理论，以帮助比较不同拓扑下的上同调。给定第
\ref{sites-cohomology-section-compare} 节中的 $\mathcal{C}$、$\tau$、$\tau'$，以及
$\mathcal{C}$ 中的态射 $f : X \to Y$，得到拓扑斯态射的交换图
\begin{equation}
\label{sites-cohomology-equation-commutative-epsilon}
\vcenter{
\xymatrix{
\Sh(\mathcal{C}_\tau/X) \ar[r]_{f_\tau} \ar[d]_{\epsilon_X} &
\Sh(\mathcal{C}_\tau/Y) \ar[d]^{\epsilon_Y} \\
\Sh(\mathcal{C}_{\tau'}/X) \ar[r]^{f_{\tau'}} &
\Sh(\mathcal{C}_{\tau'}/X)
}
}
\end{equation}
这里的态射 $\epsilon_X$（分别为 $\epsilon_Y$）是第
\ref{sites-cohomology-section-compare} 节中为带有 $\tau$ 和 $\tau'$ 两个拓扑的范畴
$\mathcal{C}/X$ 定义的比较态射。态射 $f_\tau$ 和 $f_{\tau'}$
是“重新局部化”态射（站点，引理 \ref{sites-lemma-relocalize}）。
图的交换性是站点，引理 \ref{sites-lemma-localize-morphism} 的特例
（取 $\mathcal{C} = \mathcal{C}_\tau/Y$、$\mathcal{D} = \mathcal{C}_{\tau'}/Y$、
$u = \text{id}$、$U = X$ 及 $V = X$）。此外还得到
$\epsilon_{X, *} \circ f_\tau^{-1} = f_{\tau'}^{-1} \circ \epsilon_{Y, *}$
这既可由该引理得到，也可因其显然性得到。

\begin{situation}
\label{sites-cohomology-situation-compare}
设 $\mathcal{C}$、$\tau$、$\tau'$ 如第 \ref{sites-cohomology-section-compare} 节。
假设给定子集 $\mathcal{P} \subset \text{Arrows}(\mathcal{C})$，并且对
$\mathcal{C}$ 的每个对象 $X$，给定弱 Serre 子范畴
$\mathcal{A}'_X \subset \textit{Ab}(\mathcal{C}_{\tau'}/X)$。
作如下假设：
\begin{enumerate}
\item
\label{sites-cohomology-item-base-change-P}
若 $f : X \to Y$ 属于 $\mathcal{P}$，且 $Y' \to Y$ 为任意态射，
则 $X \times_Y Y'$ 存在，且 $X \times_Y Y' \to Y'$ 属于 $\mathcal{P}$；
\item
\label{sites-cohomology-item-restriction-A}
$f_{\tau'}^{-1}$ 将 $\mathcal{A}'_Y$ 映入 $\mathcal{A}'_X$，
对 $\mathcal{C}$ 的任意态射 $f : X \to Y$ 均如此；
\item
\label{sites-cohomology-item-A-sheaf}
若 $X \in \mathcal{C}$ 且 $\mathcal{F}' \in \mathcal{A}'_X$，则
$\mathcal{F}'$ 满足 $\tau$-覆盖的层条件，即
$\mathcal{F}' = \epsilon_{X, *}\epsilon_X^{-1}\mathcal{F}'$,
\item
\label{sites-cohomology-item-A-and-P}
若 $f : X \to Y$ 属于 $\mathcal{P}$ 且
$\mathcal{F}' \in \Ob(\mathcal{A}'_X)$，则当 $i \geq 0$ 时
$R^if_{\tau', *}\mathcal{F}' \in \Ob(\mathcal{A}'_Y)$；
\item
\label{sites-cohomology-item-refine-tau-by-P}
若 $\{U_i \to U\}_{i \in I}$ 是 $\tau$-覆盖，则存在
\begin{enumerate}
\item $\tau'$-覆盖 $\{V_j \to U\}_{j \in J}$；
\item $\tau$-覆盖 $\{f_j : W_j \to V_j\}$，其中每个覆盖只含有一个
$f_j \in \mathcal{P}$；
\item $\tau'$-覆盖 $\{W_{jk} \to W_j\}_{k \in K_j}$
\end{enumerate}
使得 $\{W_{jk} \to U\}_{j \in J, k \in K_j}$ 是
$\{U_i \to U\}_{i \in I}$ 的细化。
\end{enumerate}
\end{situation}

\begin{lemma}
\label{sites-cohomology-lemma-A}
在情形 \ref{sites-cohomology-situation-compare} 中，对 $X \in \mathcal{C}$，记
$\mathcal{A}_X$ 为 $\textit{Ab}(\mathcal{C}_\tau/X)$ 中形如
$\epsilon_X^{-1}\mathcal{F}'$ 的对象，其中 $\mathcal{F}' \in \mathcal{A}'_X$。
则
\begin{enumerate}
\item 对 $\mathcal{F} \in \textit{Ab}(\mathcal{C}_\tau/X)$，有
$\mathcal{F} \in \mathcal{A}_X \Leftrightarrow
\epsilon_{X, *}\mathcal{F} \in \mathcal{A}'_X$；
\item 对 $\mathcal{C}$ 的任意态射 $f : X \to Y$，$f_\tau^{-1}$ 将
$\mathcal{A}_Y$ 映入 $\mathcal{A}_X$。
\end{enumerate}
\end{lemma}

\begin{proof}
第 (1) 部分由 (\ref{sites-cohomology-item-A-sheaf}) 得出；第 (2) 部分由
(\ref{sites-cohomology-item-restriction-A}) 以及 (\ref{sites-cohomology-equation-commutative-epsilon}) 的交换性得出，
其中
$\epsilon_X^{-1} \circ f_{\tau'}^{-1} = f_\tau^{-1} \circ \epsilon_Y^{-1}$.
\end{proof}

\noindent
我们的下一个目标是证明引理 \ref{sites-cohomology-lemma-compare-cohomology-general} 和
\ref{sites-cohomology-lemma-cohomological-descent-general}。我们将使用下面的归纳假设，
通过归纳论证完成证明。

\medskip\noindent
$(V_n)$ 对 $X \in \mathcal{C}$ 及 $\mathcal{F} \in \mathcal{A}_X$，有
$R^i\epsilon_{X, *}\mathcal{F} = 0$（$1 \leq i \leq n$）。

\begin{lemma}
\label{sites-cohomology-lemma-V-implies-C-general}
在情形 \ref{sites-cohomology-situation-compare} 中假设 $(V_n)$ 成立。
若 $f : X \to Y$ 属于 $\mathcal{P}$ 且 $\mathcal{F} \in \mathcal{A}_X$，
则当 $i \leq n$ 时有
$R^if_{\tau', *}\epsilon_{X, *}\mathcal{F} =
\epsilon_{Y, *}R^if_{\tau, *}\mathcal{F}$。
\end{lemma}

\begin{proof}
下文不再特别说明，我们将使用交换图
(\ref{sites-cohomology-equation-commutative-epsilon})。特别地，有
$$
Rf_{\tau', *}R\epsilon_{X, *}\mathcal{F} =
R\epsilon_{Y, *}Rf_{\tau, *}\mathcal{F}
$$
假设 $(V_n)$ 告诉我们
$\epsilon_{X, *}\mathcal{F} \to R\epsilon_{X, *}\mathcal{F}$
在次数 $\leq n$ 上是同构。因此
$Rf_{\tau', *}\epsilon_{X, *}\mathcal{F} \to
Rf_{\tau', *}R\epsilon_{X, *}\mathcal{F}$ 在次数 $\leq n$ 上是同构。
于是
$$
R^if_{\tau', *}\epsilon_{X, *}\mathcal{F} \to
H^i(R\epsilon_{Y, *}Rf_{\tau, *}\mathcal{F})
$$
当 $i \leq n$ 时是同构。我们考察引理
\ref{sites-cohomology-lemma-relative-Leray} 中关于
$R\epsilon_{Y, *}Rf_{\tau, *}\mathcal{F}$ 的谱序列的第二页来证明本引理。
图示如下：
$$
\begin{matrix}
\ldots &
\ldots &
\ldots &
\ldots \\
\epsilon_{Y, *}R^2f_{\tau, *}\mathcal{F} &
R^1\epsilon_{Y, *}R^2f_{\tau, *}\mathcal{F} &
R^2\epsilon_{Y, *}R^2f_{\tau, *}\mathcal{F} &
\ldots \\
\epsilon_{Y, *}R^1f_{\tau, *}\mathcal{F} &
R^1\epsilon_{Y, *}R^1f_{\tau, *}\mathcal{F} &
R^2\epsilon_{Y, *}R^1f_{\tau, *}\mathcal{F} &
\ldots \\
\epsilon_{Y, *}f_{\tau, *}\mathcal{F} &
R^1\epsilon_{Y, *}f_{\tau, *}\mathcal{F} &
R^2\epsilon_{Y, *}f_{\tau, *}\mathcal{F} &
\ldots
\end{matrix}
$$
令 $(C_m)$ 表示假设：
$R^if_{\tau', *}\epsilon_{X, *}\mathcal{F} =
\epsilon_{Y, *}R^if_{\tau, *}\mathcal{F}$（$i \leq m$）。
注意 $(C_0)$ 成立。我们将证明当 $m < n$ 时
$(C_{m - 1}) \Rightarrow (C_m)$。具体地，若 $(C_{m - 1})$ 成立，则当
$n \geq p > 0$ 且 $q \leq m - 1$ 时有
\begin{align*}
R^p\epsilon_{Y, *}R^qf_{\tau, *}\mathcal{F}
& =
R^p\epsilon_{Y, *}
\epsilon_Y^{-1} \epsilon_{Y, *} R^qf_{\tau, *}\mathcal{F} \\
& =
R^p\epsilon_{Y, *}
\epsilon_Y^{-1}R^qf_{\tau', *}\epsilon_{X, *}\mathcal{F} = 0
\end{align*}
第一个等号来自 $\epsilon_Y^{-1}\epsilon_{Y, *} = \text{id}$，
第二个来自 $(C_{m - 1})$，最后一个来自 $(V_n)$，因为由
(\ref{sites-cohomology-item-A-and-P})，$\epsilon_Y^{-1}R^qf_{\tau', *}\epsilon_{X, *}\mathcal{F}$
属于 $\mathcal{A}_Y$。考察谱序列可知
$E_2^{0, m} = \epsilon_{Y, *}R^mf_{\tau, *}\mathcal{F}$
是满足 $p + q = m$ 的 $E_2^{p, q}$ 中唯一非零项。记住
$\text{d}_r^{p, q} : E_r^{p, q} \to E_r^{p + r, q - r + 1}$。
因此不存在 $r \geq 2$ 的非零微分
$\text{d}_r^{p, q}$ 从该位置出发或进入该位置。于是
$H^m(R\epsilon_{Y, *}Rf_{\tau, *}\mathcal{F}) =
\epsilon_{Y, *}R^mf_{\tau, *}\mathcal{F}$，由上面的讨论得到
$(C_m)$。

\medskip\noindent
最后，假设 $(C_{n - 1})$。同样的分析表明
$E_2^{0, n} = \epsilon_{Y, *}R^nf_{\tau, *}\mathcal{F}$
是满足 $p + q = n$ 的 $E_2^{p, q}$ 中唯一非零项。
仍然没有非零微分进入该位置，但可能有非零微分从该位置出发。
具体地，映射
$d_{n + 1}^{0, n} : \epsilon_{Y, *}R^nf_{\tau, *}\mathcal{F} \to
R^{n + 1}\epsilon_{Y, *}f_{\tau, *}\mathcal{F}$.
因此存在正合序列
$$
0 \to R^nf_{\tau', *}\epsilon_{X, *}\mathcal{F} \to 
\epsilon_{Y, *}R^nf_{\tau, *}\mathcal{F} \to
R^{n + 1}\epsilon_{Y, *}f_{\tau, *}\mathcal{F}
$$
由 (\ref{sites-cohomology-item-A-and-P}) 和 (\ref{sites-cohomology-item-A-sheaf})，层
$R^nf_{\tau', *}\epsilon_{X, *}\mathcal{F}$
满足 $\tau$-覆盖的层性质，$\epsilon_{Y, *}R^nf_{\tau, *}\mathcal{F}$
也满足（使用第 \ref{sites-cohomology-section-compare} 节中对 $\epsilon_*$ 的描述）。
然而，$\tau'$-层
$R^{n + 1}\epsilon_{Y, *}f_{\tau, *}\mathcal{F}$
的 $\tau$-层化为零（由上同调的局部性；使用引理
\ref{sites-cohomology-lemma-kill-cohomology-class-on-covering} 和
\ref{sites-cohomology-lemma-higher-direct-images}）。因此
$R^nf_{\tau', *}\epsilon_{X, *}\mathcal{F} \to
\epsilon_{Y, *}R^nf_{\tau, *}\mathcal{F}$
必为同构，证明完成。
\end{proof}

\noindent
若 $E'$（分别地，$E$）是
$D(\mathcal{C}_{\tau'}/X)$（分别地，$D(\mathcal{C}_\tau/X)$）中的对象，
则对 $\mathcal{C}/X$ 中的对象 $U$，记
$H^n_{\tau'}(U, E')$（分别地，$H^n_\tau(U, E)$）为
$E'$（分别地，$E$）在 $U$ 上的上同调。

\begin{lemma}
\label{sites-cohomology-lemma-V-implies-cohomology-general}
在情形 \ref{sites-cohomology-situation-compare} 中，若 $(V_n)$ 成立，则对
$X \in \mathcal{C}$ 及 $L \in D(\mathcal{C}_{\tau'}/X)$，若
$i < 0$ 时 $H^i(L) = 0$，且 $0 \leq i \leq n$ 时
$H^i(L) \in \mathcal{A}'_X$，则有
$H^n_{\tau'}(X, L) = H^n_\tau(X, \epsilon_X^{-1}L)$.
\end{lemma}

\begin{proof}
由引理 \ref{sites-cohomology-lemma-Leray-unbounded}，有
$H^n_\tau(X, \epsilon_X^{-1}L) =
H^n_{\tau'}(X, R\epsilon_{X, *}\epsilon_X^{-1}L)$.
存在谱序列
$$
E_2^{p, q} = R^p\epsilon_{X, *}\epsilon_X^{-1}H^q(L)
$$
收敛到 $H^{p + q}(R\epsilon_{X, *}\epsilon_X^{-1}L)$。
由 $(V_n)$，当 $0 < p \leq n$ 且 $0 \leq q \leq n$ 时
$E_2^{p, q}$ 消失。因此
$E_2^{0, q} = \epsilon_{X, *}\epsilon_X^{-1}H^q(L) = H^q(L)$
是满足 $p + q \leq n$ 的 $E_2^{p, q}$ 中唯一可能的非零项。
于是映射
$$
L \longrightarrow R\epsilon_{X, *}\epsilon_X^{-1}L
$$
在次数 $\leq n$ 上是同构（细节略去）。因此
$H^i_{\tau'}(X, L) = H^i_{\tau'}(X, R\epsilon_{X, *}\epsilon_X^{-1}L)$
对 $i \leq n$ 成立。引理得证。
\end{proof}

\begin{lemma}
\label{sites-cohomology-lemma-V-implies-cohomology-extra-general}
在情形 \ref{sites-cohomology-situation-compare} 中，若 $(V_n)$ 成立，则对
$X \in \mathcal{C}$ 及 $\mathcal{F} \in \mathcal{A}_X$，映射
$H^{n + 1}_{\tau'}(X, \epsilon_{X, *}\mathcal{F}) \to
H^{n + 1}_\tau(X, \mathcal{F})$
是单射，其像正是那些在 $X$ 的某个 $\tau'$-覆盖上变为零的类。
\end{lemma}

\begin{proof}
记住 $\epsilon_X^{-1}\epsilon_{X, *}\mathcal{F} = \mathcal{F}$，
故该映射由沿 $\epsilon_X$ 拉回上同调类给出。Leray 谱序列（引理
\ref{sites-cohomology-lemma-Leray}）
$$
E_2^{p, q} = H^p_{\tau'}(X, R^q\epsilon_{X, *}\mathcal{F})
\Rightarrow
H^{p + q}_\tau(X, \mathcal{F})
$$
结合已假设的消失，得到正合序列
$$
0 \to
H^{n + 1}_{\tau'}(X, \epsilon_{X, *}\mathcal{F}) \to
H^{n + 1}_\tau(X, \mathcal{F}) \to
H^0_{\tau'}(X, R^{n + 1}\epsilon_{X, *}\mathcal{F})
$$
这正是引理的另一种表述。
\end{proof}

\begin{lemma}
\label{sites-cohomology-lemma-make-class-zero-general}
在情形 \ref{sites-cohomology-situation-compare} 中，设 $f : X \to Y$ 属于 $\mathcal{P}$，
且 $\{X \to Y\}$ 是 $\tau$-覆盖。设 $\mathcal{F}' \in \mathcal{A}'_Y$。
若 $n \geq 0$ 且
$$
\theta \in
\text{Equalizer}\left(
\xymatrix{
H^{n + 1}_{\tau'}(X, \mathcal{F}')
\ar@<1ex>[r] \ar@<-1ex>[r] &
H^{n + 1}_{\tau'}(X \times_Y X, \mathcal{F}')
}
\right)
$$
则存在 $\tau'$-覆盖 $\{Y_i \to Y\}$，使得 $\theta$ 在
$H^{n + 1}_{\tau'}(Y_i \times_Y X, \mathcal{F}')$ 中的限制为零。
\end{lemma}

\begin{proof}
由 (\ref{sites-cohomology-item-base-change-P})，$X \times_Y X$ 存在。对
$Z \in \mathcal{C}/Y$，记 $\mathcal{F}'|_Z$ 为 $\mathcal{F}'$ 到
$\mathcal{C}_{\tau'}/Z$ 的限制。记住
$H^{n + 1}_{\tau'}(X, \mathcal{F}') =
H^{n + 1}(\mathcal{C}_{\tau'}/X, \mathcal{F}'|_X)$，参见引理
\ref{sites-cohomology-lemma-cohomology-of-open}。本引理断言 $\theta$ 的像
$\overline{\theta} \in H^0(Y, R^{n + 1}f_{\tau', *}\mathcal{F}'|_X)$
为零。考虑笛卡尔图
$$
\xymatrix{
X \times_Y X \ar[d]_{\text{pr}_1} \ar[r]_{\text{pr}_2} &
X \ar[d]^f \\
X \ar[r]^f & Y
}
$$
由平凡基变换（引理 \ref{sites-cohomology-lemma-localize-cartesian-square}），有
$$
f_{\tau'}^{-1}R^{n + 1}f_{\tau', *}(\mathcal{F}'|_X) =
R^{n + 1}\text{pr}_{1, \tau', *}(\mathcal{F}'|_{X \times_Y X})
$$
若 $\text{pr}_1^{-1}\theta = \text{pr}_2^{-1}\theta$，则
$f_{\tau'}^{-1}\overline{\theta}$ 作为
$f_{\tau'}^{-1}R^{n + 1}f_{\tau', *}(\mathcal{F}'|_X)$ 的截面为零，
因为显然 $\text{pr}_1^{-1}\theta$ 映到
$H^0(X, R^{n + 1}\text{pr}_{1, \tau', *}(\mathcal{F}'|_{X \times_Y X}))$
中的零元素。由 (\ref{sites-cohomology-item-restriction-A})，$\mathcal{F}'|_X \in \mathcal{A}'_X$。
因此由 (\ref{sites-cohomology-item-A-and-P})，
$\mathcal{G}' = R^{n + 1}f_{\tau', *}(\mathcal{F}'|_X)$
是 $\mathcal{A}'_Y$ 中的对象。由 (\ref{sites-cohomology-item-A-sheaf})，
$\mathcal{G}'$ 满足 $\tau$-覆盖的层性质。
由于 $\{X \to Y\}$ 是 $\tau$-覆盖，限制
$\mathcal{G}'(Y) \to \mathcal{G}'(X)$ 是单射。因此
$\overline{\theta} = 0$。
\end{proof}

\begin{lemma}
\label{sites-cohomology-lemma-induction-step-V-C-general}
在情形 \ref{sites-cohomology-situation-compare} 中，有 $(V_n) \Rightarrow (V_{n + 1})$。
\end{lemma}

\begin{proof}
设 $X$ 是 $\mathcal{C}$ 中的对象，$\mathcal{F}$ 是 $\mathcal{A}_X$ 中的对象。
对某个 $U/X$，设 $\xi \in H^{n + 1}_\tau(U, \mathcal{F})$。
我们需要证明：$\xi$ 在
某个 $\tau'$-覆盖的各成员上的限制为零。参见
引理 \ref{sites-cohomology-lemma-higher-direct-images}。
由此可知，我们可以将
$U$ 替换为某个 $\tau'$-覆盖的各成员。

\medskip\noindent
由上同调的局部性
（引理 \ref{sites-cohomology-lemma-kill-cohomology-class-on-covering}），
我们可以选取一个 $\tau$-覆盖 $\{U_i \to U\}$，
使得 $\xi$ 在 $U_i$ 上的限制为零。
按照 (\ref{sites-cohomology-item-refine-tau-by-P})，选取 $\{V_j \to V\}$、$\{f_j : W_j \to V_j\}$，以及
$\{W_{jk} \to W_j\}$。
将 $U$ 替换为 $V_j$
并将 $\mathcal{F}$ 替换为其在
$\mathcal{C}_\tau/V_j$ 上的限制；由
(\ref{sites-cohomology-item-base-change-P})，这是允许的。于是我们归结为
下一段所讨论的情形。

\medskip\noindent
这里 $f : X \to Y$ 是 $\mathcal{P}$ 中的元素，
使得 $\{X \to Y\}$ 是 $\tau$-覆盖，
$\mathcal{F}$ 是 $\mathcal{A}_Y$ 中的对象，并且
$\xi \in H^{n + 1}_\tau(Y, \mathcal{F})$ 满足：
存在一个 $\tau'$-覆盖 $\{X_i \to X\}_{i \in I}$，
使得对所有 $i \in I$，$\xi$ 在 $X_i$ 上的限制为零。
问题：证明 $\xi$ 在 $Y$ 的某个 $\tau'$-覆盖上的限制为零。

\medskip\noindent
由引理 \ref{sites-cohomology-lemma-V-implies-cohomology-extra-general}，存在一个
唯一的 $\tau'$-上同调类
$\theta \in H^{n + 1}_{\tau'}(X, \epsilon_{X, *}\mathcal{F})$
其像为 $\xi|_X$。
由于 $\xi|_X$ 经两个投影拉回到 $X \times_Y X$ 后得到相同的类，
由此（利用唯一性），$\theta$ 也具有同样性质。
由引理 \ref{sites-cohomology-lemma-make-class-zero-general}
可知，将 $Y$ 替换为某个 $\tau'$-覆盖的各成员后，
可以设 $\theta = 0$。
因此可以设 $\xi|_X$ 为零。

\medskip\noindent
设 $f : X \to Y$ 是 $\mathcal{P}$ 中的一个元素，
使得 $\{X \to Y\}$ 是 $\tau$-覆盖，
$\mathcal{F}$ 是 $\mathcal{A}_Y$ 中的对象，并且
$\xi \in H^{n + 1}_\tau(Y, \mathcal{F})$
在 $H^{n + 1}_\tau(X, \mathcal{F})$ 中映为零。
问题：证明 $\xi$ 在 $Y$ 的某个 $\tau'$-覆盖上的限制为零。

\medskip\noindent
这些假设说明 $\xi$ 在下述态射下的像为零：
$$
\mathcal{F} \longrightarrow Rf_{\tau, *}f_\tau^{-1}\mathcal{F}
$$
使用引理 \ref{sites-cohomology-lemma-Leray-unbounded}。
利用截断的有区别三角
（《导出范畴》，评注
\ref{derived-remark-truncation-distinguished-triangle}）可作一个简单论证，说明
$\xi$ 在下述态射下的像为零：
$$
\mathcal{F} \longrightarrow \tau_{\leq n}Rf_{\tau, *}f_\tau^{-1}\mathcal{F}
$$
我们将把它与态射 $\epsilon_{Y, *}\mathcal{F} \to K$ 比较，
其中
$$
K = \tau_{\leq n}Rf_{\tau', *}f_{\tau'}^{-1}\epsilon_{Y, *}\mathcal{F} =
\tau_{\leq n}Rf_{\tau', *}\epsilon_{X, *}f_{\tau}^{-1}\mathcal{F}
$$
等式
$\epsilon_{X, *} f_\tau^{-1} = f_{\tau'}^{-1} \epsilon_{Y, *}$
是 (\ref{sites-cohomology-equation-commutative-epsilon}) 的一个性质。考虑态射
$$
Rf_{\tau', *}\epsilon_{X, *}f_{\tau}^{-1}\mathcal{F} \longrightarrow
Rf_{\tau', *}R\epsilon_{X, *}f_{\tau}^{-1}\mathcal{F} =
R\epsilon_{Y, *}Rf_{\tau, *}f_\tau^{-1}\mathcal{F}
$$
该态射用于引理 \ref{sites-cohomology-lemma-V-implies-C-general} 的证明，
并由伴随诱导出态射
$$
\epsilon_Y^{-1} Rf_{\tau', *}\epsilon_{X, *}f_{\tau}^{-1}\mathcal{F} \to
Rf_{\tau, *}f_\tau^{-1}\mathcal{F}
$$
对其作截断，得到态射
$$
\epsilon_Y^{-1}K
\longrightarrow
\tau_{\leq n}Rf_{\tau, *}f_\tau^{-1}\mathcal{F}
$$
由引理 \ref{sites-cohomology-lemma-V-implies-C-general}，该态射是同构；
该引理适用，因为 $f_\tau^{-1}\mathcal{F}$ 属于 $\mathcal{A}_X$，
这是引理 \ref{sites-cohomology-lemma-A} 的结论。选取一个有区别三角
$$
\epsilon_{Y, *}\mathcal{F} \to K \to L \to \epsilon_{Y, *}\mathcal{F}[1]
$$
态射 $\mathcal{F} \to f_{\tau, *}f_\tau^{-1}\mathcal{F}$
由于 $\{X \to Y\}$ 是 $\tau$-覆盖，该态射是单射。因此
$\epsilon_{Y, *}\mathcal{F} \to
\epsilon_{Y, *}f_{\tau, *}f_\tau^{-1}\mathcal{F} =
f_{\tau', *}f_{\tau'}^{-1}\epsilon_{Y, *}\mathcal{F}$
也为单射。
因此 $L$ 的非零上同调层只可能位于次数 $0, \ldots, n$。
由于 $f_{\tau', *}f_{\tau'}^{-1}\epsilon_{Y, *}\mathcal{F}$
根据 (\ref{sites-cohomology-item-restriction-A}) 和
(\ref{sites-cohomology-item-A-and-P}) 属于 $\mathcal{A}'_Y$，故
$$
H^0(L) = \Coker(\epsilon_{Y, *}\mathcal{F} \to
f_{\tau', *}f_{\tau'}^{-1}\epsilon_{Y, *}\mathcal{F})
$$
属于弱 Serre 子范畴 $\mathcal{A}'_Y$。当 $1 \leq i \leq n$ 时，
我们看到 $H^i(L) = R^if_{\tau', *}f_{\tau'}^{-1}\epsilon_{Y, *}\mathcal{F}$
根据 (\ref{sites-cohomology-item-restriction-A}) 和
(\ref{sites-cohomology-item-A-and-P}) 属于 $\mathcal{A}'_Y$。
用 $\epsilon_Y$ 拉回上面的有区别三角，得到有区别三角
$$
\mathcal{F} \to \tau_{\leq n}Rf_{\tau, *}f_\tau^{-1}\mathcal{F}
\to \epsilon_Y^{-1}L \to \mathcal{F}[1]
$$
由于 $\xi$ 在中间项中的像为零，我们得到
$\xi$ 是某个元素的像：
$\xi' \in H^n_\tau(Y, \epsilon_Y^{-1}L)$.
由引理 \ref{sites-cohomology-lemma-V-implies-cohomology-general}，有
$$
H^n_{\tau'}(Y, L) = H^n_\tau(Y, \epsilon_Y^{-1}L),
$$
因此可以将 $\xi'$ 提升为 $H^n_{\tau'}(Y, L)$ 中的元素，
并取其边界映到
$H^{n + 1}_{\tau'}(Y, \epsilon_{Y, *}\mathcal{F})$
从而 $\xi$ 属于规范态射的像：
$H^{n + 1}_{\tau'}(Y, \epsilon_{Y, *}\mathcal{F}) \to
H^{n + 1}_\tau(Y, \mathcal{F})$.
由上同调的局部性，对于
$H^{n + 1}_{\tau'}(Y, \epsilon_{Y, *}\mathcal{F})$，参见
引理 \ref{sites-cohomology-lemma-kill-cohomology-class-on-covering}，
得证。
\end{proof}

\begin{lemma}
\label{sites-cohomology-lemma-V-C-all-n-general}
在情形 \ref{sites-cohomology-situation-compare} 中，有
$(V_n)$ 对所有 $n$ 都成立。此外：
\begin{enumerate}
\item 对 $X$ 是 $\mathcal{C}$ 中的对象，且
$K' \in D^+_{\mathcal{A}'_X}(\mathcal{C}_{\tau'}/X)$ 上的态射
$K' \to R\epsilon_{X, *}(\epsilon_X^{-1}K')$ 是同构。
\item 对 $f : X \to Y$ 属于 $\mathcal{P}$，且
$K' \in D^+_{\mathcal{A}'_X}(\mathcal{C}_{\tau'}/X)$，有
$Rf_{\tau', *}K' \in D^+_{\mathcal{A}'_X}(\mathcal{C}_{\tau'}/Y)$，且
$\epsilon_Y^{-1}(Rf_{\tau', *}K') = Rf_{\tau, *}(\epsilon_X^{-1}K')$.
\end{enumerate}
\end{lemma}

\begin{proof}
注意 $(V_0)$ 成立，因为这是空条件。
于是由
引理 \ref{sites-cohomology-lemma-induction-step-V-C-general} 得到所有 $n$ 的 $(V_n)$。

\medskip\noindent
证明 (1)。对象 $K = \epsilon_X^{-1}K'$ 的上同调层
$H^i(K) = \epsilon_X^{-1}H^i(K')$ 属于 $\mathcal{A}_X$。
因此有谱序列
$$
E_2^{p, q} = R^p\epsilon_{X, *} H^q(K) \Rightarrow
H^{p + q}(R\epsilon_{X, *}K)
$$
由 $(V_n)$ 对所有 $n$ 成立，该谱序列退化，从而
$$
H^n(R\epsilon_{X, *}K) = \epsilon_{X, *}H^n(K) =
\epsilon_{X, *}\epsilon_X^{-1}H^i(K') = H^i(K').
$$
再次利用 $H^i(K')$ 属于 $\mathcal{A}'_X$ 这一事实。
因此规范态射 $K' \to R\epsilon_{X, *}(\epsilon_X^{-1}K')$
是同构。

\medskip\noindent
证明 (2)。利用谱序列
$$
E_2^{p, q} = R^pf_{\tau', *}H^q(K') \Rightarrow R^{p + q}f_{\tau', *}K'
$$
$R^pf_{\tau', *}H^q(K')$ 属于
$\mathcal{A}'_Y$（由 (\ref{sites-cohomology-item-A-and-P})），
以及 $\mathcal{A}'_Y$ 是
$\textit{Ab}(\mathcal{C}_{\tau'}/Y)$ 的弱 Serre 子范畴这一事实，
再由同调，引理 \ref{homology-lemma-biregular-ss-converges}
可知
$Rf_{\tau', *}K' \in D^+_{\mathcal{A}'_X}(\mathcal{C}_{\tau'}/X)$.
为完成证明，需证明基变换态射
$$
\epsilon_Y^{-1}(Rf_{\tau', *}K')
\longrightarrow
Rf_{\tau, *}(\epsilon_X^{-1}K')
$$
它是同构。将上面的谱序列与
谱序列
$$
E_2^{p, q} = R^pf_{\tau, *}H^q(\epsilon_X^{-1}K')
\Rightarrow R^{p + q}f_{\tau, *}\epsilon_X^{-1}K'
$$
由此可将问题化归为 $K'$ 只有一个非零
上同调层 $\mathcal{F}' \in \mathcal{A}'_X$ 的情形；细节略去。
然后由引理 \ref{sites-cohomology-lemma-V-implies-C-general} 得到
$\epsilon_Y^{-1}R^if_{\tau', *}\mathcal{F}' =
R^if_{\tau, *}\epsilon_X^{-1}\mathcal{F}'$ 对所有 $i$ 均成立，
证明完成。
\end{proof}

\begin{lemma}
\label{sites-cohomology-lemma-cohomological-descent-general}
在情形 \ref{sites-cohomology-situation-compare} 中，对于任意 $X$，
$\mathcal{C}$ 中的范畴
$\mathcal{A}_X \subset \textit{Ab}(\mathcal{C}_\tau/X)$
是弱 Serre 子范畴，且函子
$$
R\epsilon_{X, *} :
D^+_{\mathcal{A}_X}(\mathcal{C}_\tau/X)
\longrightarrow
D^+_{\mathcal{A}'_X}(\mathcal{C}_{\tau'}/X)
$$
是等价；其拟逆由 $\epsilon_X^{-1}$ 给出。
\end{lemma}

\begin{proof}
我们需要检查同调引理
\ref{homology-lemma-characterize-weak-serre-subcategory} 中列出的
关于 $\mathcal{A}_X$ 的条件。若 $\varphi : \mathcal{F} \to \mathcal{G}$ 是
$\mathcal{A}_X$ 中的态射，则 $\epsilon_{X, *}\varphi$：
$\epsilon_{X, *}\mathcal{F} \to \epsilon_{X, *}\mathcal{G}$
是 $\mathcal{A}'_X$ 中的态射。因此 $\Ker(\epsilon_{X, *}\varphi)$ 与
$\Coker(\epsilon_{X, *}\varphi)$ 是 $\mathcal{A}'_X$ 中的对象，
因为 $\mathcal{A}'_X$ 是 $\textit{Ab}(\mathcal{C}_{\tau'}/X)$ 的弱 Serre 子范畴。
应用 $\epsilon_X^{-1}$，得到正合序列
$$
0 \to
\epsilon_X^{-1}\Ker(\epsilon_{X, *}\varphi) \to
\mathcal{F} \to \mathcal{G} \to
\epsilon_X^{-1}\Coker(\epsilon_{X, *}\varphi) \to 0
$$
由此可见 $\Ker(\varphi)$ 和 $\Coker(\varphi)$ 属于
$\mathcal{A}_X$。最后，设
$$
0 \to \mathcal{F}_1 \to \mathcal{F}_2 \to \mathcal{F}_3 \to 0
$$
这是 $\textit{Ab}(\mathcal{C}_\tau/X)$ 中的短正合序列，
且 $\mathcal{F}_1$ 与 $\mathcal{F}_3$ 属于 $\mathcal{A}_X$。
随后应用 $\epsilon_{X, *}$，得到正合序列
$$
0 \to 
\epsilon_{X, *}\mathcal{F}_1 \to
\epsilon_{X, *}\mathcal{F}_2 \to
\epsilon_{X, *}\mathcal{F}_3 \to
R^1\epsilon_{X, *}\mathcal{F}_1 = 0
$$
由引理 \ref{sites-cohomology-lemma-V-C-all-n-general}，该项消失。
因此 $\epsilon_{X, *}\mathcal{F}_2$ 属于 $\mathcal{A}'_X$
因为它是 $\textit{Ab}(\mathcal{C}_{\tau'}/X)$ 的弱 Serre 子范畴。
用 $\epsilon_X$ 拉回，得到
$\mathcal{F}_2$ 属于 $\mathcal{A}_X$。

\medskip\noindent
因此 $\mathcal{A}_X$ 是
$\textit{Ab}(\mathcal{C}_\tau/X)$ 的弱 Serre 子范畴，于是考虑
范畴 $D^+_{\mathcal{A}_X}(\mathcal{C}_\tau/X)$ 是有意义的。
注意 $\epsilon_X^{-1} : \mathcal{A}'_X \to \mathcal{A}_X$
是等价，并且
$\mathcal{F}' \to R\epsilon_{X, *}\epsilon_X^{-1}\mathcal{F}'$
对 $\mathcal{F}' \in \mathcal{A}'_X$，这是同构，因为我们
根据引理 \ref{sites-cohomology-lemma-V-C-all-n-general} 已知 $(V_n)$ 对所有 $n$ 成立。
因此由引理 \ref{sites-cohomology-lemma-equivalence-bounded} 得出结论。
\end{proof}

\begin{lemma}
\label{sites-cohomology-lemma-compare-cohomology-general}
在情形 \ref{sites-cohomology-situation-compare} 中，设 $X$ 是 $\mathcal{C}$ 中的对象。
\begin{enumerate}
\item 对 $\mathcal{F}' \in \mathcal{A}'_X$，有
$H^n_{\tau'}(X, \mathcal{F}') = H^n_\tau(X, \epsilon_X^{-1}\mathcal{F}')$,
\item 对 $K' \in D^+_{\mathcal{A}'_X}(\mathcal{C}_{\tau'}/X)$，
有 $H^n_{\tau'}(X, K') = H^n_\tau(X, \epsilon_X^{-1}K')$。
\end{enumerate}
\end{lemma}

\begin{proof}
由引理 \ref{sites-cohomology-lemma-V-C-all-n-general}
以及评注 \ref{sites-cohomology-remark-before-Leray} 得出。
\end{proof}
























\section{豪斯多夫且局部拟紧空间上的上同调}
\label{sites-cohomology-section-cohomology-LC}

\noindent
我们沿用“豪斯多夫且局部拟紧”的说法，
而不采用文献中常见的“局部紧”说法。
记 $\textit{LC}$ 为如下范畴：其对象是豪斯多夫且
局部拟紧的拓扑空间，
其态射是连续映射。


\begin{lemma}
\label{sites-cohomology-lemma-LC-basic}
范畴 $\textit{LC}$ 有纤维积和终对象，因而有任意有限极限。若态射 $X \to Z$ 与 $Y \to Z$
在 $\textit{LC}$ 中，若
$X$ 与 $Y$ 拟紧，则
$X \times_Z Y$ 拟紧。
\end{lemma}

\begin{proof}
终对象是单点空间。给定 $X \to Z$ 与
$Y \to Z$ 这两个 $\textit{LC}$ 中的态射，纤维积 $X \times_Z Y$ 是
$X \times Y$ 的子空间。因此 $X \times_Z Y$ 是豪斯多夫的，因为
$X \times Y$ 由
拓扑，第 \ref{topology-section-Hausdorff} 节知是豪斯多夫的。

\medskip\noindent
若 $X$ 与 $Y$ 拟紧，则由
拓扑，定理 \ref{topology-theorem-tychonov}，$X \times Y$ 拟紧。
由于 $X \times_Z Y$ 是 $X \times Y$ 的闭子集
（拓扑，引理 \ref{topology-lemma-fibre-product-closed}），
由此根据
拓扑，引理 \ref{topology-lemma-closed-in-quasi-compact}，$X \times_Z Y$ 拟紧。

\medskip\noindent
最后回到一般情形。若 $x \in X$ 且 $y \in Y$，
可取拟紧邻域 $x \in E \subset X$ 以及
$y \in F \subset Y$，由上述结果可知 $E \times_Z F$ 是
$(x, y)$ 的拟紧邻域。因此 $X \times_Z Y$
由
拓扑，引理 \ref{topology-lemma-locally-quasi-compact-Hausdorff}，是 $\textit{LC}$ 中的对象。
\end{proof}

\noindent
我们可以在 $\textit{LC}$ 上赋予比通常拓扑更强的拓扑。

\begin{definition}
\label{sites-cohomology-definition-covering-LC}
设 $\{f_i : X_i \to X\}$ 是一族目标固定的态射，
其目标位于范畴 $\textit{LC}$ 中。我们称此族为
{\it qc 覆盖}\footnote{这是非标准记号。
我们采用它是为了使读者联想到概形的 fpqc 覆盖。}
若对每个 $x \in X$，存在 $i_1, \ldots, i_n \in I$ 以及
拟紧子集 $E_j \subset X_{i_j}$，使得
$\bigcup f_{i_j}(E_j)$ 是 $x$ 的一个邻域。
\end{definition}

\noindent
注意，$X = \bigcup U_i$ 作为 $\textit{LC}$ 中对象的开覆盖
给出 qc 覆盖 $\{U_i \to X\}$，因为 $X$ 是局部拟紧的。
我们先给出必要的引理。

\begin{lemma}
\label{sites-cohomology-lemma-qc}
设 $X$ 是豪斯多夫且局部拟紧的空间，换言之，
是 $\textit{LC}$ 中的对象。
\begin{enumerate}
\item 若 $X' \to X$ 是 $\textit{LC}$ 中的同构，则
$\{X' \to X\}$ 是 qc 覆盖。
\item 若 $\{f_i : X_i \to X\}_{i\in I}$ 是 qc 覆盖，且对每个
$i$ 有 qc 覆盖 $\{g_{ij} : X_{ij} \to X_i\}_{j\in J_i}$，则
$\{X_{ij} \to X\}_{i \in I, j\in J_i}$ 是 qc 覆盖。
\item 若 $\{X_i \to X\}_{i\in I}$ 是 qc 覆盖，
且 $X' \to X$ 是 $\textit{LC}$ 中的态射，则
$\{X' \times_X X_i \to X'\}_{i\in I}$ 是 qc 覆盖。
\end{enumerate}
\end{lemma}

\begin{proof}
第 (1) 项由上面关于开覆盖是 qc 覆盖的评注立即得到。

\medskip\noindent
证明 (2)。设 $x \in X$。选取 $i_1, \ldots, i_n \in I$ 以及
$E_a \subset X_{i_a}$ 为拟紧集，使得 $\bigcup f_{i_a}(E_a)$
是 $x$ 的邻域。对每个 $e \in E_a$，可找到
有限集 $J_e \subset J_{i_a}$ 以及拟紧集
$F_{e, j} \subset X_{ij}$（$j \in J_e$），使得 $\bigcup g_{ij}(F_{e, j})$
是 $e$ 的邻域。由于 $E_a$ 拟紧，可找到
有限个点 $e_1, \ldots, e_{m_a}$，使得
$$
E_a \subset
\bigcup\nolimits_{k = 1, \ldots, m_a}
\bigcup\nolimits_{j \in J_{e_k}} g_{ij}(F_{e_k, j})
$$
于是得到
$$
\bigcup\nolimits_{a = 1, \ldots, n}
\bigcup\nolimits_{k = 1, \ldots, m_a}
\bigcup\nolimits_{j \in J_{e_k}} f_i(g_{ij}(F_{e_k, j}))
$$
是 $x$ 的邻域。

\medskip\noindent
证明 (3)。设 $x' \in X'$ 是一个点，设其像 $x \in X$。
选取 $i_1, \ldots, i_n \in I$ 以及拟紧子集
$E_j \subset X_{i_j}$，使得 $\bigcup f_{i_j}(E_j)$ 是一个
$x$ 的邻域。选取 $X'$ 中 $x'$ 的拟紧邻域 $F \subset X'$，
其映入 $x$ 的拟紧邻域 $\bigcup f_{i_j}(E_j)$。于是
$F \times_X E_j \subset X' \times_X X_{i_j}$ 是
拟紧子集，而 $F$ 是态射
$\coprod F \times_X E_j \to F$ 的像。因此基变换是
qc 覆盖，证明完成。
\end{proof}

\noindent
由于 $\textit{LC}$ 的所有对象都是豪斯多夫空间，$\textit{LC}$ 中的任一态射
$f : X \to Y$ 都是拓扑空间之间的分离连续映射。于是 $f$ 是拓扑空间中的 proper 映射
当且仅当它是泛闭的。参见
拓扑，第 \ref{topology-section-proper} 节的讨论。


\begin{lemma}
\label{sites-cohomology-lemma-proper-surjective-is-qc-covering}
设 $f : X \to Y$ 是 $\textit{LC}$ 中的态射。
若 $f$ proper 且满射，则 $\{f : X \to Y\}$
是 qc 覆盖。
\end{lemma}

\begin{proof}
设 $y \in Y$ 是一个点。对每个 $x \in X_y$，选取一个拟紧
邻域 $E_x \subset X$。选取开集 $x \in U_x \subset E_x$。
由于 $f$ proper，纤维 $X_y$ 拟紧，因此可找到
$x_1, \ldots, x_n \in X_y$ 使得
$X_y \subset U_{x_1} \cup \ldots \cup U_{x_n}$。
我们断言 $f(E_{x_1}) \cup \ldots \cup f(E_{x_n})$ 是 $y$ 的邻域。
事实上，由于 $f$ 是闭映射
（拓扑，定理 \ref{topology-theorem-characterize-proper}）
可知 $Z = f(X \setminus U_{x_1} \cup \ldots \cup U_{x_n})$
是 $Y$ 中不含 $y$ 的闭子集。由于 $f$ 满射，
可知 $Y \setminus Z$ 包含于
$f(E_{x_1}) \cup \ldots \cup f(E_{x_n})$，正如所需。
\end{proof}
\noindent
除集合论方面的一些问题外，引理 \ref{sites-cohomology-lemma-qc} 表明
$\textit{LC}$ 与 qc 覆盖的集合构成一个站点。
我们将把这个站点（适当修改以克服
集合论问题）记为
$\textit{LC}_{qc}$。

\begin{remark}[集合论问题]
\label{sites-cohomology-remark-set-theoretic-LC}
范畴 $\textit{LC}$ 是一个“大的”范畴，因为其对象构成
真类；覆盖同样构成真类。
我们定义拓扑空间 $X$ 的 {\it size} 为
其点集的基数。选取基数上的函数
$Bound$，例如见集合，
等式 (\ref{sets-equation-bound})。
最后，令 $S_0$ 是 $\textit{LC}$ 中对象的初始集合，
例如 $S_0 = \{(\mathbf{R}, \text{欧几里得拓扑})\}$。
完全类似于集合，引理 \ref{sets-lemma-construct-category}
可选取极限序数 $\alpha$，使得
$\textit{LC}_\alpha = \textit{LC} \cap V_\alpha$
包含 $S_0$，并且在 $\textit{LC}$ 中存在的所有可数极限与
余极限下保持不变。此外，若 $X \in \textit{LC}_\alpha$
且 $Y \in \textit{LC}$，并且
$\text{size}(Y) \leq Bound(\text{size}(X))$，则 $Y$ 同构于
$\textit{LC}_\alpha$ 中的某个对象。
接着应用集合，引理 \ref{sets-lemma-coverings-site}
在 $\textit{LC}_\alpha$ 上选取 qc 覆盖的集合 $\text{Cov}$，
使得 $\textit{LC}_\alpha$ 中的每个 qc 覆盖都
在组合意义下等价于该集合中的某个覆盖。
这样得到站点 $(\textit{LC}_\alpha, \text{Cov})$，
记为 $\textit{LC}_{qc}$。
\end{remark}

\noindent
$\textit{LC}_{qc}$ 站点（评注 \ref{sites-cohomology-remark-set-theoretic-LC}）上还有第二个拓扑。
具体地，对给定对象
$X$，考虑 $\textit{LC}_{qc}$ 的所有覆盖 $\{X_i \to X\}$，
其中 $X_i \to X$ 是开浸入。我们将此站点记为
$\textit{LC}_{Zar}$。恒等函子
$\textit{LC}_{Zar} \to \textit{LC}_{qc}$ 连续，并定义
一个站点态射
$$
\epsilon : \textit{LC}_{qc} \longrightarrow \textit{LC}_{Zar}
$$
参见第 \ref{sites-cohomology-section-compare} 节。
对于豪斯多夫且局部拟紧的拓扑空间 $X$，更准确地说，
对于 $X \in \Ob(\textit{LC}_{qc})$，记诱导的态射为
$$
\epsilon_X : \textit{LC}_{qc}/X \longrightarrow \textit{LC}_{Zar}/X
$$
（参见站点，引理 \ref{sites-lemma-localize-morphism}）。
令 $X_{Zar}$ 为对象是 $X$ 的开集的站点，参见
站点，例 \ref{sites-example-site-topological}。
存在站点态射
$$
\pi_X : \textit{LC}_{Zar}/X \longrightarrow X_{Zar}
$$
由连续函子
$X_{Zar} \to \textit{LC}_{Zar}/X$，$U \mapsto U$ 给出。
具体地，$X_{Zar}$ 有纤维积和终对象，而上述
函子与它们交换，因此站点，命题 \ref{sites-proposition-get-morphism} 适用。
我们常将 $\pi$ 看作拓扑斯态射
$$
\pi_X : \Sh(\textit{LC}_{Zar}/X) \longrightarrow \Sh(X)
$$
利用等式 $\Sh(X_{Zar}) = \Sh(X)$。

\begin{lemma}
\label{sites-cohomology-lemma-describe-pullback-pi}
设 $X$ 是 $\textit{LC}_{qc}$ 中的对象，设 $\mathcal{F}$ 是
$X$ 上的层。该规则
$$
\textit{LC}_{qc}/X \longrightarrow \textit{Sets},\quad
(f : Y \to X) \longmapsto \Gamma(Y, f^{-1}\mathcal{F})
$$
是一个层，因而当然也是 $\textit{LC}_{Zar}/X$ 上的层。
该层在 $\textit{LC}_{Zar}/X$ 上等于
$\pi_X^{-1}\mathcal{F}$，并在 $\textit{LC}_{qc}/X$ 上等于
$\epsilon_X^{-1}\pi_X^{-1}\mathcal{F}$。
\end{lemma}

\begin{proof}
记由引理中的公式定义的预层为 $\mathcal{G}$。
当然，公式中的拉回 $f^{-1}$ 表示拓扑空间上层的通常
拉回。由定义立即可知，$\mathcal{G}$ 是 Zar
拓扑下的层。

\medskip\noindent
设 $Y \to X$ 是 $\textit{LC}_{qc}$ 中的态射。设
$\mathcal{V} = \{g_i : Y_i \to Y\}_{i \in I}$ 是 qc 覆盖。
要证明 $\mathcal{G}$ 是 qc 拓扑下的层，
只需证明
$\mathcal{G}(Y) \to H^0(\mathcal{V}, \mathcal{G})$
是同构（参见站点，第 \ref{sites-section-sheafification} 节）。
首先注意该态射是单射，因为 qc 覆盖
是满射，并且可以在茎上检验截面的相等性；
（使用层，引理 \ref{sheaves-lemma-sheaf-subset-stalks} 以及
\ref{sheaves-lemma-stalk-pullback-presheaf}）。因此
$\mathcal{G}$ 是 $\textit{LC}_{qc}$ 上的分离预层，
所以只需证明任意元素
$(s_i) \in H^0(\mathcal{V}, \mathcal{G})$
在将 $\mathcal{V}$ 替换为某个加细后，
映到 $\mathcal{G}(Y)$ 的像中
（站点，定理 \ref{sites-theorem-plus}）。

\medskip\noindent
将 $Y_{i, Zar}$ 上的层与 $Y_i$ 上的层等同，我们看到
$\mathcal{G}|_{Y_{i, Zar}}$ 是连续映射 $g_i : Y_i \to Y$
下的 $f^{-1}\mathcal{F}$ 的拉回。因此可以选取一个开覆盖
$Y_i = \bigcup V_{ij}$，使得对每个 $j$，存在开集
$W_{ij} \subset Y$ 以及截面 $t_{ij} \in \mathcal{G}(W_{ij})$，
满足 $V_{ij}$ 映入 $W_{ij}$，并且
$s|_{V_{ij}}$ 是 $t_{ij}$ 的拉回。换言之，将覆盖
$\{Y_i \to Y\}$ 加细后，可以假设存在开集 $W_i \subset Y$，
使得 $Y_i \to Y$ 经由 $W_i$ 分解，并且在 $W_i$ 上有
$\mathcal{G}$ 的截面 $t_i$，其限制是给定的截面 $s_i$。
此外，若 $y \in Y$ 同时属于 $Y_i \to Y$ 和 $Y_j \to Y$ 的像，
则茎 $f^{-1}\mathcal{F}_y$ 中的像 $t_{i, y}$ 与 $t_{j, y}$ 相等
（因为 $s_i$ 与 $s_j$ 在 $Y_i \times_Y Y_j$ 上相等）。
因此，对 $y \in Y$，存在一个定义良好的元素 $t_y \in
f^{-1}\mathcal{F}_y$；当 $y$ 属于 $Y_i \to Y$ 的像时，
它与 $t_{i, y}$ 相等。我们将证明元素 $(t_y)$ 来自 $Y$ 上
$f^{-1}\mathcal{F}$ 的一个全局截面，这就完成引理的证明。

\medskip\noindent
只需证明该结论在 $Y$ 上局部成立；参见层，第
\ref{sheaves-section-sheafification} 节。令 $y_0 \in Y$。
选取 $i_1, \ldots, i_n \in I$ 以及拟紧子集
$E_j \subset Y_{i_j}$，使得 $\bigcup g_{i_j}(E_j)$ 是 $y_0$
的一个邻域。令 $V \subset Y$ 是 $y_0$ 的一个开邻域，且包含于
$\bigcup g_{i_j}(E_j)$ 以及 $W_{i_1} \cap \ldots \cap W_{i_n}$。
由于 $t_{i_1, y_0} = \ldots = t_{i_n, y_0}$，缩小 $V$ 后可假设
$f^{-1}\mathcal{F}$ 的截面 $t_{i_j}|_V$（$j = 1, \ldots, n$）
彼此相等。由于 $V \subset \bigcup g_{i_j}(E_j)$，可知
$(t_y)_{y \in V}$ 来自这个截面。

\medskip\noindent
还需证明，在 $\textit{LC}_{qc}$ 上，$\mathcal{G}$ 等于
$\epsilon_X^{-1}\pi_X^{-1}\mathcal{F}$；分别地，在
$\textit{LC}_{Zar}$ 上等于 $\pi_X^{-1}\mathcal{F}$。
在这两种情形中，拉回都由取下述预层来定义：
$$
(f : Y \to X)
\longmapsto
\colim_{f(Y) \subset U \subset X} \mathcal{F}(U)
$$
然后作层化。在 Zar 拓扑下层化恰好得到我们的层
$\mathcal{G}$；又因为 $\mathcal{G}$ 是 qc 层，故在 qc 拓扑下
同样成立。
\end{proof}

\noindent
令 $X \in \Ob(\textit{LC}_{Zar})$，令 $\mathcal{H}$
是 $\textit{LC}_{Zar}/X$ 上的阿贝尔层。
对于 $\textit{LC}_{Zar}/X$ 的对象 $U$，我们将用
$H^n_{Zar}(U, \mathcal{H})$ 表示 $\mathcal{H}$ 在 $U$ 上的上同调。

\begin{lemma}
\label{sites-cohomology-lemma-collect-true-things-Zar}
令 $X$ 是 $\textit{LC}_{Zar}$ 中的对象。则
\begin{enumerate}
\item 对于 $\mathcal{F} \in \textit{Ab}(X)$，有
$H^n_{Zar}(X, \pi_X^{-1}\mathcal{F}) = H^n(X, \mathcal{F})$,
\item $\pi_{X, *} : \textit{Ab}(\textit{LC}_{Zar}/X) \to \textit{Ab}(X)$
是正合的，
\item 伴随的单位 $\text{id} \to \pi_{X, *} \circ \pi_X^{-1}$
是同构，
\item 对于 $K \in D(X)$，规范映射
$K \to R\pi_{X, *} \pi_X^{-1}K$ 是同构。
\end{enumerate}
令 $f : X \to Y$ 是 $\textit{LC}_{Zar}$ 中的态射。则
\begin{enumerate}
\item[(5)] 存在交换图
$$
\xymatrix{
\Sh(\textit{LC}_{Zar}/X) \ar[r]_{f_{Zar}} \ar[d]_{\pi_X} &
\Sh(\textit{LC}_{Zar}/Y) \ar[d]^{\pi_Y} \\
\Sh(X_{Zar}) \ar[r]^f &
\Sh(Y_{Zar})
}
$$
其中对象是拓扑斯，
\item[(6)] 对于 $L \in D^+(Y)$，有
$H^n_{Zar}(X, \pi_Y^{-1}L) = H^n(X, f^{-1}L)$,
\item[(7)] 若 $f$ 是 proper 映射，则有
\begin{enumerate}
\item $\pi_Y^{-1} \circ f_* = f_{Zar, *} \circ \pi_X^{-1}$ 是从
$\Sh(X)$ 到 $\Sh(\textit{LC}_{Zar}/Y)$ 的函子，
\item $\pi_Y^{-1} \circ Rf_* = Rf_{Zar, *} \circ \pi_X^{-1}$ 是从
$D^+(X)$ 到 $D^+(\textit{LC}_{Zar}/Y)$ 的函子。
\end{enumerate}
\end{enumerate}
\end{lemma}

\begin{proof}
证明 (1)。等式 $H^n_{Zar}(X, \pi_X^{-1}\mathcal{F}) = H^n(X, \mathcal{F})$
是一个一般事实，来自如下显然观察：$\textit{LC}_{Zar}$ 中 $X$ 的覆盖
与 $X$ 的开覆盖是同一回事。若希望看到详细证明，
应将引理 \ref{sites-cohomology-lemma-cohomology-bigger-site} 应用于函子
$X_{Zar} \to \textit{LC}_{Zar}$。

\medskip\noindent
证明 (2)。这是因为存在拓扑斯态射
$\tau_X : \Sh(X_{Zar}) \to \Sh(\textit{LC}_{Zar})$，使得
$\pi_{X, *} = \tau_X^{-1}$；这一点由将站点，引理
\ref{sites-lemma-bigger-site} 应用于定义 $\pi_X$ 所用的函子
$X_{Zar} \to \textit{LC}_{Zar}/X$ 得到。

\medskip\noindent
证明 (3)。这是因为根据站点，引理
\ref{sites-lemma-bigger-site}，$\tau_X^{-1} \circ \pi_X^{-1}$
是恒等函子。也可以从引理
\ref{sites-cohomology-lemma-describe-pullback-pi} 中对 $\pi_X^{-1}$ 的显式描述推出。

\medskip\noindent
证明 (4)。将 (3) 应用于表示 $K$ 的阿贝尔层复形。

\medskip\noindent
证明 (5)。拓扑斯态射 $f_{Zar}$ 来自应用站点，引理
\ref{sites-lemma-relocalize}；在此情形中，根据站点，引理
\ref{sites-lemma-relocalize-given-fibre-products}，它来自连续函子
$Z/Y \mapsto Z \times_Y X/X$。该图交换只是因为相应的连续函子
正确复合（参见站点，引理
\ref{sites-lemma-composition-morphisms-sites}）。

\medskip\noindent
证明 (6)。对于 $\mathcal{G} \in \textit{Ab}(Y)$，由引理
\ref{sites-cohomology-lemma-cohomology-of-open}，有
$H^n_{Zar}(X, \pi_Y^{-1}\mathcal{G}) =
H^n_{Zar}(X, f_{Zar}^{-1}\pi_Y^{-1}\mathcal{G})$。
根据 (5) 中图的交换性，这等于
$H^n_{Zar}(X, \pi_X^{-1}f^{-1}\mathcal{G})$。
当 $L$ 是集中在次数 $0$ 的单个层时，由 (1) 得到结论。
一般情形则通过用有下界的阿贝尔层复形表示
$L$ 得到。

\medskip\noindent
证明 (7a)。令 $\mathcal{F}$ 是 $X$ 上的层。
令 $g : Z \to Y$ 是 $\textit{LC}_{Zar}/Y$ 的对象。考虑纤维积
$$
\xymatrix{
Z' \ar[r]_{f'} \ar[d]_{g'} & Z \ar[d]^g \\
X \ar[r]^f & Y
}
$$
则有
$$
(f_{Zar, *}\pi_X^{-1}\mathcal{F})(Z/Y) =
(\pi_X^{-1}\mathcal{F})(Z'/X)  =
\Gamma(Z', (g')^{-1}\mathcal{F})  =
\Gamma(Z, f'_*(g')^{-1}\mathcal{F})
$$
第二个等号来自引理 \ref{sites-cohomology-lemma-describe-pullback-pi}。
另一方面，
$$
(\pi_Y^{-1}f_*\mathcal{F})(Z/Y) = \Gamma(Z, g^{-1}f_*\mathcal{F})
$$
同样由引理 \ref{sites-cohomology-lemma-describe-pullback-pi} 得到。
因此，根据集合层的真基变换定理
（上同调，引理 \ref{cohomology-lemma-proper-base-change-sheaves-of-sets}），
可知这两个集合规范同构。该同构与限制映射相容，
并定义出同构
$\pi_Y^{-1}f_*\mathcal{F} = f_{Zar, *}\pi_X^{-1}\mathcal{F}$.
从而得到函子同构
$\pi_Y^{-1} \circ f_* = f_{Zar, *} \circ \pi_X^{-1}$.

\medskip\noindent
证明 (7b)。令 $K \in D^+(X)$。根据引理
\ref{sites-cohomology-lemma-unbounded-describe-higher-direct-images}，
$Rf_{Zar, *}\pi_X^{-1}K$ 的第 $n$ 个上同调层是与下述预层
相伴的层：
$$
(g : Z \to Y) \longmapsto H^n_{Zar}(Z', \pi_X^{-1}K)
$$
记号如上。注意
\begin{align*}
H^n_{Zar}(Z', \pi_X^{-1}K)
& =
H^n(Z', (g')^{-1}K) \\
& =
H^n(Z, Rf'_*(g')^{-1}K) \\
& =
H^n(Z, g^{-1}Rf_*K) \\
& =
H^n_{Zar}(Z, \pi_Y^{-1}Rf_*K)
\end{align*}
第一个等式是将 (6) 应用于 $K$ 和 $g' : Z' \to X$ 的结果。
第二个等式是对 $f' : Z' \to Z$ 应用 Leray
（上同调，引理 \ref{cohomology-lemma-before-Leray}）。
第三个等式是真基变换定理
（上同调，定理 \ref{cohomology-theorem-proper-base-change}）。
第四个等式是将 (6) 应用于 $g : Z \to Y$ 和 $Rf_*K$ 的结果。
因此 $Rf_{Zar, *}\pi_X^{-1}K$ 与 $\pi_Y^{-1}Rf_*K$ 具有相同的
上同调层。我们省略对规范基变换映射
$\pi_Y^{-1}Rf_*K \to Rf_{Zar, *}\pi_X^{-1}K$
诱导出该同构的验证。
\end{proof}

\noindent
在引理 \ref{sites-cohomology-lemma-describe-pullback-pi} 的情形中，
$\epsilon$ 与 $\pi$ 的复合以及等式
$\Sh(X) = \Sh(X_{Zar})$ 确定一个拓扑斯态射
$$
a_X : \Sh(\textit{LC}_{qc}/X) \longrightarrow \Sh(X)
$$

\begin{lemma}
\label{sites-cohomology-lemma-push-pull-LC}
令 $f : X \to Y$ 是 $\textit{LC}_{qc}$ 中的态射。
则存在拓扑斯的交换图
$$
\vcenter{
\xymatrix{
\Sh(\textit{LC}_{qc}/X) \ar[r]_{f_{qc}} \ar[d]_{\epsilon_X} &
\Sh(\textit{LC}_{qc}/Y) \ar[d]^{\epsilon_Y} \\
\Sh(\textit{LC}_{Zar}/X) \ar[r]^{f_{Zar}} &
\Sh(\textit{LC}_{Zar}/Y)
}
}
\quad\text{和}\quad
\vcenter{
\xymatrix{
\Sh(\textit{LC}_{qc}/X) \ar[r]_{f_{qc}} \ar[d]_{a_X} &
\Sh(\textit{LC}_{qc}/Y) \ar[d]^{a_Y} \\
\Sh(X) \ar[r]^f &
\Sh(Y)
}
}
$$
其中 $a_X = \pi_X \circ \epsilon_X$，$a_Y = \pi_X \circ \epsilon_X$。
若 $f$ 是 proper 映射，则
$a_Y^{-1} \circ f_* = f_{qc, *} \circ a_X^{-1}$。
\end{lemma}

\begin{proof}
拓扑斯态射 $f_{qc}$ 是站点，引理
\ref{sites-lemma-relocalize} 中的那个态射；在此情形中，
它来自连续函子 $Z/Y \mapsto Z \times_Y X/X$，参见站点，
引理 \ref{sites-lemma-relocalize-given-fibre-products}。
左图交换是因为相应的连续函子正确复合
（参见站点，引理 \ref{sites-lemma-composition-morphisms-sites}）。
右图交换是因为左图交换，并且因为引理
\ref{sites-cohomology-lemma-collect-true-things-Zar} 的 (5) 部分。

\medskip\noindent
证明最后的断言。读者可以重复引理
\ref{sites-cohomology-lemma-collect-true-things-Zar} 的 (7a) 部分的证明；
这里我们改为从该部分推出结论。由于 $\epsilon_{Y, *}$ 是底层预层
上的恒等函子，它反映同构。引理
\ref{sites-cohomology-lemma-describe-pullback-pi} 中的描述表明
$\epsilon_{Y, *} \circ a_Y^{-1} = \pi_Y^{-1}$，对 $X$ 也同样成立。
要证明规范映射
$a_Y^{-1}f_*\mathcal{F} \to f_{qc, *}a_X^{-1}\mathcal{F}$
是同构，只需证明
$$
\pi_Y^{-1}f_*\mathcal{F} =
\epsilon_{Y, *}a_Y^{-1}f_*\mathcal{F} \to
\epsilon_{Y, *}f_{qc, *}a_X^{-1}\mathcal{F} =
f_{Zar, *}\epsilon_{X, *}a_X^{-1}\mathcal{F} =
f_{Zar, *}\pi_X^{-1}\mathcal{F}
$$
是同构。这正是引理
\ref{sites-cohomology-lemma-collect-true-things-Zar} 的 (7a) 部分。
\end{proof}

\begin{lemma}
\label{sites-cohomology-lemma-compare-qc-zar}
考虑比较态射
$\epsilon : \textit{LC}_{qc} \to \textit{LC}_{Zar}$。
令 $\mathcal{P}$ 表示拓扑空间的 proper 映射类。
对 $\textit{LC}_{Zar}$ 中的 $X$，记
$\mathcal{A}'_X \subset \textit{Ab}(\textit{LC}_{Zar}/X)$
为由形如 $\pi_X^{-1}\mathcal{F}$ 的层组成的全子范畴，
其中 $\mathcal{F} \in \textit{Ab}(X)$。
则
(\ref{sites-cohomology-item-base-change-P}),
(\ref{sites-cohomology-item-restriction-A}),
(\ref{sites-cohomology-item-A-sheaf}),
(\ref{sites-cohomology-item-A-and-P}) 和
(\ref{sites-cohomology-item-refine-tau-by-P})
这五个条件（情形 \ref{sites-cohomology-situation-compare}）均成立。
\end{lemma}

\begin{proof}
首先，通过检查同调，引理
\ref{homology-lemma-characterize-weak-serre-subcategory} 的条件
(1)、(2)、(3)、(4)，证明
$\mathcal{A}'_X \subset \textit{Ab}(\textit{LC}_{Zar}/X)$
是弱 Serre 子范畴。(1)、(2)、(3) 立即成立，因为根据引理
\ref{sites-cohomology-lemma-collect-true-things-Zar} 的 (3) 部分，$\pi_X^{-1}$
是正合且全忠实的。若
$0 \to \pi_X^{-1}\mathcal{F} \to \mathcal{G} \to \pi_X^{-1}\mathcal{F}' \to 0$
是 $\textit{Ab}(\textit{LC}_{Zar}/X)$ 中的短正合序列，则根据引理
\ref{sites-cohomology-lemma-collect-true-things-Zar} 的 (2) 部分，
$0 \to \mathcal{F} \to \pi_{X, *}\mathcal{G} \to \mathcal{F}' \to 0$
正合。因此 $\mathcal{G} = \pi_X^{-1}\pi_{X, *}\mathcal{G}$
属于 $\mathcal{A}'_X$，最后一个条件也得到验证。

\medskip\noindent
性质 (\ref{sites-cohomology-item-base-change-P}) 由引理 \ref{sites-cohomology-lemma-LC-basic} 以及
proper 映射的基变换仍是 proper 映射这一事实成立（参见拓扑，
定理 \ref{topology-theorem-characterize-proper} 和引理
\ref{topology-lemma-base-change-separated}）。

\medskip\noindent
性质 (\ref{sites-cohomology-item-restriction-A}) 来自引理
\ref{sites-cohomology-lemma-collect-true-things-Zar} 中的交换图 (5)。

\medskip\noindent
性质 (\ref{sites-cohomology-item-A-sheaf}) 就是引理 \ref{sites-cohomology-lemma-describe-pullback-pi}。

\medskip\noindent
性质 (\ref{sites-cohomology-item-A-and-P}) 是引理
\ref{sites-cohomology-lemma-collect-true-things-Zar} 的 (7)(b) 部分。

\medskip\noindent
证明 (\ref{sites-cohomology-item-refine-tau-by-P})。设给定一个 qc 覆盖
$\{U_i \to U\}$。对 $u \in U$，选取 $i_1, \ldots, i_m \in I$
以及拟紧子集 $E_j \subset U_{i_j}$，使得
$\bigcup f_{i_j}(E_j)$ 是 $u$ 的邻域。
注意，$Y = \coprod_{j = 1, \ldots, m} E_j \to U$
作为豪斯多夫拟紧空间之间的连续映射是 proper 的
（拓扑，引理 \ref{topology-lemma-closed-map}）。
选取包含于 $\bigcup f_{i_j}(E_j)$ 的开邻域 $u \in V$。
则 $Y \times_U V \to V$ 是满射 proper 态射，因而根据引理
\ref{sites-cohomology-lemma-proper-surjective-is-qc-covering} 是 $qc$ 覆盖。
对每个 $u \in U$ 都能如此处理，因此
(\ref{sites-cohomology-item-refine-tau-by-P}) 成立。
\end{proof}

\begin{lemma}
\label{sites-cohomology-lemma-V-C-all-n}
记号如上。
\begin{enumerate}
\item 对于 $X \in \Ob(\textit{LC}_{qc})$ 以及 $X$ 上的阿贝尔层
$\mathcal{F}$，有 $\epsilon_{X, *}a_X^{-1}\mathcal{F} = \pi_X^{-1}\mathcal{F}$，
并且当 $i > 0$ 时 $R^i\epsilon_{X, *}(a_X^{-1}\mathcal{F}) = 0$。
\item 对于 $\textit{LC}_{qc}$ 中的 proper 态射 $f : X \to Y$
以及 $X$ 上的阿贝尔层 $\mathcal{F}$，对所有 $i$ 有
$a_Y^{-1}(R^if_*\mathcal{F}) = R^if_{qc, *}(a_X^{-1}\mathcal{F})$。
\item 对于 $X \in \Ob(\textit{LC}_{qc})$ 以及 $D^+(X)$ 中的 $K$，
映射 $\pi_X^{-1}K \to R\epsilon_{X, *}(a_X^{-1}K)$ 是同构。
\item 对于 $\textit{LC}_{qc}$ 中的 proper 态射 $f : X \to Y$
以及 $D^+(X)$ 中的 $K$，有
$a_Y^{-1}(Rf_*K) = Rf_{qc, *}(a_X^{-1}K)$。
\end{enumerate}
\end{lemma}

\begin{proof}
根据引理 \ref{sites-cohomology-lemma-compare-qc-zar}，第
\ref{sites-cohomology-section-compare-general} 节中的引理都适用于当前情形。
为了将这些结果翻译过来，注意引理 \ref{sites-cohomology-lemma-A} 中的范畴
$\mathcal{A}_X$ 是下述函子的本质像：
$a_X^{-1} : \textit{Ab}(X) \to \textit{Ab}(\textit{LC}_{qc}/X)$。

\medskip\noindent
第 (1) 部分等价于对所有 $n$ 的 $(V_n)$，而这由引理
\ref{sites-cohomology-lemma-V-C-all-n-general} 成立。

\medskip\noindent
第 (2) 部分由将 $\epsilon_Y^{-1}$ 应用于引理
\ref{sites-cohomology-lemma-V-implies-C-general} 的结论得到。

\medskip\noindent
第 (3) 部分由引理 \ref{sites-cohomology-lemma-V-C-all-n-general} 的 (1) 部分得到，
因为 $\pi_X^{-1}K$ 属于
$D^+_{\mathcal{A}'_X}(\textit{LC}_{Zar}/X)$，且
$a_X^{-1} = \epsilon_X^{-1} \circ a_X^{-1}$。

\medskip\noindent
第 (4) 部分出于同样理由由引理
\ref{sites-cohomology-lemma-V-C-all-n-general} 的 (2) 部分得到。
\end{proof}

\begin{lemma}
\label{sites-cohomology-lemma-cohomological-descent-LC}
令 $X$ 是 $\textit{LC}_{qc}$ 中的对象。对于 $K \in D^+(X)$，映射
$$
K \longrightarrow Ra_{X, *}a_X^{-1}K
$$
是同构，其中 $a_X : \Sh(\textit{LC}_{qc}/X) \to \Sh(X)$ 如上所述。
\end{lemma}

\begin{proof}
首先将断言约化到 $K$ 由单个阿贝尔层给出的情形。具体地，
用有下界的复形 $\mathcal{F}^\bullet$ 表示 $K$。由单个层的情形可知
$\mathcal{F}^n = a_{X, *} a_X^{-1} \mathcal{F}^n$
并且当 $q > 0$ 时，层 $R^qa_{X, *}a_X^{-1}\mathcal{F}^n$
为零。将 Leray 无环引理（导出范畴，引理
\ref{derived-lemma-leray-acyclicity}）应用于 $a_X^{-1}\mathcal{F}^\bullet$
和函子 $a_{X, *}$，即可得到结论。以下设 $K = \mathcal{F}$。

\medskip\noindent
根据引理 \ref{sites-cohomology-lemma-describe-pullback-pi}，有
$a_{X, *}a_X^{-1}\mathcal{F} = \mathcal{F}$。因此只需证明当
$q > 0$ 时 $R^qa_{X, *}a_X^{-1}\mathcal{F} = 0$。
为此可使用 $a_X = \epsilon_X \circ \pi_X$ 以及 Leray 谱序列，
引理 \ref{sites-cohomology-lemma-relative-Leray}。根据引理
\ref{sites-cohomology-lemma-V-C-all-n}，当 $i > 0$ 时
$R^i\epsilon_{X, *}(a_X^{-1}\mathcal{F}) = 0$，并且
$\epsilon_{X, *}a_X^{-1}\mathcal{F} = \pi_X^{-1}\mathcal{F}$。
根据引理 \ref{sites-cohomology-lemma-collect-true-things-Zar}，当 $j > 0$ 时
$R^j\pi_{X, *}(\pi_X^{-1}\mathcal{F}) = 0$。
证明完成。
\end{proof}

\begin{lemma}
\label{sites-cohomology-lemma-compare-cohomology-LC}
令 $X \in \Ob(\textit{LC}_{qc})$，并令
$a_X : \Sh(\textit{LC}_{qc}/X) \to \Sh(X)$ 如上所述：
\begin{enumerate}
\item 对于 $X$ 上的阿贝尔层 $\mathcal{F}$，有
$H^n(X, \mathcal{F}) = H^n_{qc}(X, a_X^{-1}\mathcal{F})$,
\item 对于 $K \in D^+(X)$，有 $H^n(X, K) = H^n_{qc}(X, a_X^{-1}K)$。
\end{enumerate}
例如，若 $A$ 是阿贝尔群，则
$H^n(X, \underline{A}) = H^n_{qc}(X, \underline{A})$.
\end{lemma}

\begin{proof}
由引理 \ref{sites-cohomology-lemma-cohomological-descent-LC} 和评注
\ref{sites-cohomology-remark-before-Leray} 得到。
\end{proof}











\section{Ext 的谱序列}
\label{sites-cohomology-section-spectral-sequence-ext}

\noindent
本节汇集考虑 Ext 函子时出现的各种谱序列。对于环化站点
$(\mathcal{C}, \mathcal{O})$ 上的任意一对模复形
$\mathcal{G}^\bullet, \mathcal{F}^\bullet$，我们记
$$
\Ext^n_\mathcal{O}(\mathcal{G}^\bullet, \mathcal{F}^\bullet)
=
\Hom_{D(\mathcal{O})}(\mathcal{G}^\bullet, \mathcal{F}^\bullet[n])
$$
按照导出范畴第 \ref{derived-section-ext} 节中的一般约定。

\begin{example}
\label{sites-cohomology-example-hom-complex-into-sheaf}
令 $(\mathcal{C}, \mathcal{O})$ 是环化站点。
令 $\mathcal{K}^\bullet$ 是有上界的 $\mathcal{O}$-模复形，
令 $\mathcal{F}$ 是 $\mathcal{O}$-模。则存在以 $E_2$ 页为
$$
E_2^{i, j} =
\Ext_\mathcal{O}^i(H^{-j}(\mathcal{K}^\bullet), \mathcal{F})
\Rightarrow
\Ext_\mathcal{O}^{i + j}(\mathcal{K}^\bullet, \mathcal{F})
$$
以及另一个以 $E_1$ 页为
$$
E_1^{i, j} =
\Ext_\mathcal{O}^j(\mathcal{K}^{-i}, \mathcal{F})
\Rightarrow
\Ext_\mathcal{O}^{i + j}(\mathcal{K}^\bullet, \mathcal{F}).
$$
为构造这些谱序列，选取内射分解
$\mathcal{F} \to \mathcal{I}^\bullet$，并考虑来自双复形
$\Hom_\mathcal{O}(\mathcal{K}^\bullet, \mathcal{I}^\bullet)$
的两个谱序列，参见同调第 \ref{homology-section-double-complex} 节。
\end{example}








\section{杯积}
\label{sites-cohomology-section-cup-product}

\noindent
令 $(\mathcal{C}, \mathcal{O})$ 是环化站点。令
$K, M$ 是 $D(\mathcal{O})$ 中的对象。置
$A = \Gamma(\mathcal{C}, \mathcal{O})$。此情形中的（全局）杯积是
映射
$$
R\Gamma(\mathcal{C}, K) \otimes_A^\mathbf{L}
R\Gamma(\mathcal{C}, M)
\longrightarrow
R\Gamma(\mathcal{C}, K \otimes_\mathcal{O}^\mathbf{L} M)
$$
在 $D(A)$ 中。我们将其定义为环化拓扑斯态射
$(\Sh(\mathcal{C}), \mathcal{O}) \to (\Sh(pt), A)$
的相对杯积，见评注 \ref{sites-cohomology-remark-cup-product}。

\medskip\noindent
下面表述并证明相对杯积的一个自然相容性。设有环化拓扑斯态射
$f : (\Sh(\mathcal{C}), \mathcal{O}_\mathcal{C}) \to
(\Sh(\mathcal{D}), \mathcal{O}_\mathcal{D})$。
令 $\mathcal{K}^\bullet$ 和 $\mathcal{M}^\bullet$
是 $\mathcal{O}_\mathcal{C}$-模复形。存在朴素杯积
$$
\text{Tot}(
f_*\mathcal{K}^\bullet
\otimes_{\mathcal{O}_\mathcal{D}}
f_*\mathcal{M}^\bullet)
\longrightarrow
f_*\text{Tot}(\mathcal{K}^\bullet
\otimes_{\mathcal{O}_\mathcal{C}}
\mathcal{M}^\bullet)
$$
我们断言这与相对杯积有关。

\begin{lemma}
\label{sites-cohomology-lemma-cup-compatible-with-naive}
在上述情形，下图交换：
$$
\xymatrix{
f_*\mathcal{K}^\bullet
\otimes_{\mathcal{O}_\mathcal{D}}^\mathbf{L}
f_*\mathcal{M}^\bullet \ar[r] \ar[d]
&
Rf_*\mathcal{K}^\bullet
\otimes_{\mathcal{O}_\mathcal{D}}^\mathbf{L}
Rf_*\mathcal{M}^\bullet \ar[d]^{\text{评注 \ref{sites-cohomology-remark-cup-product}}} \\
\text{Tot}(
f_*\mathcal{K}^\bullet
\otimes_{\mathcal{O}_\mathcal{D}}
f_*\mathcal{M}^\bullet) \ar[d]_{\text{朴素杯积}} &
Rf_*(\mathcal{K}^\bullet
\otimes_{\mathcal{O}_\mathcal{C}}^\mathbf{L}
\mathcal{M}^\bullet) \ar[d] \\
f_*\text{Tot}(\mathcal{K}^\bullet
\otimes_{\mathcal{O}_\mathcal{C}}
\mathcal{M}^\bullet) \ar[r] &
Rf_*\text{Tot}(\mathcal{K}^\bullet
\otimes_{\mathcal{O}_\mathcal{C}}
\mathcal{M}^\bullet)
}
$$
\end{lemma}

\begin{proof}
根据评注 \ref{sites-cohomology-remark-cup-product} 中的构造可知，沿图顺时针走一周，
映射
$$
f_*\mathcal{K}^\bullet
\otimes_{\mathcal{O}_\mathcal{D}}^\mathbf{L}
f_*\mathcal{M}^\bullet 
\longrightarrow
Rf_*\text{Tot}(\mathcal{K}^\bullet
\otimes_{\mathcal{O}_\mathcal{C}}
\mathcal{M}^\bullet)
$$
是下述映射的伴随：
\begin{align*}
Lf^*(f_*\mathcal{K}^\bullet
\otimes_{\mathcal{O}_\mathcal{D}}^\mathbf{L}
f_*\mathcal{M}^\bullet)
& =
Lf^*f_*\mathcal{K}^\bullet
\otimes_{\mathcal{O}_\mathcal{D}}^\mathbf{L}
Lf^*f_*\mathcal{M}^\bullet \\
& \to
Lf^*Rf_*\mathcal{K}^\bullet
\otimes_{\mathcal{O}_\mathcal{D}}^\mathbf{L}
Lf^*Rf_*\mathcal{M}^\bullet \\
& \to
\mathcal{K}^\bullet
\otimes_{\mathcal{O}_\mathcal{D}}^\mathbf{L}
\mathcal{M}^\bullet \\
& \to
\text{Tot}(\mathcal{K}^\bullet
\otimes_{\mathcal{O}_\mathcal{C}}
\mathcal{M}^\bullet)
\end{align*}
根据引理 \ref{sites-cohomology-lemma-adjoints-push-pull-compatibility}，这也等于
\begin{align*}
Lf^*(f_*\mathcal{K}^\bullet
\otimes_{\mathcal{O}_\mathcal{D}}^\mathbf{L}
f_*\mathcal{M}^\bullet)
& =
Lf^*f_*\mathcal{K}^\bullet
\otimes_{\mathcal{O}_\mathcal{D}}^\mathbf{L}
Lf^*f_*\mathcal{M}^\bullet \\
& \to
f^*f_*\mathcal{K}^\bullet
\otimes_{\mathcal{O}_\mathcal{D}}^\mathbf{L}
f^*f_*\mathcal{M}^\bullet \\
& \to
\mathcal{K}^\bullet
\otimes_{\mathcal{O}_\mathcal{D}}^\mathbf{L}
\mathcal{M}^\bullet \\
& \to
\text{Tot}(\mathcal{K}^\bullet
\otimes_{\mathcal{O}_\mathcal{C}}
\mathcal{M}^\bullet)
\end{align*}
沿图逆时针走一周，我们得到下述映射的伴随：
\begin{align*}
Lf^*(f_*\mathcal{K}^\bullet
\otimes_{\mathcal{O}_\mathcal{D}}^\mathbf{L}
f_*\mathcal{M}^\bullet)
& \to
Lf^*\text{Tot}(
f_*\mathcal{K}^\bullet
\otimes_{\mathcal{O}_\mathcal{D}}
f_*\mathcal{M}^\bullet) \\
& \to
Lf^*f_*\text{Tot}(\mathcal{K}^\bullet
\otimes_{\mathcal{O}_\mathcal{C}}
\mathcal{M}^\bullet) \\
& \to
Lf^*Rf_*\text{Tot}(\mathcal{K}^\bullet
\otimes_{\mathcal{O}_\mathcal{C}}
\mathcal{M}^\bullet) \\
& \to
\text{Tot}(\mathcal{K}^\bullet
\otimes_{\mathcal{O}_\mathcal{C}}
\mathcal{M}^\bullet)
\end{align*}
根据引理 \ref{sites-cohomology-lemma-adjoints-push-pull-compatibility}，这也等于
\begin{align*}
Lf^*(f_*\mathcal{K}^\bullet
\otimes_{\mathcal{O}_\mathcal{D}}^\mathbf{L}
f_*\mathcal{M}^\bullet)
& \to
Lf^*\text{Tot}(
f_*\mathcal{K}^\bullet
\otimes_{\mathcal{O}_\mathcal{D}}
f_*\mathcal{M}^\bullet) \\
& \to
Lf^*f_*\text{Tot}(\mathcal{K}^\bullet
\otimes_{\mathcal{O}_\mathcal{C}}
\mathcal{M}^\bullet) \\
& \to
f^*f_*\text{Tot}(\mathcal{K}^\bullet
\otimes_{\mathcal{O}_\mathcal{C}}
\mathcal{M}^\bullet) \\
& \to
\text{Tot}(\mathcal{K}^\bullet
\otimes_{\mathcal{O}_\mathcal{C}}
\mathcal{M}^\bullet)
\end{align*}
现在考察下图即可完成证明：
$$
\xymatrix{
Lf^*(f_*\mathcal{K}^\bullet
\otimes_{\mathcal{O}_\mathcal{D}}^\mathbf{L}
f_*\mathcal{M}^\bullet) \ar[d] \ar[rr] & &
Lf^*f_*\mathcal{K}^\bullet \otimes_{\mathcal{O}_\mathcal{C}}^\mathbf{L}
Lf^*f_*\mathcal{M}^\bullet \ar[d] \\
Lf^*\text{Tot}(
f_*\mathcal{K}^\bullet
\otimes_{\mathcal{O}_\mathcal{D}}
f_*\mathcal{M}^\bullet) \ar[d]_{\text{朴素}} \ar[r] &
f^*\text{Tot}(
f_*\mathcal{K}^\bullet
\otimes_{\mathcal{O}_\mathcal{D}}
f_*\mathcal{M}^\bullet) \ar[ldd]^{\text{朴素}} \ar[dd] &
f^*f_*\mathcal{K}^\bullet \otimes_{\mathcal{O}_\mathcal{C}}^\mathbf{L}
f^*f_*\mathcal{M}^\bullet \ar[dd] \ar[ldd] \\
Lf^*f_*\text{Tot}(\mathcal{K}^\bullet
\otimes_{\mathcal{O}_\mathcal{C}}
\mathcal{M}^\bullet) \ar[d] \\
f^*f_*\text{Tot}(\mathcal{K}^\bullet \otimes_{\mathcal{O}_\mathcal{C}}
\mathcal{M}^\bullet) \ar[rd] &
\text{Tot}(f^*f_*\mathcal{K}^\bullet \otimes_{\mathcal{O}_\mathcal{C}}
f^*f_*\mathcal{M}^\bullet) \ar[d] &
\mathcal{K}^\bullet \otimes_{\mathcal{O}_\mathcal{C}}^\mathbf{L}
\mathcal{M}^\bullet \ar[ld] \\
& \text{Tot}(\mathcal{K}^\bullet
\otimes_{\mathcal{O}_\mathcal{C}}
\mathcal{M}^\bullet)
}
$$
该图中的所有多边形都交换。顶部的多边形由引理
\ref{sites-cohomology-lemma-tensor-pull-compatibility} 交换。
含有两个朴素杯积的方形交换，是因为
$Lf^* \to f^*$ 对模复形函子性。
含有下述两个映射的方形同理：
$\mathcal{A}^\bullet \otimes^\mathbf{L} \mathcal{B}^\bullet \to
\text{Tot}(\mathcal{A}^\bullet \otimes \mathcal{B}^\bullet)$.
最后，剩余方形的交换性在复形层面成立，并可视为朴素杯积的定义
（利用 $f^*$ 与 $f_*$ 的伴随性）。证明完成，因为沿图外圈走一周
所得的正是上面给出的两个映射。
\end{proof}

\begin{lemma}
\label{sites-cohomology-lemma-cup-product-associative}
令 $f : (\Sh(\mathcal{C}), \mathcal{O}) \to (\Sh(\mathcal{C}'), \mathcal{O}')$
是环化拓扑斯态射。评注 \ref{sites-cohomology-remark-cup-product} 中的相对杯积按下述
意义结合：对 $D(\mathcal{O})$ 中任意 $K,L,M$，下图
$$
\xymatrix{
Rf_*K \otimes_{\mathcal{O}'}^\mathbf{L}
Rf_*L \otimes_{\mathcal{O}'}^\mathbf{L}
Rf_*M \ar[r] \ar[d] &
Rf_*(K \otimes_\mathcal{O}^\mathbf{L} L)
\otimes_{\mathcal{O}'}^\mathbf{L} Rf_*M \ar[d] \\
Rf_*K \otimes_{\mathcal{O}'}^\mathbf{L}
Rf_*(L \otimes_\mathcal{O}^\mathbf{L} M) \ar[r] &
Rf_*(K \otimes_\mathcal{O}^\mathbf{L} 
L \otimes_\mathcal{O}^\mathbf{L} M)
}
$$
在 $D(\mathcal{O}')$ 中交换。
\end{lemma}

\begin{proof}
沿任一侧走一周，我们得到下述显然映射的伴随：
\begin{align*}
Lf^*(Rf_*K \otimes_{\mathcal{O}'}^\mathbf{L}
Rf_*L \otimes_{\mathcal{O}'}^\mathbf{L}
Rf_*M) & =
Lf^*(Rf_*K) \otimes_\mathcal{O}^\mathbf{L}
Lf^*(Rf_*L) \otimes_\mathcal{O}^\mathbf{L}
Lf^*(Rf_*M) \\
& \to
K \otimes_\mathcal{O}^\mathbf{L} 
L \otimes_\mathcal{O}^\mathbf{L} M
\end{align*}
该映射位于 $D(\mathcal{O})$ 中。
\end{proof}

\begin{lemma}
\label{sites-cohomology-lemma-cup-product-commutative}
令 $f : (\Sh(\mathcal{C}), \mathcal{O}) \to (\Sh(\mathcal{C}'), \mathcal{O}')$
是环化拓扑斯态射。评注 \ref{sites-cohomology-remark-cup-product} 中的相对杯积按下述
意义交换：下图
$$
\xymatrix{
Rf_*K \otimes_{\mathcal{O}'}^\mathbf{L} Rf_*L \ar[r] \ar[d]_\psi &
Rf_*(K \otimes_\mathcal{O}^\mathbf{L} L) \ar[d]^{Rf_*\psi} \\
Rf_*L \otimes_{\mathcal{O}'}^\mathbf{L} Rf_*K \ar[r] &
Rf_*(L \otimes_\mathcal{O}^\mathbf{L} K)
}
$$
对 $D(\mathcal{O})$ 中所有 $K,L$，在 $D(\mathcal{O}')$ 中交换。
这里 $\psi$ 是导出范畴中的交换约束
（引理 \ref{sites-cohomology-lemma-symmetric-monoidal-derived}）。
\end{lemma}

\begin{proof}
略去证明。
\end{proof}

\begin{lemma}
\label{sites-cohomology-lemma-compose-cup-product}
令 $f : (\Sh(\mathcal{C}), \mathcal{O}) \to (\Sh(\mathcal{C}'), \mathcal{O}')$
以及 $f' : (\Sh(\mathcal{C}'), \mathcal{O}') \to
(\Sh(\mathcal{C}''), \mathcal{O}'')$ 是环化拓扑斯态射。
评注 \ref{sites-cohomology-remark-cup-product} 中的相对杯积与复合相容，
即下图
$$
\xymatrix{
R(f' \circ f)_*K \otimes_{\mathcal{O}''}^\mathbf{L} R(f' \circ f)_*L
\ar@{=}[rr] \ar[d] & &
Rf'_*Rf_*K \otimes_{\mathcal{O}''}^\mathbf{L} Rf'_*Rf_*L \ar[d] \\
R(f' \circ f)_*(K \otimes_\mathcal{O}^\mathbf{L} L) \ar@{=}[r] &
Rf'_*Rf_*(K \otimes_\mathcal{O}^\mathbf{L} L) &
Rf'_*(Rf_*K \otimes_{\mathcal{O}'}^\mathbf{L}  Rf_*L) \ar[l]
}
$$
对 $D(\mathcal{O})$ 中所有 $K,L$，在 $D(\mathcal{O}'')$ 中交换。
\end{lemma}

\begin{proof}
这是因为沿图的任一方向走一周，都得到下述映射的伴随：
\begin{align*}
& L(f' \circ f)^*\left(R(f' \circ f)_*K
\otimes_{\mathcal{O}''}^\mathbf{L}
R(f' \circ f)_*L\right) \\
& =
L(f' \circ f)^*R(f' \circ f)_*K
\otimes_\mathcal{O}^\mathbf{L}
L(f' \circ f)^*R(f' \circ f)_*L) \\
& \to
K \otimes_\mathcal{O}^\mathbf{L} L
\end{align*}
该映射位于 $D(\mathcal{O})$ 中。为此使用如下余单位的复合：
$$
L(f' \circ f)^*R(f' \circ f)_* =
Lf^* L(f')^* Rf'_* Rf_*  \to
Lf^* Rf_* \to \text{id}
$$
它是 $L(f' \circ f)^*$ 与 $R(f' \circ f)_*$ 的余单位。
参见范畴，引理 \ref{categories-lemma-compose-counits}。
\end{proof}

\begin{lemma}
\label{sites-cohomology-lemma-base-change-cup-product}
考虑交换方形
$$
\xymatrix{
(\Sh(\mathcal{C}'), \mathcal{O}_{\mathcal{C}'}) \ar[r]_{g'} \ar[d]_{f'} &
(\Sh(\mathcal{C}, \mathcal{O}_{\mathcal{C}}) \ar[d]^f \\
(\Sh(\mathcal{D}'), \mathcal{O}_{\mathcal{D}'}) \ar[r]^g &
(\Sh(\mathcal{D}), \mathcal{O}_{\mathcal{D}}) 
}
$$
它由环化拓扑斯组成。令 $K,L \in D(\mathcal{O}_{\mathcal{C}})$。
相对杯积与该方形相容，意味着下图
$$
\xymatrix{
Lg^*(Rf_*K \otimes_{\mathcal{O}_{\mathcal{D}}}^\mathbf{L} Rf_*L)
\ar[r] \ar@{=}[d] &
Lg^*(Rf_*(K \otimes_{\mathcal{O}_{\mathcal{C}}}^\mathbf{L} L)) \ar[d] \\
Lg^*Rf_*K \otimes_{\mathcal{O}_{\mathcal{D}'}}^\mathbf{L} Lg^*Rf_*L \ar[d] &
R(f')_*L(g')^*(K \otimes_{\mathcal{O}_{\mathcal{C}}}^\mathbf{L} L) \ar@{=}[d] \\
R(f')_*(L(g')^*K \otimes_{\mathcal{O}_{\mathcal{D}'}} R(f')_*(L(g')^*L \ar[r] &
R(f')_*(L(g')^*K \otimes_{\mathcal{O}_{\mathcal{C}'}}^\mathbf{L} L(g')^*L)
}
$$
在 $D(\mathcal{O}_{\mathcal{D}'})$ 中交换。
水平箭头由相对杯积（评注 \ref{sites-cohomology-remark-cup-product}）给出，
竖直箭头由基变换映射（评注 \ref{sites-cohomology-remark-base-change}）
和引理 \ref{sites-cohomology-lemma-pullback-tensor-product} 给出。
\end{lemma}

\begin{proof}
略去证明。
\end{proof}






\section{Hom 复形}
\label{sites-cohomology-section-hom-complexes}

\noindent
令 $(\mathcal{C}, \mathcal{O})$ 是环化站点。令
$\mathcal{L}^\bullet$ 和 $\mathcal{M}^\bullet$ 是两个
$\mathcal{O}$-模复形。我们构造一个 $\mathcal{O}$-模复形
$\SheafHom^\bullet(\mathcal{L}^\bullet, \mathcal{M}^\bullet)$。
具体地，对每个 $n$ 置
$$
\SheafHom^n(\mathcal{L}^\bullet, \mathcal{M}^\bullet) =
\prod\nolimits_{n = p + q}
\SheafHom_\mathcal{O}(\mathcal{L}^{-q}, \mathcal{M}^p)
$$
可以将 $\SheafHom^n$ 理解为如下 $\mathcal{O}$-模层：
所有从 $\mathcal{L}^\bullet$ 到 $\mathcal{M}^\bullet$ 的
$\mathcal{O}$-线性映射（视为分次 $\mathcal{O}$-模）中，
次数齐次为 $n$ 的映射。在此术语下，按下式定义微分：
$$
\text{d}(f) =
\text{d}_\mathcal{M} \circ f - (-1)^n f \circ \text{d}_\mathcal{L}
$$
其中
$f \in \SheafHom^n_\mathcal{O}(\mathcal{L}^\bullet, \mathcal{M}^\bullet)$。
我们省略 $\text{d}^2 = 0$ 的验证。
该构造是微分分次代数，例 \ref{dga-example-category-complexes} 的特例。
由构造立即得到
\begin{equation}
\label{sites-cohomology-equation-cohomology-hom-complex}
H^n(\Gamma(U, \SheafHom^\bullet(\mathcal{L}^\bullet, \mathcal{M}^\bullet))) =
\Hom_{K(\mathcal{O}_U)}(\mathcal{L}^\bullet|_U, \mathcal{M}^\bullet[n]|_U)
\end{equation}
对所有 $n \in \mathbf{Z}$ 以及每个 $U \in \Ob(\mathcal{C})$ 均成立。
类似地，对全局截面复形有
\begin{equation}
\label{sites-cohomology-equation-global-cohomology-hom-complex}
H^n(\Gamma(\mathcal{C},
\SheafHom^\bullet(\mathcal{L}^\bullet, \mathcal{M}^\bullet))) =
\Hom_{K(\mathcal{O})}(\mathcal{L}^\bullet, \mathcal{M}^\bullet[n])
\end{equation}
。

\begin{lemma}
\label{sites-cohomology-lemma-compose}
令 $(\mathcal{C}, \mathcal{O})$ 是环化站点。给定
$\mathcal{O}$-模复形 $\mathcal{K}^\bullet, \mathcal{L}^\bullet, \mathcal{M}^\bullet$，
存在同构
$$
\SheafHom^\bullet(\mathcal{K}^\bullet,
\SheafHom^\bullet(\mathcal{L}^\bullet, \mathcal{M}^\bullet))
=
\SheafHom^\bullet(\text{Tot}(\mathcal{K}^\bullet \otimes_\mathcal{O}
\mathcal{L}^\bullet), \mathcal{M}^\bullet)
$$
它是关于 $\mathcal{K}^\bullet, \mathcal{L}^\bullet, \mathcal{M}^\bullet$
函子的 $\mathcal{O}$-模复形同构。
\end{lemma}

\begin{proof}
略去证明。提示：证明方式与代数进阶，引理
\ref{more-algebra-lemma-compose} 完全相同。
\end{proof}

\begin{lemma}
\label{sites-cohomology-lemma-composition}
令 $(\mathcal{C}, \mathcal{O})$ 是环化站点。给定
$\mathcal{O}$-模复形 $\mathcal{K}^\bullet, \mathcal{L}^\bullet, \mathcal{M}^\bullet$，
存在规范态射
$$
\text{Tot}\left(
\SheafHom^\bullet(\mathcal{L}^\bullet, \mathcal{M}^\bullet)
\otimes_\mathcal{O}
\SheafHom^\bullet(\mathcal{K}^\bullet, \mathcal{L}^\bullet)
\right)
\longrightarrow
\SheafHom^\bullet(\mathcal{K}^\bullet, \mathcal{M}^\bullet)
$$
它是 $\mathcal{O}$-模复形的态射。
\end{lemma}

\begin{proof}
略去证明。提示：证明方式与代数进阶，引理
\ref{more-algebra-lemma-composition} 完全相同。
\end{proof}

\begin{lemma}
\label{sites-cohomology-lemma-diagonal-better}
令 $(\mathcal{C}, \mathcal{O})$ 是环化站点。给定
$\mathcal{O}$-模复形 $\mathcal{K}^\bullet, \mathcal{L}^\bullet, \mathcal{M}^\bullet$，
存在规范态射
$$
\text{Tot}\left(
\mathcal{K}^\bullet \otimes_\mathcal{O}
\SheafHom^\bullet(\mathcal{M}^\bullet, \mathcal{L}^\bullet)
\right)
\longrightarrow
\SheafHom^\bullet(\mathcal{M}^\bullet,
\text{Tot}(\mathcal{K}^\bullet \otimes_\mathcal{O} \mathcal{L}^\bullet))
$$
它是关于这三个复形函子的 $\mathcal{O}$-模复形态射。
\end{lemma}

\begin{proof}
略去证明。提示：证明方式与代数进阶，引理
\ref{more-algebra-lemma-diagonal-better} 完全相同。
\end{proof}

\begin{lemma}
\label{sites-cohomology-lemma-diagonal}
令 $(\mathcal{C}, \mathcal{O})$ 是环化站点。给定
$\mathcal{O}$-模复形 $\mathcal{K}^\bullet, \mathcal{L}^\bullet, \mathcal{M}^\bullet$，
存在规范态射
$$
\mathcal{K}^\bullet
\longrightarrow
\SheafHom^\bullet(\mathcal{L}^\bullet,
\text{Tot}(\mathcal{K}^\bullet \otimes_\mathcal{O} \mathcal{L}^\bullet))
$$
它是关于两个复形函子的 $\mathcal{O}$-模复形态射。
\end{lemma}

\begin{proof}
略去证明。提示：证明方式与代数进阶，引理
\ref{more-algebra-lemma-diagonal} 完全相同。
\end{proof}

\begin{lemma}
\label{sites-cohomology-lemma-evaluate-and-more}
令 $(\mathcal{C}, \mathcal{O})$ 是环化站点。给定
$\mathcal{O}$-模复形 $\mathcal{K}^\bullet, \mathcal{L}^\bullet, \mathcal{M}^\bullet$，
存在规范态射
$$
\text{Tot}(\SheafHom^\bullet(\mathcal{L}^\bullet,
\mathcal{M}^\bullet) \otimes_\mathcal{O} \mathcal{K}^\bullet)
\longrightarrow
\SheafHom^\bullet(\SheafHom^\bullet(\mathcal{K}^\bullet,
\mathcal{L}^\bullet), \mathcal{M}^\bullet)
$$
它是关于这三个复形函子的 $\mathcal{O}$-模复形态射。
\end{lemma}

\begin{proof}
略去证明。提示：证明方式与代数进阶，引理
\ref{more-algebra-lemma-evaluate-and-more} 完全相同。
\end{proof}

\begin{lemma}
\label{sites-cohomology-lemma-RHom-into-K-injective}
令 $(\mathcal{C}, \mathcal{O})$ 是环化站点。令 $L,M$
是 $D(\mathcal{O})$ 中的对象。令 $\mathcal{I}^\bullet$
是表示 $M$ 的 K-内射 $\mathcal{O}$-模复形。令
$\mathcal{L}^\bullet$ 是表示 $L$ 的 $\mathcal{O}$-模复形。
则
$$
H^0(\Gamma(U, \SheafHom^\bullet(\mathcal{L}^\bullet, \mathcal{I}^\bullet))) =
\Hom_{D(\mathcal{O}_U)}(L|_U, M|_U)
$$
对所有 $U \in \Ob(\mathcal{C})$ 成立。类似地，
$H^0(\Gamma(\mathcal{C},
\SheafHom^\bullet(\mathcal{L}^\bullet, \mathcal{I}^\bullet))) =
\Hom_{D(\mathcal{O})}(L, M)$.
\end{lemma}

\begin{proof}
我们有
\begin{align*}
H^0(\Gamma(U, \SheafHom^\bullet(\mathcal{L}^\bullet, \mathcal{I}^\bullet)))
& =
\Hom_{K(\mathcal{O}_U)}(\mathcal{L}^\bullet|_U, \mathcal{I}^\bullet|_U) \\
& =
\Hom_{D(\mathcal{O}_U)}(L|_U, M|_U)
\end{align*}
第一个等式由 (\ref{sites-cohomology-equation-cohomology-hom-complex}) 得到。
第二个等式成立，是因为根据引理
\ref{sites-cohomology-lemma-restrict-K-injective-to-open}，$\mathcal{I}^\bullet|_U$
是 K-内射的。最后一个等式的证明类似，只是使用
(\ref{sites-cohomology-equation-global-cohomology-hom-complex})。
\end{proof}

\begin{lemma}
\label{sites-cohomology-lemma-RHom-well-defined}
令 $(\mathcal{C}, \mathcal{O})$ 是环化站点。令
$(\mathcal{I}')^\bullet \to \mathcal{I}^\bullet$
是 K-内射 $\mathcal{O}$-模复形之间的拟同构。
令 $(\mathcal{L}')^\bullet \to \mathcal{L}^\bullet$
是 $\mathcal{O}$-模复形之间的拟同构。则
$$
\SheafHom^\bullet(\mathcal{L}^\bullet, (\mathcal{I}')^\bullet)
\longrightarrow
\SheafHom^\bullet((\mathcal{L}')^\bullet, \mathcal{I}^\bullet)
$$
是拟同构。
\end{lemma}

\begin{proof}
令 $M$ 是由 $\mathcal{I}^\bullet$ 和 $(\mathcal{I}')^\bullet$
表示的 $D(\mathcal{O})$ 对象。令 $L$ 是由
$\mathcal{L}^\bullet$ 和 $(\mathcal{L}')^\bullet$ 表示的
$D(\mathcal{O})$ 对象。根据引理 \ref{sites-cohomology-lemma-RHom-into-K-injective}，
层
$$
H^0(\SheafHom^\bullet(\mathcal{L}^\bullet, (\mathcal{I}')^\bullet))
\quad\text{和}\quad
H^0(\SheafHom^\bullet((\mathcal{L}')^\bullet, \mathcal{I}^\bullet))
$$
都等于与下述预层相伴的层：
$$
U \longmapsto \Hom_{D(\mathcal{O}_U)}(L|_U, M|_U)
$$
因此该映射是拟同构。
\end{proof}

\begin{lemma}
\label{sites-cohomology-lemma-RHom-from-K-flat-into-K-injective}
令 $(\mathcal{C}, \mathcal{O})$ 是环化站点。令 $\mathcal{I}^\bullet$
是 K-内射的 $\mathcal{O}$-模复形，令 $\mathcal{L}^\bullet$ 是
K-平坦的 $\mathcal{O}$-模复形。则
$\SheafHom^\bullet(\mathcal{L}^\bullet, \mathcal{I}^\bullet)$
是 K-内射的 $\mathcal{O}$-模复形。
\end{lemma}

\begin{proof}
具体地，若 $\mathcal{K}^\bullet$ 是无环的 $\mathcal{O}$-模复形，则
\begin{align*}
\Hom_{K(\mathcal{O})}(\mathcal{K}^\bullet,
\SheafHom^\bullet(\mathcal{L}^\bullet, \mathcal{I}^\bullet))
& =
H^0(\Gamma(\mathcal{C},
\SheafHom^\bullet(\mathcal{K}^\bullet,
\SheafHom^\bullet(\mathcal{L}^\bullet, \mathcal{I}^\bullet)))) \\
& =
H^0(\Gamma(\mathcal{C}, \SheafHom^\bullet(\text{Tot}(
\mathcal{K}^\bullet \otimes_\mathcal{O} \mathcal{L}^\bullet),
\mathcal{I}^\bullet))) \\
& =
\Hom_{K(\mathcal{O})}(
\text{Tot}(\mathcal{K}^\bullet \otimes_\mathcal{O} \mathcal{L}^\bullet),
\mathcal{I}^\bullet) \\
& =
0
\end{align*}
第一个等式由 (\ref{sites-cohomology-equation-global-cohomology-hom-complex}) 得到。
第二个等式由引理 \ref{sites-cohomology-lemma-compose} 得到。
第三个等式由 (\ref{sites-cohomology-equation-global-cohomology-hom-complex}) 得到。
最后一个等式是因为
$\text{Tot}(\mathcal{K}^\bullet \otimes_\mathcal{O} \mathcal{L}^\bullet)$
是无环的：$\mathcal{L}^\bullet$ 是 K-平坦的
（定义 \ref{sites-cohomology-definition-K-flat}），而 $\mathcal{I}^\bullet$ 是 K-内射的。
\end{proof}








\section{导出范畴中的内部 Hom}
\label{sites-cohomology-section-internal-hom}

\noindent
令 $(\mathcal{C}, \mathcal{O})$ 是环化站点。令 $L,M$ 是
$D(\mathcal{O})$ 中的对象。我们希望构造 $D(\mathcal{O})$ 中的对象
$R\SheafHom(L,M)$，使得对 $D(\mathcal{O})$ 中任意第三个对象 $K$，
存在规范双射
\begin{equation}
\label{sites-cohomology-equation-internal-hom}
\Hom_{D(\mathcal{O})}(K, R\SheafHom(L, M))
=
\Hom_{D(\mathcal{O})}(K \otimes_\mathcal{O}^\mathbf{L} L, M)
\end{equation}
注意，根据 Yoneda 引理（范畴，引理
\ref{categories-lemma-yoneda}），该公式将 $R\SheafHom(L,M)$
确定到唯一同构。

\medskip\noindent
为构造这样的对象，选取表示 $M$ 的 K-内射
$\mathcal{O}$-模复形 $\mathcal{I}^\bullet$，以及任意表示 $L$ 的
$\mathcal{O}$-模复形 $\mathcal{L}^\bullet$。置
$$
R\SheafHom(L, M) = \SheafHom^\bullet(\mathcal{L}^\bullet, \mathcal{I}^\bullet)
$$
其中右端是第 \ref{sites-cohomology-section-hom-complexes} 节构造的
$\mathcal{O}$-模复形。根据引理 \ref{sites-cohomology-lemma-RHom-well-defined}，
该定义良好。得到函子
$$
D(\mathcal{O})^{opp} \times D(\mathcal{O}) \longrightarrow D(\mathcal{O}),
\quad
(K, L) \longmapsto R\SheafHom(K, L)
$$
作为证明 (\ref{sites-cohomology-equation-internal-hom}) 的准备，我们计算
$R\SheafHom(K,L)$ 的上同调群。

\begin{lemma}
\label{sites-cohomology-lemma-section-RHom-over-U}
令 $(\mathcal{C}, \mathcal{O})$ 是环化站点。令 $L,M$ 是
$D(\mathcal{O})$ 中的对象。对 $\mathcal{C}$ 的每个对象 $U$，有
$$
H^0(U, R\SheafHom(L, M)) =
\Hom_{D(\mathcal{O}_U)}(L|_U, M|_U)
$$
并且有 $H^0(\mathcal{C}, R\SheafHom(L,M)) =
\Hom_{D(\mathcal{O})}(L, M)$.
\end{lemma}

\begin{proof}
选取表示 $M$ 的 K-内射 $\mathcal{O}$-模复形
$\mathcal{I}^\bullet$，以及表示 $L$ 的 K-平坦复形
$\mathcal{L}^\bullet$。则
$\SheafHom^\bullet(\mathcal{L}^\bullet, \mathcal{I}^\bullet)$
根据引理 \ref{sites-cohomology-lemma-RHom-from-K-flat-into-K-injective}，
它是 K-内射的。因此只需取 $U$ 上的截面即可计算 $U$ 上的上同调，
结论由引理 \ref{sites-cohomology-lemma-RHom-into-K-injective} 得到。
\end{proof}

\begin{lemma}
\label{sites-cohomology-lemma-internal-hom}
令 $(\mathcal{C}, \mathcal{O})$ 是环化站点。令 $K,L,M$ 是
$D(\mathcal{O})$ 中的对象。利用上述构造，存在规范同构
$$
R\SheafHom(K, R\SheafHom(L, M)) =
R\SheafHom(K \otimes_\mathcal{O}^\mathbf{L} L, M)
$$
它位于 $D(\mathcal{O})$ 中，并且关于 $K,L,M$ 函子性；取
$H^0(\mathcal{C},-)$ 后恢复 (\ref{sites-cohomology-equation-internal-hom})。
\end{lemma}

\begin{proof}
选取表示 $M$ 的 K-内射复形 $\mathcal{I}^\bullet$，以及表示 $L$
的 K-平坦 $\mathcal{O}$-模复形 $\mathcal{L}^\bullet$。
对任意 $\mathcal{O}$-模复形 $\mathcal{K}^\bullet$，有
$$
\SheafHom^\bullet(\mathcal{K}^\bullet,
\SheafHom^\bullet(\mathcal{L}^\bullet, \mathcal{I}^\bullet))
=
\SheafHom^\bullet(
\text{Tot}(\mathcal{K}^\bullet \otimes_\mathcal{O} \mathcal{L}^\bullet),
\mathcal{I}^\bullet)
$$
由引理 \ref{sites-cohomology-lemma-compose} 得到。注意左端表示
$R\SheafHom(K,R\SheafHom(L,M))$（使用引理
\ref{sites-cohomology-lemma-RHom-from-K-flat-into-K-injective}），右端表示
$R\SheafHom(K \otimes_\mathcal{O}^\mathbf{L} L,M)$。
这就证明了引理中的显示公式。取全局截面并使用引理
\ref{sites-cohomology-lemma-section-RHom-over-U}，得到
(\ref{sites-cohomology-equation-internal-hom})。
\end{proof}

\begin{lemma}
\label{sites-cohomology-lemma-restriction-RHom-to-U}
令 $(\mathcal{C}, \mathcal{O})$ 是环化站点。令 $K,L$ 是
$D(\mathcal{O})$ 中的对象。$R\SheafHom(K,L)$ 的构造
与限制相容，即对 $\mathcal{C}$ 的每个对象 $U$，有
$R\SheafHom(K|_U, L|_U) = R\SheafHom(K, L)|_U$.
\end{lemma}

\begin{proof}
这由构造和引理 \ref{sites-cohomology-lemma-restrict-K-injective-to-open} 立即得到。
\end{proof}

\begin{lemma}
\label{sites-cohomology-lemma-RHom-triangulated}
令 $(\mathcal{C}, \mathcal{O})$ 是环化站点。双函子
$R\SheafHom(-,-)$ 将两个变量中的有区别三角变为有区别三角。
\end{lemma}

\begin{proof}
这由下述对应的性质得到：
$$
(\mathcal{L}^\bullet, \mathcal{M}^\bullet) \longmapsto
\SheafHom^\bullet(\mathcal{L}^\bullet, \mathcal{M}^\bullet)
$$
它将任一变量中逐项分裂的复形短正合序列变为逐项分裂的短正合序列。
细节略去。
\end{proof}

\begin{lemma}
\label{sites-cohomology-lemma-internal-hom-evaluate}
令 $(\mathcal{C}, \mathcal{O})$ 是环化站点。令 $K,L,M$ 是
$D(\mathcal{O})$ 中的对象。存在规范态射
$$
R\SheafHom(L, M) \otimes_\mathcal{O}^\mathbf{L} K
\longrightarrow
R\SheafHom(R\SheafHom(K, L), M)
$$
它位于 $D(\mathcal{O})$ 中，并且关于 $K,L,M$ 函子性。
\end{lemma}

\begin{proof}
选取表示 $M$ 的 K-内射复形 $\mathcal{I}^\bullet$，
表示 $L$ 的 K-内射复形 $\mathcal{J}^\bullet$，以及表示 $K$ 的
K-平坦复形 $\mathcal{K}^\bullet$。该映射利用引理
$$
\text{Tot}(\SheafHom^\bullet(\mathcal{J}^\bullet,
\mathcal{I}^\bullet) \otimes_\mathcal{O} \mathcal{K}^\bullet)
\longrightarrow
\SheafHom^\bullet(\SheafHom^\bullet(\mathcal{K}^\bullet,
\mathcal{J}^\bullet), \mathcal{I}^\bullet)
$$
\ref{sites-cohomology-lemma-evaluate-and-more} 中的映射定义。由所选复形，
左端表示 $R\SheafHom(L,M) \otimes_\mathcal{O}^\mathbf{L}K$，
右端表示 $R\SheafHom(R\SheafHom(K,L),M)$。
关于 $D(\mathcal{O})$ 中三个对象的函子性证明略去。
\end{proof}

\begin{lemma}
\label{sites-cohomology-lemma-internal-hom-composition}
\begin{slogan}
在 $R\SheafHom$ 上复合。
\end{slogan}
令 $(\mathcal{C}, \mathcal{O})$ 是环化站点。给定
$D(\mathcal{O})$ 中的 $K,L,M$，存在规范态射
$$
R\SheafHom(L, M) \otimes_\mathcal{O}^\mathbf{L} R\SheafHom(K, L)
\longrightarrow R\SheafHom(K, M)
$$
它位于 $D(\mathcal{O})$ 中。
\end{lemma}

\begin{proof}
选取表示 $M$ 的 K-内射复形 $\mathcal{I}^\bullet$、表示 $L$ 的
K-内射复形 $\mathcal{J}^\bullet$，以及任意表示 $K$ 的
$\mathcal{O}$-模复形 $\mathcal{K}^\bullet$。根据引理
\ref{sites-cohomology-lemma-composition}，存在复形映射
$$
\text{Tot}\left(
\SheafHom^\bullet(\mathcal{J}^\bullet, \mathcal{I}^\bullet)
\otimes_\mathcal{O}
\SheafHom^\bullet(\mathcal{K}^\bullet, \mathcal{J}^\bullet)
\right)
\longrightarrow
\SheafHom^\bullet(\mathcal{K}^\bullet, \mathcal{I}^\bullet)
$$
下列 $\mathcal{O}$-模复形
$\SheafHom^\bullet(\mathcal{J}^\bullet, \mathcal{I}^\bullet)$,
$\SheafHom^\bullet(\mathcal{K}^\bullet, \mathcal{J}^\bullet)$，以及
$\SheafHom^\bullet(\mathcal{K}^\bullet, \mathcal{I}^\bullet)$
分别表示 $R\SheafHom(L,M)$、$R\SheafHom(K,L)$ 和
$R\SheafHom(K,M)$。若选取 K-平坦复形 $\mathcal{H}^\bullet$ 以及拟同构
$\mathcal{H}^\bullet \to
\SheafHom^\bullet(\mathcal{K}^\bullet, \mathcal{J}^\bullet)$,
则存在映射
$$
\text{Tot}\left(
\SheafHom^\bullet(\mathcal{J}^\bullet, \mathcal{I}^\bullet)
\otimes_\mathcal{O} \mathcal{H}^\bullet
\right)
\longrightarrow
\text{Tot}\left(
\SheafHom^\bullet(\mathcal{J}^\bullet, \mathcal{I}^\bullet)
\otimes_\mathcal{O}
\SheafHom^\bullet(\mathcal{K}^\bullet, \mathcal{J}^\bullet)
\right)
$$
其源表示
$R\SheafHom(L,M) \otimes_\mathcal{O}^\mathbf{L} R\SheafHom(K,L)$。
复合显示的两个箭头即得到所需映射。构造的函子性证明略去。
\end{proof}

\begin{lemma}
\label{sites-cohomology-lemma-internal-hom-diagonal-better}
令 $(\mathcal{C}, \mathcal{O})$ 是环化站点。给定
$D(\mathcal{O})$ 中的 $K,L,M$，存在规范态射
$$
K \otimes_\mathcal{O}^\mathbf{L} R\SheafHom(M, L)
\longrightarrow
R\SheafHom(M, K \otimes_\mathcal{O}^\mathbf{L} L)
$$
它位于 $D(\mathcal{O})$ 中，并且关于 $K,L,M$ 函子性。
\end{lemma}

\begin{proof}
选取表示 $K$ 的 K-平坦复形 $\mathcal{K}^\bullet$，
表示 $L$ 的 K-内射复形 $\mathcal{I}^\bullet$，以及任意表示 $M$ 的
$\mathcal{O}$-模复形 $\mathcal{M}^\bullet$。选取拟同构
$\text{Tot}(\mathcal{K}^\bullet \otimes_{\mathcal{O}_X} \mathcal{I}^\bullet)
\to \mathcal{J}^\bullet$
其中 $\mathcal{J}^\bullet$ 是 K-内射的。然后使用映射
$$
\text{Tot}\left(
\mathcal{K}^\bullet \otimes_\mathcal{O}
\SheafHom^\bullet(\mathcal{M}^\bullet, \mathcal{I}^\bullet)
\right)
\to
\SheafHom^\bullet(\mathcal{M}^\bullet,
\text{Tot}(\mathcal{K}^\bullet \otimes_\mathcal{O} \mathcal{I}^\bullet))
\to
\SheafHom^\bullet(\mathcal{M}^\bullet, \mathcal{J}^\bullet)
$$
其中第一个映射来自引理 \ref{sites-cohomology-lemma-diagonal-better}。
\end{proof}

\begin{lemma}
\label{sites-cohomology-lemma-internal-hom-diagonal}
令 $(\mathcal{C}, \mathcal{O})$ 是环化站点。
给定 $D(\mathcal{O})$ 中的 $K,L$，存在规范态射
$$
K \longrightarrow R\SheafHom(L, K \otimes_\mathcal{O}^\mathbf{L} L)
$$
它位于 $D(\mathcal{O})$ 中，并且关于 $K,L$ 函子性。
\end{lemma}

\begin{proof}
选取表示 $K$ 的 K-平坦复形 $\mathcal{K}^\bullet$，以及表示 $L$ 的
任意 $\mathcal{O}$-模复形 $\mathcal{L}^\bullet$。选取 K-内射复形
$\mathcal{J}^\bullet$ 和拟同构
$\text{Tot}(\mathcal{K}^\bullet \otimes_\mathcal{O} \mathcal{L}^\bullet)
\to \mathcal{J}^\bullet$。然后使用
$$
\mathcal{K}^\bullet \to
\SheafHom^\bullet(\mathcal{L}^\bullet,
\text{Tot}(\mathcal{K}^\bullet \otimes_\mathcal{O} \mathcal{L}^\bullet))
\to
\SheafHom^\bullet(\mathcal{L}^\bullet, \mathcal{J}^\bullet)
$$
其中第一个映射来自引理 \ref{sites-cohomology-lemma-diagonal}。
\end{proof}

\begin{lemma}
\label{sites-cohomology-lemma-dual}
令 $(\mathcal{C}, \mathcal{O})$ 是环化站点。令 $L$ 是
$D(\mathcal{O})$ 中的对象。置 $L^\vee = R\SheafHom(L,\mathcal{O})$。
对于 $D(\mathcal{O})$ 中的 $M$，存在规范映射
\begin{equation}
\label{sites-cohomology-equation-eval}
M \otimes^\mathbf{L}_\mathcal{O} L^\vee \longrightarrow R\SheafHom(L, M)
\end{equation}
它诱导规范映射
$$
H^0(\mathcal{C}, M \otimes_\mathcal{O}^\mathbf{L} L^\vee)
\longrightarrow
\Hom_{D(\mathcal{O})}(L, M)
$$
且关于 $D(\mathcal{O})$ 中的 $M$ 函子性。
\end{lemma}

\begin{proof}
映射 (\ref{sites-cohomology-equation-eval}) 是引理
\ref{sites-cohomology-lemma-internal-hom-composition} 的特例，其中使用了识别
$M = R\SheafHom(\mathcal{O},M)$。
\end{proof}

\begin{remark}
\label{sites-cohomology-remark-projection-formula-for-internal-hom}
令 $f : (\Sh(\mathcal{C}), \mathcal{O}_\mathcal{C}) \to
(\Sh(\mathcal{D}), \mathcal{O}_\mathcal{D})$ 是环化拓扑斯态射。
令 $K,L$ 是 $D(\mathcal{O}_\mathcal{C})$ 中的对象。断言存在规范映射
$$
Rf_*R\SheafHom(L, K) \longrightarrow R\SheafHom(Rf_*L, Rf_*K)
$$
具体地，根据 (\ref{sites-cohomology-equation-internal-hom})，这等价于映射
$Rf_*R\SheafHom(L, K) \otimes_{\mathcal{O}_\mathcal{D}}^\mathbf{L} Rf_*L
\to Rf_*K$.
为此可使用复合
$$
Rf_*R\SheafHom(L, K) \otimes_{\mathcal{O}_\mathcal{D}}^\mathbf{L} Rf_*L \to
Rf_*(R\SheafHom(L, K) \otimes_{\mathcal{O}_\mathcal{C}}^\mathbf{L} L) \to
Rf_*K
$$
其中第一个箭头是相对杯积（评注 \ref{sites-cohomology-remark-cup-product}），
第二个箭头是将 $Rf_*$ 应用于来自引理
\ref{sites-cohomology-lemma-internal-hom-composition} 的规范映射
$R\SheafHom(L,K) \otimes_{\mathcal{O}_\mathcal{C}}^\mathbf{L} L \to K$
（其中一个位置取 $\mathcal{O}_\mathcal{C}$）。
\end{remark}

\begin{remark}
\label{sites-cohomology-remark-prepare-fancy-base-change}
令 $h : (\Sh(\mathcal{C}), \mathcal{O}) \to (\Sh(\mathcal{C}'), \mathcal{O}')$
是环化拓扑斯态射。令 $K,L$ 是 $D(\mathcal{O}')$ 中的对象。
断言存在规范映射
$$
Lh^*R\SheafHom(K, L) \longrightarrow R\SheafHom(Lh^*K, Lh^*L)
$$
它位于 $D(\mathcal{O})$ 中。具体地，根据
(\ref{sites-cohomology-equation-internal-hom})（该式由引理 \ref{sites-cohomology-lemma-internal-hom} 证明），
这样的映射等价于映射
$$
Lh^*R\SheafHom(K, L) \otimes^\mathbf{L} Lh^*K \longrightarrow Lh^*L
$$
根据引理 \ref{sites-cohomology-lemma-pullback-tensor-product}，该箭头的源是
$Lh^*(\SheafHom(K,L) \otimes^\mathbf{L} K)$，
因此只需构造规范映射
$$
R\SheafHom(K, L) \otimes^\mathbf{L} K \longrightarrow L.
$$
为此取与下式对应的箭头：
$$
\text{id} :
R\SheafHom(K, L)
\longrightarrow
R\SheafHom(K, L)
$$
通过 (\ref{sites-cohomology-equation-internal-hom})。
\end{remark}

\begin{remark}
\label{sites-cohomology-remark-fancy-base-change}
设
$$
\xymatrix{
(\Sh(\mathcal{C}'), \mathcal{O}_{\mathcal{C}'})
\ar[r]_h \ar[d]_{f'} &
(\Sh(\mathcal{C}), \mathcal{O}_\mathcal{C}) \ar[d]^f \\
(\Sh(\mathcal{D}'), \mathcal{O}_{\mathcal{D}'})
\ar[r]^g &
(\Sh(\mathcal{D}), \mathcal{O}_\mathcal{D})
}
$$
是环化拓扑斯的交换图。令 $K,L$ 是
$D(\mathcal{O}_\mathcal{C})$ 中的对象。断言存在规范基变换映射
$$
Lg^*Rf_*R\SheafHom(K, L)
\longrightarrow
R(f')_*R\SheafHom(Lh^*K, Lh^*L)
$$
它位于 $D(\mathcal{O}_{\mathcal{D}'})$ 中。具体地，取下述复合的伴随映射：
\begin{align*}
L(f')^*Lg^*Rf_*R\SheafHom(K, L)
& =
Lh^*Lf^*Rf_*R\SheafHom(K, L) \\
& \to
Lh^*R\SheafHom(K, L) \\
& \to
R\SheafHom(Lh^*K, Lh^*L)
\end{align*}
其中第一个箭头使用伴随映射 $Lf^*Rf_* \to \text{id}$，
第二个箭头是评注 \ref{sites-cohomology-remark-prepare-fancy-base-change} 中构造的规范映射。
\end{remark}




\section{全局导出 Hom}
\label{sites-cohomology-section-global-RHom}

\noindent
令 $(\Sh(\mathcal{C}), \mathcal{O})$ 是环化拓扑斯。
令 $K,L \in D(\mathcal{O})$。利用导出范畴中内部 Hom 的构造，
得到定义良好的对象
$$
R\Hom_\mathcal{O}(K, L) = R\Gamma(\mathcal{C}, R\SheafHom(K, L))
$$
它位于 $D(\Gamma(\mathcal{C}, \mathcal{O}))$ 中。根据引理
\ref{sites-cohomology-lemma-section-RHom-over-U}，有
$$
H^0(R\Hom_\mathcal{O}(K, L)) = \Hom_{D(\mathcal{O})}(K, L)
$$
并且
$$
H^p(R\Hom_\mathcal{O}(K, L)) = \Ext_{D(\mathcal{O})}^p(K, L)
$$
若 $f : (\mathcal{C}', \mathcal{O}') \to (\mathcal{C}, \mathcal{O})$
是环化拓扑斯态射，则存在规范映射
$$
R\Hom_\mathcal{O}(K, L) \longrightarrow R\Hom_{\mathcal{O}'}(Lf^*K, Lf^*L)
$$
它位于 $D(\Gamma(\mathcal{O}))$ 中，是取评注
\ref{sites-cohomology-remark-prepare-fancy-base-change} 所定义映射的全局截面得到的。









\section{导出下 shriek}
\label{sites-cohomology-section-derived-lower-shriek}

\noindent
本节研究环化拓扑斯的态射 $g$：除 $Lg^*$ 和 $Rg_*$ 外，
还存在导出函子 $Lg_!$。

\begin{lemma}
\label{sites-cohomology-lemma-pullback-injective-pre-limp}
令 $u : \mathcal{C} \to \mathcal{D}$ 是站点的连续且余连续函子。
令 $g : \Sh(\mathcal{C}) \to \Sh(\mathcal{D})$ 是相应的拓扑斯态射。
令 $\mathcal{O}_\mathcal{D}$ 是环层，令 $\mathcal{I}$ 是内射的
$\mathcal{O}_\mathcal{D}$-模。则
$H^p(U, g^{-1}\mathcal{I}) = 0$ 对所有 $p > 0$ 以及
$U \in \Ob(\mathcal{C})$ 成立。
\end{lemma}

\begin{proof}
若能证明对 $\mathcal{C}$ 的任意覆盖
$\mathcal{U} = \{U_i \to U\}$，所有高阶 Čech 上同调群
$\check H^p(\mathcal{U},g^{-1}\mathcal{I})$ 都消失，
则由引理 \ref{sites-cohomology-lemma-cech-vanish-collection} 得到本引理的消失结论。
由于 $u$ 连续，$u(\mathcal{U}) = \{u(U_i) \to u(U)\}$
是 $\mathcal{D}$ 的覆盖，并且
$u(U_{i_0} \times_U \ldots \times_U U_{i_n}) =
u(U_{i_0}) \times_{u(U)} \ldots \times_{u(U)} u(U_{i_n})$.
因此有
$$
\check H^p(\mathcal{U}, g^{-1}\mathcal{I}) =
\check H^p(u(\mathcal{U}), \mathcal{I})
$$
因为根据站点，引理 \ref{sites-lemma-when-shriek}，$g^{-1} = u^p$。
由于 $\mathcal{I}$ 是内射的 $\mathcal{O}_\mathcal{D}$-模，这些 Čech
上同调群消失，参见引理 \ref{sites-cohomology-lemma-injective-module-trivial-cech}。
\end{proof}

\begin{lemma}
\label{sites-cohomology-lemma-existence-derived-lower-shriek}
令 $u : \mathcal{C} \to \mathcal{D}$ 是站点的连续且余连续函子。
令 $g : \Sh(\mathcal{C}) \to \Sh(\mathcal{D})$ 是相应的拓扑斯态射。
令 $\mathcal{O}_\mathcal{D}$ 是环层，并置
$\mathcal{O}_\mathcal{C} = g^{-1}\mathcal{O}_\mathcal{D}$。
函子 $g_! : \textit{Mod}(\mathcal{O}_\mathcal{C}) \to
\textit{Mod}(\mathcal{O}_\mathcal{D})$
（参见站点上的模，引理
\ref{sites-modules-lemma-lower-shriek-modules}）具有左导出函子
$$
Lg_! : D(\mathcal{O}_\mathcal{C}) \longrightarrow D(\mathcal{O}_\mathcal{D})
$$
它是 $g^*$ 的左伴随。此外，对 $U \in \Ob(\mathcal{C})$，有
$$
Lg_!(j_{U!}\mathcal{O}_U) =
g_!j_{U!}\mathcal{O}_U =
j_{u(U)!} \mathcal{O}_{u(U)}.
$$
其中 $j_{U!}$ 和 $j_{u(U)!}$ 是与局部化态射相关的零延拓：
$j_U : \mathcal{C}/U \to \mathcal{C}$ 以及
$j_{u(U)} : \mathcal{D}/u(U) \to \mathcal{D}$。
\end{lemma}

\begin{proof}
我们将使用导出范畴，命题
\ref{derived-proposition-left-derived-exists} 构造 $Lg_!$。
为此需验证该命题的假设 (1)、(2)、(3)、(4)、(5)。首先，由于
$g_!$ 是左伴随，它右正合且与所有余极限交换，故 (5) 成立。
(3)、(4) 成立，因为环化站点上的模范畴是 Grothendieck 阿贝尔范畴。
令 $\mathcal{P} \subset \Ob(\textit{Mod}(\mathcal{O}_\mathcal{C}))$
为由形如 $j_{U!}\mathcal{O}_U$ 的模作直和得到的
$\mathcal{O}_\mathcal{C}$-模集合。注意
$g_!j_{U!}\mathcal{O}_U = j_{u(U)!}\mathcal{O}_{u(U)}$，参见站点上的模，
引理 \ref{sites-modules-lemma-lower-shriek-modules} 的证明。
每个 $\mathcal{O}_\mathcal{C}$-模都是 $\mathcal{P}$ 中某对象的商，参见
站点上的模，引理 \ref{sites-modules-lemma-module-quotient-flat}。
因此 (1) 成立。最后需证明 (2)。
令 $\mathcal{K}^\bullet$ 是有上界的无环
$\mathcal{O}_\mathcal{C}$-模复形，且所有 $n$ 都有 $\mathcal{K}^n \in \mathcal{P}$。
需证明 $g_!\mathcal{K}^\bullet$ 正合。为此，只需证明对每个内射的
$\mathcal{O}_\mathcal{D}$-模 $\mathcal{I}$，
$$
\Hom_{D(\mathcal{O}_\mathcal{D})}(
g_!\mathcal{K}^\bullet, \mathcal{I}[n]) = 0
$$
对所有 $n \in \mathbf{Z}$ 均成立。由于 $\mathcal{I}$ 内射，有
\begin{align*}
\Hom_{D(\mathcal{O}_\mathcal{D})}(
g_!\mathcal{K}^\bullet, \mathcal{I}[n])
& =
\Hom_{K(\mathcal{O}_\mathcal{D})}(
g_!\mathcal{K}^\bullet, \mathcal{I}[n]) \\
& =
H^n(\Hom_{\mathcal{O}_\mathcal{D}}(
g_!\mathcal{K}^\bullet, \mathcal{I})) \\
& =
H^n(\Hom_{\mathcal{O}_\mathcal{C}}(
\mathcal{K}^\bullet, g^{-1}\mathcal{I}))
\end{align*}
最后一个等式来自 $g_!$ 与 $g^{-1}$ 的伴随性。

\medskip\noindent
若 $g^{-1}\mathcal{I}$ 是内射的 $\mathcal{O}_\mathcal{C}$-模，
该群的消失显然成立。但 $g_!$ 一般不正合，所以 $g^{-1}\mathcal{I}$
不一定是内射的 $\mathcal{O}_\mathcal{C}$-模。不过我们知道
$$
\Ext^p_{\mathcal{O}_\mathcal{C}}(
j_{U!}\mathcal{O}_U, g^{-1}\mathcal{I}) =
H^p(U, g^{-1}\mathcal{I}) = 0 \text{，其中 }p \geq 1
$$
这里第一个等式来自
$\Hom_{\mathcal{O}_\mathcal{C}}(j_{U!}\mathcal{O}_U, \mathcal{H}) =
\mathcal{H}(U)$ 并取导出函子；而根据引理
\ref{sites-cohomology-lemma-pullback-injective-pre-limp}，对 $p > 0$ 和
$U \in \Ob(\mathcal{C})$，$H^p(U,g^{-1}\mathcal{I})$ 消失。
由于每个 $\mathcal{K}^{-q}$ 都是形如 $j_{U!}\mathcal{O}_U$ 的模的直和，
可知
$$
\Ext^p_{\mathcal{O}_\mathcal{C}}(\mathcal{K}^{-q}, g^{-1}\mathcal{I}) = 0
\text{ 对所有 }p \geq 1\text{ 和 }q
$$
使用谱序列（参见例 \ref{sites-cohomology-example-hom-complex-into-sheaf}）
$$
E_1^{p, q} = \Ext^q_{\mathcal{O}_\mathcal{C}}(
\mathcal{K}^{-p}, g^{-1}\mathcal{I})
\Rightarrow
\Ext^{p + q}_{\mathcal{O}_\mathcal{C}}(
\mathcal{K}^\bullet, g^{-1}\mathcal{I}) = 0.
$$
注意该谱序列收敛到零，因为 $\mathcal{K}^\bullet$ 无环
（因此在导出范畴中为零，从而 Ext 群消失）。由上面证明的高阶 Ext
消失，$E_1$ 页上唯一可能非零的项是
$E_1^{p, 0} = \Hom_{\mathcal{O}_\mathcal{C}}(
\mathcal{K}^{-p}, g^{-1}\mathcal{I})$.
因此复形
$\Hom_{\mathcal{O}_\mathcal{C}}(
\mathcal{K}^\bullet, g^{-1}\mathcal{I})$
是无环的，正如所需。

\medskip\noindent
因此左导出函子 $Lg_!$ 存在。它是
$g^{-1} = g^* = Rg^* = Lg^*$ 的左伴随，即
\begin{equation}
\label{sites-cohomology-equation-to-prove}
\Hom_{D(\mathcal{O}_\mathcal{C})}(K, g^*L) =
\Hom_{D(\mathcal{O}_\mathcal{D})}(Lg_!K, L)
\end{equation}
由导出范畴，引理 \ref{derived-lemma-derived-adjoint-functors} 得到。
证明完成。
\end{proof}

\begin{remark}
\label{sites-cohomology-remark-when-derived-shriek-equal}
警告！令 $u : \mathcal{C} \to \mathcal{D}$、$g$、$\mathcal{O}_\mathcal{D}$
和 $\mathcal{O}_\mathcal{C}$ 如引理
\ref{sites-cohomology-lemma-existence-derived-lower-shriek} 中所述。
一般而言，下图并不交换：
$$
\xymatrix{
D(\mathcal{O}_\mathcal{C}) \ar[r]_{Lg_!} \ar[d]_{forget} &
D(\mathcal{O}_\mathcal{D}) \ar[d]^{forget} \\
D(\mathcal{C}) \ar[r]^{Lg^{Ab}_!} &
D(\mathcal{D})
}
$$
其中函子 $Lg_!^{Ab}$ 是引理
\ref{sites-cohomology-lemma-existence-derived-lower-shriek} 中构造的函子，
但在 $\mathcal{C}$ 和 $\mathcal{D}$ 上都使用常值层 $\mathbf{Z}$ 作为结构层。
一般甚至不成立 $g_! = g_!^{Ab}$（参见站点上的模，评注
\ref{sites-modules-remark-when-shriek-equal}）；但即使
$g_! = g_!^{Ab}$，上述现象也可能发生！具体地，
引理 \ref{sites-cohomology-lemma-existence-derived-lower-shriek} 证明中 $Lg_!$ 的构造表明，
当且仅当下述规范映射
$$
Lg^{Ab}_!j_{U!}\mathcal{O}_U \longrightarrow j_{u(U)!}\mathcal{O}_{u(U)}
$$
对 $\mathcal{C}$ 中所有对象 $U$ 在 $D(\mathcal{D})$ 中是同构时，
$Lg_!$ 才与 $Lg_!^{\textit{Ab}}$ 一致。一般我们只能说存在自然变换
$$
Lg_!^{Ab} \circ forget \longrightarrow forget \circ Lg_!
$$
\end{remark}

\begin{lemma}
\label{sites-cohomology-lemma-pullback-injective-limp}
令 $u : \mathcal{C} \to \mathcal{D}$ 是站点的连续且余连续函子。
令 $g : \Sh(\mathcal{C}) \to \Sh(\mathcal{D})$ 是相应的拓扑斯态射。
令 $\mathcal{O}_\mathcal{D}$ 是环层，令 $\mathcal{I}$ 是内射的
$\mathcal{O}_\mathcal{D}$-模。若
$g_!^{Sh} : \Sh(\mathcal{C}) \to \Sh(\mathcal{D})$
与纤维积交换\footnote{若 $\mathcal{C}$ 有有限连通极限且 $u$ 与之交换则成立，
参见站点，引理 \ref{sites-lemma-preserve-equalizers}。}，则
$g^{-1}\mathcal{I}$ 完全无环。
\end{lemma}

\begin{proof}
我们将使用引理 \ref{sites-cohomology-lemma-characterize-limp} 的判别准则。
条件 (1) 由引理 \ref{sites-cohomology-lemma-pullback-injective-pre-limp} 成立。
令 $K' \to K$ 是 $\mathcal{C}$ 上集合层的满射。
由于 $g_!^{Sh}$ 是左伴随，$g_!^{Sh}K' \to g_!^{Sh}K$ 满射。
注意
\begin{align*}
H^0(K' \times_K \ldots \times_K K', g^{-1}\mathcal{I}) 
& =
H^0(g_!^{Sh}(K' \times_K \ldots \times_K K'), \mathcal{I}) \\
& =
H^0(g_!^{Sh}K' \times_{g_!^{Sh}K} \ldots \times_{g_!^{Sh}K} g_!^{Sh}K',
\mathcal{I})
\end{align*}
由我们对 $g_!^{Sh}$ 的假设成立。由于 $\mathcal{I}$ 是内射模，
根据引理 \ref{sites-cohomology-lemma-direct-image-injective-sheaf}（应用于恒等态射），
它完全无环。因此可使用引理 \ref{sites-cohomology-lemma-characterize-limp} 的逆命题，
得到复形
$$
0 \to H^0(K, g^{-1}\mathcal{I}) \to H^0(K', g^{-1}\mathcal{I}) \to
H^0(K' \times_K K', g^{-1}\mathcal{I}) \to \ldots
$$
正合，证毕。
\end{proof}

\begin{lemma}
\label{sites-cohomology-lemma-pullback-same-cohomology}
令 $u : \mathcal{C} \to \mathcal{D}$ 是站点的连续且余连续函子。
令 $g : \Sh(\mathcal{C}) \to \Sh(\mathcal{D})$ 是相应的拓扑斯态射。
令 $U \in \Ob(\mathcal{C})$。
\begin{enumerate}
\item 对于 $D(\mathcal{D})$ 中的 $M$，有
$R\Gamma(U, g^{-1}M) = R\Gamma(u(U), M)$.
\item 若 $\mathcal{O}_\mathcal{D}$ 是环层且
$\mathcal{O}_\mathcal{C} = g^{-1}\mathcal{O}_\mathcal{D}$，则对于
$D(\mathcal{O}_\mathcal{D})$ 中的 $M$，有
$R\Gamma(U, g^*M) = R\Gamma(u(U), M)$.
\end{enumerate}
\end{lemma}

\begin{proof}
在有下界情形中，用有下界的内射复形表示 $K$，并使用引理
\ref{sites-cohomology-lemma-pullback-injective-pre-limp} 以及 Leray 无环引理，
即可得到 (1)、(2)。在无界情形中，先注意 (1) 是 (2) 的特例。
对于 (2)，可使用
$$
R\Gamma(U, g^*M) =
R\Hom_{\mathcal{O}_\mathcal{C}}(j_{U!}\mathcal{O}_U, g^*M) =
R\Hom_{\mathcal{O}_\mathcal{D}}(j_{u(U)!}\mathcal{O}_{u(U)}, M) =
R\Gamma(u(U), M)
$$
其中中间等式是引理
\ref{sites-cohomology-lemma-existence-derived-lower-shriek} 的结果。
\end{proof}

\begin{lemma}
\label{sites-cohomology-lemma-special-square-cocontinuous}
假设给定交换图
$$
\xymatrix{
(\Sh(\mathcal{C}'), \mathcal{O}_{\mathcal{C}'})
\ar[r]_{(g', (g')^\sharp)} \ar[d]_{(f', (f')^\sharp)} &
(\Sh(\mathcal{C}), \mathcal{O}_\mathcal{C}) \ar[d]^{(f, f^\sharp)} \\
(\Sh(\mathcal{D}'), \mathcal{O}_{\mathcal{D}'}) \ar[r]^{(g, g^\sharp)} &
(\Sh(\mathcal{D}), \mathcal{O}_\mathcal{D})
}
$$
它由环化拓扑斯组成。假设
\begin{enumerate}
\item $f$、$f'$、$g$ 和 $g'$ 分别对应于余连续函子
$u$、$u'$、$v$ 和 $v'$，如《站点》引理
\ref{sites-lemma-cocontinuous-morphism-topoi} 所述；
\item $v \circ u' = u \circ v'$,
\item $v$ 和 $v'$ 既连续又余连续；
\item 对 $\mathcal{D}'$ 的任意对象 $V'$，由 $v$ 给出的函子
${}^{u'}_{V'}\mathcal{I} \to {}^{\ \ \ u}_{v(V')}\mathcal{I}$
是共终的；
\item $g^{-1}\mathcal{O}_{\mathcal{D}} = \mathcal{O}_{\mathcal{D}'}$
且 $(g')^{-1}\mathcal{O}_{\mathcal{C}} = \mathcal{O}_{\mathcal{C}'}$；
\item $g'_! : \textit{Ab}(\mathcal{C}') \to \textit{Ab}(\mathcal{C})$
是正合的\footnote{若 $\mathcal{C}'$ 中存在纤维积和等化子，且 $v'$ 与它们交换，
则此条件成立，见《站点上的模》引理
\ref{sites-modules-lemma-exactness-lower-shriek}。}。
\end{enumerate}
则作为函子
$D(\mathcal{O}_\mathcal{C}) \to D(\mathcal{O}_{\mathcal{D}'})$.
有 $Rf'_* \circ (g')^* = g^* \circ Rf_*$。
\end{lemma}

\begin{proof}
由条件 (5)，有
$g^* = Lg^* = g^{-1}$ 且 $(g')^* = L(g')^* = (g')^{-1}$。
根据引理 \ref{sites-cohomology-lemma-modules-abelian-unbounded}，只需在阿贝尔层的
导出范畴 $D(\mathcal{C})$ 上证明结论。选取对象
$K \in D(\mathcal{C})$。令 $\mathcal{I}^\bullet$ 是 $\mathcal{C}$ 上
表示 $K$ 的阿贝尔层 K-内射复形。由《导出范畴》引理
\ref{derived-lemma-adjoint-preserve-K-injectives} 和假设 (6)，
$(g')^{-1}\mathcal{I}^\bullet$ 是 $\mathcal{C}'$ 上阿贝尔层的 K-内射复形。
由《站点上的模》引理
\ref{sites-modules-lemma-special-square-cocontinuous}，有
$f'_*(g')^{-1}\mathcal{I}^\bullet = g^{-1}f_*\mathcal{I}^\bullet$。
由于 $f_*\mathcal{I}^\bullet$ 表示 $Rf_*K$，而
$f'_*(g')^{-1}\mathcal{I}^\bullet$ 表示 $Rf'_*(g')^{-1}K$，结论成立。
\end{proof}

\begin{lemma}
\label{sites-cohomology-lemma-special-square-continuous}
考虑交换图
$$
\xymatrix{
(\Sh(\mathcal{C}'), \mathcal{O}_{\mathcal{C}'})
\ar[r]_{(g', (g')^\sharp)} \ar[d]_{(f', (f')^\sharp)} &
(\Sh(\mathcal{C}), \mathcal{O}_\mathcal{C}) \ar[d]^{(f, f^\sharp)} \\
(\Sh(\mathcal{D}'), \mathcal{O}_{\mathcal{D}'}) \ar[r]^{(g, g^\sharp)} &
(\Sh(\mathcal{D}), \mathcal{O}_\mathcal{D})
}
$$
它由环化拓扑斯组成，并假设有函子
$$
\xymatrix{
\mathcal{C}' \ar[r]_{v'} &
\mathcal{C} \\
\mathcal{D}' \ar[r]^v \ar[u]^{u'} &
\mathcal{D} \ar[u]_u
}
$$
使得（采用《站点》第 \ref{sites-section-morphism-sites} 节和第
\ref{sites-section-cocontinuous-morphism-topoi} 节的记号）有
\begin{enumerate}
\item $u$ 和 $u'$ 连续，并分别给出态射 $f$ 和 $f'$；
\item $v$ 和 $v'$ 余连续，并分别给出态射 $g$ 和 $g'$；
\item $u \circ v = v' \circ u'$,
\item $v$ 和 $v'$ 既连续又余连续；
\item $g^{-1}\mathcal{O}_{\mathcal{D}} = \mathcal{O}_{\mathcal{D}'}$
且 $(g')^{-1}\mathcal{O}_{\mathcal{C}} = \mathcal{O}_{\mathcal{C}'}$。
\end{enumerate}
则作为函子 $D^+(\mathcal{O}_\mathcal{C}) \to
D^+(\mathcal{O}_{\mathcal{D}'})$，有
$Rf'_* \circ (g')^* = g^* \circ Rf_*$。若此外
\begin{enumerate}
\item[(6)] $g'_! : \textit{Ab}(\mathcal{C}') \to \textit{Ab}(\mathcal{C})$
是正合的\footnote{若 $\mathcal{C}'$ 中存在纤维积和等化子，且 $v'$ 与它们交换，
则此条件成立，见《站点上的模》引理
\ref{sites-modules-lemma-exactness-lower-shriek}。}，
\end{enumerate}
则作为函子 $D(\mathcal{O}_\mathcal{C}) \to
D(\mathcal{O}_{\mathcal{D}'})$，有
$Rf'_* \circ (g')^* = g^* \circ Rf_*$。
\end{lemma}

\begin{proof}
由条件 (5)，有
$g^* = Lg^* = g^{-1}$ 且 $(g')^* = L(g')^* = (g')^{-1}$。
根据引理 \ref{sites-cohomology-lemma-modules-abelian-unbounded}，只需在阿贝尔层的
导出范畴 $D^+(\mathcal{C})$ 或 $D(\mathcal{C})$ 上证明结论。

\medskip\noindent
选取对象 $K \in D^+(\mathcal{C})$。令 $\mathcal{I}^\bullet$ 是
$\mathcal{C}$ 上表示 $K$ 的下有界内射阿贝尔层复形。由引理
\ref{sites-cohomology-lemma-pullback-injective-pre-limp}，对所有 $p > 0$、任意 $q$
以及任意 $U' \in \Ob(\mathcal{C}')$，都有
$H^p(U', (g')^{-1}\mathcal{I}^q) = 0$。回忆
$R^pf'_*(g')^{-1}\mathcal{I}^q$ 是与预层
$V' \mapsto H^p(u'(V'), (g')^{-1}\mathcal{I}^q)$ 关联的层，见引理
\ref{sites-cohomology-lemma-higher-direct-images}。因此 $(g')^{-1}\mathcal{I}^q$
对函子 $f'_*$ 右无环。由 Leray 无环性引理（《导出范畴》引理
\ref{derived-lemma-leray-acyclicity}），
$f'_*(g')^*\mathcal{I}^\bullet$ 表示 $Rf'_*(g')^{-1}K$。
由《站点上的模》引理
\ref{sites-modules-lemma-special-square-continuous}，有
$f'_*(g')^{-1}\mathcal{I}^\bullet = g^{-1}f_*\mathcal{I}^\bullet$。
由于 $g^{-1}f_*\mathcal{I}^\bullet$ 表示 $g^{-1}Rf_*K$，结论成立。

\medskip\noindent
选取对象 $K \in D(\mathcal{C})$。令 $\mathcal{I}^\bullet$ 是
$\mathcal{C}$ 上表示 $K$ 的阿贝尔层 K-内射复形。由《导出范畴》引理
\ref{derived-lemma-adjoint-preserve-K-injectives} 和假设 (6)，
$(g')^{-1}\mathcal{I}^\bullet$ 是 $\mathcal{C}'$ 上阿贝尔层的
K-内射复形。由《站点上的模》引理
\ref{sites-modules-lemma-special-square-continuous}，有
$f'_*(g')^{-1}\mathcal{I}^\bullet = g^{-1}f_*\mathcal{I}^\bullet$。
由于 $f_*\mathcal{I}^\bullet$ 表示 $Rf_*K$，而
$f'_*(g')^{-1}\mathcal{I}^\bullet$ 表示 $Rf'_*(g')^{-1}K$，结论成立。
\end{proof}








\section{纤维化范畴的导出下 shriek}
\label{sites-cohomology-section-derived-lower-shriek-fibred}

\noindent
本节详细处理第 \ref{sites-cohomology-section-derived-lower-shriek} 节所讨论情形的若干特殊情况。
我们将确保模上的下 shriek 与阿贝尔群层上的下 shriek 相等。
建议读者初次阅读时跳过本节。

\begin{situation}
\label{sites-cohomology-situation-fibred-category}
设 $(\mathcal{D}, \mathcal{O}_\mathcal{D})$ 为环化站点，且
$p : \mathcal{C} \to \mathcal{D}$ 为纤维化范畴。赋予 $\mathcal{C}$
从 $\mathcal{D}$ 继承的拓扑（《叠》第
\ref{stacks-section-topology} 节）。以
$\pi : \Sh(\mathcal{C}) \to \Sh(\mathcal{D})$ 表示与 $p$ 关联的
拓扑斯态射（《叠》引理
\ref{stacks-lemma-topology-inherited-functorial}）。置
$\mathcal{O}_\mathcal{C} = \pi^{-1}\mathcal{O}_\mathcal{D}$，
从而得到环化拓扑斯的态射
$$
\pi :
(\Sh(\mathcal{C}), \mathcal{O}_\mathcal{C})
\longrightarrow
(\Sh(\mathcal{D}), \mathcal{O}_\mathcal{D})
$$
\end{situation}

\begin{lemma}
\label{sites-cohomology-lemma-fibred-category-with-object}
假设和记号如情形 \ref{sites-cohomology-situation-fibred-category}。对
$U \in \Ob(\mathcal{C})$，考虑诱导的拓扑斯态射
$$
\pi_U : \Sh(\mathcal{C}/U) \longrightarrow \Sh(\mathcal{D}/p(U))
$$
则存在拓扑斯态射
$$
\sigma : \Sh(\mathcal{D}/p(U)) \to \Sh(\mathcal{C}/U)
$$
使得 $\pi_U \circ \sigma = \text{id}$ 且 $\sigma^{-1} = \pi_{U, *}$。
\end{lemma}

\begin{proof}
注意，$\pi_U$ 是 $\pi$ 在局部化上的限制，见《站点》引理
\ref{sites-lemma-localize-cocontinuous}。对 $\mathcal{D}/p(U)$ 的对象
$V \to p(U)$，以 $V \times_{p(U)} U \to U$ 表示 $\mathcal{C}$ 中
$\mathcal{D}$ 上的强笛卡儿态射；它因 $p$ 是纤维化范畴而存在。函子
$$
v : \mathcal{D}/p(U) \to \mathcal{C}/U,\quad
V/p(U) \mapsto V \times_{p(U)} U/U
$$
依照 $\mathcal{C}$ 上拓扑的定义是连续的。此外，由强笛卡儿态射的定义，
它是 $p$ 的右伴随。因此我们处于《站点》第
\ref{sites-section-cocontinuous-adjoint} 节所讨论的情形，并且层
$\pi_{U, *}\mathcal{F}$ 等于
$V \mapsto \mathcal{F}(V \times_{p(U)} U)$（尤其见《站点》引理
\ref{sites-lemma-have-functor-other-way-morphism}）。

\medskip\noindent
但这里还有更多结论。事实上，函子 $v$ 也余连续（因为 $\mathcal{C}$
的覆盖中的所有态射都是强笛卡儿的）。因此 $v$ 定义了引理所述的态射
$\sigma$。等式 $\sigma^{-1} = \pi_{U, *}$ 由定义立即得到。由于
$\pi_U^{-1}\mathcal{G}$ 由规则
$U'/U \mapsto \mathcal{G}(p(U')/p(U))$ 给出，遂有
$\sigma^{-1} \circ \pi_U^{-1} = \text{id}$，这证明了等式
$\pi_U \circ \sigma = \text{id}$。
\end{proof}

\begin{situation}
\label{sites-cohomology-situation-morphism-fibred-categories}
设 $(\mathcal{D}, \mathcal{O}_\mathcal{D})$ 为环化站点。设
$u : \mathcal{C}' \to \mathcal{C}$ 为 $\mathcal{D}$ 上纤维化范畴的
$1$-态射（《范畴》定义
\ref{categories-definition-fibred-categories-over-C}）。赋予
$\mathcal{C}$ 和 $\mathcal{C}'$ 各自的继承拓扑（《叠》定义
\ref{stacks-definition-topology-inherited}），并令
$\pi : \Sh(\mathcal{C}) \to \Sh(\mathcal{D})$,
$\pi' : \Sh(\mathcal{C}') \to \Sh(\mathcal{D})$，以及
$g : \Sh(\mathcal{C}') \to \Sh(\mathcal{C})$
为相应的拓扑斯态射（《叠》引理
\ref{stacks-lemma-topology-inherited-functorial}）。置
$\mathcal{O}_\mathcal{C} = \pi^{-1}\mathcal{O}_\mathcal{D}$ 且
$\mathcal{O}_{\mathcal{C}'} = (\pi')^{-1}\mathcal{O}_\mathcal{D}$。
注意 $g^{-1}\mathcal{O}_\mathcal{C} = \mathcal{O}_{\mathcal{C}'}$，故
$$
\xymatrix{
(\Sh(\mathcal{C}'), \mathcal{O}_{\mathcal{C}'}) \ar[rd]_{\pi'} \ar[rr]_g & &
(\Sh(\mathcal{C}), \mathcal{O}_\mathcal{C}) \ar[ld]^\pi \\
& (\Sh(\mathcal{D}), \mathcal{O}_\mathcal{D})
}
$$
是环化拓扑斯态射的交换图。
\end{situation}

\begin{lemma}
\label{sites-cohomology-lemma-morphism-fibred-categories-with-object}
假设和记号如情形 \ref{sites-cohomology-situation-morphism-fibred-categories}。
对 $U' \in \Ob(\mathcal{C}')$，置 $U = u(U')$ 且 $V = p'(U')$，
并考虑诱导的环化拓扑斯态射
$$
\xymatrix{
(\Sh(\mathcal{C}'/U'), \mathcal{O}_{U'}) \ar[rd]_{\pi'_{U'}} \ar[rr]_{g'} & &
(\Sh(\mathcal{C}), \mathcal{O}_U) \ar[ld]^{\pi_U} \\
& (\Sh(\mathcal{D}/V), \mathcal{O}_V)
}
$$
则存在拓扑斯态射
$$
\sigma' : \Sh(\mathcal{D}/V) \to \Sh(\mathcal{C}'/U'),
$$
使得置 $\sigma = g' \circ \sigma'$ 后，有
$\pi'_{U'} \circ \sigma' = \text{id}$、$\pi_U \circ \sigma = \text{id}$、
$(\sigma')^{-1} = \pi'_{U', *}$ 且 $\sigma^{-1} = \pi_{U, *}$。
\end{lemma}

\begin{proof}
令 $v' : \mathcal{D}/V \to \mathcal{C}'/U'$ 为引理
\ref{sites-cohomology-lemma-fibred-category-with-object} 的证明中从
$p' : \mathcal{C}' \to \mathcal{D}'$ 和对象 $U'$ 出发所构造的函子。
由于 $u$ 是 $\mathcal{D}$ 上纤维化范畴的 $1$-态射，它把强笛卡儿态射
变为强笛卡儿态射，故函子 $v = u \circ v'$ 正是引理
\ref{sites-cohomology-lemma-fibred-category-with-object} 的证明中关于
$p : \mathcal{C} \to \mathcal{D}$ 和 $U$ 的函子。因此本引理由该引理得到。
\end{proof}

\begin{lemma}
\label{sites-cohomology-lemma-properties-lower-shriek-fibred-category}
假设和记号如情形 \ref{sites-cohomology-situation-morphism-fibred-categories}。
\begin{enumerate}
\item 在模和阿贝尔层上，$g^* = g^{-1}$ 分别有左伴随
$g_! : \textit{Mod}(\mathcal{O}_{\mathcal{C}'}) \to
\textit{Mod}(\mathcal{O}_\mathcal{C})$ 和
$g_!^{\textit{Ab}} : \textit{Ab}(\mathcal{C}') \to \textit{Ab}(\mathcal{C})$
。
\item 图
$$
\xymatrix{
\textit{Mod}(\mathcal{O}_{\mathcal{C}'}) \ar[d] \ar[r]_{g_!} &
\textit{Mod}(\mathcal{O}_\mathcal{C}) \ar[d] \\
\textit{Ab}(\mathcal{C}') \ar[r]^{g_!^{\textit{Ab}}} &
\textit{Ab}(\mathcal{C})
}
$$
交换。
\item 在模和阿贝尔层的导出范畴上，$g^* = g^{-1}$ 分别有左伴随
$Lg_! : D(\mathcal{O}_{\mathcal{C}'}) \to D(\mathcal{O}_\mathcal{C})$
以及
$Lg_!^{\textit{Ab}} : D(\mathcal{C}') \to D(\mathcal{C})$
。
\item 图
$$
\xymatrix{
D(\mathcal{O}_{\mathcal{C}'}) \ar[d] \ar[r]_{Lg_!} &
D(\mathcal{O}_\mathcal{C}) \ar[d] \\
D(\mathcal{C}') \ar[r]^{Lg_!^{\textit{Ab}}} &
D(\mathcal{C})
}
$$
交换。
\end{enumerate}
\end{lemma}

\begin{proof}
函子 $u$ 连续且余连续，见《叠》引理
\ref{stacks-lemma-topology-inherited-functorial}。因此函子
$g_!$、$g_!^{\textit{Ab}}$、$Lg_!$ 和 $Lg_!^{\textit{Ab}}$ 的存在性见
《站点上的模》第 \ref{sites-modules-section-exactness-lower-shriek} 节、
第 \ref{sites-modules-section-lower-shriek-modules} 节以及第
\ref{sites-cohomology-section-derived-lower-shriek} 节。

\medskip\noindent
为证明 (2)，只需证明规范映射
$$
g_!^{\textit{Ab}}j_{U'!}\mathcal{O}_{U'} \to j_{u(U')!}\mathcal{O}_{u(U')}
$$
对 $\mathcal{C}'$ 的所有对象 $U'$ 均为同构，见《站点上的模》评注
\ref{sites-modules-remark-when-shriek-equal}。类似地，为证明 (4)，
只需证明规范映射
$$
Lg_!^{\textit{Ab}}j_{U'!}\mathcal{O}_{U'} \to j_{u(U')!}\mathcal{O}_{u(U')}
$$
对 $\mathcal{C}'$ 的所有对象 $U'$ 在 $D(\mathcal{C})$ 中均为同构，
见评注 \ref{sites-cohomology-remark-when-derived-shriek-equal}。这也会推出前一个公式，
因此下面证明这一点。

\medskip\noindent
我们将使用如下事实：对局部化态射 $j$，函子 $j_!$ 与
$j_!^{\textit{Ab}}$ 一致（见《站点上的模》评注
\ref{sites-modules-remark-localize-shriek-equal}），并且 $j_!$ 正合
（《站点上的模》引理
\ref{sites-modules-lemma-extension-by-zero-exact}）。采用引理
\ref{sites-cohomology-lemma-morphism-fibred-categories-with-object} 的记号。由于
$Lg_!^{\textit{Ab}} \circ j_{U'!} = j_{U!} \circ L(g')^{\textit{Ab}}_!$
（由《站点》引理 \ref{sites-lemma-localize-cocontinuous} 的交换性
以及伴随函子的唯一性），只需证明
$L(g')^{\textit{Ab}}_!\mathcal{O}_{U'} = \mathcal{O}_U$。利用引理
\ref{sites-cohomology-lemma-morphism-fibred-categories-with-object} 的结论，对
$D(\mathcal{C}/u(U'))$ 的任意对象 $E$，有如下等式链：
\begin{align*}
\Hom_{D(\mathcal{C}/U)}(L(g')_!^{\textit{Ab}}\mathcal{O}_{U'}, E)
& =
\Hom_{D(\mathcal{C}'/U')}(\mathcal{O}_{U'}, (g')^{-1}E) \\
& =
\Hom_{D(\mathcal{C}'/U')}((\pi'_{U'})^{-1}\mathcal{O}_V, (g')^{-1}E) \\
& =
\Hom_{D(\mathcal{D}/V)}(\mathcal{O}_V, R\pi'_{U', *}(g')^{-1}E) \\
& =
\Hom_{D(\mathcal{D}/V)}(\mathcal{O}_V, (\sigma')^{-1}(g')^{-1}E) \\
& =
\Hom_{D(\mathcal{D}/V)}(\mathcal{O}_V, \sigma^{-1}E) \\
& =
\Hom_{D(\mathcal{D}/V)}(\mathcal{O}_V, \pi_{U, *}E) \\
& =
\Hom_{D(\mathcal{C}/U)}(\pi_U^{-1}\mathcal{O}_V, E) \\
& =
\Hom_{D(\mathcal{C}/U)}(\mathcal{O}_U, E)
\end{align*}
由 Yoneda 引理即得结论。
\end{proof}

\begin{remark}
\label{sites-cohomology-remark-fibred-category}
假设和记号如情形 \ref{sites-cohomology-situation-fibred-category}。注意，置
$\mathcal{C}' = \mathcal{D}$，并令 $u$ 等于 $\mathcal{C}$ 的结构函子，
就得到情形 \ref{sites-cohomology-situation-morphism-fibred-categories} 中的情形。因此引理
\ref{sites-cohomology-lemma-properties-lower-shriek-fibred-category} 表明，我们有函子
$\pi_!$、$\pi_!^{\textit{Ab}}$、$L\pi_!$ 和 $L\pi_!^{\textit{Ab}}$，满足
$forget \circ \pi_! = \pi_!^{\textit{Ab}} \circ forget$ 且
$forget \circ L\pi_! = L\pi_!^{\textit{Ab}} \circ forget$.
\end{remark}

\begin{remark}
\label{sites-cohomology-remark-morphism-fibred-categories}
假设和记号如情形 \ref{sites-cohomology-situation-morphism-fibred-categories}。设
$\mathcal{F}$ 为 $\mathcal{C}$ 上的阿贝尔层，$\mathcal{F}'$ 为
$\mathcal{C}'$ 上的阿贝尔层，并设
$t : \mathcal{F}' \to g^{-1}\mathcal{F}$ 为映射。则得到规范映射
$$
L\pi'_!(\mathcal{F}') \longrightarrow L\pi_!(\mathcal{F})
$$
它由 $t$ 的伴随 $g_!\mathcal{F}' \to \mathcal{F}$、映射
$Lg_!(\mathcal{F}') \to g_!\mathcal{F}'$ 以及等式
$L\pi'_! = L\pi_! \circ Lg_!$ 给出。
\end{remark}

\begin{lemma}
\label{sites-cohomology-lemma-compute-pi-shriek}
假设和记号如情形 \ref{sites-cohomology-situation-fibred-category}。对
$\textit{Ab}(\mathcal{C})$ 中的 $\mathcal{F}$，层 $\pi_!\mathcal{F}$
是与预层
$$
V \longmapsto \colim_{\mathcal{C}_V^{opp}} \mathcal{F}|_{\mathcal{C}_V}
$$
关联的层，其限制映射如证明中所示。
\end{lemma}

\begin{proof}
以 $\mathcal{H}$ 表示引理中的规则。对 $\mathcal{D}$ 的态射
$h : V' \to V$，有纤维化范畴的拉回函子
$h^* : \mathcal{C}_V \to \mathcal{C}_{V'}$（《范畴》定义
\ref{categories-definition-pullback-functor-fibred-category}）。此外，对
$U \in \Ob(\mathcal{C}_V)$，存在覆盖 $h$ 的强笛卡儿态射
$h^*U \to U$。沿这些强笛卡儿态射作限制，定义了函子的变换
$$
\mathcal{F}|_{\mathcal{C}_V}
\longrightarrow
\mathcal{F}|_{\mathcal{C}_{V'}} \circ h^*.
$$
从而得到余极限之间的映射
$\mathcal{H}(V) \to \mathcal{H}(V')$，见《范畴》引理
\ref{categories-lemma-functorial-colimit}。

\medskip\noindent
为证明本引理，我们证明
$$
\Mor_{\textit{PSh}(\mathcal{D})}(\mathcal{H}, \mathcal{G}) =
\Mor_{\Sh(\mathcal{C})}(\mathcal{F}, \pi^{-1}\mathcal{G})
$$
对每个 $\mathcal{C}$ 上的层 $\mathcal{G}$ 成立。左端的一个元素是对
$\mathcal{C}$ 中所有 $U$ 的相容映射系
$\mathcal{F}(U) \to \mathcal{G}(p(U))$。根据我们对 $\mathcal{C}$ 上
拓扑的选择，$\pi^{-1}\mathcal{G}(U) = \mathcal{G}(p(U))$，故右端也描述
同一对象，结论成立。
\end{proof}





\section{范畴上的同调}
\label{sites-cohomology-section-homology}

\noindent
在点上的范畴这一情形，我们把左导出下 shriek 函子称为同调函子。

\begin{example}[点上的范畴]
\label{sites-cohomology-example-category-to-point}
设 $\mathcal{C}$ 为范畴。赋予 $\mathcal{C}$ 混沌拓扑（《站点》例
\ref{sites-example-indiscrete}）。于是 $\mathcal{C}$ 上预层与层相同。
令 $*$ 为仅有一个对象和一个态射的范畴，则函子
$p : \mathcal{C} \to *$ 余连续且连续。令
$\pi : \Sh(\mathcal{C}) \to \Sh(*)$ 为相应的拓扑斯态射。设 $B$ 为环。
赋予 $*$ 环层 $B$，并赋予 $\mathcal{C}$ 结构层
$\mathcal{O}_\mathcal{C} = \pi^{-1}B$，将其记为 $\underline{B}$。这样
$$
\pi : (\Sh(\mathcal{C}), \underline{B}) \to (\Sh(*), B)
$$
就是情形 \ref{sites-cohomology-situation-fibred-category} 的一个例子。由评注
\ref{sites-cohomology-remark-fibred-category}，无需区分模上的 $\pi_!$ 与阿贝尔层上的
$\pi_!$。由引理 \ref{sites-cohomology-lemma-compute-pi-shriek}，有
$\pi_!\mathcal{F} = \colim_{\mathcal{C}^{opp}} \mathcal{F}$.
因此 $L_n\pi_!$ 是取余极限函子的第 $n$ 个左导出函子。以下记
$$
H_n(\mathcal{C}, \mathcal{F}) = L_n\pi_!(\mathcal{F})
$$
并称其为 $\mathcal{F}$ 在 $\mathcal{C}$ 上的{\it 第 $n$ 个同调群}。
\end{example}

\begin{example}[计算同调]
\label{sites-cohomology-example-left-derived-colimits}
在例 \ref{sites-cohomology-example-category-to-point} 中，可如下计算函子
$H_n(\mathcal{C}, -)$。设
$\mathcal{F} \in \Ob(\textit{Ab}(\mathcal{C}))$。考虑链复形
$$
K_\bullet(\mathcal{F}) :
\ \ldots \to
\bigoplus\nolimits_{U_2 \to U_1 \to U_0} \mathcal{F}(U_0)
\to
\bigoplus\nolimits_{U_1 \to U_0} \mathcal{F}(U_0)
\to
\bigoplus\nolimits_{U_0} \mathcal{F}(U_0)
$$
其中过渡映射由下式给出：
$$
(U_2 \to U_1 \to U_0, s)
\longmapsto
(U_1 \to U_0, s) - (U_2 \to U_0, s) + (U_2 \to U_1, s|_{U_1})
$$
其他次数类似。由构造，
$$
H_0(\mathcal{C}, \mathcal{F}) =
\colim_{\mathcal{C}^{opp}} \mathcal{F} =
H_0(K_\bullet(\mathcal{F})),
$$
见《范畴》引理
\ref{categories-lemma-colimits-coproducts-coequalizers}。
$K_\bullet(\mathcal{F})$ 的构造关于 $\mathcal{F}$ 是函子性的，并把
$\textit{Ab}(\mathcal{C})$ 的短正合列变为复形的短正合列。因此函子列
$\mathcal{F} \mapsto H_n(K_\bullet(\mathcal{F}))$ 构成 $\delta$-函子，见
《同调》定义 \ref{homology-definition-cohomological-delta-functor} 和引理
\ref{homology-lemma-long-exact-sequence-cochain}。当
$\mathcal{F} = j_{U!}\mathbf{Z}_U$ 时，复形 $K_\bullet(\mathcal{F})$
是单纯集 $X_\bullet$ 上自由 $\mathbf{Z}$-模所关联的复形，其各项为
$$
X_n = \coprod\nolimits_{U_n \to \ldots \to U_1 \to U_0}
\Mor_\mathcal{C}(U_0, U)
$$
该单纯集同伦等价于单点集 $\{*\}$ 上的常值单纯集。事实上，映射
$X_\bullet \to \{*\}$ 是显然的，映射 $\{*\} \to X_n$ 把 $*$ 映到
$(U \to \ldots \to U, \text{id}_U)$，而映射
$$
h_{n, i} : X_n \longrightarrow X_n
$$
（《单纯形》引理 \ref{simplicial-lemma-relations-homotopy}）定义两个映射
$X_\bullet \to X_\bullet$ 之间的同伦，并由规则
$$
h_{n, i} :
(U_n \to \ldots \to U_0, f)
\longmapsto
(U_n \to \ldots \to U_i \to U \to \ldots \to U, \text{id})
$$
给出，其中 $i > 0$，且 $h_{n, 0} = \text{id}$。验证略去。
这表明 $K_\bullet(j_{U!}\mathbf{Z}_U)$ 在负次数的上同调为零
（由《单纯形》评注 \ref{simplicial-remark-homotopy-better} 的函子性
以及《单纯形》引理 \ref{simplicial-lemma-homotopy-s-N} 的结论）。
因此，例如由《同调》引理
\ref{homology-lemma-efface-implies-universal} 和《导出范畴》引理
\ref{derived-lemma-right-derived-delta-functor} 的对偶，
$K_\bullet(\mathcal{F})$ 计算 $H_0(\mathcal{C}, -)$ 的左导出函子
$H_n(\mathcal{C}, -)$。
\end{example}

\begin{example}
\label{sites-cohomology-example-morphism-categories}
设 $u : \mathcal{C}' \to \mathcal{C}$ 为函子。如例
\ref{sites-cohomology-example-category-to-point}，赋予 $\mathcal{C}'$ 和 $\mathcal{C}$
混沌拓扑。令 $*$ 为仅有一个对象和一个态射的范畴，则函子
$u$、$\mathcal{C}' \to *$ 以及 $\mathcal{C} \to *$ 都余连续且连续。令
$g : \Sh(\mathcal{C}') \to \Sh(\mathcal{C})$、
$\pi' : \Sh(\mathcal{C}') \to \Sh(*)$ 和
$\pi : \Sh(\mathcal{C}) \to \Sh(*)$
为相应的拓扑斯态射。设 $B$ 为环。赋予 $*$ 环层 $B$，并赋予
$\mathcal{C}'$、$\mathcal{C}$ 常值层 $\underline{B}$。这样
$$
\xymatrix{
(\Sh(\mathcal{C}'), \underline{B}) \ar[rd]_{\pi'} \ar[rr]_g & &
(\Sh(\mathcal{C}), \underline{B}) \ar[ld]^\pi \\
& (\Sh(*), B)
}
$$
就是情形 \ref{sites-cohomology-situation-morphism-fibred-categories} 的一个例子。因此引理
\ref{sites-cohomology-lemma-properties-lower-shriek-fibred-category} 可应用于 $g$，
故无需区分模上的 $g_!$ 与阿贝尔层上的 $g_!$。特别地，评注
\ref{sites-cohomology-remark-morphism-fibred-categories} 给出规范映射
$$
H_n(\mathcal{C}', \mathcal{F}')
\longrightarrow
H_n(\mathcal{C}, \mathcal{F})
$$
只要给定 $\textit{Ab}(\mathcal{C})$ 中的 $\mathcal{F}$、
$\textit{Ab}(\mathcal{C}')$ 中的 $\mathcal{F}'$ 以及映射
$t : \mathcal{F}' \to g^{-1}\mathcal{F}$。用例
\ref{sites-cohomology-example-left-derived-colimits} 给出的同调计算表示时，
这些映射来自复形映射
$$
K_\bullet(\mathcal{F}') \longrightarrow K_\bullet(\mathcal{F})
$$
它由规则
$$
(U'_n \to \ldots \to U'_0, s') \longmapsto
(u(U'_n) \to \ldots \to u(U'_0), t(s'))
$$
给出，记号含义显然。
\end{example}

\begin{remark}
\label{sites-cohomology-remark-map-evaluation-to-derived}
记号和假设如例 \ref{sites-cohomology-example-category-to-point}。设
$\mathcal{F}^\bullet$ 为 $\mathcal{C}$ 上阿贝尔层的有界复形。
对 $\mathcal{C}$ 的任意对象 $U$，存在规范映射
$$
\mathcal{F}^\bullet(U) \longrightarrow L\pi_!(\mathcal{F}^\bullet)
$$
它位于 $D(\textit{Ab})$ 中。若 $\mathcal{F}^\bullet$ 是
$\underline{B}$-模复形，则该映射位于 $D(B)$ 中。为证明这一点，注意
$L\pi_!(\mathcal{F}^\bullet)$ 可通过取拟同构
$\mathcal{P}^\bullet \to \mathcal{F}^\bullet$ 来计算，其中
$\mathcal{P}^\bullet$ 是投射对象的复形。然而，由于拓扑是混沌的，
$\mathcal{P}^\bullet(U) \to \mathcal{F}^\bullet(U)$ 是拟同构，
故可在 $D(\textit{Ab})$（相应地在 $D(B)$）中求逆。再与从把 $\pi_!$
计算为余极限而来的规范映射
$\mathcal{P}^\bullet(U) \to \pi_!(\mathcal{P}^\bullet)$ 复合，
即得所需箭头。
\end{remark}

\begin{lemma}
\label{sites-cohomology-lemma-initial-final}
记号和假设如例 \ref{sites-cohomology-example-category-to-point}。若 $\mathcal{C}$ 有始对象
或终对象，则在 $D(\textit{Ab})$（相应地在 $D(B)$）上有
$L\pi_! \circ \pi^{-1} = \text{id}$。
\end{lemma}

\begin{proof}
若 $\mathcal{C}$ 有始对象，则 $\pi_!$ 由在该对象上取值来计算，
断言显然。若 $\mathcal{C}$ 有终对象，则 $R\pi_*$ 由在该对象上取值来计算，
故在 $D(\textit{Ab})$（相应地在 $D(B)$）上有
$R\pi_* \circ \pi^{-1} \cong \text{id}$。这表明
$\pi^{-1} : D(\textit{Ab}) \to D(\mathcal{C})$,
相应地，$\pi^{-1} : D(B) \to D(\underline{B})$ 完全忠实，见《范畴》引理
\ref{categories-lemma-adjoint-fully-faithful}。同一引理于是推出
$L\pi_! \circ \pi^{-1} = \text{id}$，正如所需。
\end{proof}

\begin{lemma}
\label{sites-cohomology-lemma-change-of-rings}
记号和假设如例 \ref{sites-cohomology-example-category-to-point}。设 $B \to B'$ 为环同态。
考虑环化拓扑斯的交换图
$$
\xymatrix{
(\Sh(\mathcal{C}), \underline{B}) \ar[d]_\pi &
(\Sh(\mathcal{C}), \underline{B'}) \ar[d]^{\pi'} \ar[l]^h \\
(*, B) & (*, B') \ar[l]_f
}
$$
则 $L\pi_! \circ Lh^* = Lf^* \circ L\pi'_!$。
\end{lemma}

\begin{proof}
两个函子都是显然函子 $D(B') \to D(\underline{B})$ 的右伴随。
\end{proof}

\begin{lemma}
\label{sites-cohomology-lemma-compute-by-cosimplicial-resolution}
记号和假设如例 \ref{sites-cohomology-example-category-to-point}。设 $U_\bullet$ 为
$\mathcal{C}$ 中的余单纯对象，使得对每个
$U \in \Ob(\mathcal{C})$，单纯集 $\Mor_\mathcal{C}(U_\bullet, U)$
同伦等价于单点集上的常值单纯集。则
$$
L\pi_!(\mathcal{F}) = \mathcal{F}(U_\bullet)
$$
在 $D(\textit{Ab})$（相应地在 $D(B)$）中成立，并且关于
$\textit{Ab}(\mathcal{C})$（相应地 $\textit{Mod}(\underline{B})$）中的
$\mathcal{F}$ 是函子性的。
\end{lemma}

\begin{proof}
由引理 \ref{sites-cohomology-lemma-properties-lower-shriek-fibred-category}，模上的
$L\pi_!$ 与阿贝尔层上的 $L\pi_!$ 一致，故只需在 $\mathcal{F}$ 为
阿贝尔层时证明。对 $U \in \Ob(\mathcal{C})$，阿贝尔层
$j_{U!}\mathbf{Z}_U$ 是 $\textit{Ab}(\mathcal{C})$ 的投射对象，因为
$\Hom(j_{U!}\mathbf{Z}_U, \mathcal{F}) = \mathcal{F}(U)$，而在混沌拓扑下
取截面是正合函子。由《站点上的模》引理
\ref{sites-modules-lemma-module-quotient-flat}，每个阿贝尔层都是若干
$j_{U!}\mathbf{Z}_U$ 的直和的商。因此可选取分辨
$$
\ldots \to \mathcal{G}^{-1} \to \mathcal{G}^0 \to \mathcal{F} \to 0
$$
其各项均为上述形式的层的直和，并取
$L\pi_!(\mathcal{F}) = \pi_!(\mathcal{G}^\bullet)$，以此计算
$L\pi_!(\mathcal{F})$。考虑双复形
$A^{\bullet, \bullet} = \mathcal{G}^\bullet(U_\bullet)$。映射
$\mathcal{G}^0 \to \mathcal{F}$ 给出复形映射
$A^{0, \bullet} \to \mathcal{F}(U_\bullet)$。由于 $\pi_!$ 通过取
$\mathcal{C}^{opp}$ 上的余极限来计算（引理
\ref{sites-cohomology-lemma-compute-pi-shriek}），可知两个复合
$\mathcal{G}^m(U_1) \to \mathcal{G}^m(U_0) \to \pi_!\mathcal{G}^m$
相等。因此得到规范的复形映射
$$
\text{Tot}(A^{\bullet, \bullet})
\longrightarrow
\pi_!(\mathcal{G}^\bullet) = L\pi_!(\mathcal{F})
$$
为证明本引理，只需证明复形
$$
\ldots \to \mathcal{G}^m(U_1) \to \mathcal{G}^m(U_0) \to
\pi_!\mathcal{G}^m \to 0
$$
正合，见《同调》引理
\ref{homology-lemma-double-complex-gives-resolution}。由于层
$\mathcal{G}^m$ 是若干 $j_{U!}\mathbf{Z}_U$ 的直和，归结为
$\mathcal{G} = j_{U!}\mathbf{Z}_U$。复形
$j_{U!}\mathbf{Z}_U(U_\bullet)$ 是单纯集
$\Mor_\mathcal{C}(U_\bullet, U)$ 上自由 $\mathbf{Z}$-模所关联的
阿贝尔群复形；我们假设该单纯集同伦等价于单点集。因此
$$
j_{U!}\mathbf{Z}_U(U_\bullet) \to \mathbf{Z}
$$
是阿贝尔群的同伦等价，因而是拟同构（《单纯形》评注
\ref{simplicial-remark-homotopy-better} 和引理
\ref{simplicial-lemma-homotopy-s-N}）。又因引理
\ref{sites-cohomology-lemma-properties-lower-shriek-fibred-category} 的证明已表明
$\pi_!j_{U!}\mathbf{Z}_U = \mathbf{Z}$，证明完成。
\end{proof}

\begin{lemma}
\label{sites-cohomology-lemma-get-it-now}
记号和假设如例 \ref{sites-cohomology-example-morphism-categories}。若存在
$\mathcal{C}'$ 的余单纯对象 $U'_\bullet$，使得引理
\ref{sites-cohomology-lemma-compute-by-cosimplicial-resolution} 同时适用于
$\mathcal{C}'$ 中的 $U'_\bullet$ 和 $\mathcal{C}$ 中的
$u(U'_\bullet)$，则作为函子 $D(\mathcal{C}) \to D(\textit{Ab})$
（相应地 $D(\mathcal{C}, \underline{B}) \to D(B)$），有
$L\pi'_! \circ g^{-1} = L\pi_!$。
\end{lemma}

\begin{proof}
这立即由引理 \ref{sites-cohomology-lemma-compute-by-cosimplicial-resolution} 以及
$g^{-1}$ 由与 $u$ 预复合给出这一事实得到。
\end{proof}

\begin{lemma}
\label{sites-cohomology-lemma-product-categories}
设 $\mathcal{C}_i$（$i = 1, 2$）为范畴，
$u_i : \mathcal{C}_1 \times \mathcal{C}_2 \to \mathcal{C}_i$ 为投影函子。
设 $B$ 为环，并设
$g_i : (\Sh(\mathcal{C}_1 \times \mathcal{C}_2), \underline{B}) \to
(\Sh(\mathcal{C}_i), \underline{B})$ 为相应的环化拓扑斯态射，见例
\ref{sites-cohomology-example-morphism-categories}。对 $K_i \in D(\mathcal{C}_i, B)$，有
$$
L(\pi_1 \times \pi_2)_!(
g_1^{-1}K_1 \otimes_{\underline{B}}^\mathbf{L} g_2^{-1}K_2)
=
L\pi_{1, !}(K_1) \otimes_B^\mathbf{L} L\pi_{2, !}(K_2)
$$
该等式在 $D(B)$ 中成立，记号含义显然。
\end{lemma}

\begin{proof}
由于两端都与余极限交换，只需对
$K_1 = j_{U!}\underline{B}_U$ 和 $K_2 = j_{V!}\underline{B}_V$
证明，其中 $U \in \Ob(\mathcal{C}_1)$ 且
$V \in \Ob(\mathcal{C}_2)$。见引理
\ref{sites-cohomology-lemma-existence-derived-lower-shriek} 中 $L\pi_!$ 的构造。
在此情形，
$$
g_1^{-1}K_1 \otimes_{\underline{B}}^\mathbf{L} g_2^{-1}K_2 =
g_1^{-1}K_1 \otimes_{\underline{B}} g_2^{-1}K_2 =
j_{(U, V)!}\underline{B}_{(U, V)}
$$
验证略去。引理公式的左右两端都取值为 $B$，故结论成立，见引理
\ref{sites-cohomology-lemma-existence-derived-lower-shriek} 中 $L\pi_!$ 的构造。
\end{proof}

\begin{lemma}
\label{sites-cohomology-lemma-eilenberg-zilber}
记号和假设如例 \ref{sites-cohomology-example-category-to-point}。若存在
$\mathcal{C}$ 的余单纯对象 $U_\bullet$，使得引理
\ref{sites-cohomology-lemma-compute-by-cosimplicial-resolution} 适用，则
$$
L\pi_!(K_1 \otimes^\mathbf{L}_{\underline{B}} K_2) =
L\pi_!(K_1) \otimes^\mathbf{L}_B L\pi_!(K_2)
$$
对所有 $K_i \in D(\underline{B})$ 成立。
\end{lemma}

\begin{proof}
考虑范畴和函子的图
$$
\xymatrix{
& & \mathcal{C} \\
\mathcal{C} \ar[r]^-u &
\mathcal{C} \times \mathcal{C} \ar[rd]^{u_2} \ar[ru]_{u_1} \\
& & \mathcal{C}
}
$$
其中 $u$ 为对角函子，$u_i$ 为投影函子。这给出环化拓扑斯态射
$g$、$g_1$、$g_2$。对 $\mathcal{C}$ 的任意对象 $(U_1, U_2)$，有
$$
\Mor_{\mathcal{C} \times \mathcal{C}}(u(U_\bullet), (U_1, U_2)) =
\Mor_\mathcal{C}(U_\bullet, U_1) \times \Mor_\mathcal{C}(U_\bullet, U_2)
$$
由《单纯形》引理 \ref{simplicial-lemma-products-homotopy}，它同伦等价于一点。
因此引理 \ref{sites-cohomology-lemma-get-it-now} 给出：对
$D(\mathcal{C} \times \mathcal{C}, B)$ 中任意 $K$，有
$L\pi_!(g^{-1}K) = L(\pi \times \pi)_!(K)$。取
$K = g_1^{-1}K_1 \otimes_B^\mathbf{L} g_2^{-1}K_2$。由于
$g^{-1} = g^* = Lg^*$ 与导出张量积交换（引理
\ref{sites-cohomology-lemma-pullback-tensor-product}），有
$g^{-1}K = K_1 \otimes^\mathbf{L}_{\underline{B}} K_2$。
最后应用引理 \ref{sites-cohomology-lemma-product-categories}。
\end{proof}

\begin{remark}[单纯模]
\label{sites-cohomology-remark-simplicial-modules}
令 $\mathcal{C} = \Delta$，并令 $B$ 为任意环。这是例
\ref{sites-cohomology-example-category-to-point} 的一个特殊情形，其中引理
\ref{sites-cohomology-lemma-compute-by-cosimplicial-resolution} 的假设成立。事实上，
令 $U_\bullet$ 为由恒等函子给出的 $\Delta$ 的余单纯对象。为验证条件，
必须证明对 $[m] \in \Ob(\Delta)$，单纯集
$\Delta[m] : n \mapsto \Mor_\Delta([n], [m])$ 同伦等价于一点。
这在《单纯形》例 \ref{simplicial-example-simplex-contractible} 中说明。

\medskip\noindent
在此情形，范畴 $\textit{Mod}(\underline{B})$ 就是单纯 $B$-模范畴，
而函子 $L\pi_!$ 把单纯 $B$-模 $M_\bullet$ 映到其关联的 $B$-模复形
$s(M_\bullet)$。因此上述结论可重新解释为关于单纯模的结论。例如，引理
\ref{sites-cohomology-lemma-eilenberg-zilber} 的一个特殊情形如下：若
$M_\bullet$、$M'_\bullet$ 是平坦单纯 $B$-模，则复形
$s(M_\bullet \otimes_B M'_\bullet)$ 拟同构于双复形
$s(M_\bullet) \otimes_B s(M'_\bullet)$ 所关联的全复形。
（提示：利用平坦性把导出张量积化为通常张量积。）
这是 Eilenberg--Zilber 定理的一个特殊情形，见 \cite{Eilenberg-Zilber}。
\end{remark}

\begin{lemma}
\label{sites-cohomology-lemma-O-homology-qis}
设 $\mathcal{C}$ 为范畴（赋予混沌拓扑）。设
$\mathcal{O} \to \mathcal{O}'$ 为 $\mathcal{C}$ 上环层的映射。假设
\begin{enumerate}
\item 存在引理 \ref{sites-cohomology-lemma-compute-by-cosimplicial-resolution} 所述的
$\mathcal{C}$ 中余单纯对象 $U_\bullet$；
\item $L\pi_!\mathcal{O} \to L\pi_!\mathcal{O}'$ 是同构。
\end{enumerate}
对 $D(\mathcal{O})$ 中的 $K$，有
$$
L\pi_!(K) = L\pi_!(K \otimes_\mathcal{O}^\mathbf{L} \mathcal{O}')
$$
该等式在 $D(\textit{Ab})$ 中成立。
\end{lemma}

\begin{proof}
注意：在本证明中，$L\pi_!$ 表示阿贝尔层上 $\pi_!$ 的左导出函子。
由于 $L\pi_!$ 与余极限交换，只需对 $\mathcal{O}$-模的上有界复形
证明（与《导出范畴》命题
\ref{derived-proposition-left-derived-exists} 的论证比较，或直接始终限于
上有界复形）。每个这样的复形都拟同构于一个上有界复形，其各项是
$j_{U!}\mathcal{O}_U$（$U \in \Ob(\mathcal{C})$）的直和，见
《站点上的模》引理 \ref{sites-modules-lemma-module-quotient-flat}。
因此只需对 $j_{U!}\mathcal{O}_U$ 证明本引理。由假设，
$$
S_\bullet = \Mor_\mathcal{C}(U_\bullet, U)
$$
是同伦等价于单点集上的常值单纯集的单纯集。置
$P_n = \mathcal{O}(U_n)$ 且 $P'_n = \mathcal{O}'(U_n)$。注意，
单纯阿贝尔群
$$
X_\bullet : n \longmapsto \bigoplus\nolimits_{s \in S_n} P_n
$$
所关联的复形由引理 \ref{sites-cohomology-lemma-compute-by-cosimplicial-resolution} 计算
$L\pi_!(j_{U!}\mathcal{O}_U)$。由于 $j_{U!}\mathcal{O}_U$ 是平坦
$\mathcal{O}$-模，有
$j_{U!}\mathcal{O}_U \otimes^\mathbf{L}_\mathcal{O} \mathcal{O}' =
j_{U!}\mathcal{O}'_U$，而它的 $L\pi_!$ 由单纯阿贝尔群
$$
X'_\bullet : n \longmapsto \bigoplus\nolimits_{s \in S_n} P'_n
$$
所关联的复形计算。把单纯集 $T_\bullet$ 映到各项为
$\bigoplus_{t \in T_n} P_n$ 的单纯阿贝尔群的规则是函子，故
$X_\bullet \to P_\bullet$ 是单纯阿贝尔群的同伦等价。类似地，把单纯集
$T_\bullet$ 映到各项为 $\bigoplus_{t \in T_n} P'_n$ 的单纯阿贝尔群的
规则也是函子。因此 $X'_\bullet \to P'_\bullet$ 是单纯阿贝尔群的
同伦等价。由假设，$P_\bullet \to P'_\bullet$ 是拟同构（因为由引理
\ref{sites-cohomology-lemma-compute-by-cosimplicial-resolution}，$P_\bullet$ 计算
$L\pi_!\mathcal{O}$，相应地 $P'_\bullet$ 计算 $L\pi_!\mathcal{O}'$）。
所以 $X_\bullet$ 与 $X'_\bullet$ 拟同构，正如所需。
\end{proof}

\begin{remark}
\label{sites-cohomology-remark-O-homology-B-homology-general}
设 $\mathcal{C}$ 和 $B$ 如例 \ref{sites-cohomology-example-category-to-point}。假设存在
引理 \ref{sites-cohomology-lemma-compute-by-cosimplicial-resolution} 所述的余单纯对象。
设 $\mathcal{O} \to \underline{B}$ 为 $\mathcal{C}$ 上环层的映射，
它诱导同构 $L\pi_!\mathcal{O} \to L\pi_!\underline{B}$。在此情形，
得到三角范畴的正合函子
$$
L\pi_! : D(\mathcal{O}) \longrightarrow D(B)
$$
事实上，对 $D(\mathcal{O})$ 的任意对象 $K$，由引理
\ref{sites-cohomology-lemma-O-homology-qis} 有
$L\pi^{\textit{Ab}}_!(K) =
L\pi^{\textit{Ab}}_!(K \otimes_{\mathcal{O}}^\mathbf{L} \underline{B})$
。因此可把所显示的函子定义为
$- \otimes^\mathbf{L}_\mathcal{O} \underline{B}$ 与函子
$L\pi_! : D(\underline{B}) \to D(B)$ 的复合。换言之，通过引理给出的
$L\pi_!(K)$ 与
$L\pi_!(K \otimes_\mathcal{O}^\mathbf{L} \underline{B})$ 的规范函子性识别，
得到 $L\pi_!(K)$ 上的 $B$-模结构。
\end{remark}







\section{计算导出下 shriek}
\label{sites-cohomology-section-calculate}

\noindent
本节应用第 \ref{sites-cohomology-section-homology} 节的结论，计算情形
\ref{sites-cohomology-situation-fibred-category} 中的 $L\pi_!$ 和情形
\ref{sites-cohomology-situation-morphism-fibred-categories} 中的 $Lg_!$。

\begin{lemma}
\label{sites-cohomology-lemma-compute-left-derived-pi-shriek-pre}
假设和记号如情形 \ref{sites-cohomology-situation-fibred-category}。对
$\textit{PAb}(\mathcal{C})$ 中的 $\mathcal{F}$ 和 $n \geq 0$，
考虑 $\mathcal{D}$ 上的阿贝尔层 $L_n(\mathcal{F})$，它是与预层
$$
V \longmapsto H_n(\mathcal{C}_V, \mathcal{F}|_{\mathcal{C}_V})
$$
关联的层，其限制映射如证明中所示。则
$L_n(\mathcal{F}) = L_n(\mathcal{F}^\#)$。
\end{lemma}

\begin{proof}
对 $\mathcal{D}$ 的态射 $h : V' \to V$，有纤维化范畴的拉回函子
$h^* : \mathcal{C}_V \to \mathcal{C}_{V'}$（《范畴》定义
\ref{categories-definition-pullback-functor-fibred-category}）。此外，对
$U \in \Ob(\mathcal{C}_V)$，存在覆盖 $h$ 的强笛卡儿态射
$h^*U \to U$。沿这些强笛卡儿态射作限制，定义了函子的变换
$$
\mathcal{F}|_{\mathcal{C}_V}
\longrightarrow
\mathcal{F}|_{\mathcal{C}_{V'}} \circ h^*.
$$
由例 \ref{sites-cohomology-example-morphism-categories}，得到所需的限制映射
$$
H_n(\mathcal{C}_V, \mathcal{F}|_{\mathcal{C}_V})
\longrightarrow
H_n(\mathcal{C}_{V'}, \mathcal{F}|_{\mathcal{C}_{V'}})
$$
以 $L_{n, p}(\mathcal{F})$ 表示该预层，于是
$L_n(\mathcal{F}) = L_{n, p}(\mathcal{F})^\#$。规范映射
$\gamma : \mathcal{F} \to \mathcal{F}^+$（《站点》定理
\ref{sites-theorem-plus}）定义规范映射
$L_{n, p}(\mathcal{F}) \to L_{n, p}(\mathcal{F}^+)$。
我们必须证明该映射在层化后成为同构。

\medskip\noindent
使用例 \ref{sites-cohomology-example-left-derived-colimits} 给出的同调计算。以
$K_\bullet(\mathcal{F}|_{\mathcal{C}_V})$ 表示 $\mathcal{F}$ 到纤维范畴
$\mathcal{C}_V$ 的限制所关联的复形。由上述说明，得到复形预层
$K_\bullet(\mathcal{F})$：
$$
V \longmapsto K_\bullet(\mathcal{F}|_{\mathcal{C}_V})
$$
它的上同调预层就是预层 $L_{n, p}(\mathcal{F})$。因此只需证明
$$
K_\bullet(\mathcal{F}) \longrightarrow K_\bullet(\mathcal{F}^+)
$$
在层化后成为同构。

\medskip\noindent
单射性。设 $V$ 为 $\mathcal{D}$ 的对象，且
$\xi \in K_n(\mathcal{F})(V)$ 为映到
$K_n(\mathcal{F}^+)(V)$ 中零元的元素。必须证明存在覆盖
$\{V_j \to V\}$，使得 $\xi|_{V_j}$ 在
$K_n(\mathcal{F})(V_j)$ 中为零。写成
$$
\xi = \sum (U_{i, n + 1} \to \ldots \to U_{i, 0}, \sigma_i)
$$
其中 $\sigma_i \in \mathcal{F}(U_{i, 0})$。可以安排为
$\mathcal{C}_V$ 的每个态射列 $U_n \to \ldots \to U_0$ 至多出现一次。
由于复形 $K_\bullet$ 的定义中的和是直和，该元素映到
$K_\bullet(\mathcal{F}^+)(V)$ 中零元的唯一可能是所有 $\sigma_i$ 都映到
$\mathcal{F}^+(U_{i, 0})$ 中的零元。由 $\mathcal{F}^+$ 的构造，存在覆盖
$\{U_{i, 0, j} \to U_{i, 0}\}$，使得
$\sigma_i|_{U_{i, 0, j}}$ 为零。由我们对 $\mathcal{C}$ 上拓扑的构造，
$U_{i, 0, j} \to U_{i, 0}$ 可写成某些态射 $V_{i, j} \to V$ 的拉回
（《范畴》定义 \ref{categories-definition-pullback-functor-fibred-category}），
而且每个 $\{V_{i, j} \to V\}$ 都是覆盖。选取同时加细所有覆盖
$\{V_{i, j} \to V\}$ 的覆盖 $\{V_j \to V\}$。显然
$\xi|_{V_j} = 0$。

\medskip\noindent
满射性。证明略去。提示：仿照单射性的证明论证。
\end{proof}

\begin{lemma}
\label{sites-cohomology-lemma-compute-left-derived-pi-shriek}
假设和记号如情形 \ref{sites-cohomology-situation-fibred-category}。对
$\textit{Ab}(\mathcal{C})$ 中的 $\mathcal{F}$ 和 $n \geq 0$，
层 $L_n\pi_!(\mathcal{F})$ 等于引理
\ref{sites-cohomology-lemma-compute-left-derived-pi-shriek-pre} 中构造的层
$L_n(\mathcal{F})$。
\end{lemma}

\begin{proof}
考虑从 $\textit{PAb}(\mathcal{C}) \to \textit{Ab}(\mathcal{C})$ 的函子列
$\mathcal{F} \mapsto L_n(\mathcal{F})$。由于对每个
$V \in \Ob(\mathcal{D})$，函子列 $H_n(\mathcal{C}_V, - )$ 构成
$\delta$-函子，函子 $\mathcal{F} \mapsto L_n(\mathcal{F})$ 也如此。
我们的目标是证明它们构成万有 $\delta$-函子。为此构造一些使这些函子
消失的阿贝尔预层。

\medskip\noindent
对 $U' \in \Ob(\mathcal{C})$，考虑阿贝尔预层
$\mathcal{F}_{U'} = j_{U'!}^{\textit{PAb}}\mathbf{Z}_{U'}$
（《站点上的模》评注
\ref{sites-modules-remark-localize-presheaves}）。回忆
$$
\mathcal{F}_{U'}(U) =
\bigoplus\nolimits_{\Mor_\mathcal{C}(U, U')} \mathbf{Z}
$$
若 $U$ 位于 $\mathcal{D})$ 中的 $V = p(U)$ 之上，且 $U'$ 位于
$V' = p(U')$ 之上，则任意态射 $a : U \to U'$ 都唯一分解为
$U \to h^*U' \to U'$，其中 $h = p(a) : V \to V'$（见《范畴》定义
\ref{categories-definition-pullback-functor-fibred-category}）。因此
$$
\mathcal{F}_{U'}|_{\mathcal{C}_V}
=
\bigoplus\nolimits_{h \in \Mor_\mathcal{D}(V, V')}
j_{h^*U'!}\mathbf{Z}_{h^*U'}
$$
其中 $j_{h^*U'} : \Sh(\mathcal{C}_V/h^*U') \to \Sh(\mathcal{C}_V)$
是局部化态射。层 $j_{h^*U'!}\mathbf{Z}_{h^*U'}$ 的高阶同调群消失
（见例 \ref{sites-cohomology-example-left-derived-colimits}）。所以对所有 $n > 0$ 和所有
$U'$，有 $L_n(\mathcal{F}_{U'}) = 0$。由此，任意阿贝尔预层
$\mathcal{F}$ 都是某个阿贝尔预层 $\mathcal{G}$ 的商，且对所有
$n > 0$ 有 $L_n(\mathcal{G}) = 0$（《站点上的模》引理
\ref{sites-modules-lemma-module-quotient-flat}）。由于
$L_n(\mathcal{F}) = L_n(\mathcal{F}^\#)$，阿贝尔层也具有同一性质。
因此函子列 $L_n(-)$ 是 $\textit{Ab}(\mathcal{C})$ 上的万有
$\delta$-函子（《同调》引理
\ref{homology-lemma-efface-implies-universal}）。由引理
\ref{sites-cohomology-lemma-compute-pi-shriek}，在 $n = 0$ 时它与
$H^{-n}(L\pi_!(-))$ 一致；再由万有 $\delta$-函子的唯一性（《同调》引理
\ref{homology-lemma-uniqueness-universal-delta-functor}）以及《导出范畴》引理
\ref{derived-lemma-right-derived-delta-functor}，得到结论。
\end{proof}

\begin{lemma}
\label{sites-cohomology-lemma-compute-left-derived-g-shriek}
假设和记号如情形 \ref{sites-cohomology-situation-morphism-fibred-categories}。对
$\mathcal{C}'$ 上的阿贝尔层 $\mathcal{F}'$，层
$L_ng_!(\mathcal{F}')$ 是与预层
$$
U \longmapsto H_n(\mathcal{I}_U, \mathcal{F}'_U)
$$
关联的层。记号和限制映射见证明。
\end{lemma}

\begin{proof}
设 $p(U) = V$。范畴 $\mathcal{I}_U$ 的对象是偶对
$(U', \varphi)$，其中 $\varphi : U \to u(U')$ 是 $\mathcal{C}$ 的态射，
且 $p(\varphi) = \text{id}_V$；换言之，$\varphi$ 是纤维范畴
$\mathcal{C}_V$ 的态射。态射
$(U'_1, \varphi_1) \to (U'_2, \varphi_2)$ 由纤维范畴
$\mathcal{C}'_V$ 的态射 $a : U'_1 \to U'_2$ 给出，并满足
$\varphi_2 = u(a) \circ \varphi_1$。预层 $\mathcal{F}'_U$ 把
$(U', \varphi)$ 映到 $\mathcal{F}'(U')$。下面构造限制映射。

\medskip\noindent
选取 $u$ 的分解
$$
\xymatrix{
\mathcal{C}' \ar@<1ex>[r]^{u'} &
\mathcal{C}'' \ar[r]^{u''} \ar@<1ex>[l]^w & \mathcal{C}
}
$$
如《范畴》引理
\ref{categories-lemma-ameliorate-morphism-fibred-categories} 所述。
则 $g_! = g''_! \circ g'_!$，导出函子也同样成立。另一方面，函子
$g'_!$ 正合，见《站点上的模》引理
\ref{sites-modules-lemma-have-left-adjoint}。因此
$Lg_!(\mathcal{F}') = Lg''_!(\mathcal{F}'')$，其中
$\mathcal{F}'' = g'_!\mathcal{F}'$。注意
$\mathcal{F}'' = h^{-1}\mathcal{F}'$，其中
$h : \Sh(\mathcal{C}'') \to \Sh(\mathcal{C}')$ 是与 $w$ 关联的
拓扑斯态射，见《站点》引理 \ref{sites-lemma-have-left-adjoint}。
函子 $u''$ 使 $\mathcal{C}''$ 成为 $\mathcal{C}$ 上的纤维化范畴，
故引理 \ref{sites-cohomology-lemma-compute-left-derived-pi-shriek} 可用于计算
$L_ng''_!$。又因为《范畴》引理
\ref{categories-lemma-ameliorate-morphism-fibred-categories} 的证明中
$\mathcal{C}''$ 的构造表明纤维范畴 $\mathcal{C}''_U$ 等于
$\mathcal{I}_U$，结论成立。此外，
$h^{-1}\mathcal{F}'|_{\mathcal{C}''_U}$ 由上述规则给出（由《叠》引理
\ref{stacks-lemma-topology-inherited-functorial}，$w$ 连续且余连续，
故可应用《站点》引理 \ref{sites-lemma-when-shriek}）。
\end{proof}








\section{单纯模}
\label{sites-cohomology-section-simplicial-modules}

\noindent
设 $A_\bullet$ 是一个单纯环。回忆我们可以将 $A_\bullet$
看作 $\Delta$（赋予混沌拓扑）上的层，见
《单纯对象》第 \ref{simplicial-section-simplicial-presheaves} 节。
于是，$A_\bullet$ 上的单纯模 $M_\bullet$ 正是 $\Delta$ 上的
$A_\bullet$-模层。换言之，对每个 $n \geq 0$，都有一个
$A_n$-模 $M_n$；对每个映射 $\varphi : [n] \to [m]$，都有相应的映射
$$
M_\bullet(\varphi) : M_m \longrightarrow M_n
$$
它是 $A_\bullet(\varphi)$-线性的，并且这些映射按通常方式复合。

\medskip\noindent
设 $\mathcal{C}$ 是一个站点。$\mathcal{C}$ 上的{\it 单纯环层}
$\mathcal{A}_\bullet$，是 $\mathcal{C}$ 上环层范畴中的单纯对象。
在这种情况下，指派 $U \mapsto \mathcal{A}_\bullet(U)$ 是一个
单纯环层；事实上，这两个概念等价。对于单纯阿贝尔层、
单纯 Lie 代数层等，也有类似的讨论。

\medskip\noindent
不过，与上面的单纯环情形一样，还可用另一种方式理解单纯层。
具体地，考虑投影
$$
p : \Delta \times \mathcal{C} \longrightarrow \mathcal{C}
$$
它定义了一个纤维化范畴，其中强笛卡儿态射恰为形如
$([n], U) \to ([n], V)$ 的态射。我们在范畴
$\Delta \times \mathcal{C}$ 上赋予从 $\mathcal{C}$ 继承的拓扑
（见《栈》第 \ref{stacks-section-topology} 节）。
$\Delta \times \mathcal{C}$ 中覆盖族的简单描述
（《栈》引理 \ref{stacks-lemma-topology-inherited}）立即表明：
$\mathcal{C}$ 上的单纯环层与 $\Delta \times \mathcal{C}$ 上的环层是同一回事。

\medskip\noindent
类比单纯环上的单纯模，我们如下定义单纯环层上的单纯模。

\begin{definition}
\label{sites-cohomology-definition-simplicial-module}
设 $\mathcal{C}$ 是一个站点，$\mathcal{A}_\bullet$ 是
$\mathcal{C}$ 上的单纯环层。一个
{\it 单纯 $\mathcal{A}_\bullet$-模} $\mathcal{F}_\bullet$
（有时称为{\it 单纯 $\mathcal{A}_\bullet$-模层}）是
$\Delta \times \mathcal{C}$ 上与 $\mathcal{A}_\bullet$ 相联系的环层之上的模层。
\end{definition}

\noindent
由此得到单纯模范畴 $\textit{Mod}(\mathcal{A}_\bullet)$ 以及相应的导出范畴
$D(\mathcal{A}_\bullet)$。给定单纯环层的映射
$\mathcal{A}_\bullet \to \mathcal{B}_\bullet$，可得函子
$$
- \otimes^\mathbf{L}_{\mathcal{A}_\bullet} \mathcal{B}_\bullet :
D(\mathcal{A}_\bullet)
\longrightarrow
D(\mathcal{B}_\bullet)
$$
此外，前几节的内容给出函子
$$
L\pi_! : D(\mathcal{A}_\bullet) \longrightarrow D(\mathcal{C})
$$
给定单纯模 $\mathcal{F}_\bullet$，对象
$L\pi_!(\mathcal{F}_\bullet)$ 由相联系的链复形
$s(\mathcal{F}_\bullet)$ 表示
（《单纯对象》第 \ref{simplicial-section-complexes} 节）。
这由引理 \ref{sites-cohomology-lemma-compute-left-derived-pi-shriek} 和
\ref{sites-cohomology-lemma-compute-by-cosimplicial-resolution} 得出。

\begin{lemma}
\label{sites-cohomology-lemma-base-change-by-qis}
设 $\mathcal{C}$ 是一个站点，并设
$\mathcal{A}_\bullet \to \mathcal{B}_\bullet$ 是 $\mathcal{C}$ 上单纯环层的同态。
如果 $L\pi_!\mathcal{A}_\bullet \to L\pi_!\mathcal{B}_\bullet$ 是
$D(\mathcal{C})$ 中的同构，那么
$$
L\pi_!(K) =
L\pi_!(K \otimes^\mathbf{L}_{\mathcal{A}_\bullet} \mathcal{B}_\bullet)
$$
对 $D(\mathcal{A}_\bullet)$ 中每个 $K$ 都成立。
\end{lemma}

\begin{proof}
设 $([n], U)$ 是 $\Delta \times \mathcal{C}$ 的一个对象。由于 $L\pi_!$
与余极限交换，只需对 $\mathcal{O}$-模的上有界复形证明此结论
（可与《导出范畴》命题 \ref{derived-proposition-left-derived-exists}
的论证比较，或者径直限于上有界复形）。每个这样的复形都拟同构于
一个各项均为平坦模的上有界复形，见《站点上的模》
引理 \ref{sites-modules-lemma-module-quotient-flat}。
因此，只需对平坦 $\mathcal{A}_\bullet$-模 $\mathcal{F}$ 证明该引理。
此时导出张量积就是通常的张量积，并且也是一个层。于是由
引理 \ref{sites-cohomology-lemma-compute-left-derived-pi-shriek}，可按
引理 \ref{sites-cohomology-lemma-compute-left-derived-pi-shriek-pre} 的步骤计算等式两边的同调层。
故只需对 $\mathcal{F}$ 在各纤维范畴（即 $\Delta \times U$）上的限制证明结论。
在此情形下，结论由引理 \ref{sites-cohomology-lemma-O-homology-qis} 得出。
\end{proof}

\begin{remark}
\label{sites-cohomology-remark-homology-augmentation}
设 $\mathcal{C}$ 是一个站点。设
$\epsilon : \mathcal{A}_\bullet \to \mathcal{O}$ 是环层范畴中的一个增广
（《单纯对象》定义 \ref{simplicial-definition-augmentation}）。
假设 $\epsilon$ 诱导拟同构
$s(\mathcal{A}_\bullet) \to \mathcal{O}$。
在此情形下，我们得到三角范畴间的正合函子
$$
L\pi_! : D(\mathcal{A}_\bullet) \longrightarrow D(\mathcal{O})
$$
事实上，对 $D(\mathcal{A}_\bullet)$ 的任意对象 $K$，由
引理 \ref{sites-cohomology-lemma-base-change-by-qis} 有
$L\pi_!(K) = L\pi_!(K \otimes_{\mathcal{A}_\bullet}^\mathbf{L} \mathcal{O})$
。因此，可以把所显示的函子定义为
$- \otimes^\mathbf{L}_{\mathcal{A}_\bullet} \mathcal{O}$ 与
注记 \ref{sites-cohomology-remark-fibred-category} 中函子
$L\pi_! : D(\Delta \times \mathcal{C}, \pi^{-1}\mathcal{O}) \to
D(\mathcal{O})$ 的复合。换言之，由该引理给出的
$L\pi_!(K)$ 与
$L\pi_!(K \otimes_{\mathcal{A}_\bullet}^\mathbf{L} \mathcal{O})$
之间规范且函子性的等同，我们在 $L\pi_!(K)$ 上得到一个
$\mathcal{O}$-模结构。
\end{remark}






\section{范畴上的上同调}
\label{sites-cohomology-section-cohomology}

\noindent
在例 \ref{sites-cohomology-example-category-to-point} 的情形下，除导出函子
$L\pi_!$ 外，我们还有函子 $R\pi_*$。对 $\mathcal{C}$ 上的阿贝尔层
$\mathcal{F}$，有
$H_n(\mathcal{C}, \mathcal{F}) = H^{-n}(L\pi_!\mathcal{F})$
和 $H^n(\mathcal{C}, \mathcal{F}) = H^n(R\pi_*\mathcal{F})$。

\begin{example}[计算上同调]
\label{sites-cohomology-example-right-derived-limits}
在例 \ref{sites-cohomology-example-category-to-point} 中，可如下计算函子
$H^n(\mathcal{C}, -)$。设
$\mathcal{F} \in \Ob(\textit{Ab}(\mathcal{C}))$，考虑余链复形
$$
K^\bullet(\mathcal{F}) :
\prod\nolimits_{U_0} \mathcal{F}(U_0)
\to
\prod\nolimits_{U_0 \to U_1} \mathcal{F}(U_0)
\to
\prod\nolimits_{U_0 \to U_1 \to U_2} \mathcal{F}(U_0)
\to \ldots
$$
其中转移映射为
$$
(s_{U_0 \to U_1})
\longmapsto
((U_0 \to U_1 \to U_2) \mapsto s_{U_0 \to U_1} - s_{U_0 \to U_2}
+ s_{U_1 \to U_2}|_{U_0})
$$
，其他次数中的情形类似。由构造，
$$
H^0(\mathcal{C}, \mathcal{F}) =
\lim_{\mathcal{C}^{opp}} \mathcal{F} =
H^0(K^\bullet(\mathcal{F})),
$$
见《范畴》引理 \ref{categories-lemma-limits-products-equalizers}。
$K^\bullet(\mathcal{F})$ 的构造关于 $\mathcal{F}$ 是函子性的，
并把 $\textit{Ab}(\mathcal{C})$ 的短正合列变为复形的短正合列。
因此，函子列
$\mathcal{F} \mapsto H^n(K^\bullet(\mathcal{F}))$ 构成一个 $\delta$-函子，见
《同调》定义 \ref{homology-definition-cohomological-delta-functor} 和
引理 \ref{homology-lemma-long-exact-sequence-cochain}。
对 $\mathcal{C}$ 的对象 $U$，以
$p_U : \Sh(*) \to \Sh(\mathcal{C})$ 表示相应的点，其中
$p_U^{-1}$ 等于在 $U$ 处取值，见
《站点》例 \ref{sites-example-indiscrete-points}。
设 $A$ 是一个阿贝尔群，并令 $\mathcal{F} = p_{U, *}A$。
此时，复形 $K^\bullet(\mathcal{F})$ 的各项为
$\text{Map}(X_n, A)$，其中
$$
X_n = \coprod\nolimits_{U_0 \to \ldots \to U_{n - 1} \to U_n}
\Mor_\mathcal{C}(U, U_0)
$$
这个单纯集同伦等价于单点集 $\{*\}$ 上的常单纯集。事实上，
映射 $X_\bullet \to \{*\}$ 是显然的；映射
$\{*\} \to X_n$ 把 $*$ 映到
$(U \to \ldots \to U, \text{id}_U)$；而定义两个映射
$X_\bullet \to X_\bullet$ 之间同伦的映射
$$
h_{n, i} : X_n \longrightarrow X_n
$$
（《单纯对象》引理 \ref{simplicial-lemma-relations-homotopy}）
由下述规则给出：
$$
h_{n, i} :
(U_0 \to \ldots \to U_n, f)
\longmapsto
(U \to \ldots \to U \to U_i \to \ldots \to U_n, \text{id})
$$
这里 $i > 0$，且 $h_{n, 0} = \text{id}$。验证从略。
由于 $\text{Map}(-, A)$ 是反变函子，这表明
$K^\bullet(p_{U, *}A)$ 在正次数中的上同调为零
（这使用《单纯对象》注记 \ref{simplicial-remark-homotopy-better}
的函子性以及《单纯对象》引理
\ref{simplicial-lemma-homotopy-s-Q} 的结论）。
由此还可知：若 $\mathcal{F}$ 是形如
$p_{U, *}A$ 的层的乘积，则 $K^\bullet(\mathcal{F})$ 在正次数中无环。
由于 $\mathcal{C}$ 上每个阿贝尔层都嵌入这样的乘积，故例如由
《同调》引理 \ref{homology-lemma-efface-implies-universal} 和
《导出范畴》引理 \ref{derived-lemma-right-derived-delta-functor}，
可知 $K^\bullet(\mathcal{F})$ 计算 $H^0(\mathcal{C}, -)$ 的左导出函子
$H^n(\mathcal{C}, -)$。
\end{example}

\begin{example}[计算 Ext]
\label{sites-cohomology-example-computing-exts}
在例 \ref{sites-cohomology-example-category-to-point} 中，进一步假设给定了
$\mathcal{C}$ 上的环层 $\mathcal{O}$。设
$\mathcal{F}$、$\mathcal{G}$ 是 $\mathcal{O}$-模。
考虑复形 $K^\bullet(\mathcal{G}, \mathcal{F})$，其 $n$ 次项为
$$
\prod\nolimits_{U_0 \to U_1 \to \ldots \to U_n}
\Hom_{\mathcal{O}(U_n)}(\mathcal{G}(U_n), \mathcal{F}(U_0))
$$
，转移映射为
$$
(\varphi_{U_0 \to U_1})
\longmapsto
((U_0 \to U_1 \to U_2) \mapsto
\varphi_{U_0 \to U_1} \circ \rho^{U_2}_{U_1}
- \varphi_{U_0 \to U_2}
+ \rho^{U_1}_{U_0} \circ \varphi_{U_1 \to U_2}
$$
，其他次数中的情形类似。这里的 $\rho$ 表示限制映射。由构造，
$$
\Hom_\mathcal{O}(\mathcal{G}, \mathcal{F}) =
H^0(K^\bullet(\mathcal{G}, \mathcal{F}))
$$
对任意一对 $\mathcal{O}$-模 $\mathcal{F}, \mathcal{G}$ 都成立。指派
$(\mathcal{G}, \mathcal{F}) \mapsto K^\bullet(\mathcal{G}, \mathcal{F})$
是一个双函子；它把第一变元中的直和变为乘积，并与第二变元中的乘积交换。
我们断言
$$
\Ext^i_\mathcal{O}(\mathcal{G}, \mathcal{F}) =
H^i(K^\bullet(\mathcal{G}, \mathcal{F}))
$$
对 $i \geq 0$ 成立，只要下列条件之一成立：
\begin{enumerate}
\item 对每个 $U \in \Ob(\mathcal{C})$，
$\mathcal{G}(U)$ 都是投射 $\mathcal{O}(U)$-模；或
\item 对每个 $U \in \Ob(\mathcal{C})$，
$\mathcal{F}(U)$ 都是内射 $\mathcal{O}(U)$-模。
\end{enumerate}
在情形 (1) 中，函子 $K^\bullet(\mathcal{G}, -)$ 是从
$\mathcal{O}$-模范畴到阿贝尔群余链复形范畴的正合函子。
因此，仿照例 \ref{sites-cohomology-example-right-derived-limits} 的论证，
只需证明：当 $\mathcal{F}=p_{U, *}A$，其中 $A$ 是一个
$\mathcal{O}(U)$-模时，$K^\bullet(\mathcal{G}, \mathcal{F})$
在正次数中无环。选取短正合列
\begin{equation}
\label{sites-cohomology-equation-split}
0 \to \mathcal{G}' \to \bigoplus j_{U_i!}\mathcal{O}_{U_i} \to
\mathcal{G} \to 0
\end{equation}
，见《站点上的模》引理
\ref{sites-modules-lemma-module-quotient-flat}。
由于 (1) 对中间的层和右端的层成立，它也对 $\mathcal{G}'$ 成立；
而在 $\mathcal{C}$ 的任一对象上对
(\ref{sites-cohomology-equation-split}) 取值，得到模的分裂正合列。
于是，对任意 $\mathcal{F}$，特别是对
$\mathcal{F} = p_{U, *}A$，得到复形的短正合列
$$
0 \to
K^\bullet(\mathcal{G}, \mathcal{F}) \to
\prod K^\bullet(j_{U_i!}\mathcal{O}_{U_i}, \mathcal{F}) \to
K^\bullet(\mathcal{G}', \mathcal{F}) \to 0
$$
。在 $H^0$ 上得到
$$
0 \to \Hom(\mathcal{G}, p_{U, *}A) \to
\Hom(\prod j_{U_i!}\mathcal{O}_{U_i}, p_{U, *}A) \to
\Hom(\mathcal{G}', p_{U, *}A) \to 0
$$
，它是正合的，因为
$\Hom(\mathcal{H}, p_{U, *}A) = \Hom_{\mathcal{O}(U)}(\mathcal{H}(U), A)$
，并且 (\ref{sites-cohomology-equation-split}) 在 $U$ 上的截面列分裂正合。
因此，利用维数转移可知，只需证明对所有
$U, U' \in \Ob(\mathcal{C})$，
$K^\bullet(j_{U'!}\mathcal{O}_{U'}, p_{U, *}A)$ 在正次数中无环。
在此情形下，
$K^n(j_{U'!}\mathcal{O}_{U'}, p_{U, *}A)$ 等于
$$
\prod\nolimits_{U \to U_0 \to U_1 \to \ldots \to U_n \to U'} A
$$
换言之，$K^\bullet(j_{U'!}\mathcal{O}_{U'}, p_{U, *}A)$
是各项为 $\text{Map}(X_\bullet, A)$ 的复形，其中
$$
X_n = \coprod\nolimits_{U_0 \to \ldots \to U_{n - 1} \to U_n}
\Mor_\mathcal{C}(U, U_0) \times \Mor_\mathcal{C}(U_n, U')
$$
这个单纯集同伦等价于单点集 $\{*\}$ 上的常单纯集；其证明与
例 \ref{sites-cohomology-example-right-derived-limits} 中相应陈述的证明完全相同。
这就完成了该断言的证明。

\medskip\noindent
情形 (2) 中的论证类似（但与之对偶）。
\end{example}






\section{范畴上的模}
\label{sites-cohomology-section-modules-cohomology}

\noindent
本节内容将用于定义概形范畴上的群胚栈之拟凝聚模导出范畴的一种变体。
见《栈上的层》第 \ref{stacks-sheaves-section-QC} 节。

\medskip\noindent
设 $\mathcal{C}$ 是一个范畴。我们把 $\mathcal{C}$ 看作赋予混沌拓扑的站点。
如例 \ref{sites-cohomology-example-computing-exts}，设 $\mathcal{O}$ 是
$\mathcal{C}$ 上的环层。换言之，$\mathcal{O}$ 是范畴
$\mathcal{C}$ 上的环预层，见《范畴》定义
\ref{categories-definition-presheaf}。

\begin{definition}
\label{sites-cohomology-definition-cartesian}
在上述情形下，以 $\mathit{QC}(\mathcal{C}, \mathcal{O})$，或简记为
{\it $\mathit{QC}(\mathcal{O})$}，表示
$D(\mathcal{O}) = D(\mathcal{C}, \mathcal{O})$ 的如下全子范畴：
其对象 $K$ 满足，对 $\mathcal{C}$ 中每个 $U \to V$，规范映射
$$
R\Gamma(V, K) \otimes_{\mathcal{O}(V)}^\mathbf{L} \mathcal{O}(U)
\longrightarrow
R\Gamma(U, K)
$$
是 $D(\mathcal{O}(U))$ 中的同构。
\end{definition}

\begin{lemma}
\label{sites-cohomology-lemma-cartesian-good-sub}
在上述情形下，子范畴 $\mathit{QC}(\mathcal{O})$ 是
$D(\mathcal{O})$ 的严格全、饱和三角子范畴，并且在任意直和下保持不变。
\end{lemma}

\begin{proof}
设 $U$ 是 $\mathcal{C}$ 的一个对象。由于 $\mathcal{C}$ 上的拓扑是混沌拓扑，
函子 $\mathcal{F} \mapsto \mathcal{F}(U)$ 正合且与直和交换。
因此，要计算正合函子 $K \mapsto R\Gamma(U, K)$，可用任意
$\mathcal{O}$-模复形 $\mathcal{F}^\bullet$ 表示 $K$，再取
$\mathcal{F}^\bullet(U)$。故 $R\Gamma(U, -)$ 与直和交换，见
《内射对象》引理 \ref{injectives-lemma-derived-products}。
类似地，给定 $\mathcal{C}$ 的态射 $U \to V$，导出张量积函子
$- \otimes_{\mathcal{O}(V)}^\mathbf{L} \mathcal{O}(U) :
D(\mathcal{O}(V)) \to D(\mathcal{O}(U))$
正合且与直和交换。由这些观察可直接得到该引理；细节从略。
\end{proof}

\begin{lemma}
\label{sites-cohomology-lemma-QC-bounded-above}
在上述情形下，假设 $M$ 是 $\mathit{QC}(\mathcal{O})$ 的一个对象，
且 $b \in \mathbf{Z}$ 满足对所有 $i > b$ 都有 $H^i(M) = 0$。
那么，$H^b(M)$ 是 $(\mathcal{C}, \mathcal{O})$ 上的拟凝聚模，
这里采用《站点上的模》定义 \ref{sites-modules-definition-site-local} 的含义。
\end{lemma}

\begin{proof}
由《站点上的模》引理
\ref{sites-modules-lemma-chaotic-quasi-coherent}，只需证明对
$\mathcal{C}$ 的每个态射 $U \to V$，映射
$$
H^p(M)(V) \otimes_{\mathcal{O}(V)} \mathcal{O}(U) \to H^b(M)(U)
$$
是同构。已知映射
$$
R\Gamma(V, M) \otimes_{\mathcal{O}(V)}^\mathbf{L} \mathcal{O}(U)
\to R\Gamma(U, M)
$$
是同构。于是，例如由 Tor 谱序列可得结论。细节从略。
\end{proof}

\begin{lemma}
\label{sites-cohomology-lemma-QC-final-object}
在上述情形下，假设 $\mathcal{C}$ 有终对象 $X$。
令 $R = \mathcal{O}(X)$，并以
$f : (\mathcal{C}, \mathcal{O}) \to (pt, R)$ 表示显然的站点态射。
那么，由 $Lf^*$ 和 $Rf_*$ 给出 $\mathit{QC}(\mathcal{O}) = D(R)$。
\end{lemma}

\begin{proof}
从略。
\end{proof}

\begin{lemma}
\label{sites-cohomology-lemma-QC-hom-out-of}
在上述情形下，假设 $K$ 是 $\mathit{QC}(\mathcal{O})$ 的一个对象，
而 $M$ 是 $D(\mathcal{O})$ 中的任意对象。对 $\mathcal{C}$ 的每个对象
$U$，都有
$$
\Hom_{D(\mathcal{O}_U)}(K|_U, M|_U) =
R\Hom_{\mathcal{O}(U)}(R\Gamma(U, K), R\Gamma(U, M))
$$
\end{lemma}

\begin{proof}
可用 $\mathcal{C}/U$ 代替 $\mathcal{C}$。因而可设
$U = X$ 是 $\mathcal{C}$ 的终对象。由
引理 \ref{sites-cohomology-lemma-QC-final-object}，可知 $K = Lf^*P$，其中
$P = R\Gamma(U, K) = R\Gamma(X, K) = Rf_*K$。
由于 $Lf^*$ 是 $Rf_*(-) = R\Gamma(U, -)$ 的左伴随，结论随即成立。
\end{proof}

\noindent
设 $(\mathcal{C}, \mathcal{O})$ 如上。对 $\mathcal{O}$-模复形
$\mathcal{F}^\bullet$，把 $\mathcal{F}^\bullet$ 的{\it 大小}
$|\mathcal{F}^\bullet|$ 定义为
$$
|\mathcal{F}^\bullet| =
\left|
\coprod\nolimits_{i \in \mathbf{Z},\ U \in \Ob(\mathcal{C})} \mathcal{F}^i(U)
\right|
$$
对 $D(\mathcal{O})$ 的对象 $K$，把 $K$ 的{\it 大小} $|K|$ 定义为基数
$$
|K| = \min
\left\{
\left|
\mathcal{F}^\bullet
\right|
\text{ 其中 }\mathcal{F}^\bullet\text{ 表示 }K
\right\}
$$
由基数的性质，这个最小值存在。

\begin{lemma}
\label{sites-cohomology-lemma-build-from}
在上述情形下，存在一个基数 $\kappa$，具有如下性质：
给定 $\mathcal{O}$-模复形 $\mathcal{F}^\bullet$ 以及子集
$\Omega^i_U \subset \mathcal{F}^i(U)$，存在子复形
$\mathcal{H}^\bullet \subset \mathcal{F}^\bullet$，使得
$\Omega^i_U \subset \mathcal{H}^i(U)$，并且
$|\mathcal{H}^\bullet| \leq \max(\kappa, |\bigcup \Omega^i_U|)$。
\end{lemma}

\begin{proof}
把 $\mathcal{H}^i(U)$ 定义为 $\mathcal{F}^i(U)$ 的如下
$\mathcal{O}(U)$-子模：它由 $\Omega^i_V$ 在沿任意态射
$f : U \to V$ 限制所得的像以及 $\text{d}(\Omega^{i - 1}_U)$ 生成。
$\mathcal{H}^i(U)$ 的基数不超过下列各基数的最大值：
$\aleph_0$、$\mathcal{O}(U)$ 的基数、
$\text{Arrows}(\mathcal{C})$ 的基数以及
$|\bigcup \Omega^i_U|$。细节从略。
\end{proof}

\begin{lemma}
\label{sites-cohomology-lemma-qis-small}
在上述情形下，存在一个基数 $\kappa$，具有如下性质：
给定表示 $D(\mathcal{O})$ 的对象 $K$ 的 $\mathcal{O}$-模复形
$\mathcal{F}^\bullet$，存在子复形
$\mathcal{H}^\bullet \subset \mathcal{F}^\bullet$，使得
$\mathcal{H}^\bullet$ 表示 $K$，且
$|\mathcal{H}^\bullet| \leq \max(\kappa, |K|)$。
\end{lemma}

\begin{proof}
首先，对每个 $i$ 和 $U$，选取子集
$\Omega^i_U \subset \Ker(\text{d} :
\mathcal{F}^i(U) \to \mathcal{F}^{i + 1}(U))$
，使它双射到
$H^i(K)(U) = H^i(\mathcal{F}^\bullet(U))$ 上。
由于可用大小为 $|K|$ 的复形表示 $K$，故
$|\Omega^i_U| \leq |K|$。应用引理 \ref{sites-cohomology-lemma-build-from}，
得到包含 $\Omega^i_U$ 的子复形
$\mathcal{S}^\bullet \subset \mathcal{F}^\bullet$，其大小至多为
$\max(\kappa, |K|)$；因而
$H^i(\mathcal{S}^\bullet) \to H^i(\mathcal{F}^\bullet)$
是层的满射。

\medskip\noindent
我们将归纳构造子复形
$$
\mathcal{S}^\bullet =
\mathcal{S}_0^\bullet \subset
\mathcal{S}_1^\bullet \subset
\mathcal{S}_2^\bullet \subset \ldots \subset \mathcal{F}^\bullet
$$
，使其大小均不超过 $\max(\kappa, |K|)$，并且
$H^i(\mathcal{S}_n^\bullet) \to H^i(\mathcal{F}^\bullet)$ 的核
与 $H^i(\mathcal{S}_n^\bullet) \to
H^i(\mathcal{S}_{n + 1}^\bullet)$ 的核相同。
构造完成后，可取
$\mathcal{H}^\bullet = \bigcup \mathcal{S}_n^\bullet$ 作为所求。

\medskip\noindent
现由 $\mathcal{S}_n^\bullet$ 构造
$\mathcal{S}_{n + 1}^\bullet$。对每个 $U$ 和 $i$，取子集
$\Omega^{i - 1}_U \subset \mathcal{F}^{i - 1}(U)$，使得
$\text{d} : \mathcal{F}^{i - 1}(U) \to \mathcal{F}^i(U)$
把 $\Omega^{i - 1}_U$ 双射到
$$
\mathcal{S}_n^i(U) \cap
\text{Im}(\text{d} : \mathcal{F}^{i - 1}(U)\to \mathcal{F}^i(U))
$$
上。注意 $|\Omega^i_U| \leq |K|$，因为
$\mathcal{S}_n^i(U)$ 具有这个界。随后对下列子集应用
引理 \ref{sites-cohomology-lemma-build-from}：
$$
\mathcal{S}^i(U) \cup \Omega^i_U \subset \mathcal{F}^i(U)
$$
即可得到 $\mathcal{S}_{n + 1}^\bullet$；其余均显然。
\end{proof}

\begin{lemma}
\label{sites-cohomology-lemma-cartesian-colimit-kappa}
在上述情形下，存在一个具有如下性质的基数 $\kappa$：
\begin{enumerate}
\item 对 $\mathit{QC}(\mathcal{O})$ 的每个非零对象 $K$，
存在 $\mathit{QC}(\mathcal{O})$ 中的非零态射 $E \to K$，
使得 $|E| \leq \kappa$；
\item 对 $\mathit{QC}(\mathcal{O})$ 中每个满足
$|E| \leq \kappa$ 的态射
$\alpha : E \to \bigoplus_n K_n$，存在
$\mathit{QC}(\mathcal{O})$ 中的态射 $E_n \to K_n$，其中
$|E_n| \leq \kappa$，使得 $\alpha$ 经由
$\bigoplus E_n \to \bigoplus K_n$ 分解。
\end{enumerate}
\end{lemma}

\begin{proof}
取 $\kappa$ 为下列基数集合的一个上界：
\begin{enumerate}
\item 对所有 $U \in \Ob(\mathcal{C})$ 的
$|\coprod_V j_{U!}\mathcal{O}_U(V)|$；
\item 对 $\mathcal{C}$ 中所有态射 $U \to V$，在《代数进阶》
引理 \ref{more-algebra-lemma-enlarge} 中得到的基数
$\kappa(\mathcal{O}(V) \to \mathcal{O}(U))$；
\item 引理 \ref{sites-cohomology-lemma-qis-small} 中得到的基数。
\end{enumerate}
我们断言：对表示 $\mathit{QC}(\mathcal{O})$ 中一个对象的任意复形
$\mathcal{F}^\bullet$，以及任意满足
$|\mathcal{S}^\bullet| \leq \kappa$ 的子复形
$\mathcal{S}^\bullet \subset \mathcal{F}^\bullet$，
存在 $\mathcal{F}^\bullet$ 的子复形 $\mathcal{H}^\bullet$，
它包含 $\mathcal{S}^\bullet$，表示
$\mathit{QC}(\mathcal{O})$ 中的一个对象，并且
$|\mathcal{H}^\bullet| \leq \kappa$。
下面两段说明该断言蕴含此引理。

\medskip\noindent
如 (1) 所述，设 $K$ 是 $\mathit{QC}(\mathcal{O})$ 的非零对象。
设 $K$ 由 $\mathcal{O}$-模复形 $\mathcal{F}^\bullet$ 表示。
那么对某个 $i$，$H^i(\mathcal{F}^\bullet)$ 非零。因此，
存在 $\mathcal{C}$ 的对象 $U$ 以及满足 $d(s) = 0$ 的截面
$s \in \mathcal{F}^i(U)$，它确定
$H^i(\mathcal{F}^\bullet)$ 在 $U$ 上的一个非零截面。
于是，$s : j_{U!}\mathcal{O}_U[-i] \to \mathcal{F}^\bullet$
的像是满足 $|\mathcal{S}^\bullet| \leq \kappa$ 的子复形
$\mathcal{S}^\bullet \subset \mathcal{F}^\bullet$。
应用该断言，得到 $\mathit{QC}(\mathcal{O})$ 中的非零态射
$\mathcal{H}^\bullet \to \mathcal{F}^\bullet$，其中
$|\mathcal{H}^\bullet| \leq \kappa$。故 (1) 成立。

\medskip\noindent
设 $\alpha : E \to \bigoplus K_n$ 如 (2)。
任取表示 $K_n$ 的复形 $\mathcal{K}_n^\bullet$。
于是 $\bigoplus \mathcal{K}_n^\bullet$ 表示 $\bigoplus K_n$。
由导出范畴的构造，可用复形 $\mathcal{E}^\bullet$ 表示 $E$，
使 $\alpha$ 由复形态射
$a : \mathcal{E}^\bullet \to \bigoplus \mathcal{K}_n^\bullet$
表示。由引理 \ref{sites-cohomology-lemma-qis-small} 以及上面对 $\kappa$ 的选取，
可设 $|\mathcal{E}^\bullet| \leq \kappa$。
应用该断言，得到子复形
$\mathcal{E}_n^\bullet \subset \mathcal{K}_n^\bullet$；
它们表示 $\mathit{QC}(\mathcal{O})$ 中满足
$|E_n| \leq \kappa$ 的对象 $E_n$，并包含
$a_n : \mathcal{E}^\bullet \to \mathcal{K}_n^\bullet$ 的像，正合所需。

\medskip\noindent
下面证明该断言。设 $\mathcal{F}^\bullet$ 是表示
$\mathit{QC}(\mathcal{O})$ 中一个对象的复形，并设
$\mathcal{S}^\bullet \subset \mathcal{F}^\bullet$ 是大小不超过
$\kappa$ 的子复形。我们将归纳构造子复形
$$
\mathcal{S}^\bullet = \mathcal{S}_0^\bullet \subset
\mathcal{S}_1^\bullet \subset
\mathcal{S}_2^\bullet \subset \ldots \subset \mathcal{F}^\bullet
$$
，使其大小均不超过 $\kappa$，且对 $\mathcal{C}$ 的每个态射
$f : U \to V$ 和每个 $i \in \mathbf{Z}$，满足
\begin{enumerate}
\item 映射
$H^i(\mathcal{S}_n^\bullet(V)
\otimes_{\mathcal{O}(V)}^\mathbf{L} \mathcal{O}(U)) \to
H^i(\mathcal{S}_n^\bullet(U))$
的核在
$H^i(\mathcal{S}_{n + 1}^\bullet(V)
\otimes_{\mathcal{O}(V)}^\mathbf{L} \mathcal{O}(U))$ 中的像为零；
\item 映射
$H^i(\mathcal{S}_n^\bullet(U)) \to H^i(\mathcal{S}_{n + 1}^\bullet(U))$
的像包含于下述映射的像：
$H^i(\mathcal{S}_{n + 1}^\bullet(V)
\otimes_{\mathcal{O}(V)}^\mathbf{L} \mathcal{O}(U)) \to
H^i(\mathcal{S}_{n + 1}^\bullet(U))$。
\end{enumerate}
构造完成后，可令
$\mathcal{H}^\bullet = \bigcup \mathcal{S}_n^\bullet$。
事实上，由于导出张量积以及对环上模复形取上同调都与滤过余极限交换，
条件 (1) 和 (2) 合起来保证
$$
\mathcal{H}^\bullet(V) \otimes_{\mathcal{O}(V)}^\mathbf{L} \mathcal{O}(U)
\longrightarrow
\mathcal{H}^\bullet(U)
$$
对所有 $\mathcal{C}$ 中的 $f : U \to V$，在每个次数的上同调上都是同构，
因而是 $D(\mathcal{O}(U))$ 中的同构。
所以 $\mathcal{H}^\bullet$ 如所需地表示
$\mathit{QC}(\mathcal{O})$ 中的一个对象。

\medskip\noindent
现由 $\mathcal{S}_n$ 构造 $\mathcal{S}_{n + 1}$。
对 $\mathcal{C}$ 的每个态射 $f : U \to V$，考虑交换图
$$
\xymatrix{
\mathcal{S}_n^\bullet(V)
\ar[r] \ar[d] &
\mathcal{S}_n^\bullet(U) \ar[d] \\
\mathcal{F}^\bullet(V)
\ar[r] &
\mathcal{F}^\bullet(U)
}
$$
对于环映射 $\mathcal{O}(V) \to \mathcal{O}(U)$，这是
《代数进阶》引理 \ref{more-algebra-lemma-enlarge} 中那样的图；
也就是说，底行诱导同构
$$
\mathcal{F}^\bullet(V) \otimes_{\mathcal{O}(V)}^\mathbf{L} \mathcal{O}(U)
\longrightarrow
\mathcal{F}^\bullet(U)
$$
于 $D(\mathcal{O}(U))$ 中。因此，可以按照《代数进阶》
引理 \ref{more-algebra-lemma-enlarge} 选取子复形
$$
\mathcal{S}_n^\bullet(V) \subset M^\bullet_f \subset \mathcal{F}^\bullet(V)
\quad\text{且}\quad
\mathcal{S}_n^\bullet(U) \subset N^\bullet_f \subset \mathcal{F}^\bullet(U)
$$
，特别有 $|N^i_f|, |M^i_f| \leq \kappa$。接着，利用子集
$$
\mathcal{S}_n^i(U) \amalg \coprod\nolimits_{f : U \to V} N^i_f
\amalg \coprod\nolimits_{g : W \to U} M^i_g
\subset
\mathcal{F}^i(U)
$$
应用引理 \ref{sites-cohomology-lemma-build-from}，得到子复形
$$
\mathcal{S}_n^\bullet \subset
\mathcal{S}_{n + 1}^\bullet \subset
\mathcal{F}^\bullet
$$
，它包含这些子集，并且
$|\mathcal{S}_{n + 1}^\bullet| \leq \kappa$。
条件 (1) 和 (2) 成立，因为由《代数进阶》引理
\ref{more-algebra-lemma-enlarge} 中的构造，相应陈述对
$\mathcal{S}_n^\bullet(V) \subset M^\bullet_f$ 和
$\mathcal{S}_n^\bullet(U) \subset N^\bullet_f$ 成立。
证明完毕。
\end{proof}

\begin{proposition}
\label{sites-cohomology-proposition-cartesian-brown}
设 $\mathcal{C}$ 是一个范畴，并将其视为赋予混沌拓扑的站点。
设 $\mathcal{O}$ 是 $\mathcal{C}$ 上的环层。令
$\mathit{QC}(\mathcal{O})$ 如定义 \ref{sites-cohomology-definition-cartesian}，则
\begin{enumerate}
\item $\mathit{QC}(\mathcal{O})$ 是 $D(\mathcal{O})$ 的严格全、饱和三角子范畴，
并且在任意直和下保持不变；
\item 任何把直和变为乘积的反变上同调函子
$H : \mathit{QC}(\mathcal{O}) \to \textit{Ab}$
都是可表的；
\item 任何把直和变为直和的三角范畴间正合函子
$F : \mathit{QC}(\mathcal{O}) \to \mathcal{D}$
都有正合右伴随；
\item 包含函子 $\mathit{QC}(\mathcal{O}) \to D(\mathcal{O})$
有正合右伴随。
\end{enumerate}
\end{proposition}

\begin{proof}
(1) 是引理 \ref{sites-cohomology-lemma-cartesian-good-sub}。
(2) 由引理 \ref{sites-cohomology-lemma-cartesian-colimit-kappa} 和
《导出范畴》引理 \ref{derived-lemma-brown-bis} 得出。
(3) 由引理 \ref{sites-cohomology-lemma-cartesian-colimit-kappa} 和
《导出范畴》命题 \ref{derived-proposition-brown-bis} 得出。
(4) 是 (3) 的特例。
\end{proof}

\noindent
设 $u : \mathcal{C}' \to \mathcal{C}$ 是范畴间的函子。
若把 $\mathcal{C}$ 和 $\mathcal{C}'$ 视为赋予混沌拓扑的站点，
则 $u$ 是连续且余连续的函子。因此得到拓扑斯态射
$g : \Sh(\mathcal{C}') \to \Sh(\mathcal{C})$，见
《站点》引理 \ref{sites-lemma-cocontinuous-morphism-topoi}。
此外，假设给定 $\mathcal{C}$ 上的环层 $\mathcal{O}$、
$\mathcal{C}'$ 上的环层 $\mathcal{O}'$，以及映射
$g^\sharp : g^{-1}\mathcal{O} \to \mathcal{O}'$。
把相应的环化拓扑斯态射简记为
$g : (\Sh(\mathcal{C}'), \mathcal{O}') \to (\Sh(\mathcal{C}), \mathcal{O})$,
见《站点上的模》第 \ref{sites-modules-section-ringed-topoi} 节。

\begin{lemma}
\label{sites-cohomology-lemma-cartesian-functorial}
设
$g : (\Sh(\mathcal{C}'), \mathcal{O}') \to
(\Sh(\mathcal{C}), \mathcal{O})$ 如上。那么，函子
$Lg^* : D(\mathcal{O}) \to D(\mathcal{O}')$
把 $\mathit{QC}(\mathcal{O})$ 映入 $\mathit{QC}(\mathcal{O}')$。
\end{lemma}

\begin{proof}
设 $U' \in \Ob(\mathcal{C}')$，其在 $\mathcal{C}$ 中的像为
$U = u(U')$。以 $pt$ 表示只有一个对象和一个态射的范畴。
以 $(\Sh(pt), \mathcal{O}'(U'))$ 和
$(\Sh(pt), \mathcal{O}(U))$ 表示所指的环化拓扑斯。
当然，我们把这些环化拓扑斯上模的导出范畴分别等同于
$D(\mathcal{O}'(U'))$ 和 $D(\mathcal{O}(U))$。
于是有环化拓扑斯的交换图
$$
\xymatrix{
(\Sh(pt), \mathcal{O}'(U')) \ar[rr]_{U'} \ar[d] & &
(\Sh(\mathcal{C}'), \mathcal{O}') \ar[d]^g \\
(\Sh(pt), \mathcal{O}(U)) \ar[rr]^U & &
(\Sh(\mathcal{C}), \mathcal{O})
}
$$
沿下方水平态射拉回，把 $D(\mathcal{O})$ 中的 $K$ 送到
$R\Gamma(U, K)$。沿左侧竖直箭头拉回，把 $M$ 送到
$M \otimes_{\mathcal{O}(U)}^\mathbf{L} \mathcal{O}'(U')$。
沿图的两个方向复合所得结果相同
（引理 \ref{sites-cohomology-lemma-derived-pullback-composition}），故
$$
R\Gamma(U', Lg^*K) =
R\Gamma(U, K) \otimes_{\mathcal{O}(U)}^\mathbf{L} \mathcal{O}'(U')
$$
最后，设 $f' : U' \to V'$ 是 $\mathcal{C}'$ 中的态射，并以
$f = u(f') : U = u(U') \to V = u(V')$ 表示其在 $\mathcal{C}$ 中的像。
若 $K$ 属于 $\mathit{QC}(\mathcal{O})$，则
\begin{align*}
R\Gamma(V', Lg^*K) \otimes_{\mathcal{O}'(V')}^\mathbf{L} \mathcal{O}'(U')
& =
R\Gamma(V, K) \otimes_{\mathcal{O}(V)}^\mathbf{L} \mathcal{O}'(V')
\otimes_{\mathcal{O}'(V')}^\mathbf{L} \mathcal{O}'(U') \\
& =
R\Gamma(V, K) \otimes_{\mathcal{O}(V)}^\mathbf{L} \mathcal{O}'(U') \\
& =
R\Gamma(V, K) \otimes_{\mathcal{O}(V)}^\mathbf{L} \mathcal{O}(U)
\otimes_{\mathcal{O}(U)}^\mathbf{L} \mathcal{O}'(U') \\
& =
R\Gamma(U, K) \otimes_{\mathcal{O}(U)}^\mathbf{L} \mathcal{O}'(U') \\
& =
R\Gamma(U', Lg^*K)
\end{align*}
，正合所需。这里对 $U'$ 和 $V'$ 都使用了上面的观察。
\end{proof}

\begin{lemma}
\label{sites-cohomology-lemma-cartesian-flat-quasi-coherent}
设 $\mathcal{C}$ 是一个范畴，并将其视为赋予混沌拓扑的站点。
设 $\mathcal{O}$ 是 $\mathcal{C}$ 上的环层。
假设对 $\mathcal{C}$ 中所有 $U \to V$，限制映射
$\mathcal{O}(V) \to \mathcal{O}(U)$ 都是平坦环映射。
那么，$\mathit{QC}(\mathcal{O})$ 等于
$D(\mathcal{O})$ 中由上同调层均拟凝聚的复形组成的子范畴
$D_\QCoh(\mathcal{O}) \subset D(\mathcal{O})$。
\end{lemma}

\begin{proof}
回忆在我们的假设下，
$\QCoh(\mathcal{O}) \subset \textit{Mod}(\mathcal{O})$
是弱 Serre 子范畴，见《站点上的模》引理
\ref{sites-modules-lemma-chaotic-flat-quasi-coherent}。
因此，可取 $D(\mathcal{O})$ 的全子范畴
$$
D_\QCoh(\mathcal{O}) = D_{\QCoh(\mathcal{O})}(\textit{Mod}(\mathcal{O}))
$$
，见《导出范畴》第 \ref{derived-section-triangulated-sub} 节。
（严格说来，在该引理的证明中并不需要这一点。）

\medskip\noindent
设 $M$ 是 $\mathit{QC}(\mathcal{O})$ 的一个对象。
由于对 $\mathcal{C}$ 的每个态射 $U \to V$，限制映射
$\mathcal{O}(V) \to \mathcal{O}(U)$ 都是平坦的，故
\begin{align*}
H^i(M)(U)
& =
H^i(R\Gamma(U, M)) \\
& =
H^i(R\Gamma(V, M) \otimes_{\mathcal{O}(V)}^\mathbf{L} \mathcal{O}(U)) \\
& =
H^i(R\Gamma(V, M)) \otimes_{\mathcal{O}(V)} \mathcal{O}(U) \\
& =
H^i(M)(V) \otimes_{\mathcal{O}(V)} \mathcal{O}(U)
\end{align*}
，从而由《站点上的模》引理
\ref{sites-modules-lemma-chaotic-quasi-coherent}，
$H^i(M)$ 是拟凝聚的。上面的第一个和最后一个等式源于如下事实：
由于 $\mathcal{C}$ 上的拓扑是混沌拓扑，在 $\mathcal{C}$ 的一个对象上取截面是正合函子。

\medskip\noindent
反过来，若 $M$ 是 $D_\QCoh(\mathcal{O})$ 的一个对象，则由
《站点上的模》引理
\ref{sites-modules-lemma-chaotic-quasi-coherent}，
映射 $R\Gamma(V, M) \to R\Gamma(U, M)$ 诱导同构
$H^i(M)(U) \to H^i(M)(V) \otimes_{\mathcal{O}(V)} \mathcal{O}(U)$。
于是，由 $\mathcal{O}(V) \to \mathcal{O}(U)$ 的平坦性，
$R\Gamma(V, K) \otimes_{\mathcal{O}(V)}^\mathbf{L} \mathcal{O}(U)
\to R\Gamma(U, K)$ 是 $D(\mathcal{O}(U))$ 中的同构，
故 $M$ 属于 $\mathit{QC}(\mathcal{O})$。
\end{proof}

\begin{lemma}
\label{sites-cohomology-lemma-cartesion-plus-topology}
设 $\epsilon : (\mathcal{C}_\tau, \mathcal{O}_\tau) \to
(\mathcal{C}_{\tau'}, \mathcal{O}_{\tau'})$ 如
第 \ref{sites-cohomology-section-compare} 节。假设
\begin{enumerate}
\item $\tau'$ 是范畴 $\mathcal{C}$ 上的混沌拓扑；
\item 对所有 $U \in \Ob(\mathcal{C})$ 以及所有
$\mathcal{O}(U)$-模的 K-平坦复形 $M^\bullet$，映射
$$
M^\bullet \longrightarrow
R\Gamma((\mathcal{C}/U)_\tau,
(M^\bullet \otimes_{\mathcal{O}(U)} \mathcal{O}_U)^\#)
$$
是拟同构（解释见证明）。
\end{enumerate}
那么，$\epsilon^*$ 和 $R\epsilon_*$ 在
$\mathit{QC}(\mathcal{O})$ 与
$D(\mathcal{C}_\tau, \mathcal{O}_\tau)$ 的如下全子范畴之间给出互为拟逆的等价：
该全子范畴由满足 $R\epsilon_*K$ 属于
$\mathit{QC}(\mathcal{O})$ 的对象 $K$ 组成\footnote{这意味着，对
$\mathcal{C}$ 中所有 $U \to V$，映射
$R\Gamma(V, K) \otimes_{\mathcal{O}(V)}^\mathbf{L} \mathcal{O}(U)
\to R\Gamma(U, K)$ 是同构。}。
\end{lemma}

\begin{proof}
我们将使用第 \ref{sites-cohomology-section-compare} 节中的观察而不再特别说明。
由于 $R\epsilon_*$ 完全忠实且
$\epsilon^* \circ R\epsilon_* = \text{id}$，为证明该引理，
只需证明对 $\mathit{QC}(\mathcal{O})$ 中的 $M$ 有
$R\epsilon_*(\epsilon^*M) = M$。条件 (2) 恰好是得到此结论所需的条件。
具体地，选取一个各项平坦、表示 $M$ 的
$\mathcal{O}$-模 K-平坦复形 $\mathcal{M}^\bullet$。
于是，$\epsilon^*M$ 由 $\mathcal{M}^\bullet$ 的
$\tau$-层化 $(\mathcal{M}^\bullet)^\#$ 表示。
设 $U \in \Ob(\mathcal{C})$。由 Leray 得
$$
R\Gamma(U, R\epsilon_*(\epsilon^*M)) =
R\Gamma((\mathcal{C}/U)_\tau, (\mathcal{M}^\bullet)^\#|_{\mathcal{C}/U}) =
R\Gamma((\mathcal{C}/U)_\tau, (\mathcal{M}^\bullet|_{\mathcal{C}/U})^\#)
$$
最后一个等式成立是因为层化与限制到 $\mathcal{C}/U$ 交换。
照例，以 $\mathcal{O}_U$ 表示 $\mathcal{O}$ 到
$\mathcal{C}/U$ 的限制。考虑 $\mathcal{O}_U$-模复形的映射
$$
\mathcal{M}^\bullet(U) \otimes_{\mathcal{O}(U)} \mathcal{O}_U
\longrightarrow
\mathcal{M}^\bullet|_{\mathcal{C}/U}
$$
（在 $\tau'$-拓扑中）。由我们对 $\mathcal{M}^\bullet$ 的选取，
复形 $\mathcal{M}^\bullet(U)$ 是 $\mathcal{O}(U)$-模的 K-平坦复形；
见引理 \ref{sites-cohomology-lemma-pullback-K-flat}，并利用 $U$ 到
$\mathcal{C}$ 的包含定义环化拓扑斯态射
$(\Sh(pt), \mathcal{O}(U)) \to (\Sh(\mathcal{C}_{\tau'}), \mathcal{O})$.
由于 $M$ 属于 $\mathit{QC}(\mathcal{O})$，可知所显示的箭头是拟同构。
又因层化正合，层化后仍然如此。因此
$$
R\Gamma(U, R\epsilon_*(\epsilon^*M)) =
R\Gamma((\mathcal{C}/U)_\tau,
(M^\bullet \otimes_{\mathcal{O}(U)} \mathcal{O}_U)^\#)
$$
，而假设 (2) 表明它等于
$R\Gamma(U, M) = \mathcal{M}^\bullet(U)$，正合所需。
\end{proof}

\begin{lemma}
\label{sites-cohomology-lemma-QC-hom-out-of-plus-topology}
采用引理 \ref{sites-cohomology-lemma-cartesion-plus-topology} 的记号和假设。
假设 $K$ 是 $\mathit{QC}(\mathcal{O})$ 的一个对象，而
$M$ 是 $D(\mathcal{O}_\tau)$ 中的任意对象。
对 $\mathcal{C}$ 的每个对象 $U$，都有
$$
\Hom_{D((\mathcal{O}_U)_\tau)}(\epsilon^*K|_U, M|_U) =
R\Hom_{\mathcal{O}(U)}(R\Gamma(U, K), R\Gamma(U, M))
$$
\end{lemma}

\begin{proof}
由伴随性，有
$$
\Hom_{D((\mathcal{O}_U)_\tau)}(\epsilon^*K|_U, M|_U) =
\Hom_{D((\mathcal{O}_U)_{\tau'})}(K|_U, R\epsilon_*M|_U)
$$
故由引理 \ref{sites-cohomology-lemma-QC-hom-out-of} 以及下式得到结论：
$$
R\Gamma(U, M) = R\Gamma(U, R\epsilon_*M)
$$
，其中使用了 Leray。
\end{proof}







\section{严格完美复形}
\label{sites-cohomology-section-strictly-perfect}

\noindent
本节是《上同调》第 \ref{cohomology-section-strictly-perfect} 节的对应版本。

\begin{definition}
\label{sites-cohomology-definition-strictly-perfect}
设 $(\mathcal{C}, \mathcal{O})$ 是一个环化站点，
$\mathcal{E}^\bullet$ 是 $\mathcal{O}$-模复形。
如果除有限多个 $i$ 外 $\mathcal{E}^i$ 均为零，并且对每个 $i$，
$\mathcal{E}^i$ 都是某个有限自由 $\mathcal{O}$-模的直和项，
则称 $\mathcal{E}^\bullet$ 是{\it 严格完美的}。
\end{definition}

\noindent
设 $U$ 是 $\mathcal{C}$ 的一个对象。我们常说
“设 $\mathcal{E}^\bullet$ 是 $\mathcal{O}_U$-模的严格完美复形”，
其含义是 $\mathcal{E}^\bullet$ 是环化站点
$(\mathcal{C}/U, \mathcal{O}_U)$ 上模的严格完美复形，见
《站点上的模》定义 \ref{sites-modules-definition-localize-ringed-site}。

\begin{lemma}
\label{sites-cohomology-lemma-cone}
严格完美复形之间态射的锥是严格完美的。
\end{lemma}

\begin{proof}
由定义立即可得。
\end{proof}

\begin{lemma}
\label{sites-cohomology-lemma-tensor}
与两个严格完美复形的张量积相联系的全复形是严格完美的。
\end{lemma}

\begin{proof}
从略。
\end{proof}

\begin{lemma}
\label{sites-cohomology-lemma-strictly-perfect-pullback}
设 $(f, f^\sharp) : (\mathcal{C}, \mathcal{O}_\mathcal{C}) \to
(\mathcal{D}, \mathcal{O}_\mathcal{D})$
是环化拓扑斯的态射。如果 $\mathcal{F}^\bullet$ 是
$\mathcal{O}_\mathcal{D}$-模的严格完美复形，那么
$f^*\mathcal{F}^\bullet$ 是
$\mathcal{O}_\mathcal{C}$-模的严格完美复形。
\end{lemma}

\begin{proof}
我们已在《站点上的模》引理
\ref{sites-modules-lemma-global-pullback} 中看到，有限自由模的拉回仍为有限自由模。
函子 $f^*$ 是加性函子，因而保持直和项。引理得证。
\end{proof}

\begin{lemma}
\label{sites-cohomology-lemma-local-lift-map}
设 $(\mathcal{C}, \mathcal{O})$ 是一个环化站点，$U$ 是
$\mathcal{C}$ 的一个对象。给定 $\mathcal{O}_U$-模的实线图
$$
\xymatrix{
\mathcal{E} \ar@{..>}[dr] \ar[r] & \mathcal{F} \\
& \mathcal{G} \ar[u]_p
}
$$
，其中 $\mathcal{E}$ 是某个有限自由 $\mathcal{O}_U$-模的直和项，
且 $p$ 为满射，则存在覆盖 $\{U_i \to U\}$，使得在每个 $U_i$ 上
都存在令该图交换的虚线箭头。
\end{lemma}

\begin{proof}
可设对某个 $n$ 有
$\mathcal{E} = \mathcal{O}_U^{\oplus n}$。
此时，找到虚线箭头等价于提升 $\Gamma(U, \mathcal{F})$ 中各基元素的像。
由层的满射之刻画，这在局部可以做到
（《站点》第 \ref{sites-section-sheaves-injective} 节）。
\end{proof}

\begin{lemma}
\label{sites-cohomology-lemma-local-homotopy}
设 $(\mathcal{C}, \mathcal{O})$ 是一个环化站点，$U$ 是
$\mathcal{C}$ 的一个对象。
\begin{enumerate}
\item 设 $\alpha : \mathcal{E}^\bullet \to \mathcal{F}^\bullet$
是 $\mathcal{O}_U$-模复形的态射，其中 $\mathcal{E}^\bullet$
严格完美，而 $\mathcal{F}^\bullet$ 无环。那么，存在覆盖
$\{U_i \to U\}$，使得每个 $\alpha|_{U_i}$ 都同伦于零。
\item 设 $\alpha : \mathcal{E}^\bullet \to \mathcal{F}^\bullet$
是 $\mathcal{O}_U$-模复形的态射，其中 $\mathcal{E}^\bullet$
严格完美，对 $i < a$ 有 $\mathcal{E}^i = 0$，并且对
$i \geq a$ 有 $H^i(\mathcal{F}^\bullet) = 0$。
那么，存在覆盖 $\{U_i \to U\}$，使得每个
$\alpha|_{U_i}$ 都同伦于零。
\end{enumerate}
\end{lemma}

\begin{proof}
第一条由第二条得出，故只证明 (2)。
我们对复形 $\mathcal{E}^\bullet$ 的长度作归纳。
如果对某个有限自由 $\mathcal{O}$-模的直和项 $\mathcal{E}$
以及整数 $n \geq a$，有
$\mathcal{E}^\bullet \cong \mathcal{E}[-n]$，则结论由
引理 \ref{sites-cohomology-lemma-local-lift-map} 以及下述事实得出：
$\mathcal{F}^{n - 1} \to \Ker(\mathcal{F}^n \to \mathcal{F}^{n + 1})$
因假设 $H^n(\mathcal{F}^\bullet)$ 为零而是满射。
若除 $i \in [a,b]$ 外 $\mathcal{E}^i$ 均为零，则有复形的分裂正合列
$$
0 \to \mathcal{E}^b[-b] \to \mathcal{E}^\bullet \to
\sigma_{\leq b - 1}\mathcal{E}^\bullet \to 0
$$
，它在 $K(\mathcal{O}_U)$ 中确定一个有区别三角，因而由三角范畴公理得到正合列
$$
\Hom_{K(\mathcal{O}_U)}(
\sigma_{\leq b - 1}\mathcal{E}^\bullet, \mathcal{F}^\bullet)
\to
\Hom_{K(\mathcal{O}_U)}(\mathcal{E}^\bullet, \mathcal{F}^\bullet)
\to
\Hom_{K(\mathcal{O}_U)}(\mathcal{E}^b[-b], \mathcal{F}^\bullet)
$$
。由上面的论证，复合
$\mathcal{E}^b[-b] \to \mathcal{F}^\bullet$
在 $U$ 的某个覆盖的各成员上同伦于零。因此，可设我们的映射来自
上面所显示正合列左端的一个元素。由归纳假设，该元素在
$U$ 的某个覆盖的各成员上为零。
\end{proof}

\begin{lemma}
\label{sites-cohomology-lemma-lift-through-quasi-isomorphism}
设 $(\mathcal{C}, \mathcal{O})$ 是一个环化站点，$U$ 是
$\mathcal{C}$ 的一个对象。给定 $\mathcal{O}_U$-模复形的实线图
$$
\xymatrix{
\mathcal{E}^\bullet \ar@{..>}[dr] \ar[r]_\alpha & \mathcal{F}^\bullet \\
& \mathcal{G}^\bullet \ar[u]_f
}
$$
，其中 $\mathcal{E}^\bullet$ 严格完美，对 $j < a$ 有
$\mathcal{E}^j = 0$，并且 $H^j(f)$ 对 $j > a$ 是同构、
对 $j = a$ 是满射，则存在覆盖 $\{U_i \to U\}$，并且对每个 $i$，
在 $U_i$ 上都存在使该图交换至同伦的虚线箭头。
\end{lemma}

\begin{proof}
关于 $f$ 的假设蕴含锥 $C(f)^\bullet$ 在次数 $\geq a$ 中的上同调层为零。
因此，引理 \ref{sites-cohomology-lemma-local-homotopy} 保证存在覆盖
$\{U_i \to U\}$，使得复合
$\mathcal{E}^\bullet \to \mathcal{F}^\bullet \to C(f)^\bullet$
在 $U_i$ 上同伦于零。由于
$$
\mathcal{G}^\bullet \to \mathcal{F}^\bullet \to C(f)^\bullet \to
\mathcal{G}^\bullet[1]
$$
限制为 $K(\mathcal{O}_{U_i})$ 中的有区别三角，可知
$\alpha|_{U_i}$ 能够提升至同伦为映射
$\alpha_i : \mathcal{E}^\bullet|_{U_i} \to
\mathcal{G}^\bullet|_{U_i}$，正合所需。
\end{proof}

\begin{lemma}
\label{sites-cohomology-lemma-local-actual}
设 $(\mathcal{C}, \mathcal{O})$ 是一个环化站点，$U$ 是
$\mathcal{C}$ 的一个对象。设 $\mathcal{E}^\bullet$、
$\mathcal{F}^\bullet$ 是 $\mathcal{O}_U$-模复形，其中
$\mathcal{E}^\bullet$ 严格完美。
\begin{enumerate}
\item 对任意元素
$\alpha \in \Hom_{D(\mathcal{O}_U)}(\mathcal{E}^\bullet, \mathcal{F}^\bullet)$
，存在覆盖 $\{U_i \to U\}$，使得
$\alpha|_{U_i}$ 由复形态射
$\alpha_i : \mathcal{E}^\bullet|_{U_i} \to
\mathcal{F}^\bullet|_{U_i}$ 给出。
\item 给定复形态射
$\alpha : \mathcal{E}^\bullet \to \mathcal{F}^\bullet$
，若它在群
$\Hom_{D(\mathcal{O}_U)}(\mathcal{E}^\bullet, \mathcal{F}^\bullet)$
中的像为零，则存在覆盖 $\{U_i \to U\}$，使得
$\alpha|_{U_i}$ 同伦于零。
\end{enumerate}
\end{lemma}

\begin{proof}
证明 (1)。由导出范畴的构造，可以找到拟同构
$f : \mathcal{F}^\bullet \to \mathcal{G}^\bullet$ 和复形映射
$\beta : \mathcal{E}^\bullet \to \mathcal{G}^\bullet$，
使得 $\alpha = f^{-1}\beta$。于是结论由
引理 \ref{sites-cohomology-lemma-lift-through-quasi-isomorphism} 得出。
(2) 的证明从略。
\end{proof}

\begin{lemma}
\label{sites-cohomology-lemma-Rhom-strictly-perfect}
设 $(\mathcal{C}, \mathcal{O})$ 是一个环化站点。
设 $\mathcal{E}^\bullet$、$\mathcal{F}^\bullet$ 是
$\mathcal{O}$-模复形，其中 $\mathcal{E}^\bullet$ 严格完美。
那么，内部 Hom
$R\SheafHom(\mathcal{E}^\bullet, \mathcal{F}^\bullet)$
由复形 $\mathcal{H}^\bullet$ 表示，其各项为
$$
\mathcal{H}^n =
\bigoplus\nolimits_{n = p + q}
\SheafHom_\mathcal{O}(\mathcal{E}^{-q}, \mathcal{F}^p)
$$
，微分如第 \ref{sites-cohomology-section-internal-hom} 节所述。
\end{lemma}

\begin{proof}
选取到 K-内射复形的拟同构
$\mathcal{F}^\bullet \to \mathcal{I}^\bullet$。
设 $(\mathcal{H}')^\bullet$ 是各项如下的复形：
$$
(\mathcal{H}')^n =
\prod\nolimits_{n = p + q}
\SheafHom_\mathcal{O}(\mathcal{L}^{-q}, \mathcal{I}^p)
$$
由第 \ref{sites-cohomology-section-internal-hom} 节中的构造，它表示
$R\SheafHom(\mathcal{E}^\bullet, \mathcal{F}^\bullet)$。
只需证明映射
$$
\mathcal{H}^\bullet \longrightarrow (\mathcal{H}')^\bullet
$$
是拟同构。给定 $\mathcal{C}$ 的对象 $U$，直接检验可得
$$
H^0(\mathcal{H}^\bullet(U)) =
\Hom_{K(\mathcal{O}_U)}(\mathcal{E}^\bullet|_U, \mathcal{I}^\bullet|_U)
\to
H^0((\mathcal{H}')^\bullet(U)) =
\Hom_{D(\mathcal{O}_U)}(\mathcal{E}^\bullet|_U, \mathcal{I}^\bullet|_U)
$$
由引理 \ref{sites-cohomology-lemma-local-actual}，
$U \mapsto H^0(\mathcal{H}^\bullet(U))$ 的层化等于
$U \mapsto H^0((\mathcal{H}')^\bullet(U))$ 的层化。
对其他上同调层可作类似论证。因此，
$\mathcal{H}^\bullet$ 与 $(\mathcal{H}')^\bullet$ 拟同构，引理得证。
\end{proof}

\begin{lemma}
\label{sites-cohomology-lemma-Rhom-complex-of-direct-summands-finite-free}
设 $(\mathcal{C}, \mathcal{O})$ 是一个环化站点。
设 $\mathcal{E}^\bullet$、$\mathcal{F}^\bullet$ 是
$\mathcal{O}$-模复形，并且
\begin{enumerate}
\item 对 $n \ll 0$ 有 $\mathcal{F}^n = 0$；
\item 对 $n \gg 0$ 有 $\mathcal{E}^n = 0$；
\item $\mathcal{E}^n$ 同构于某个有限自由 $\mathcal{O}$-模的直和项。
\end{enumerate}
那么，内部 Hom
$R\SheafHom(\mathcal{E}^\bullet, \mathcal{F}^\bullet)$
由复形 $\mathcal{H}^\bullet$ 表示，其各项为
$$
\mathcal{H}^n =
\bigoplus\nolimits_{n = p + q}
\SheafHom_\mathcal{O}(\mathcal{E}^{-q}, \mathcal{F}^p)
$$
，微分如第 \ref{sites-cohomology-section-internal-hom} 节所述。
\end{lemma}

\begin{proof}
选取拟同构 $\mathcal{F}^\bullet \to \mathcal{I}^\bullet$，
其中 $\mathcal{I}^\bullet$ 是内射对象的下有界复形。
注意 $\mathcal{I}^\bullet$ 是 K-内射的
（《导出范畴》引理
\ref{derived-lemma-bounded-below-injectives-K-injective}）。
因此，第 \ref{sites-cohomology-section-internal-hom} 节中的构造表明，
$R\SheafHom(\mathcal{E}^\bullet, \mathcal{F}^\bullet)$
由复形 $(\mathcal{H}')^\bullet$ 表示，其各项为
$$
(\mathcal{H}')^n =
\prod\nolimits_{n = p + q}
\SheafHom_\mathcal{O}(\mathcal{E}^{-q}, \mathcal{I}^p) =
\bigoplus\nolimits_{n = p + q}
\SheafHom_\mathcal{O}(\mathcal{E}^{-q}, \mathcal{I}^p)
$$
（等号成立是因为只有有限多个非零项）。
注意 $\mathcal{H}^\bullet$ 是与各项为
$\SheafHom_\mathcal{O}(\mathcal{E}^{-q}, \mathcal{F}^p)$
的双复形相联系的全复形；$(\mathcal{H}')^\bullet$ 也类似。
自然映射 $(\mathcal{H}')^\bullet \to \mathcal{H}^\bullet$
来自双复形的一个映射。
因此，为证明该映射是拟同构，可以使用双复形的谱序列
（《同调》引理 \ref{homology-lemma-first-quadrant-ss}）
$$
{}'E_1^{p, q} =
H^p(\SheafHom_\mathcal{O}(\mathcal{E}^{-q}, \mathcal{F}^\bullet))
$$
，它收敛到 $H^{p + q}(\mathcal{H}^\bullet)$；
对 $(\mathcal{H}')^\bullet$ 也有类似谱序列。
为完成引理的证明，只需证明：每当 $\mathcal{E}$ 是某个有限自由
$\mathcal{O}$-模的直和项时，
$\mathcal{F}^\bullet \to \mathcal{I}^\bullet$ 在上同调层上诱导同构
$$
H^p(\SheafHom_\mathcal{O}(\mathcal{E}, \mathcal{F}^\bullet))
\longrightarrow
H^p(\SheafHom_\mathcal{O}(\mathcal{E}, \mathcal{I}^\bullet))
$$
。当 $\mathcal{E}$ 有限自由时这是显然的，故一般结论也随之成立。
\end{proof}









\section{伪凝聚模}
\label{sites-cohomology-section-pseudo-coherent}

\noindent
本节讨论伪凝聚复形。

\begin{definition}
\label{sites-cohomology-definition-pseudo-coherent}
设 $(\mathcal{C}, \mathcal{O})$ 是一个环化站点，
$\mathcal{E}^\bullet$ 是 $\mathcal{O}$-模复形，并设
$m \in \mathbf{Z}$。
\begin{enumerate}
\item 如果对 $\mathcal{C}$ 的每个对象 $U$，都存在覆盖
$\{U_i \to U\}$，并且对每个 $i$ 都存在复形态射
$\alpha_i : \mathcal{E}_i^\bullet \to \mathcal{E}^\bullet|_{U_i}$，
其中 $\mathcal{E}_i^\bullet$ 是 $\mathcal{O}_{U_i}$-模的严格完美复形，
$H^j(\alpha_i)$ 对 $j > m$ 是同构，而
$H^m(\alpha_i)$ 是满射，则称 $\mathcal{E}^\bullet$
是{\it $m$-伪凝聚的}。
\item 如果 $\mathcal{E}^\bullet$ 对所有 $m$ 都是 $m$-伪凝聚的，
则称它是{\it 伪凝聚的}。
\item 当且仅当 $D(\mathcal{O})$ 的对象 $E$ 能由一个
$m$-伪凝聚（相应地，伪凝聚）的 $\mathcal{O}$-模复形表示时，
称 $E$ 是{\it $m$-伪凝聚的}（相应地，{\it 伪凝聚的}）。
\end{enumerate}
\end{definition}

\noindent
如果 $\mathcal{C}$ 有拟紧的终对象 $X$
（例如，$X$ 的每个覆盖都可由有限覆盖加细），那么
$D(\mathcal{O})$ 的 $m$-伪凝聚对象属于
$D^-(\mathcal{O})$。但一般而言未必如此。

\begin{lemma}
\label{sites-cohomology-lemma-pseudo-coherent-independent-representative}
设 $(\mathcal{C}, \mathcal{O})$ 是一个环化站点，
$E$ 是 $D(\mathcal{O})$ 的一个对象。
\begin{enumerate}
\item 如果 $\mathcal{C}$ 有终对象 $X$，并且存在覆盖
$\{U_i \to X\}$、$\mathcal{O}_{U_i}$-模的严格完美复形
$\mathcal{E}_i^\bullet$，以及 $D(\mathcal{O}_{U_i})$ 中的映射
$\alpha_i : \mathcal{E}_i^\bullet \to E|_{U_i}$，使得
$H^j(\alpha_i)$ 对 $j > m$ 是同构，且
$H^m(\alpha_i)$ 是满射，那么 $E$ 是 $m$-伪凝聚的。
\item 如果 $E$ 是 $m$-伪凝聚的，那么表示 $E$ 的任意
$\mathcal{O}$-模复形都是 $m$-伪凝聚的。
\item 如果对 $\mathcal{C}$ 的每个对象 $U$，都存在覆盖
$\{U_i \to U\}$，使得 $E|_{U_i}$ 是 $m$-伪凝聚的，
那么 $E$ 是 $m$-伪凝聚的。
\end{enumerate}
\end{lemma}

\begin{proof}
设 $\mathcal{F}^\bullet$ 是表示 $E$ 的任意复形，并设
$X$、$\{U_i \to X\}$ 以及
$\alpha_i : \mathcal{E}_i \to E|_{U_i}$ 如 (1)。
我们将证明复形 $\mathcal{F}^\bullet$ 是 $m$-伪凝聚的；
这将证明 (1)，并在 $\mathcal{C}$ 有终对象时证明 (2)。
由引理 \ref{sites-cohomology-lemma-local-actual}，加细覆盖
$\{U_i \to X\}$ 后，可以用复形映射
$\alpha_i : \mathcal{E}_i^\bullet \to
\mathcal{F}^\bullet|_{U_i}$ 表示这些映射。
根据假设，$H^j(\alpha_i)$ 对 $j > m$ 是同构，而
$H^m(\alpha_i)$ 是满射，故 $\mathcal{F}^\bullet$ 是 $m$-伪凝聚的。

\medskip\noindent
证明 (2)。由上可知，对 $\mathcal{C}$ 的所有对象 $U$，
$\mathcal{F}^\bullet|_U$ 作为 $\mathcal{O}_U$-模复形是
$m$-伪凝聚的。由定义形式地推出
$\mathcal{F}^\bullet$ 是 $m$-伪凝聚的。

\medskip\noindent
证明 (3)。由定义以及《站点》定义
\ref{sites-definition-site} 的 (2) 得出。
\end{proof}

\begin{lemma}
\label{sites-cohomology-lemma-pseudo-coherent-pullback}
设 $(f, f^\sharp) : (\mathcal{C}, \mathcal{O}_\mathcal{C}) \to
(\mathcal{D}, \mathcal{O}_\mathcal{D})$
是环化站点的态射。设 $E$ 是
$D(\mathcal{O}_\mathcal{C})$ 的一个对象。如果 $E$ 是
$m$-伪凝聚的，那么 $Lf^*E$ 是 $m$-伪凝聚的。
\end{lemma}

\begin{proof}
设 $f$ 由函子 $u : \mathcal{D} \to \mathcal{C}$ 给出。
设 $U$ 是 $\mathcal{C}$ 的一个对象。由《站点》引理
\ref{sites-lemma-morphism-of-sites-covering}，
可以找到覆盖 $\{U_i \to U\}$，并且对每个 $i$，都有从
$U_i$ 到某个 $\mathcal{D}$ 中对象 $V_i$ 的像 $u(V_i)$ 的态射。
由引理 \ref{sites-cohomology-lemma-pseudo-coherent-independent-representative}，
只需证明 $Lf^*E|_{U_i}$ 是 $m$-伪凝聚的。
为此，只需证明 $Lf^*E|_{u(V_i)}$ 是 $m$-伪凝聚的，
因为 $Lf^*E|_{U_i}$ 是 $Lf^*E|_{u(V_i)}$ 到
$\mathcal{C}/U_i$ 的限制（使用《站点上的模》引理
\ref{sites-modules-lemma-relocalize}）。
由《站点上的模》引理
\ref{sites-modules-lemma-localize-morphism-ringed-sites} 中的交换图，
只需对环化站点态射
$(\mathcal{C}/u(V_i), \mathcal{O}_{u(V_i)}) \to
(\mathcal{D}/V_i, \mathcal{O}_{V_i})$ 证明该引理。
因此，可设 $\mathcal{D}$ 有终对象 $Y$，且
$X = u(Y)$ 是 $\mathcal{C}$ 的终对象。

\medskip\noindent
设 $\{V_i \to Y\}$ 是一个覆盖，使得对每个 $i$，都存在
$\mathcal{O}_{V_i}$-模的严格完美复形
$\mathcal{F}_i^\bullet$ 以及 $D(\mathcal{O}_{V_i})$ 中的态射
$\alpha_i : \mathcal{F}_i^\bullet \to E|_{V_i}$，并且
$H^j(\alpha_i)$ 对 $j > m$ 是同构，而
$H^m(\alpha_i)$ 是满射。仿照上面的论证，只需对
$(\mathcal{C}/u(V_i), \mathcal{O}_{u(V_i)}) \to
(\mathcal{D}/V_i, \mathcal{O}_{V_i})$ 证明结论。
因此，可设存在 $\mathcal{O}_\mathcal{D}$-模的严格完美复形
$\mathcal{F}^\bullet$ 以及 $D(\mathcal{O}_\mathcal{D})$ 中的态射
$\alpha : \mathcal{F}^\bullet \to E$，使得
$H^j(\alpha)$ 对 $j > m$ 是同构，而
$H^m(\alpha)$ 是满射。在此情形下，选取有区别三角
$$
\mathcal{F}^\bullet \to E \to C \to \mathcal{F}^\bullet[1]
$$
关于 $\alpha$ 的假设恰好意味着对所有 $j \geq m$，
上同调层 $H^j(C)$ 都为零。应用 $Lf^*$，得到有区别三角
$$
Lf^*\mathcal{F}^\bullet \to Lf^*E \to Lf^*C \to Lf^*\mathcal{F}^\bullet[1]
$$
由 $Lf^*$ 作为左导出函子的构造，可知对 $j \geq m$ 有
$H^j(Lf^*C) = 0$（使用《导出范畴》引理
\ref{derived-lemma-negative-vanishing} 的对偶）。
因此，$H^j(Lf^*\alpha)$ 对 $j > m$ 是同构，而
$H^m(Lf^*\alpha)$ 是满射。另一方面，由于
$\mathcal{F}^\bullet$ 是平坦
$\mathcal{O}_\mathcal{D}$-模的上有界复形，有
$Lf^*\mathcal{F}^\bullet = f^*\mathcal{F}^\bullet$。
应用引理 \ref{sites-cohomology-lemma-strictly-perfect-pullback} 即得结论。
\end{proof}

\begin{lemma}
\label{sites-cohomology-lemma-cone-pseudo-coherent}
设 $(\mathcal{C}, \mathcal{O})$ 是一个环化站点，
$m \in \mathbf{Z}$。设
$(K, L, M, f, g, h)$ 是 $D(\mathcal{O})$ 中的有区别三角。
\begin{enumerate}
\item 如果 $K$ 是 $(m + 1)$-伪凝聚的，而 $L$ 是 $m$-伪凝聚的，
那么 $M$ 是 $m$-伪凝聚的。
\item 如果 $K$ 和 $M$ 都是 $m$-伪凝聚的，那么 $L$ 是 $m$-伪凝聚的。
\item 如果 $L$ 是 $(m + 1)$-伪凝聚的，而 $M$ 是 $m$-伪凝聚的，
那么 $K$ 是 $(m + 1)$-伪凝聚的。
\end{enumerate}
\end{lemma}

\begin{proof}
证明 (1)。设 $U$ 是 $\mathcal{C}$ 的一个对象。选取覆盖
$\{U_i \to U\}$ 以及 $D(\mathcal{O}_{U_i})$ 中的映射
$\alpha_i : \mathcal{K}_i^\bullet \to K|_{U_i}$，其中
$\mathcal{K}_i^\bullet$ 严格完美，$H^j(\alpha_i)$ 对
$j > m + 1$ 是同构，并且对 $j = m + 1$ 是满射。
可用 $\sigma_{\geq m + 1}\mathcal{K}_i^\bullet$ 代替
$\mathcal{K}_i^\bullet$，从而可设对 $j < m + 1$ 有
$\mathcal{K}_i^j = 0$。加细覆盖后，可以选取
$D(\mathcal{O}_{U_i})$ 中的映射
$\beta_i : \mathcal{L}_i^\bullet \to L|_{U_i}$，
其中 $\mathcal{L}_i^\bullet$ 严格完美，并且
$H^j(\beta)$ 对 $j > m$ 是同构、对 $j = m$ 是满射。
由引理 \ref{sites-cohomology-lemma-lift-through-quasi-isomorphism}，
再次加细覆盖后，可找到复形映射
$\gamma_i : \mathcal{K}^\bullet \to \mathcal{L}^\bullet$，
使得图
$$
\xymatrix{
K|_{U_i} \ar[r] & L|_{U_i} \\
\mathcal{K}_i^\bullet \ar[u]^{\alpha_i} \ar[r]^{\gamma_i} &
\mathcal{L}_i^\bullet \ar[u]_{\beta_i}
}
$$
在 $D(\mathcal{O}_{U_i})$ 中交换（这要求用实际的复形映射表示
$\alpha_i$、$\beta_i$ 以及 $K|_{U_i} \to L|_{U_i}$；
一些细节从略）。锥 $C(\gamma_i)^\bullet$ 严格完美
（引理 \ref{sites-cohomology-lemma-cone}）。
该图的交换性蕴含存在有区别三角的态射
$$
(\mathcal{K}_i^\bullet, \mathcal{L}_i^\bullet, C(\gamma_i)^\bullet)
\longrightarrow
(K|_{U_i}, L|_{U_i}, M|_{U_i}).
$$
由上同调长正合列上诱导的映射以及《同调》引理
\ref{homology-lemma-four-lemma} 和
\ref{homology-lemma-five-lemma}，可知
$C(\gamma_i)^\bullet \to M|_{U_i}$ 在次数 $>m$ 的上同调上诱导同构，
在次数 $m$ 上诱导满射。因此，由引理
\ref{sites-cohomology-lemma-pseudo-coherent-independent-representative}，
$M$ 是 $m$-伪凝聚的。

\medskip\noindent
旋转有区别三角，即由 (1) 得到断言 (2) 和 (3)。
\end{proof}

\begin{lemma}
\label{sites-cohomology-lemma-tensor-pseudo-coherent}
设 $(\mathcal{C}, \mathcal{O})$ 是一个环化站点，
$K,L$ 是 $D(\mathcal{O})$ 的对象。
\begin{enumerate}
\item 如果 $K$ 是 $n$-伪凝聚的，且对 $i > a$ 有
$H^i(K) = 0$；而 $L$ 是 $m$-伪凝聚的，且对 $j > b$ 有
$H^j(L) = 0$，那么
$K \otimes_\mathcal{O}^\mathbf{L} L$ 是 $t$-伪凝聚的，其中
$t = \max(m + a, n + b)$。
\item 如果 $K$ 和 $L$ 都是伪凝聚的，那么
$K \otimes_\mathcal{O}^\mathbf{L} L$ 是伪凝聚的。
\end{enumerate}
\end{lemma}

\begin{proof}
证明 (1)。设 $U$ 是 $\mathcal{C}$ 的一个对象。
用某个覆盖的各成员代替 $U$，并用局部化
$\mathcal{C}/U$ 代替 $\mathcal{C}$ 后，可设存在严格完美复形
$\mathcal{K}^\bullet$、$\mathcal{L}^\bullet$ 以及映射
$\alpha : \mathcal{K}^\bullet \to K$ 和
$\beta : \mathcal{L}^\bullet \to L$，使得
$H^i(\alpha)$ 对 $i > n$ 是同构、对 $i = n$ 是满射，并且
$H^i(\beta)$ 对 $i > m$ 是同构、对 $i = m$ 是满射。
于是映射
$$
\alpha \otimes^\mathbf{L} \beta :
\text{Tot}(\mathcal{K}^\bullet \otimes_\mathcal{O} \mathcal{L}^\bullet)
\to K \otimes_\mathcal{O}^\mathbf{L} L
$$
在 $i > t$ 时的第 $i$ 个上同调层上诱导同构，并在
$i=t$ 时诱导满射。这由 Tor 谱序列得出（细节从略）。

\medskip\noindent
证明 (2)。设 $U$ 是 $\mathcal{C}$ 的一个对象。
先用某个覆盖的各成员代替 $U$，并用局部化
$\mathcal{C}/U$ 代替 $\mathcal{C}$，可归结到 $K$ 和 $L$
均上有界的情形。此时断言立即由 (1) 得出。
\end{proof}

\begin{lemma}
\label{sites-cohomology-lemma-summands-pseudo-coherent}
设 $(\mathcal{C}, \mathcal{O})$ 是一个环化站点，
$m \in \mathbf{Z}$。如果 $D(\mathcal{O})$ 中的
$K \oplus L$ 是 $m$-伪凝聚的（相应地，伪凝聚的），
那么 $K$ 和 $L$ 也都是如此。
\end{lemma}

\begin{proof}
假设 $K \oplus L$ 是 $m$-伪凝聚的。设 $U$ 是
$\mathcal{C}$ 的一个对象。用某个覆盖的各成员代替 $U$ 后，
可设 $K \oplus L \in D^-(\mathcal{O}_U)$，从而
$L \in D^-(\mathcal{O}_U)$。注意有有区别三角
$$
(K \oplus L, K \oplus L, L \oplus L[1]) =
(K, K, 0) \oplus (L, L, L \oplus L[1])
$$
，见《导出范畴》引理
\ref{derived-lemma-direct-sum-triangles}。
由引理 \ref{sites-cohomology-lemma-cone-pseudo-coherent}，
$L \oplus L[1]$ 是 $m$-伪凝聚的。因此
$L[1] \oplus L[2]$ 也是 $m$-伪凝聚的。
归纳可知 $L[n] \oplus L[n + 1]$ 是 $m$-伪凝聚的。
由于 $L$ 上有界，对充分大的 $n$，$L[n]$ 是 $m$-伪凝聚的。
于是，利用有区别三角
$$
(L[n], L[n] \oplus L[n - 1], L[n - 1])
$$
反向推导，可知 $L[n - 1], L[n - 2], \ldots, L$ 都是
$m$-伪凝聚的，正合所需。
\end{proof}

\begin{lemma}
\label{sites-cohomology-lemma-finite-cohomology}
设 $(\mathcal{C}, \mathcal{O})$ 是一个环化站点，
$K$ 是 $D(\mathcal{O})$ 的一个对象，并设
$m \in \mathbf{Z}$。
\begin{enumerate}
\item 如果 $K$ 是 $m$-伪凝聚的，且对 $i > m$ 有
$H^i(K) = 0$，那么 $H^m(K)$ 是有限型 $\mathcal{O}$-模。
\item 如果 $K$ 是 $m$-伪凝聚的，且对 $i > m + 1$ 有
$H^i(K) = 0$，那么 $H^{m + 1}(K)$ 是有限表示
$\mathcal{O}$-模。
\end{enumerate}
\end{lemma}

\begin{proof}
证明 (1)。设 $U$ 是 $\mathcal{C}$ 的一个对象。需要证明：
$H^m(K)$ 在 $U$ 的某个覆盖的各成员上可由有限多个截面生成
（见《站点上的模》定义
\ref{sites-modules-definition-site-local}）。
因此，在证明过程中，可以有限多次选取覆盖
$\{U_i \to U\}$，并用 $\mathcal{C}/U_i$ 代替
$\mathcal{C}$、用 $U_i$ 代替 $U$。特别地，由我们的定义，
可设存在严格完美复形 $\mathcal{E}^\bullet$ 和映射
$\alpha : \mathcal{E}^\bullet \to K$，它在次数 $>m$ 的上同调上诱导同构，
并在次数 $m$ 上诱导满射。只需对 $\mathcal{E}^\bullet$ 证明结论。
设 $n$ 是使 $\mathcal{E}^n \not = 0$ 的最大整数。
若 $n=m$，则 $H^m(\mathcal{E}^\bullet)$ 是
$\mathcal{E}^n$ 的商，结论显然。
若 $n>m$，则由于 $H^n(E^\bullet)=0$，
$\mathcal{E}^{n - 1} \to \mathcal{E}^n$ 是满射。
由引理 \ref{sites-cohomology-lemma-local-lift-map}，用某个覆盖的各成员代替
$U$ 后，可找到该满射的截面，并写成
$\mathcal{E}^{n - 1} = \mathcal{E}' \oplus \mathcal{E}^n$。
因而，只需对复形 $(\mathcal{E}')^\bullet$ 证明结论；
它与 $\mathcal{E}^\bullet$ 相同，但在次数 $n-1$ 处为
$\mathcal{E}'$，在次数 $n$ 处为 $0$。对 $n$ 归纳即得结论。

\medskip\noindent
证明 (2)。取 $\mathcal{C}$ 的一个对象 $U$。
如 (1) 的证明，可在 $U$ 上局部工作。
因而可设存在严格完美复形 $\mathcal{E}^\bullet$ 和映射
$\alpha : \mathcal{E}^\bullet \to K$，它在次数 $>m$ 的上同调上诱导同构，
并在次数 $m$ 上诱导满射。如 (1) 的证明，可归结到对
$i > m+1$ 有 $\mathcal{E}^i=0$ 的情形。于是
$H^{m + 1}(K) \cong H^{m + 1}(\mathcal{E}^\bullet) =
\Coker(\mathcal{E}^m \to \mathcal{E}^{m + 1})$
是有限表示的。
\end{proof}








\section{Tor 维数}
\label{sites-cohomology-section-tor}

\noindent
本节更细致地考察平坦模分辨。

\begin{definition}
\label{sites-cohomology-definition-tor-amplitude}
设 $(\mathcal{C}, \mathcal{O})$ 是一个环化站点，
$E$ 是 $D(\mathcal{O})$ 的一个对象。
设 $a,b \in \mathbf{Z}$ 且 $a \leq b$。
\begin{enumerate}
\item 如果对所有 $\mathcal{O}$-模 $\mathcal{F}$ 以及所有
$i \not \in [a,b]$，都有
$H^i(E \otimes_\mathcal{O}^\mathbf{L} \mathcal{F}) = 0$，
则称 $E$ 的 {\it Tor 振幅在 $[a,b]$ 中}。
\item 如果对某些 $a,b$，$E$ 的 Tor 振幅在 $[a,b]$ 中，
则称 $E$ 具有{\it 有限 Tor 维数}。
\item 如果对 $\mathcal{C}$ 的任意对象 $U$，都存在覆盖
$\{U_i \to U\}$，使得所有 $E|_{U_i}$ 都具有有限 Tor 维数，
则称 $E$ {\it 局部具有有限 Tor 维数}。
\end{enumerate}
如果把 $\mathcal{F}[0]$ 看作 $D(\mathcal{O})$ 的对象时，
其 Tor 振幅在 $[-d,0]$ 中，则称 $\mathcal{O}$-模
$\mathcal{F}$ 的 {\it Tor 维数 $\leq d$}。
\end{definition}

\noindent
注意，如果定义中的 $E$ 具有有限 Tor 维数，那么
$E$ 是 $D^b(\mathcal{O})$ 的对象；这可由在上述定义中取
$\mathcal{F} = \mathcal{O}$ 看出。

\begin{lemma}
\label{sites-cohomology-lemma-last-one-flat}
设 $(\mathcal{C}, \mathcal{O})$ 是一个环化站点。
设 $\mathcal{E}^\bullet$ 是平坦 $\mathcal{O}$-模的上有界复形，
其 Tor 振幅在 $[a,b]$ 中。那么
$\Coker(d_{\mathcal{E}^\bullet}^{a - 1})$ 是平坦
$\mathcal{O}$-模。
\end{lemma}

\begin{proof}
由于 $\mathcal{E}^\bullet$ 是平坦模的上有界复形，对任意
$\mathcal{O}$-模 $\mathcal{F}$ 都有
$\mathcal{E}^\bullet \otimes_\mathcal{O} \mathcal{F} =
\mathcal{E}^\bullet \otimes_\mathcal{O}^{\mathbf{L}} \mathcal{F}$
。因此，对每个 $\mathcal{O}$-模 $\mathcal{F}$，序列
$$
\mathcal{E}^{a - 2} \otimes_\mathcal{O} \mathcal{F} \to
\mathcal{E}^{a - 1} \otimes_\mathcal{O} \mathcal{F} \to
\mathcal{E}^a \otimes_\mathcal{O} \mathcal{F}
$$
在中间正合。由于
$\mathcal{E}^{a - 2} \to \mathcal{E}^{a - 1} \to \mathcal{E}^a \to
\Coker(d^{a - 1}) \to 0$
是平坦分辨，这蕴含对所有 $\mathcal{O}$-模 $\mathcal{F}$ 都有
$\text{Tor}_1^\mathcal{O}(\Coker(d^{a - 1}), \mathcal{F}) = 0$
。这意味着 $\Coker(d^{a - 1})$ 平坦，见
引理 \ref{sites-cohomology-lemma-flat-tor-zero}。
\end{proof}

\begin{lemma}
\label{sites-cohomology-lemma-tor-amplitude}
设 $(\mathcal{C}, \mathcal{O})$ 是一个环化站点，
$E$ 是 $D(\mathcal{O})$ 的一个对象。
设 $a,b \in \mathbf{Z}$ 且 $a \leq b$。下列条件等价：
\begin{enumerate}
\item $E$ 的 Tor 振幅在 $[a,b]$ 中。
\item $E$ 由平坦 $\mathcal{O}$-模复形
$\mathcal{E}^\bullet$ 表示，并且对
$i \not \in [a,b]$ 有 $\mathcal{E}^i=0$。
\end{enumerate}
\end{lemma}

\begin{proof}
如果 (2) 成立，则可按下式计算：
$E \otimes_\mathcal{O}^\mathbf{L} \mathcal{F} =
\mathcal{E}^\bullet \otimes_\mathcal{O} \mathcal{F}$
，显然 (1) 成立。

\medskip\noindent
假设 (1) 成立。可用平坦 $\mathcal{O}$-模的上有界复形
$\mathcal{K}^\bullet$ 表示 $E$，见第 \ref{sites-cohomology-section-flat} 节。
设 $n$ 是使 $\mathcal{K}^n \not = 0$ 的最大整数。
若 $n>b$，则因 $H^n(\mathcal{K}^\bullet)=0$，
$\mathcal{K}^{n - 1} \to \mathcal{K}^n$ 是满射。
由于 $\mathcal{K}^n$ 平坦，
$\Ker(\mathcal{K}^{n - 1} \to \mathcal{K}^n)$ 平坦
（《站点上的模》引理 \ref{sites-modules-lemma-flat-ses}）。
因此，可以用 $\tau_{\leq n - 1}\mathcal{K}^\bullet$
代替 $\mathcal{K}^\bullet$。于是，对 $n$ 归纳，可归结到
$K^\bullet$ 是平坦 $\mathcal{O}$-模复形，并且对 $i>b$ 有
$\mathcal{K}^i=0$ 的情形。

\medskip\noindent
令 $\mathcal{E}^\bullet = \tau_{\geq a}\mathcal{K}^\bullet$。
除 $\mathcal{E}^a$ 的平坦性外，一切均显然；而这一点立即由
引理 \ref{sites-cohomology-lemma-last-one-flat} 和定义得出。
\end{proof}

\begin{lemma}
\label{sites-cohomology-lemma-bounded-below-tor-amplitude}
设 $(\mathcal{C}, \mathcal{O})$ 是一个环化站点，
$E$ 是 $D(\mathcal{O})$ 的一个对象。设
$a \in \mathbf{Z}$。下列条件等价：
\begin{enumerate}
\item $E$ 的 Tor 振幅在 $[a,\infty]$ 中。
\item $E$ 可由平坦 $\mathcal{O}$-模的 K-平坦复形
$\mathcal{E}^\bullet$ 表示，并且对
$i \not \in [a,\infty]$ 有 $\mathcal{E}^i=0$。
\end{enumerate}
此外，可以选取 $\mathcal{E}^\bullet$，使得沿任意环化站点态射的拉回
都是各项平坦的 K-平坦复形。
\end{lemma}

\begin{proof}
(2) $\Rightarrow$ (1) 是立即的。假设 (1) 成立。
首先选取表示 $E$ 且各项平坦的 K-平坦复形
$\mathcal{K}^\bullet$，见引理 \ref{sites-cohomology-lemma-K-flat-resolution}。
对任意 $\mathcal{O}$-模 $\mathcal{M}$，下列复形在中间的上同调
$$
\mathcal{K}^{n - 1} \otimes_\mathcal{O} \mathcal{M} \to
\mathcal{K}^n \otimes_\mathcal{O} \mathcal{M} \to
\mathcal{K}^{n + 1} \otimes_\mathcal{O} \mathcal{M}
$$
计算 $H^n(E \otimes_\mathcal{O}^\mathbf{L} \mathcal{M})$。
当 $n<a$ 时它总为零。因此，对复形
$\ldots \to \mathcal{K}^{a - 1} \to \mathcal{K}^a \to
\mathcal{K}^{a + 1}$ 应用引理 \ref{sites-cohomology-lemma-last-one-flat}，
可知
$\mathcal{N} = \Coker(\mathcal{K}^{a - 1} \to \mathcal{K}^a)$
是平坦 $\mathcal{O}$-模。令
$$
\mathcal{E}^\bullet = \tau_{\geq a}\mathcal{K}^\bullet =
(\ldots \to 0 \to \mathcal{N} \to \mathcal{K}^{a + 1} \to \ldots )
$$
映射 $\mathcal{K}^\bullet \to \mathcal{E}^\bullet$ 的核
$\mathcal{L}^\bullet$ 是复形
$$
\mathcal{L}^\bullet = (\ldots \to \mathcal{K}^{a - 1} \to
\mathcal{I} \to 0 \to \ldots)
$$
，其中 $\mathcal{I} \subset \mathcal{K}^a$ 是
$\mathcal{K}^{a - 1} \to \mathcal{K}^a$ 的像。
由于有短正合列
$0 \to \mathcal{I} \to \mathcal{K}^a \to \mathcal{N} \to 0$
，可知 $\mathcal{I}$ 是平坦 $\mathcal{O}$-模。
因此，$\mathcal{L}^\bullet$ 是平坦模的上有界复形，
故由引理 \ref{sites-cohomology-lemma-bounded-flat-K-flat} 可知它 K-平坦。
随后由引理 \ref{sites-cohomology-lemma-K-flat-two-out-of-three-ses} 可知
$\mathcal{E}^\bullet$ K-平坦。

\medskip\noindent
证明最后一个断言。设
$f : (\mathcal{C}', \mathcal{O}') \to (\mathcal{C}, \mathcal{O})$
是环化站点态射。由引理 \ref{sites-cohomology-lemma-pullback-K-flat}，
复形 $f^*\mathcal{K}^\bullet$ K-平坦且各项平坦。
复形 $f^*\mathcal{L}^\bullet$ 是平坦 $\mathcal{O}'$-模的上有界复形，
故 K-平坦。我们有 $\mathcal{O}'$-模复形的短正合列
$$
0 \to f^*\mathcal{L}^\bullet \to f^*\mathcal{K}^\bullet \to
f^*\mathcal{E}^\bullet \to 0
$$
，因为平坦模的短正合列
$0 \to \mathcal{I} \to \mathcal{K}^a \to \mathcal{N} \to 0$
拉回后仍为短正合列。由引理
\ref{sites-cohomology-lemma-K-flat-two-out-of-three-ses}，
复形 $f^*\mathcal{E}^\bullet$ K-平坦，证明完毕。
\end{proof}

\begin{lemma}
\label{sites-cohomology-lemma-tor-amplitude-pullback}
设 $(f, f^\sharp) : (\mathcal{C}, \mathcal{O}_\mathcal{C}) \to
(\mathcal{D}, \mathcal{O}_\mathcal{D})$
是环化站点的态射。设 $E$ 是
$D(\mathcal{O}_\mathcal{D})$ 的一个对象。
如果 $E$ 的 Tor 振幅在 $[a,b]$ 中，那么
$Lf^*E$ 的 Tor 振幅在 $[a,b]$ 中。
\end{lemma}

\begin{proof}
假设 $E$ 的 Tor 振幅在 $[a,b]$ 中。由
引理 \ref{sites-cohomology-lemma-tor-amplitude}，可用平坦
$\mathcal{O}$-模的复形 $\mathcal{E}^\bullet$ 表示 $E$，
其中对 $i \not \in [a,b]$ 有 $\mathcal{E}^i=0$。
于是 $Lf^*E$ 由 $f^*\mathcal{E}^\bullet$ 表示。
由《站点上的模》引理 \ref{sites-modules-lemma-pullback-flat}，
模 $f^*\mathcal{E}^i$ 是平坦的。因此，由
引理 \ref{sites-cohomology-lemma-tor-amplitude} 可知
$Lf^*E$ 的 Tor 振幅在 $[a,b]$ 中。
\end{proof}

\begin{lemma}
\label{sites-cohomology-lemma-cone-tor-amplitude}
设 $(\mathcal{C}, \mathcal{O})$ 是一个环化站点。
设 $(K,L,M,f,g,h)$ 是 $D(\mathcal{O})$ 中的有区别三角，
并设 $a,b \in \mathbf{Z}$。
\begin{enumerate}
\item 如果 $K$ 的 Tor 振幅在 $[a+1,b+1]$ 中，而
$L$ 的 Tor 振幅在 $[a,b]$ 中，那么
$M$ 的 Tor 振幅在 $[a,b]$ 中。
\item 如果 $K$ 和 $M$ 的 Tor 振幅都在 $[a,b]$ 中，那么
$L$ 的 Tor 振幅在 $[a,b]$ 中。
\item 如果 $L$ 的 Tor 振幅在 $[a+1,b+1]$ 中，而
$M$ 的 Tor 振幅在 $[a,b]$ 中，那么
$K$ 的 Tor 振幅在 $[a+1,b+1]$ 中。
\end{enumerate}
\end{lemma}

\begin{proof}
从略。提示：这直接由与有区别三角相联系的上同调长正合列以及
$- \otimes_\mathcal{O}^{\mathbf{L}} \mathcal{F}$
保持有区别三角这一事实得出。最容易证明的是 (2)，其余各条由平移得出。
\end{proof}

\begin{lemma}
\label{sites-cohomology-lemma-tensor-tor-amplitude}
设 $(\mathcal{C}, \mathcal{O})$ 是一个环化站点，
$K,L$ 是 $D(\mathcal{O})$ 的对象。如果 $K$ 的 Tor 振幅在
$[a,b]$ 中，而 $L$ 的 Tor 振幅在 $[c,d]$ 中，那么
$K \otimes_\mathcal{O}^\mathbf{L} L$ 的 Tor 振幅在
$[a+c,b+d]$ 中。
\end{lemma}

\begin{proof}
从略。提示：使用 Tor 谱序列。
\end{proof}

\begin{lemma}
\label{sites-cohomology-lemma-summands-tor-amplitude}
设 $(\mathcal{C}, \mathcal{O})$ 是一个环化站点，
$a,b \in \mathbf{Z}$。对 $D(\mathcal{O})$ 的对象 $K,L$，
如果 $K \oplus L$ 的 Tor 振幅在 $[a,b]$ 中，那么
$K$ 和 $L$ 也都如此。
\end{lemma}

\begin{proof}
由 Tor 函子的加性立即可得。
\end{proof}

\begin{lemma}
\label{sites-cohomology-lemma-bounded}
设 $(\mathcal{C}, \mathcal{O})$ 是一个环化站点，
$\mathcal{I} \subset \mathcal{O}$ 是理想层，
$K$ 是 $D(\mathcal{O})$ 的一个对象。
\begin{enumerate}
\item 如果
$K \otimes_\mathcal{O}^\mathbf{L} \mathcal{O}/\mathcal{I}$
上有界，那么对所有 $n$，
$K \otimes_\mathcal{O}^\mathbf{L} \mathcal{O}/\mathcal{I}^n$
一致上有界。
\item 如果
$K \otimes_\mathcal{O}^\mathbf{L} \mathcal{O}/\mathcal{I}$
作为 $D(\mathcal{O}/\mathcal{I})$ 的对象，其 Tor 振幅在
$[a,b]$ 中，那么对所有 $n$，
$K \otimes_\mathcal{O}^\mathbf{L} \mathcal{O}/\mathcal{I}^n$
作为 $D(\mathcal{O}/\mathcal{I}^n)$ 的对象，其 Tor 振幅也在
$[a,b]$ 中。
\end{enumerate}
\end{lemma}

\begin{proof}
证明 (1)。假设
$K \otimes_\mathcal{O}^\mathbf{L} \mathcal{O}/\mathcal{I}$
上有界，例如
$H^i(K \otimes_\mathcal{O}^\mathbf{L} \mathcal{O}/\mathcal{I}) = 0$
对 $i>b$ 成立。注意有有区别三角
$$
K \otimes_\mathcal{O}^\mathbf{L}
\mathcal{I}^n/\mathcal{I}^{n + 1} \to
K \otimes_\mathcal{O}^\mathbf{L}
\mathcal{O}/\mathcal{I}^{n + 1} \to
K \otimes_\mathcal{O}^\mathbf{L}
\mathcal{O}/\mathcal{I}^n \to
K \otimes_\mathcal{O}^\mathbf{L}
\mathcal{I}^n/\mathcal{I}^{n + 1}[1]
$$
，并且
$$
K \otimes_\mathcal{O}^\mathbf{L}
\mathcal{I}^n/\mathcal{I}^{n + 1} =
\left(
K \otimes_\mathcal{O}^\mathbf{L}
\mathcal{O}/\mathcal{I}\right)
\otimes_{\mathcal{O}/\mathcal{I}}^\mathbf{L}
\mathcal{I}^n/\mathcal{I}^{n + 1}
$$
。归纳可知，对所有 $n$ 以及所有 $i>b$，
$H^i(K \otimes_\mathcal{O}^\mathbf{L} \mathcal{O}/\mathcal{I}^n) = 0$
。

\medskip\noindent
证明 (2)。假设
$K \otimes_\mathcal{O}^\mathbf{L} \mathcal{O}/\mathcal{I}$
作为 $D(\mathcal{O}/\mathcal{I})$ 的对象，其 Tor 振幅在
$[a,b]$ 中。设 $\mathcal{F}$ 是
$\mathcal{O}/\mathcal{I}^n$-模层。则有有限滤过
$$
0 \subset \mathcal{I}^{n - 1}\mathcal{F} \subset \ldots
\subset \mathcal{I}\mathcal{F} \subset \mathcal{F}
$$
，其逐次商是 $\mathcal{O}/\mathcal{I}$-模层。因此，为证明
$K \otimes_\mathcal{O}^\mathbf{L} \mathcal{O}/\mathcal{I}^n$
的 Tor 振幅在 $[a,b]$ 中，只需证明对所有
$\mathcal{O}/\mathcal{I}$-模 $\mathcal{G}$ 以及所有
$i \not \in [a,b]$，都有
$H^i(K \otimes_\mathcal{O}^\mathbf{L} \mathcal{O}/\mathcal{I}^n
\otimes_{\mathcal{O}/\mathcal{I}^n}^\mathbf{L} \mathcal{G})$
为零。由于
$$
\left(K \otimes_\mathcal{O}^\mathbf{L} \mathcal{O}/\mathcal{I}^n\right)
\otimes_{\mathcal{O}/\mathcal{I}^n}^\mathbf{L} \mathcal{G}
=
\left(K \otimes_\mathcal{O}^\mathbf{L} \mathcal{O}/\mathcal{I}\right)
\otimes_{\mathcal{O}/\mathcal{I}}^\mathbf{L} \mathcal{G}
$$
对每个 $\mathcal{O}/\mathcal{I}$-模层 $\mathcal{G}$ 都成立，结论随即得出。
\end{proof}

\begin{lemma}
\label{sites-cohomology-lemma-tor-amplitude-stalk}
设 $(\mathcal{C}, \mathcal{O})$ 是一个环化站点，
$E$ 是 $D(\mathcal{O})$ 的一个对象，并设
$a,b \in \mathbf{Z}$。
\begin{enumerate}
\item 如果 $E$ 的 Tor 振幅在 $[a,b]$ 中，那么对站点
$\mathcal{C}$ 的每个点 $p$，
$D(\mathcal{O}_p)$ 的对象 $E_p$ 的 Tor 振幅也在 $[a,b]$ 中。
\item 如果 $\mathcal{C}$ 有足够多的点，那么逆命题成立。
\end{enumerate}
\end{lemma}

\begin{proof}
证明 (1)。这是因为在 $p$ 处取茎等同于沿环化站点态射
$(p, \mathcal{O}_p) \to (\mathcal{C}, \mathcal{O})$
拉回，因而可应用引理 \ref{sites-cohomology-lemma-tor-amplitude-pullback}。

\medskip\noindent
证明 (2)。如果 $\mathcal{C}$ 有足够多的点，则可在各茎处检验
$H^i(E \otimes_\mathcal{O}^\mathbf{L} \mathcal{F})$
的消失性，见《站点上的模》引理
\ref{sites-modules-lemma-check-exactness-stalks}。
由于
$H^i(E \otimes_\mathcal{O}^\mathbf{L} \mathcal{F})_p =
H^i(E_p \otimes_{\mathcal{O}_p}^\mathbf{L} \mathcal{F}_p)$，结论得证。
\end{proof}










\section{完美复形}
\label{sites-cohomology-section-perfect}

\noindent
本节讨论环化站点上完美复形的性质。

\begin{definition}
\label{sites-cohomology-definition-perfect}
设 $(\mathcal{C}, \mathcal{O})$ 是一个环化站点，
$\mathcal{E}^\bullet$ 是 $\mathcal{O}$-模复形。
如果对 $\mathcal{C}$ 的每个对象 $U$，都存在覆盖
$\{U_i \to U\}$，使得对每个 $i$ 都存在复形态射
$\mathcal{E}_i^\bullet \to \mathcal{E}^\bullet|_{U_i}$
，它是拟同构，且 $\mathcal{E}_i^\bullet$ 严格完美，
则称 $\mathcal{E}^\bullet$ 是{\it 完美的}。
如果 $D(\mathcal{O})$ 的对象 $E$ 能由完美
$\mathcal{O}$-模复形表示，则称 $E$ 是{\it 完美的}。
\end{definition}

\noindent
如果 $\Sh(\mathcal{C})$ 拟紧
（《站点》第 \ref{sites-section-quasi-compact} 节），
那么 $D(\mathcal{O})$ 的完美对象属于
$D^b(\mathcal{O})$。否则未必如此。

\begin{lemma}
\label{sites-cohomology-lemma-perfect-independent-representative}
设 $(\mathcal{C}, \mathcal{O})$ 是一个环化站点，
$E$ 是 $D(\mathcal{O})$ 的一个对象。
\begin{enumerate}
\item 如果 $\mathcal{C}$ 有终对象 $X$，并且存在覆盖
$\{U_i \to X\}$、$\mathcal{O}_{U_i}$-模的严格完美复形
$\mathcal{E}_i^\bullet$，以及 $D(\mathcal{O}_{U_i})$ 中的同构
$\alpha_i : \mathcal{E}_i^\bullet \to E|_{U_i}$，
那么 $E$ 是完美的。
\item 如果 $E$ 是完美的，那么表示 $E$ 的任意复形都是完美的。
\end{enumerate}
\end{lemma}

\begin{proof}
与引理 \ref{sites-cohomology-lemma-pseudo-coherent-independent-representative}
的证明完全相同。
\end{proof}

\begin{lemma}
\label{sites-cohomology-lemma-perfect-precise}
设 $(\mathcal{C}, \mathcal{O})$ 是一个环化站点，
$E$ 是 $D(\mathcal{O})$ 的一个对象。
设整数 $a \leq b$。如果 $E$ 的 Tor 振幅在 $[a,b]$ 中，
并且 $E$ 是 $(a-1)$-伪凝聚的，那么 $E$ 是完美的。
\end{lemma}

\begin{proof}
设 $U$ 是 $\mathcal{C}$ 的一个对象。用某个覆盖的各成员代替
$U$，并用局部化 $\mathcal{C}/U$ 代替 $\mathcal{C}$ 后，
可设存在严格完美复形 $\mathcal{E}^\bullet$ 和映射
$\alpha : \mathcal{E}^\bullet \to E$，使得
$H^i(\alpha)$ 对 $i \geq a$ 是同构。
我们可以并确实用
$\sigma_{\geq a - 1}\mathcal{E}^\bullet$ 代替
$\mathcal{E}^\bullet$。选取有区别三角
$$
\mathcal{E}^\bullet \to E \to C \to \mathcal{E}^\bullet[1]
$$
。由 $E$ 和 $\mathcal{E}^\bullet$ 的上同调层之消失性以及
关于 $\alpha$ 的假设，得到
$C \cong \mathcal{K}[2-a]$，其中
$\mathcal{K} = \Ker(\mathcal{E}^{a-1} \to \mathcal{E}^a)$。
设 $\mathcal{F}$ 是一个 $\mathcal{O}$-模。
应用 $- \otimes_\mathcal{O}^\mathbf{L} \mathcal{F}$，
关于 $E$ 的 Tor 振幅在 $[a,b]$ 中的假设蕴含
$\mathcal{K} \otimes_\mathcal{O} \mathcal{F} \to
\mathcal{E}^{a - 1} \otimes_\mathcal{O} \mathcal{F}$ 的像为
$\Ker(\mathcal{E}^{a - 1} \otimes_\mathcal{O} \mathcal{F}
\to \mathcal{E}^a \otimes_\mathcal{O} \mathcal{F})$.
因此
$\text{Tor}_1^\mathcal{O}(\mathcal{E}', \mathcal{F})=0$，
其中
$\mathcal{E}' = \Coker(\mathcal{E}^{a-1}\to\mathcal{E}^a)$。
故 $\mathcal{E}'$ 平坦（引理 \ref{sites-cohomology-lemma-flat-tor-zero}）。
于是，由《站点上的模》引理
\ref{sites-modules-lemma-flat-locally-finite-presentation}，
存在覆盖 $\{U_i \to U\}$，使得
$\mathcal{E}'|_{U_i}$ 是某个有限自由模的直和项。
因此复形
$$
\mathcal{E}'|_{U_i} \to \mathcal{E}^{a + 1}|_{U_i} \to \ldots \to
\mathcal{E}^b|_{U_i}
$$
与 $E|_{U_i}$ 拟同构，所以 $E$ 是完美的。
\end{proof}

\begin{lemma}
\label{sites-cohomology-lemma-perfect}
设 $(\mathcal{C}, \mathcal{O})$ 是一个环化站点，
$E$ 是 $D(\mathcal{O})$ 的一个对象。下列条件等价：
\begin{enumerate}
\item $E$ 是完美的；
\item $E$ 是伪凝聚的，并且局部具有有限 Tor 维数。
\end{enumerate}
\end{lemma}

\begin{proof}
假设 (1)。设 $U$ 是 $\mathcal{C}$ 的一个对象。
由定义，存在覆盖 $\{U_i \to U\}$，使得
$E|_{U_i}$ 由严格完美复形表示。因此，由引理
\ref{sites-cohomology-lemma-pseudo-coherent-independent-representative}，
$E$ 是伪凝聚的（即对所有 $m$ 都是 $m$-伪凝聚的）。
此外，有限自由模的直和项是平坦的，故由引理
\ref{sites-cohomology-lemma-tor-amplitude}，$E|_{U_i}$ 具有有限 Tor 维数。
所以 (2) 成立。

\medskip\noindent
假设 (2)。设 $U$ 是 $\mathcal{C}$ 的一个对象。
用某个覆盖的各成员代替 $U$ 后，可设存在整数
$a \leq b$，使得 $E|_U$ 的 Tor 振幅在 $[a,b]$ 中。
由于 $E|_U$ 对所有 $m$ 都是 $m$-伪凝聚的，
应用引理 \ref{sites-cohomology-lemma-perfect-precise} 即得结论。
\end{proof}

\begin{lemma}
\label{sites-cohomology-lemma-perfect-pullback}
设 $(f, f^\sharp) : (\mathcal{C}, \mathcal{O}_\mathcal{C}) \to
(\mathcal{D}, \mathcal{O}_\mathcal{D})$
是环化站点的态射。设 $E$ 是
$D(\mathcal{O}_\mathcal{D})$ 的一个对象。
如果 $E$ 在 $D(\mathcal{O}_\mathcal{D})$ 中是完美的，
那么 $Lf^*E$ 在 $D(\mathcal{O}_\mathcal{C})$ 中是完美的。
\end{lemma}

\begin{proof}
这由引理 \ref{sites-cohomology-lemma-perfect}、
\ref{sites-cohomology-lemma-tor-amplitude-pullback} 和
\ref{sites-cohomology-lemma-pseudo-coherent-pullback} 得出。
\end{proof}

\begin{lemma}
\label{sites-cohomology-lemma-two-out-of-three-perfect}
设 $(\mathcal{C}, \mathcal{O})$ 是一个环化站点。
设 $(K,L,M,f,g,h)$ 是 $D(\mathcal{O})$ 中的有区别三角。
如果 $K,L,M$ 中有两个是完美的，那么第三个也是完美的。
\end{lemma}

\begin{proof}
第一种证明：合并引理 \ref{sites-cohomology-lemma-perfect}、
\ref{sites-cohomology-lemma-cone-pseudo-coherent} 和
\ref{sites-cohomology-lemma-cone-tor-amplitude}。

第二种证明（概要）：设 $K$ 和 $L$ 完美。设 $U$ 是
$\mathcal{C}$ 的一个对象。用某个覆盖的各成员代替 $U$ 后，
可设 $K|_U$ 和 $L|_U$ 分别由严格完美复形
$\mathcal{K}^\bullet$ 和 $\mathcal{L}^\bullet$ 表示。
再次用某个覆盖的各成员代替 $U$ 后，可设映射
$K|_U \to L|_U$ 由复形映射
$\alpha : \mathcal{K}^\bullet \to \mathcal{L}^\bullet$ 给出，
见引理 \ref{sites-cohomology-lemma-local-actual}。
于是 $M|_U$ 同构于 $\alpha$ 的锥；由引理
\ref{sites-cohomology-lemma-cone}，该锥严格完美。
\end{proof}

\begin{lemma}
\label{sites-cohomology-lemma-tensor-perfect}
设 $(\mathcal{C}, \mathcal{O})$ 是一个环化站点。
如果 $K,L$ 是 $D(\mathcal{O})$ 的完美对象，那么
$K \otimes_\mathcal{O}^\mathbf{L} L$ 也是完美对象。
\end{lemma}

\begin{proof}
由引理 \ref{sites-cohomology-lemma-perfect}、
\ref{sites-cohomology-lemma-tensor-pseudo-coherent} 和
\ref{sites-cohomology-lemma-tensor-tor-amplitude} 得出。
\end{proof}

\begin{lemma}
\label{sites-cohomology-lemma-summands-perfect}
设 $(\mathcal{C}, \mathcal{O})$ 是一个环化站点。
如果 $K \oplus L$ 是 $D(\mathcal{O})$ 的完美对象，
那么 $K$ 和 $L$ 也都是完美对象。
\end{lemma}

\begin{proof}
由引理 \ref{sites-cohomology-lemma-perfect}、
\ref{sites-cohomology-lemma-summands-pseudo-coherent} 和
\ref{sites-cohomology-lemma-summands-tor-amplitude} 得出。
\end{proof}














\section{对偶}
\label{sites-cohomology-section-duals}

\noindent
本节刻画复形范畴和导出范畴中的可对偶对象。
特别地，我们将看到，$D(\mathcal{O})$ 的对象有对偶，当且仅当它是完美的
（这由例 \ref{sites-cohomology-example-dual-derived} 和
引理 \ref{sites-cohomology-lemma-left-dual-derived} 得出）。

\begin{lemma}
\label{sites-cohomology-lemma-symmetric-monoidal-cat-complexes}
设 $(\mathcal{C}, \mathcal{O})$ 是一个环化站点。以
$\mathcal{F}^\bullet \otimes \mathcal{G}^\bullet =
\text{Tot}(\mathcal{F}^\bullet \otimes_\mathcal{O} \mathcal{G}^\bullet)$
定义张量积的 $\mathcal{O}$-模复形范畴是一个对称幺半范畴。
\end{lemma}

\begin{proof}
略。提示：取在次数 $0$ 上为 $\mathcal{O}$、在其余次数上为零的复形
作为单位对象 $\mathbf{1}$，并取显然的同构
$\text{Tot}(\mathbf{1} \otimes_\mathcal{O} \mathcal{G}^\bullet) =
\mathcal{G}^\bullet$ 以及
$\text{Tot}(\mathcal{F}^\bullet \otimes_\mathcal{O} \mathbf{1}) =
\mathcal{F}^\bullet$。
要证明该引理，需要验证若干图表的交换性；见《范畴》定义
\ref{categories-definition-monoidal-category} 和
\ref{categories-definition-symmetric-monoidal-category}。
各项验证都很直接。
\end{proof}

\begin{example}
\label{sites-cohomology-example-dual}
设 $(\mathcal{C}, \mathcal{O})$ 是一个环化站点。设 $\mathcal{F}^\bullet$
为一个 $\mathcal{O}$-模复形，并且对每个
$U \in \Ob(\mathcal{C})$，都存在一个覆盖 $\{U_i \to U\}$，
使得 $\mathcal{F}^\bullet|_{U_i}$ 严格完美。考虑复形
$$
\mathcal{G}^\bullet = \SheafHom^\bullet(\mathcal{F}^\bullet, \mathcal{O})
$$
，如第 \ref{sites-cohomology-section-hom-complexes} 节所述。令
$$
\eta :
\mathcal{O}
\to
\text{Tot}(\mathcal{F}^\bullet \otimes_\mathcal{O} \mathcal{G}^\bullet)
\quad\text{和}\quad
\epsilon :
\text{Tot}(\mathcal{G}^\bullet \otimes_\mathcal{O} \mathcal{F}^\bullet)
\to
\mathcal{O}
$$
分别为 $\eta = \sum \eta_n$ 和 $\epsilon = \sum \epsilon_n$，
其中 $\eta_n : \mathcal{O} \to
\mathcal{F}^n \otimes_\mathcal{O} \mathcal{G}^{-n}$
以及
$\epsilon_n : \mathcal{G}^{-n} \otimes_\mathcal{O} \mathcal{F}^n
\to \mathcal{O}$ 如《站点上的模》例 \ref{sites-modules-example-dual}
中所定义。于是按照《范畴》定义 \ref{categories-definition-dual}，
$\mathcal{G}^\bullet, \eta, \epsilon$ 是 $\mathcal{F}^\bullet$ 的左对偶。
我们略去下列等式的验证：
$(1 \otimes \epsilon) \circ (\eta \otimes 1) = \text{id}_{\mathcal{F}^\bullet}$
和
$(\epsilon \otimes 1) \circ (1 \otimes \eta) =
\text{id}_{\mathcal{G}^\bullet}$。请与《代数详论》引理
\ref{more-algebra-lemma-left-dual-complex} 比较。
\end{example}

\begin{lemma}
\label{sites-cohomology-lemma-left-dual-complex}
设 $(\mathcal{C}, \mathcal{O})$ 是一个环化站点。设 $\mathcal{F}^\bullet$
为一个 $\mathcal{O}$-模复形。如果 $\mathcal{F}^\bullet$
在 $\mathcal{O}$-模复形的幺半范畴中有左对偶
（《范畴》定义 \ref{categories-definition-dual}），
那么对 $\mathcal{C}$ 的每个对象 $U$，都存在一个覆盖
$\{U_i \to U\}$，使得 $\mathcal{F}^\bullet|_{U_i}$
严格完美，并且该左对偶正是例 \ref{sites-cohomology-example-dual} 中构造的左对偶。
\end{lemma}

\begin{proof}
由左对偶的唯一性（《范畴》注记
\ref{categories-remark-left-dual-adjoint}），只要证明
$\mathcal{F}^\bullet$ 具有所述性质，就可得到最后一个断言。
设 $\mathcal{G}^\bullet, \eta, \epsilon$ 为一个左对偶。写成
$\eta = \sum \eta_n$ 和 $\epsilon = \sum \epsilon_n$，其中
$\eta_n : \mathcal{O} \to
\mathcal{F}^n \otimes_\mathcal{O} \mathcal{G}^{-n}$
以及
$\epsilon_n : \mathcal{G}^{-n} \otimes_\mathcal{O} \mathcal{F}^n
\to \mathcal{O}$。由于
$(1 \otimes \epsilon) \circ (\eta \otimes 1) = \text{id}_{\mathcal{F}^\bullet}$
和
$(\epsilon \otimes 1) \circ (1 \otimes \eta) = \text{id}_{\mathcal{G}^\bullet}$
由《范畴》定义 \ref{categories-definition-dual}，立即可见
$(1 \otimes \epsilon_n) \circ (\eta_n \otimes 1) = \text{id}_{\mathcal{F}^n}$
以及
$(\epsilon_n \otimes 1) \circ (1 \otimes \eta_n) =
\text{id}_{\mathcal{G}^{-n}}$.
换言之，$\mathcal{G}^{-n}$ 是 $\mathcal{F}^n$ 的左对偶，
而且《站点上的模》引理 \ref{sites-modules-lemma-left-dual-module}
适用于每个 $\mathcal{F}^n$。设 $U$ 是 $\mathcal{C}$ 的一个对象。
存在覆盖 $\{U_i \to U\}$，使得对每个 $i$，
只有有限多个 $\eta_n|_{U_i}$ 非零。因此，用 $U_i$ 替换 $U$ 后，
可以假设只有有限多个 $\eta_n|_U$ 非零；由所引引理，
这蕴含只有有限多个 $\mathcal{F}^n|_U$ 非零。
再次应用该引理，便可找到一个覆盖 $\{U_i \to U\}$，
使得每个 $\mathcal{F}^n|_{U_i}$ 都是有限自由
$\mathcal{O}$-模的直和项，证明完毕。
\end{proof}

\begin{lemma}
\label{sites-cohomology-lemma-dual-perfect-complex}
设 $(\mathcal{C}, \mathcal{O})$ 是一个环化站点。
设 $K$ 是 $D(\mathcal{O})$ 的一个完美对象。
则 $K^\vee = R\SheafHom(K, \mathcal{O})$ 也是完美对象，且
$(K^\vee)^\vee \cong K$。存在函子性同构
$$
M \otimes^\mathbf{L}_\mathcal{O} K^\vee = R\SheafHom_\mathcal{O}(K, M)
$$
和
$$
H^0(\mathcal{C}, M \otimes^\mathbf{L}_\mathcal{O} K^\vee) =
\Hom_{D(\mathcal{O})}(K, M)
$$
，其中 $M$ 是 $D(\mathcal{O})$ 的对象。
\end{lemma}

\begin{proof}
下文将不再特别说明地使用内部 Hom 的形成与限制可交换这一事实
（引理 \ref{sites-cohomology-lemma-restriction-RHom-to-U}）。设 $U$
是 $\mathcal{C}$ 的任意对象。要验证 $K^\vee$ 是完美的，
只需证明存在覆盖 $\{U_i \to U\}$，使得对所有 $i$，
$K^\vee|_{U_i}$ 都是完美的。存在典范映射
$$
K = R\SheafHom(\mathcal{O}_X, \mathcal{O}_X)
\otimes_{\mathcal{O}_X}^\mathbf{L} K \longrightarrow
R\SheafHom(R\SheafHom(K, \mathcal{O}_X), \mathcal{O}_X) =
(K^\vee)^\vee
$$
；见引理 \ref{sites-cohomology-lemma-internal-hom-evaluate}。只需证明存在覆盖
$\{U_i \to U\}$，使得该映射在 $\mathcal{C}/U_i$
上的限制对所有 $i$ 都是同构。由引理 \ref{sites-cohomology-lemma-dual}，
要证明最后一个断言，只需验证映射 (\ref{sites-cohomology-equation-eval})
$$
M \otimes^\mathbf{L}_\mathcal{O} K^\vee
\longrightarrow
R\SheafHom(K, M)
$$
是同构。这同样是上述意义下的局部问题。因此，只需在 $K$
由严格完美复形表示时证明该引理。

\medskip\noindent
假设 $K$ 由严格完美复形 $\mathcal{E}^\bullet$ 表示。
由引理 \ref{sites-cohomology-lemma-Rhom-strictly-perfect}，$K^\vee$
由下述复形表示：其次数 $-n$ 上的项为
$(\mathcal{E}^n)^\vee =
\SheafHom_\mathcal{O}(\mathcal{E}^n, \mathcal{O})$
。由于 $\mathcal{E}^n$ 是有限自由 $\mathcal{O}$-模的直和项，
$(\mathcal{E}^n)^\vee$ 也如此。因此，$K^\vee$
也由一个严格完美复形表示，从而 $K^\vee$ 是完美的。
映射 $K \to (K^\vee)^\vee$ 是同构，因为它在相差一个符号的意义下
由评价映射 $\mathcal{E}^n \to ((\mathcal{E}^n)^\vee)^\vee$
给出，而这些映射都是同构。要证明 (\ref{sites-cohomology-equation-eval}) 是同构，
用一个 K-平坦复形 $\mathcal{F}^\bullet$ 表示 $M$。
由引理 \ref{sites-cohomology-lemma-Rhom-strictly-perfect}，复形
$R\SheafHom(K, M)$ 由各项为
$$
\bigoplus\nolimits_{n = p + q}
\SheafHom_\mathcal{O}(\mathcal{E}^{-q}, \mathcal{F}^p)
$$
的复形表示。另一方面，对象
$M \otimes^\mathbf{L}_\mathcal{O} K^\vee$ 由各项为
$$
\bigoplus\nolimits_{n = p + q}
\mathcal{F}^p \otimes_\mathcal{O} (\mathcal{E}^{-q})^\vee
$$
的复形表示。因此，(\ref{sites-cohomology-equation-eval}) 是同构这一断言归结为：
当 $\mathcal{E}$ 是有限自由 $\mathcal{O}$-模的直和项且
$\mathcal{F}$ 是任意 $\mathcal{O}$-模时，典范映射
$$
\mathcal{F}
\otimes_\mathcal{O}
\SheafHom_\mathcal{O}(\mathcal{E}, \mathcal{O})
\longrightarrow
\SheafHom_\mathcal{O}(\mathcal{E}, \mathcal{F})
$$
是同构。这立即由 $\mathcal{E}$ 为有限自由模时的相应断言得出。
\end{proof}

\begin{lemma}
\label{sites-cohomology-lemma-symmetric-monoidal-derived}
设 $(\mathcal{C}, \mathcal{O})$ 是一个环化站点。导出范畴
$D(\mathcal{O})$ 是一个对称幺半范畴，其张量积由导出张量积给出，
并带有通常的结合性约束和交换性约束（符号规则见《代数详论》第
\ref{more-algebra-section-sign-rules} 节）。
\end{lemma}

\begin{proof}
略。与引理 \ref{sites-cohomology-lemma-symmetric-monoidal-cat-complexes} 比较。
\end{proof}

\begin{example}
\label{sites-cohomology-example-dual-derived}
设 $(\mathcal{C}, \mathcal{O})$ 是一个环化站点。设 $K$ 是
$D(\mathcal{O})$ 的一个完美对象。如引理
\ref{sites-cohomology-lemma-dual-perfect-complex} 中那样，令
$K^\vee = R\SheafHom(K, \mathcal{O})$。则映射
$$
K \otimes_\mathcal{O}^\mathbf{L} K^\vee \longrightarrow R\SheafHom(K, K)
$$
是同构（由该引理）。记映射
$$
\eta :
\mathcal{O}
\longrightarrow
K \otimes_\mathcal{O}^\mathbf{L} K^\vee
$$
为在上述同构下将 $1$ 映到与 $\text{id}_K$ 对应的截面的映射。
记
$$
\epsilon : 
K^\vee
\otimes_\mathcal{O}^\mathbf{L} K
\longrightarrow
\mathcal{O}
$$
为评价映射（例如，可以用引理
\ref{sites-cohomology-lemma-internal-hom-composition} 构造它）。于是按照《范畴》定义
\ref{categories-definition-dual}，$K^\vee, \eta, \epsilon$
是 $K$ 的左对偶。我们略去下列等式的验证：
$(1 \otimes \epsilon) \circ (\eta \otimes 1) = \text{id}_K$
和
$(\epsilon \otimes 1) \circ (1 \otimes \eta) =
\text{id}_{K^\vee}$.
\end{example}

\begin{lemma}
\label{sites-cohomology-lemma-left-dual-derived}
设 $(\mathcal{C}, \mathcal{O})$ 是一个环化站点。设 $M$ 是
$D(\mathcal{O})$ 的一个对象。如果 $M$ 在幺半范畴
$D(\mathcal{O})$ 中有左对偶（《范畴》定义
\ref{categories-definition-dual}），那么 $M$ 是完美的，
并且该左对偶正是例 \ref{sites-cohomology-example-dual-derived} 中构造的左对偶。
\end{lemma}

\begin{proof}
设 $N, \eta, \epsilon$ 为一个左对偶。注意，对 $\mathcal{C}$
的任意对象 $U$，限制 $N|_U, \eta|_U, \epsilon|_U$
都是 $M|_U$ 的一个左对偶。

\medskip\noindent
设 $U$ 是 $\mathcal{C}$ 的一个对象。只需找到 $\mathcal{C}$
的一个覆盖 $\{U_i \to U\}_{i \in I}$，使得 $M|_{U_i}$
是 $D(\mathcal{O}_{U_i})$ 的完美对象。因此，可以分别用
$\mathcal{C}/U, \mathcal{O}_U, M|_U, N|_U, \eta|_U, \epsilon|_U$
替换 $\mathcal{C}, \mathcal{O}, M, N, \eta, \epsilon$，并假设
$\mathcal{C}$ 有终对象 $X$。此外，在证明过程中，可以（有限多次）
用 $X$ 的覆盖 $\{U_i \to X\}$ 的成员替换 $X$。

\medskip\noindent
我们将多次使用下述论证。任取一个表示 $M$ 的 $\mathcal{O}$-模复形
$\mathcal{M}^\bullet$。选取一个表示 $N$ 的 K-平坦复形
$\mathcal{N}^\bullet$，其各项均为平坦 $\mathcal{O}$-模；
见引理 \ref{sites-cohomology-lemma-K-flat-resolution}。考虑映射
$$
\eta : \mathcal{O} \to
\text{Tot}(\mathcal{M}^\bullet \otimes_\mathcal{O} \mathcal{N}^\bullet)
$$
用某个覆盖的成员替换 $X$ 后，可以找到整数 $N$，并且对
$i = 1, \ldots, N$，找到整数 $n_i \in \mathbf{Z}$ 以及
$\mathcal{M}^{n_i}$ 和 $\mathcal{N}^{-n_i}$ 的截面 $f_i$ 和 $g_i$，使得
$$
\eta(1) = \sum\nolimits_i f_i \otimes g_i
$$
设 $\mathcal{K}^\bullet \subset \mathcal{M}^\bullet$
是任意一个包含截面 $f_i$（$i = 1, \ldots, N$）的
$\mathcal{O}$-模子复形。由于
$\text{Tot}(\mathcal{K}^\bullet \otimes_\mathcal{O} \mathcal{N}^\bullet)
\subset
\text{Tot}(\mathcal{M}^\bullet \otimes_\mathcal{O} \mathcal{N}^\bullet)$
（这是由模 $\mathcal{N}^n$ 的平坦性得到的），可见 $\eta$ 分解经过
$$
\tilde \eta :
\mathcal{O} \to
\text{Tot}(\mathcal{K}^\bullet \otimes_\mathcal{O} \mathcal{N}^\bullet)
$$
记 $K$ 为 $D(\mathcal{O})$ 中由 $\mathcal{K}^\bullet$
表示的对象，便得到交换图
$$
\xymatrix{
M \ar[rr]_-{\eta \otimes 1} \ar[rrd]_{\tilde \eta \otimes 1} & &
M \otimes^\mathbf{L} N \otimes^\mathbf{L} M
\ar[r]_-{1 \otimes \epsilon} &
M \\
& &
K \otimes^\mathbf{L} N \otimes^\mathbf{L} M
\ar[u] \ar[r]^-{1 \otimes \epsilon} &
K \ar[u]
}
$$
由于上行的复合是 $M$ 上的恒等映射，故 $M$
在 $D(\mathcal{O})$ 中是 $K$ 的一个直和项。

\medskip\noindent
首次应用上述论证时，可以选取子复形
$\mathcal{K}^\bullet = \sigma_{\geq a} \tau_{\leq b}\mathcal{M}^\bullet$
，其中对 $i = 1, \ldots, N$ 有 $a < n_i < b$。于是 $M$
在 $D(\mathcal{O})$ 中是一个有界复形的直和项，从而可以假设
$M$ 属于 $D^b(\mathcal{O})$。（回忆上述过程涉及用一个覆盖的成员替换
$X$。）

\medskip\noindent
由于 $M$ 属于 $D^b(\mathcal{O})$，可以选取 $\mathcal{M}^\bullet$
为一个由平坦模组成的上有界复形（由《模》引理
\ref{modules-lemma-module-quotient-flat} 和《导出范畴》引理
\ref{derived-lemma-subcategory-left-resolution}）。在上述论证中，
可以选取 $\mathcal{K}^\bullet = \sigma_{\geq a}\mathcal{M}^\bullet$，
其中对 $i = 1, \ldots, N$ 有 $a < n_i$。因此，可以假设
$M$ 在 $D(\mathcal{O})$ 中是一个由平坦模组成的有界复形的直和项。
特别地，$M$ 具有有限 Tor 振幅。

\medskip\noindent
设 $M$ 的 Tor 振幅包含于 $[a, b]$。假设 $M$ 是 $m$-伪凝聚的；
我们将证明（用一个覆盖的成员替换 $X$ 后）可以假设 $M$
是 $(m - 1)$-伪凝聚的。结合引理 \ref{sites-cohomology-lemma-perfect-precise}
以及 $M$ 无论如何都是 $(b + 1)$-伪凝聚的这一事实，这就会完成证明。
用一个覆盖的成员替换 $X$ 后，可以假设存在严格完美复形
$\mathcal{E}^\bullet$ 和 $D(\mathcal{O})$ 中的映射
$\alpha : \mathcal{E}^\bullet \to M$，使得 $H^i(\alpha)$
在 $i > m$ 时是同构，在 $i = m$ 时是满射。可以并且确实假设
$\mathcal{E}^i = 0$ 对 $i < m$ 成立。选取一个有区别三角
$$
\mathcal{E}^\bullet \to M \to L \to \mathcal{E}^\bullet[1]
$$
。注意，对 $i \geq m$ 有 $H^i(L) = 0$。因此，可以用一个满足
$\mathcal{L}^i = 0$（$i \geq m$）的复形 $\mathcal{L}^\bullet$
表示 $L$。映射 $L \to \mathcal{E}^\bullet[1]$ 由复形映射
$\mathcal{L}^\bullet \to \mathcal{E}^\bullet[1]$
给出；它除次数 $m - 1$ 外处处为零，而在次数 $m - 1$ 上得到映射
$\mathcal{L}^{m - 1} \to \mathcal{E}^m$；见《导出范畴》引理
\ref{derived-lemma-negative-exts}。于是 $M$ 由复形
$$
\mathcal{M}^\bullet :
\ldots \to
\mathcal{L}^{m - 2} \to
\mathcal{L}^{m - 1} \to
\mathcal{E}^m \to
\mathcal{E}^{m + 1} \to \ldots
$$
表示。把第二段中的论证应用于此复形，得到 $\mathcal{M}^{n_i}$
的截面 $f_i$，其中 $i = 1, \ldots, N$。对 $n < m$，令
$\mathcal{K}^n \subset \mathcal{L}^n$ 为由满足 $n_i = n$
的截面 $f_i$ 以及满足 $n_i = n - 1$ 的 $d(f_i)$ 所生成的
$\mathcal{O}$-子模。对 $n \geq m$，令
$\mathcal{K}^n = \mathcal{E}^n$。显然，存在有区别三角之间的态射
$$
\xymatrix{
\mathcal{E}^\bullet \ar[r] &
\mathcal{M}^\bullet \ar[r] &
\mathcal{L}^\bullet \ar[r] &
\mathcal{E}^\bullet[1] \\
\mathcal{E}^\bullet \ar[r] \ar[u] &
\mathcal{K}^\bullet \ar[r] \ar[u] &
\sigma_{\leq m - 1}\mathcal{K}^\bullet \ar[r] \ar[u] &
\mathcal{E}^\bullet[1] \ar[u]
}
$$
，其中所有态射均如上所示。记 $K$ 为 $D(\mathcal{O})$
中与复形 $\mathcal{K}^\bullet$ 对应的对象。
由证明第二段中的论证，得到 $D(\mathcal{O})$ 中的态射
$s : M \to K$，使得复合 $M \to K \to M$ 是 $M$ 上的恒等映射。
我们不知道图表
$$
\xymatrix{
\mathcal{E}^\bullet \ar[r] &
\mathcal{K}^\bullet \ar@{=}[r] &
K \\
\mathcal{E}^\bullet \ar[u]^{\text{id}} \ar[r]^i &
\mathcal{M}^\bullet \ar@{=}[r] &
M \ar[u]_s
}
$$
是否交换，但知道它与映射 $K \to M$ 复合后交换。由引理
\ref{sites-cohomology-lemma-local-actual}，用一个覆盖的成员替换 $X$ 后，
可以假设 $s \circ i$ 由复形映射
$\sigma : \mathcal{E}^\bullet \to \mathcal{K}^\bullet$ 给出。
由同一个引理，可以假设 $\sigma$ 与包含映射
$\mathcal{K}^\bullet \subset \mathcal{M}^\bullet$ 的复合通过某个同伦
$\{h^i : \mathcal{E}^i \to \mathcal{M}^{i - 1}\}$ 同伦于零。
因此，用 $\mathcal{K}^{m - 1} + \Im(h^m)$ 替换
$\mathcal{K}^{m - 1}$ 后（注意这样做之后，
$\mathcal{K}^{m - 1}$ 仍由有限多个整体截面生成），
可见 $\sigma$ 本身同伦于零！这意味着存在下列实线部分交换的图表
$$
\xymatrix{
\mathcal{E}^\bullet \ar[r] &
M \ar[r] &
\mathcal{L}^\bullet \ar[r] &
\mathcal{E}^\bullet[1] \\
\mathcal{E}^\bullet \ar[r] \ar[u] &
K \ar[r] \ar[u] &
\sigma_{\leq m - 1}\mathcal{K}^\bullet \ar[r] \ar[u] &
\mathcal{E}^\bullet[1] \ar[u] \\
\mathcal{E}^\bullet \ar[r] \ar[u] &
M \ar[r] \ar[u]^s &
\mathcal{L}^\bullet \ar[r] \ar@{..>}[u] &
\mathcal{E}^\bullet[1] \ar[u]
}
$$
。由三角范畴的公理，得到一条使图表成立的虚线箭头。
考察次数 $m - 1$ 上的上同调层，得到
$$
\xymatrix{
H^{m - 1}(M) \ar[r] &
H^{m - 1}(\mathcal{L}^\bullet) \ar[r] &
H^m(\mathcal{E}^\bullet) \\
H^{m - 1}(K) \ar[r] \ar[u] &
H^{m - 1}(\sigma_{\leq m - 1}\mathcal{K}^\bullet) \ar[r] \ar[u] &
H^m(\mathcal{E}^\bullet) \ar[u] \\
H^{m - 1}(M) \ar[r] \ar[u] &
H^{m - 1}(\mathcal{L}^\bullet) \ar[r] \ar[u] &
H^m(\mathcal{E}^\bullet) \ar[u]
}
$$
。由于左列和右列中的竖直复合都是恒等映射，故中列的竖直复合
$H^{m - 1}(\mathcal{L}^\bullet) \to
H^{m - 1}(\sigma_{\leq m - 1}\mathcal{K}^\bullet) \to
H^{m - 1}(\mathcal{L}^\bullet)$ 是满射！特别地，
$H^{m - 1}(\sigma_{\leq m - 1}\mathcal{K}^\bullet) \to
H^{m - 1}(\mathcal{L}^\bullet)$ 是满射。
利用上述三角之间的态射所诱导的上同调层长正合序列之间的映射，
追图可见，这蕴含 $H^i(K) \to H^i(M)$ 在 $i \geq m$
时是同构，在 $i = m - 1$ 时是满射。
由构造，可以选取 $r \geq 0$ 和一个满射
$\mathcal{O}^{\oplus r} \to \mathcal{K}^{m - 1}$。于是复合
$$
(\mathcal{O}^{\oplus r} \to \mathcal{E}^m \to
\mathcal{E}^{m + 1} \to \ldots ) \longrightarrow
K \longrightarrow M
$$
在次数 $\geq m$ 的上同调层上诱导同构，在次数 $m - 1$
上诱导满射，证明完毕。
\end{proof}

\begin{lemma}
\label{sites-cohomology-lemma-colim-and-lim-of-duals}
\begin{slogan}
完美对象系的平凡对偶性。
\end{slogan}
设 $(\mathcal{C}, \mathcal{O})$ 是一个环化站点。设
$(K_n)_{n \in \mathbf{N}}$ 是 $D(\mathcal{O})$ 中的一个完美对象系。
设 $K = \text{hocolim} K_n$ 为导出余极限
（《导出范畴》定义 \ref{derived-definition-derived-colimit}）。
那么对 $D(\mathcal{O})$ 的任意对象 $E$，有
$$
R\SheafHom(K, E) = R\lim E \otimes^\mathbf{L}_\mathcal{O} K_n^\vee
$$
，其中 $(K_n^\vee)$ 是由完美复形的对偶组成的逆系统。
\end{lemma}

\begin{proof}
由引理 \ref{sites-cohomology-lemma-dual-perfect-complex}，有
$R\lim E \otimes^\mathbf{L}_\mathcal{O} K_n^\vee =
R\lim R\SheafHom(K_n, E)$
，它嵌入有区别三角
$$
R\lim R\SheafHom(K_n, E) \to
\prod R\SheafHom(K_n, E) \to
\prod R\SheafHom(K_n, E)
$$
。由于 $K$ 类似地嵌入有区别三角
$\bigoplus K_n \to \bigoplus K_n \to K$，只需证明
$\prod R\SheafHom(K_n, E) = R\SheafHom(\bigoplus K_n, E)$。
这是 (\ref{sites-cohomology-equation-internal-hom}) 以及导出张量积与直和可交换这一事实
的形式推论。
\end{proof}






\section{导出范畴中的可逆对象}
\label{sites-cohomology-section-invertible-D-or-R}

\noindent
我们刻画环化站点的导出范畴中的可逆对象（既包括局部环化拓扑斯的情形，
也包括一般情形）。

\begin{lemma}
\label{sites-cohomology-lemma-category-summands-finite-free}
设 $(\mathcal{C}, \mathcal{O})$ 是一个环化站点。令
$R = \Gamma(\mathcal{C}, \mathcal{O})$。作为有限自由
$\mathcal{O}$-模的直和项的那些 $\mathcal{O}$-模所成的范畴，
等价于有限投射 $R$-模所成的范畴。
\end{lemma}

\begin{proof}
注意，有限投射 $R$-模恰好就是有限自由 $R$-模的直和项。
该等价由函子 $\mathcal{E} \mapsto
\Gamma(\mathcal{C}, \mathcal{E})$ 给出。
逆函子由下述构造给出。考虑拓扑斯态射
$f : \Sh(\mathcal{C}) \to \Sh(\text{pt})$，其中 $f_*$
由取整体截面给出，而 $f^{-1}$ 把集合 $S$，即
$\Sh(\text{pt})$ 的对象，映到取值为 $S$ 的常值层。
利用恒等映射 $R \to f_*\mathcal{O}$，得到环化拓扑斯的态射
$(f, f^\sharp) : (\Sh(\mathcal{C}), \mathcal{O}) \to (\Sh(\text{pt}), R)$
。于是逆函子由 $f^*$ 给出。
\end{proof}

\begin{lemma}
\label{sites-cohomology-lemma-invertible-derived}
设 $(\mathcal{C}, \mathcal{O})$ 是一个环化站点。设 $M$ 是
$D(\mathcal{O})$ 的一个对象。下列条件等价：
\begin{enumerate}
\item $M$ 在 $D(\mathcal{O})$ 中可逆；见《范畴》定义
\ref{categories-definition-invertible}；
\item 存在一个局部有限的\footnote{这意味着，对 $\mathcal{C}$
的每个对象 $U$，都存在一个覆盖 $\{U_i \to U\}$，使得对每个 $i$，
层 $\mathcal{O}_n|_{U_i}$ 只对有限多个 $n$ 非零。}直积分解
$$
\mathcal{O} = \prod\nolimits_{n \in \mathbf{Z}} \mathcal{O}_n
$$
，并且对每个 $n$，存在可逆 $\mathcal{O}_n$-模
$\mathcal{H}^n$
（《站点上的模》定义 \ref{sites-modules-definition-invertible-sheaf}），
而且在 $D(\mathcal{O})$ 中有 $M = \bigoplus \mathcal{H}^n[-n]$。
\end{enumerate}
如果 (1) 和 (2) 成立，那么 $M$ 是 $D(\mathcal{O})$ 的完美对象。
如果 $(\mathcal{C}, \mathcal{O})$ 是一个局部环化站点，
这些条件还等价于：
\begin{enumerate}
\item[(3)] 对 $\mathcal{C}$ 的每个对象 $U$，存在一个覆盖
$\{U_i \to U\}$，并且对每个 $i$ 存在整数 $n_i$，使得
$M|_{U_i}$ 由置于次数 $n_i$ 上的可逆
$\mathcal{O}_{U_i}$-模表示。
\end{enumerate}
\end{lemma}

\begin{proof}
假设 (2) 成立。考虑对象 $R\SheafHom(M, \mathcal{O})$ 和复合映射
$$
R\SheafHom(M, \mathcal{O}) \otimes_\mathcal{O}^\mathbf{L} M \to \mathcal{O}
$$
。要证明它是同构，可以局部地工作。因此，可以假设
$\mathcal{O} = \prod_{a \leq n \leq b} \mathcal{O}_n$ 且
$M = \bigoplus_{a \leq n \leq b} \mathcal{H}^n[-n]$。
只需证明
$$
R\SheafHom(\mathcal{H}^m, \mathcal{O})
\otimes_\mathcal{O}^\mathbf{L} \mathcal{H}^n
$$
在 $n \not = m$ 时为零，在 $n = m$ 时等于 $\mathcal{O}_n$。
$n \not = m$ 的情形来自如下事实：$\mathcal{O}_n$ 和
$\mathcal{O}_m$ 都是平坦 $\mathcal{O}$-代数，并且
$\mathcal{O}_n \otimes_\mathcal{O} \mathcal{O}_m = 0$。
利用可逆 $\mathcal{O}$-模的局部结构（《站点上的模》引理
\ref{sites-modules-lemma-invertible}）并局部地工作，
$n = m$ 情形中的同构可直接得出；我们略去细节。
由于 $D(\mathcal{O})$ 是对称幺半范畴，故 $M$ 可逆。

\medskip\noindent
假设 (1) 成立。(2) 中的描述给出 $\mathcal{O}_n$ 的一个候选，即
$\SheafHom_\mathcal{O}(H^n(M), H^n(M))$。
如果这是一个局部有限的环层族，并且
$\mathcal{O} = \prod \mathcal{O}_n$，那么利用来自该直积分解的
$\mathcal{O}$ 中的幂等元，立即得到直和分解
$M = \bigoplus H^n(M)[-n]$。这表明，为证明 (2)，可以在如下意义下
局部地工作。设 $U$ 是 $\mathcal{C}$ 的一个对象。必须证明存在
$U$ 的一个覆盖 $\{U_i \to U\}$，使得取上述 $\mathcal{O}_n$ 时，
限制到 $\mathcal{C}/U_i$ 后，上述各断言以及 (2) 中的断言都成立。
因此，可以假设 $\mathcal{C}$ 有终对象 $X$；在 (2) 的证明过程中，
可以有限多次用 $X$ 的某个覆盖的成员替换 $X$。

\medskip\noindent
选取 $D(\mathcal{O})$ 的一个对象 $N$ 以及一个同构
$M \otimes_\mathcal{O}^\mathbf{L} N \cong \mathcal{O}$。
于是 $N$ 是幺半范畴 $D(\mathcal{O})$ 中 $M$ 的左对偶；
由引理 \ref{sites-cohomology-lemma-left-dual-derived}，$M$ 是完美的。
由对称性，$N$ 也是完美的。用一个覆盖的成员替换 $X$ 后，
可以假设 $M$ 和 $N$ 分别由严格完美复形
$\mathcal{E}^\bullet$ 和 $\mathcal{F}^\bullet$ 表示。
于是 $M \otimes_\mathcal{O}^\mathbf{L} N$ 由
$\text{Tot}(\mathcal{E}^\bullet \otimes_\mathcal{O} \mathcal{F}^\bullet)$
表示。用 $X$ 的一个覆盖的成员替换 $X$ 后，
可以假设互逆同构
$\mathcal{O} \to M \otimes_\mathcal{O}^\mathbf{L} N$ 和
$M \otimes_\mathcal{O}^\mathbf{L} N \to \mathcal{O}$ 由复形映射
$$
\alpha : \mathcal{O} \to
\text{Tot}(\mathcal{E}^\bullet \otimes_\mathcal{O} \mathcal{F}^\bullet)
\quad\text{和}\quad
\beta :
\text{Tot}(\mathcal{E}^\bullet \otimes_\mathcal{O} \mathcal{F}^\bullet)
\to \mathcal{O}
$$
给出；见引理 \ref{sites-cohomology-lemma-local-actual}。于是，作为复形映射有
$\beta \circ \alpha = 1$，而作为 $D(\mathcal{O})$ 中的态射有
$\alpha \circ \beta = 1$。用 $X$ 的一个覆盖的成员替换 $X$ 后，
可以假设复合 $\alpha \circ \beta$ 通过某个同伦 $\theta$
同伦于 $1$，其分量为
$$
\theta^n :
\text{Tot}^n(\mathcal{E}^\bullet \otimes_\mathcal{O} \mathcal{F}^\bullet)
\to
\text{Tot}^{n - 1}(
\mathcal{E}^\bullet \otimes_\mathcal{O} \mathcal{F}^\bullet)
$$
；这是由先前同一个引理得到的。令
$R = \Gamma(\mathcal{C}, \mathcal{O})$。由引理
\ref{sites-cohomology-lemma-category-summands-finite-free}，得到：
\begin{enumerate}
\item $M^\bullet = \Gamma(X, \mathcal{E}^\bullet)$
是有限投射 $R$-模的有界复形；
\item $N^\bullet = \Gamma(X, \mathcal{F}^\bullet)$
是有限投射 $R$-模的有界复形；
\item $\alpha$ 和 $\beta$ 分别对应于复形映射
$a : R \to \text{Tot}(M^\bullet \otimes_R N^\bullet)$ 和
$b : \text{Tot}(M^\bullet \otimes_R N^\bullet) \to R$；
\item $\theta^n$ 对应于映射
$h^n : \text{Tot}^n(M^\bullet \otimes_R N^\bullet) \to
\text{Tot}^{n - 1}(M^\bullet \otimes_R N^\bullet)$；以及
\item $b \circ a = 1$ 且 $b \circ a - 1 = dh + hd$。
\end{enumerate}
由此，$M^\bullet$ 和 $N^\bullet$ 定义 $D(R)$ 中互逆的对象。
由《代数详论》引理 \ref{more-algebra-lemma-invertible-derived}，
得到直积分解 $R = \prod_{a \leq n \leq b} R_n$ 以及可逆
$R_n$-模 $H^n$，使得
$M^\bullet \cong \bigoplus_{a \leq n \leq b} H^n[-n]$。
$D(R)$ 中的这个同构可以提升为复形态射
$$
\bigoplus H^n[-n] \longrightarrow M^\bullet
$$
，因为每个 $H^n$ 作为 $R$-模都是投射的。相应地，
再次使用引理 \ref{sites-cohomology-lemma-category-summands-finite-free}，得到态射
$$
\bigoplus H^n \otimes_R \mathcal{O}[-n] \to \mathcal{E}^\bullet
$$
，它在 $D(\mathcal{O})$ 中是同构。这里 $M \otimes_R \mathcal{O}$
表示引理 \ref{sites-cohomology-lemma-category-summands-finite-free} 的证明中构造的
从有限投射 $R$-模到 $\mathcal{O}$-模的函子。
令 $\mathcal{O}_n = R_n \otimes_R \mathcal{O}$，便得出 (2) 成立。

\medskip\noindent
如果 $(\mathcal{C}, \mathcal{O})$ 是一个局部环化站点，
那么给定对象 $U$ 和有限直积分解
$\mathcal{O}|_U = \prod_{a \leq n \leq b} \mathcal{O}_n|_U$，
可以找到覆盖 $\{U_i \to U\}$，使得对每个 $i$，至多有一个 $n$
满足 $\mathcal{O}_n|_{U_i}$ 非零。这很容易由《站点上的模》引理
\ref{sites-modules-lemma-locally-ringed} 的 (2) 以及《站点上的模》定义
\ref{sites-modules-definition-locally-ringed} 中给出的局部环化站点定义
得出。由此易见 (2) $\Rightarrow$ (3)。不对环化站点作任何假设时，
(3) $\Rightarrow$ (2) 也成立。我们略去细节。
\end{proof}










\section{投影公式}
\label{sites-cohomology-section-projection-formula}

\noindent
设 $f : (\Sh(\mathcal{C}), \mathcal{O}_\mathcal{C}) \to
(\Sh(\mathcal{D}), \mathcal{O}_\mathcal{D})$ 是环化拓扑斯的一个态射。
设 $E \in D(\mathcal{O}_\mathcal{C})$ 且
$K \in D(\mathcal{O}_\mathcal{D})$。不作任何进一步假设，存在映射
\begin{equation}
\label{sites-cohomology-equation-projection-formula-map}
Rf_*E \otimes^\mathbf{L}_{\mathcal{O}_\mathcal{D}} K
\longrightarrow
Rf_*(E \otimes^\mathbf{L}_{\mathcal{O}_\mathcal{C}} Lf^*K)
\end{equation}
具体而言，它是典范映射
$$
Lf^*(Rf_*E \otimes^\mathbf{L}_{\mathcal{O}_\mathcal{D}} K) =
Lf^*Rf_*E \otimes^\mathbf{L}_{\mathcal{O}_\mathcal{C}} Lf^*K
\longrightarrow
E \otimes^\mathbf{L}_{\mathcal{O}_\mathcal{C}} Lf^*K
$$
的伴随映射；该典范映射来自映射 $Lf^*Rf_*E \to E$ 以及引理
\ref{sites-cohomology-lemma-pullback-tensor-product} 和 \ref{sites-cohomology-lemma-adjoint}。
下面是一个颇为一般的投影公式版本。

\begin{lemma}
\label{sites-cohomology-lemma-projection-formula}
设 $f : (\Sh(\mathcal{C}), \mathcal{O}_\mathcal{C}) \to
(\Sh(\mathcal{D}), \mathcal{O}_\mathcal{D})$ 是环化拓扑斯的一个态射。
设 $E \in D(\mathcal{O}_\mathcal{C})$ 且
$K \in D(\mathcal{O}_\mathcal{D})$。如果 $K$ 是完美的，那么
$$
Rf_*E \otimes^\mathbf{L}_{\mathcal{O}_\mathcal{D}} K =
Rf_*(E \otimes^\mathbf{L}_{\mathcal{O}_\mathcal{C}} Lf^*K)
$$
在 $D(\mathcal{O}_\mathcal{D})$ 中成立。
\end{lemma}

\begin{proof}
要验证 (\ref{sites-cohomology-equation-projection-formula-map}) 是同构，可以在
$\mathcal{D}$ 上局部地工作。也就是说，对 $\mathcal{D}$ 的任意对象
$V$，必须找到一个覆盖 $\{V_j \to V\}$，使得该映射在 $V_j$
上的限制是同构。由完美对象的定义，这意味着可以假设 $K$
由一个严格完美的 $\mathcal{O}_\mathcal{D}$-模复形表示。
注意，完全一般地，该断言对 $K = K_1 \oplus K_2$ 成立，当且仅当
它对 $K_1$ 和 $K_2$ 都成立。因此，可以假设 $K$
是由有限自由 $\mathcal{O}_\mathcal{D}$-模组成的有限复形。
此时，一个涉及朴素截断的简单论证把该断言归约到 $K$
由有限自由 $\mathcal{O}_\mathcal{D}$-模表示的情形。
由于该断言在变量 $K$ 的有限直和项下不变，故只需对
$K = \mathcal{O}_\mathcal{D}[n]$ 证明它；在这种情形下它是平凡的。
\end{proof}

\begin{remark}
\label{sites-cohomology-remark-compatible-with-diagram}
映射 (\ref{sites-cohomology-equation-projection-formula-map}) 在下述意义下与注记
\ref{sites-cohomology-remark-base-change} 的基变换映射相容。具体而言，假设
$$
\xymatrix{
(\Sh(\mathcal{C}'), \mathcal{O}_{\mathcal{C}'})
\ar[r]_{g'} \ar[d]_{f'} &
(\Sh(\mathcal{C}), \mathcal{O}_\mathcal{C}) \ar[d]^f \\
(\Sh(\mathcal{D}'), \mathcal{O}_{\mathcal{D}'})
\ar[r]^g &
(\Sh(\mathcal{D}), \mathcal{O}_\mathcal{D})
}
$$
是环化拓扑斯的交换图。设 $E \in D(\mathcal{O}_\mathcal{C})$
且 $K \in D(\mathcal{O}_\mathcal{D})$。则图表
$$
\xymatrix{
Lg^*(Rf_*E \otimes^\mathbf{L}_{\mathcal{O}_\mathcal{D}} K)
\ar[r]_p \ar[d]_t &
Lg^*Rf_*(E \otimes^\mathbf{L}_{\mathcal{O}_\mathcal{C}} Lf^*K)
\ar[d]_b \\
Lg^*Rf_*E \otimes^\mathbf{L}_{\mathcal{O}_{\mathcal{D}'}}
Lg^*K \ar[d]_b &
Rf'_*L(g')^*(E \otimes^\mathbf{L}_{\mathcal{O}_\mathcal{C}}
Lf^*K) \ar[d]_t \\
Rf'_*L(g')^*E \otimes^\mathbf{L}_{\mathcal{O}_{\mathcal{D}'}}
Lg^*K \ar[rd]_p &
Rf'_*(L(g')^*E \otimes^\mathbf{L}_{\mathcal{O}_{\mathcal{D}'}}
L(g')^*Lf^*K) \ar[d]_c \\
& Rf'_*(L(g')^*E \otimes^\mathbf{L}_{\mathcal{O}_{\mathcal{D}'}}
L(f')^*Lg^*K)
}
$$
交换。这里，标记为 $t$ 的箭头由应用引理
\ref{sites-cohomology-lemma-pullback-tensor-product} 得到，标记为 $b$ 的箭头由应用注记
\ref{sites-cohomology-remark-base-change} 得到，标记为 $p$ 的箭头由应用
(\ref{sites-cohomology-equation-projection-formula-map}) 得到，而 $c$ 来自
$L(g')^* \circ Lf^* = L(f')^* \circ Lg^*$。我们略去验证。
\end{remark}






\section{弱可缩对象}
\label{sites-cohomology-section-w-contractible}

\noindent
站点的对象 $U$ 称为{\it 弱可缩的}，如果集合层的每个满射
$\mathcal{F} \to \mathcal{G}$ 都诱导满射
$\mathcal{F}(U) \to \mathcal{G}(U)$；见《站点》定义
\ref{sites-definition-w-contractible}。

\begin{lemma}
\label{sites-cohomology-lemma-w-contractible}
设 $\mathcal{C}$ 是一个站点。设 $U$ 是 $\mathcal{C}$
的一个弱可缩对象。则：
\begin{enumerate}
\item 函子 $\mathcal{F} \mapsto \mathcal{F}(U)$ 是正合函子
$\textit{Ab}(\mathcal{C}) \to \textit{Ab}$；
\item $H^p(U, \mathcal{F}) = 0$
对每个阿贝尔层 $\mathcal{F}$ 以及所有 $p \geq 1$ 成立；以及
\item 对任意群层 $\mathcal{G}$，每个 $\mathcal{G}$-挠子
在 $U$ 上都有截面。
\end{enumerate}
\end{lemma}

\begin{proof}
第一个断言立即由定义得出（另见《同调》第
\ref{homology-section-functors} 节）。由《导出范畴》引理
\ref{derived-lemma-right-derived-exact-functor}，高阶导出函子消失。
设 $\mathcal{F}$ 是一个 $\mathcal{G}$-挠子。则
$\mathcal{F} \to *$ 是层的满射。因此，(3) 也由定义得出。
\end{proof}

\noindent
在此列出具有足够多弱可缩对象的一些推论会很方便。

\begin{proposition}
\label{sites-cohomology-proposition-enough-weakly-contractibles}
设 $\mathcal{C}$ 是一个站点。设 $\mathcal{B} \subset \Ob(\mathcal{C})$，
使得每个 $U \in \mathcal{B}$ 都是弱可缩的，并且 $\mathcal{C}$
的每个对象都有一个由 $\mathcal{B}$ 的元素组成的覆盖。
设 $\mathcal{O}$ 是 $\mathcal{C}$ 上的一个环层。则：
\begin{enumerate}
\item $\mathcal{O}$-模复形
$\mathcal{F}_1 \to \mathcal{F}_2 \to \mathcal{F}_3$ 正合，
当且仅当对所有 $U \in \mathcal{B}$，
$\mathcal{F}_1(U) \to \mathcal{F}_2(U) \to \mathcal{F}_3(U)$
都正合。
\item $D(\mathcal{O})$ 的每个对象 $K$ 都是其典范截断的导出极限：
$K = R\lim \tau_{\geq -n} K$。
\item 给定逆系统
$\ldots \to \mathcal{F}_3 \to \mathcal{F}_2 \to \mathcal{F}_1$
，若其过渡映射均为满射，则投影
$\lim \mathcal{F}_n \to \mathcal{F}_1$ 是满射。
\item $\textit{Mod}(\mathcal{O})$ 中的乘积是正合的。
\item $D(\mathcal{O})$ 中的乘积可以通过对任意表示复形取乘积来计算。
\item 如果 $(\mathcal{F}_n)$ 是一个 $\mathcal{O}$-模逆系统，
那么对所有 $p > 1$ 有 $R^p\lim \mathcal{F}_n = 0$，并且
$$
R^1\lim \mathcal{F}_n  =
\Coker(\prod \mathcal{F}_n \to \prod \mathcal{F}_n)
$$
，其中映射为 $(x_n) \mapsto (x_n - f(x_{n + 1}))$。
\item 如果 $(K_n)$ 是 $D(\mathcal{O})$ 的对象的一个逆系统，
那么存在短正合列
$$
0 \to R^1\lim H^{p - 1}(K_n) \to H^p(R\lim K_n) \to
\lim H^p(K_n) \to 0
$$
\end{enumerate}
\end{proposition}

\begin{proof}
证明 (1)。如果该序列正合，那么由引理 \ref{sites-cohomology-lemma-w-contractible}，
在 $\mathcal{C}$ 的任意弱可缩元素上取值都会得到正合序列。
反之，假设对所有 $U \in \mathcal{B}$，
$\mathcal{F}_1(U) \to \mathcal{F}_2(U) \to \mathcal{F}_3(U)$
都正合。设 $V$ 是 $\mathcal{C}$ 的一个对象，并设
$s \in \mathcal{F}_2(V)$ 是 $\mathcal{F}_2 \to \mathcal{F}_3$
的核中的元素。由假设，存在一个满足 $U_i \in \mathcal{B}$
的覆盖 $\{U_i \to V\}$。于是 $s|_{U_i}$ 可以提升为截面
$s_i \in \mathcal{F}_1(U_i)$。因此，$s$ 是像层
$\Im(\mathcal{F}_1 \to \mathcal{F}_2)$ 的一个截面。
换言之，序列
$\mathcal{F}_1 \to \mathcal{F}_2 \to \mathcal{F}_3$
正合。

\medskip\noindent
证明 (2)。这由引理 \ref{sites-cohomology-lemma-is-limit-dimension} 中 $d = 0$
的情形成立。

\medskip\noindent
证明 (3)。设 $(\mathcal{F}_n)$ 是如 (2) 中的系统，并令
$\mathcal{F} = \lim \mathcal{F}_n$。如果 $U \in \mathcal{B}$，那么
$\mathcal{F}(U) = \lim \mathcal{F}_n(U)$
满射到 $\mathcal{F}_1(U)$，因为所有过渡映射
$\mathcal{F}_{n + 1}(U) \to \mathcal{F}_n(U)$ 都是满射。
由《站点》定义 \ref{sites-definition-sheaves-injective-surjective}
以及每个对象都有一个由 $\mathcal{B}$ 的元素组成的覆盖这一假设，
$\mathcal{F} \to \mathcal{F}_1$ 是满射。

\medskip\noindent
证明 (4)。设
$\mathcal{F}_{i, 1} \to \mathcal{F}_{i, 2} \to \mathcal{F}_{i, 3}$
是一个 $\mathcal{O}$-模正合列族。要证明
$\prod \mathcal{F}_{i, 1} \to \prod \mathcal{F}_{i, 2} \to
\prod \mathcal{F}_{i, 3}$ 正合。使用 (1) 的判据。
设 $U \in \mathcal{B}$。则
$$
(\prod \mathcal{F}_{i, 1})(U) \to
(\prod \mathcal{F}_{i, 2})(U) \to
(\prod \mathcal{F}_{i, 3})(U)
$$
与
$$
\prod \mathcal{F}_{i, 1}(U) \to
\prod \mathcal{F}_{i, 2}(U) \to
\prod \mathcal{F}_{i, 3}(U)
$$
相同。每个序列
$\mathcal{F}_{i, 1}(U) \to \mathcal{F}_{i, 2}(U) \to \mathcal{F}_{i, 3}(U)$
均由 (1) 正合。因此，由《同调》引理
\ref{homology-lemma-product-abelian-groups-exact}，所显示的序列正合。
再次由 (1) 得出结论。

\medskip\noindent
证明 (5)。这由 (4) 以及《导出范畴》引理
\ref{derived-lemma-products} 的稍加推广版本得出。

\medskip\noindent
证明 (6) 和 (7)。关于导出极限、同伦极限及其关系的讨论，
参见第 \ref{sites-cohomology-section-derived-limits} 节。由《导出范畴》定义
\ref{derived-definition-derived-limit}，有有区别三角
$$
R\lim K_n \to \prod K_n \to \prod K_n \to R\lim K_n[1]
$$
。取上同调层的长正合序列，得到
$$
H^{p - 1}(\prod K_n) \to H^{p - 1}(\prod K_n) \to
H^p(R\lim K_n) \to H^p(\prod K_n) \to H^p(\prod K_n)
$$
。由于乘积由 (4) 正合，这变为
$$
\prod H^{p - 1}(K_n) \to \prod H^{p - 1}(K_n) \to
H^p(R\lim K_n) \to \prod H^p(K_n) \to \prod H^p(K_n)
$$
。先把它应用于 $K_n = \mathcal{F}_n[0]$ 的情形，
其中 $(\mathcal{F}_n)$ 如 (6) 中所述，便得出 (6) 成立。
再把它应用于 (7) 中的 $(K_n)$，便得出 (7) 成立。
\end{proof}






\section{紧对象}
\label{sites-cohomology-section-compact}

\noindent
本节研究环化站点上的模的导出范畴中的紧对象。回忆紧对象定义于
《导出范畴》定义 \ref{derived-definition-compact-object}。

\begin{lemma}
\label{sites-cohomology-lemma-compact-in-terms-of-generators}
设 $\mathcal{A}$ 是一个 Grothendieck 阿贝尔范畴。设
$S \subset \Ob(\mathcal{A})$ 是一个对象集，满足：
\begin{enumerate}
\item $\mathcal{A}$ 的任意对象都是 $S$ 中元素的某个直和的商；以及
\item 对任意 $E \in S$，函子 $\Hom_\mathcal{A}(E, -)$ 与直和可交换。
\end{enumerate}
那么，$D(\mathcal{A})$ 的每个紧对象在 $D(\mathcal{A})$ 中都是
一个有限复形的直和项，而该复形的每一项都是 $S$ 中元素的有限直和。
\end{lemma}

\begin{proof}
假设 $K \in D(\mathcal{A})$ 是一个紧对象。用复形 $K^\bullet$
表示 $K$，并考虑映射
$$
K^\bullet
\longrightarrow
\bigoplus\nolimits_{n \geq 0} \tau_{\geq n} K^\bullet
$$
，其中使用了典范截断；见《同调》第
\ref{homology-section-truncations} 节。这是有意义的，因为右边的直和
在每个次数上都是有限的。由假设，该映射分解经过一个有限直和。
因此，至少对一个 $n$，映射 $K \to \tau_{\geq n} K$ 为零，
即 $K$ 属于 $D^{-}(R)$。

\medskip\noindent
可以用一个上有界复形 $K^\bullet$ 表示 $K$，其每一项都是 $S$
中对象的直和；见《导出范畴》引理
\ref{derived-lemma-subcategory-left-resolution}。注意有
$$
K^\bullet = \bigcup\nolimits_{n \leq 0} \sigma_{\geq n}K^\bullet
$$
，其中使用了朴素截断；见《同调》第
\ref{homology-section-truncations} 节。因此，由《导出范畴》引理
\ref{derived-lemma-colim-hocolim} 和
\ref{derived-lemma-commutes-with-countable-sums}，可见在 $D(R)$ 中，
$1 : K^\bullet \to K^\bullet$ 分解经过
$\sigma_{\geq n}K^\bullet \to K^\bullet$。因此，
$1 : K^\bullet \to K^\bullet$ 在 $D(\mathcal{A})$ 中分解为
$$
K^\bullet \xrightarrow{\varphi} L^\bullet \xrightarrow{\psi} K^\bullet
$$
，其中复形 $L^\bullet$ 有界，且各项都是 $S$ 中元素的直和。
设 $L^i$ 对 $i \not \in [a, b]$ 为零。令 $c$ 为满足如下条件的
最大整数 $\leq b + 1$：对 $i < c$，$L^i$ 是 $S$
中元素的有限直和。断言：如果 $c < b + 1$，那么可以修改
$L^\bullet$ 以增大 $c$。对此断言作归纳，即可表明 $1_K$ 有分解
$$
K \xrightarrow{\varphi} L \xrightarrow{\psi} K
$$
，其中在 $D(\mathcal{A})$ 中，$L$ 可以由一个有限复形表示，
且其各项都是 $S$ 中元素的有限直和。注意
$e = \varphi \circ \psi \in \text{End}_{D(\mathcal{A})}(L)$
是一个幂等元。由《导出范畴》引理
\ref{derived-lemma-projectors-have-images-triangulated}，有
$L = \Ker(e) \oplus \Ker(1 - e)$。映射
$\varphi : K \to L$ 在 $D(R)$ 中诱导与 $\Ker(1 - e)$
的同构，结论得证。

\medskip\noindent
证明该断言。写成 $L^c = \bigoplus_{\lambda \in \Lambda} E_\lambda$。
由于 $L^{c - 1}$ 是 $S$ 中元素的有限直和，由假设 (2)，
可以找到有限子集 $\Lambda' \subset \Lambda$，使得
$L^{c - 1} \to L^c$ 分解经过
$\bigoplus_{\lambda \in \Lambda'} E_\lambda \subset L^c$。
考虑复形映射
$$
\pi :
L^\bullet
\longrightarrow
(\bigoplus\nolimits_{\lambda \in \Lambda \setminus \Lambda'} E_\lambda)[-i]
$$
，它在次数 $i$ 上由到与 $\Lambda \setminus \Lambda'$ 对应的因子的
投影给出。由关于 $K$ 的假设可见，必要时用一个更大的有限子集替换
$\Lambda'$ 后，可以假设 $D(\mathcal{A})$ 中
$\pi \circ \varphi = 0$。设
$(L')^\bullet \subset L^\bullet$ 为 $\pi$ 的核。
由于 $\pi$ 是满射，得到一个复形的短正合列，它在
$D(\mathcal{A})$ 中给出一个有区别三角（见《导出范畴》引理
\ref{derived-lemma-derived-canonical-delta-functor}）。
由于 $\Hom_{D(\mathcal{A})}(K, -)$ 是同调函子（见《导出范畴》引理
\ref{derived-lemma-representable-homological}）且
$\pi \circ \varphi = 0$，可以找到 $D(\mathcal{A})$ 中的态射
$\varphi' : K^\bullet \to (L')^\bullet$，使得它与
$(L')^\bullet \to L^\bullet$ 的复合等于 $\varphi$。
令 $\psi'$ 等于 $\psi$ 与 $(L')^\bullet \to L^\bullet$ 的复合，
便得到一个新的分解。由于 $(L')^\bullet$ 除次数 $c$ 外与
$L^\bullet$ 相同，并且
$(L')^c = \bigoplus_{\lambda \in \Lambda'} E_\lambda$，断言得证。
\end{proof}

\begin{lemma}
\label{sites-cohomology-lemma-compact-objects-if-enough-qc}
设 $(\mathcal{C}, \mathcal{O})$ 是一个环化站点。假设
$\mathcal{C}$ 的每个对象都有一个由拟紧对象组成的覆盖。
那么 $D(\mathcal{O})$ 的每个紧对象在 $D(\mathcal{O})$ 中都是
某个有限复形的直和项；该复形的各项都是形如
$j_!\mathcal{O}_U$ 的 $\mathcal{O}$-模的有限直和，
其中 $U$ 是 $\mathcal{C}$ 的一个拟紧对象。
\end{lemma}

\begin{proof}
应用引理 \ref{sites-cohomology-lemma-compact-in-terms-of-generators}，其中
$S \subset \Ob(\textit{Mod}(\mathcal{O}))$ 是所有形如
$j_!\mathcal{O}_U$ 的模所成的集合，且
$U \in \Ob(\mathcal{C})$ 拟紧。由《站点上的模》引理
\ref{sites-modules-lemma-module-quotient-flat} 以及每个 $U$
都可以被拟紧对象覆盖这一假设，假设 (1) 成立。假设 (2) 来自
$$
\Hom_\mathcal{O}(j_!\mathcal{O}_U, \mathcal{F}) = \mathcal{F}(U)
$$
，它由《站点》引理 \ref{sites-lemma-directed-colimits-sections}
与直和可交换。
\end{proof}

\noindent
在上述引理的情形中，模 $j_!\mathcal{O}_U$ 不一定是
$D(\mathcal{O})$ 的紧对象（即使 $U$ 是 $\mathcal{C}$
的拟紧对象也是如此）。下面两个引理处理这一问题。


\begin{lemma}
\label{sites-cohomology-lemma-when-jshriek-lower-compact}
设 $(\mathcal{C}, \mathcal{O})$ 是一个环化站点。设 $U$ 是
$\mathcal{C}$ 的一个对象。假设函子
$\mathcal{F} \mapsto H^p(U, \mathcal{F})$ 与直和可交换。
那么 $\mathcal{O}$-模 $j_!\mathcal{O}_U$ 在下述意义下是
$D^+(\mathcal{O})$ 的紧对象：如果 $D(\mathcal{O})$ 中的
$M = \bigoplus_{i \in I} M_i$ 下有界，那么
$\Hom(j_{U!}\mathcal{O}_U, M) =
\bigoplus_{i \in I} \Hom(j_{U!}\mathcal{O}_U, M_i)$。
\end{lemma}

\begin{proof}
由《站点上的模》公式
(\ref{sites-modules-equation-map-lower-shriek-OU-into-module})，
$\Hom(j_{U!}\mathcal{O}_U, -)$ 与函子
$\mathcal{F} \mapsto \mathcal{F}(U)$ 相同，因此只需证明
$H^p(U, M) = \bigoplus H^p(U, M_i)$。
设 $\mathcal{I}_i$（$i \in I$）是一个内射
$\mathcal{O}$-模族。由假设，有
$$
H^p(U, \bigoplus\nolimits_{i \in I} \mathcal{I}_i) =
\bigoplus\nolimits_{i \in I} H^p(U, \mathcal{I}_i) = 0
$$
对所有 $p$ 成立。由于 $M = \bigoplus M_i$ 下有界，
可见存在 $a \in \mathbf{Z}$，使得对 $n < a$ 有
$H^n(M_i) = 0$。因此，可以选取表示 $M_i$ 的内射
$\mathcal{O}$-模复形 $\mathcal{I}_i^\bullet$，使得对 $n < a$ 有
$\mathcal{I}_i^n = 0$；见《导出范畴》引理
\ref{derived-lemma-injective-resolutions-exist}。
由《内射对象》引理 \ref{injectives-lemma-derived-products}，
直和复形 $\bigoplus \mathcal{I}_i^\bullet$ 表示 $M$。
由 Leray 非循环性（《导出范畴》引理
\ref{derived-lemma-leray-acyclicity}），有
$$
R\Gamma(U, M) = \Gamma(U, \bigoplus \mathcal{I}_i^\bullet) =
\bigoplus \Gamma(U, \bigoplus \mathcal{I}_i^\bullet) =
\bigoplus R\Gamma(U, M_i)
$$
，正如所需。
\end{proof}

\begin{lemma}
\label{sites-cohomology-lemma-when-jshriek-lower-compact-worked-out}
设 $(\mathcal{C}, \mathcal{O})$ 是一个环化站点，
其覆盖集为 $\text{Cov}_\mathcal{C}$。设
$\mathcal{B} \subset \Ob(\mathcal{C})$ 且
$\text{Cov} \subset \text{Cov}_\mathcal{C}$ 为子集。假设：
\begin{enumerate}
\item 对每个 $\mathcal{U} \in \text{Cov}$，有
$\mathcal{U} = \{U_i \to U\}_{i \in I}$，其中 $I$ 有限，
$U, U_i \in \mathcal{B}$，并且每个
$U_{i_0} \times_U \ldots \times_U U_{i_p} \in \mathcal{B}$；
\item 对每个 $U \in \mathcal{B}$，$\text{Cov}$ 中出现的
$U$ 的覆盖构成 $U$ 的覆盖的一个共尾系统。
\end{enumerate}
那么，对 $U \in \mathcal{B}$，对象 $j_{U!}\mathcal{O}_U$
在下述意义下是 $D^+(\mathcal{O})$ 的紧对象：如果
$D(\mathcal{O})$ 中的 $M = \bigoplus_{i \in I} M_i$ 下有界，那么
$\Hom(j_{U!}\mathcal{O}_U, M) =
\bigoplus_{i \in I} \Hom(j_{U!}\mathcal{O}_U, M_i)$。
\end{lemma}

\begin{proof}
这由引理 \ref{sites-cohomology-lemma-when-jshriek-lower-compact} 和引理
\ref{sites-cohomology-lemma-colim-works-over-collection} 得出。
\end{proof}

\begin{lemma}
\label{sites-cohomology-lemma-when-jshriek-compact}
设 $(\mathcal{C}, \mathcal{O})$ 是一个环化站点。设 $U$ 是
$\mathcal{C}$ 的一个对象。如果存在整数 $d$ 使得：
\begin{enumerate}
\item 对所有 $p > d$，有 $H^p(U, \mathcal{F}) = 0$；以及
\item 函子 $\mathcal{F} \mapsto H^p(U, \mathcal{F})$ 与直和可交换，
\end{enumerate}
那么 $\mathcal{O}$-模 $j_!\mathcal{O}_U$ 是
$D(\mathcal{O})$ 的紧对象。
\end{lemma}

\begin{proof}
假设 (1) 和 (2)。回忆，对 $D(\mathcal{O})$ 中的 $K$，有
$\Hom(j_!\mathcal{O}_U, K) = R\Gamma(U, K)$。因此必须证明
$R\Gamma(U, -)$ 与直和可交换。第一个假设意味着函子
$F = H^0(U, -)$ 具有有限上同调维数。此外，第二个假设蕴含
内射模的任意直和对 $F$ 都是非循环的。设 $K_i$ 是
$D(\mathcal{O})$ 的一个对象族。选取表示 $K_i$ 且各项内射的
K-内射代表元 $I_i^\bullet$；见《内射对象》定理
\ref{injectives-theorem-K-injective-embedding-grothendieck}。
由于可以把 $F$ 应用于任意由非循环对象组成的复形来计算 $RF$
（《导出范畴》引理 \ref{derived-lemma-unbounded-right-derived}），
并且 $\bigoplus K_i$ 由 $\bigoplus I_i^\bullet$ 表示
（《内射对象》引理 \ref{injectives-lemma-derived-products}），
故 $R\Gamma(U, \bigoplus K_i)$ 由
$\bigoplus H^0(U, I_i^\bullet)$ 表示。因此，$R\Gamma(U, -)$
如所需地与直和可交换。
\end{proof}

\begin{lemma}
\label{sites-cohomology-lemma-quasi-compact-weakly-contractible-compact}
设 $(\mathcal{C}, \mathcal{O})$ 是一个环化站点。设 $U$
是 $\mathcal{C}$ 的一个拟紧且弱可缩的对象。那么
$j_!\mathcal{O}_U$ 是 $D(\mathcal{O})$ 的紧对象。
\end{lemma}

\begin{proof}
把引理 \ref{sites-cohomology-lemma-when-jshriek-compact} 和
\ref{sites-cohomology-lemma-w-contractible} 与《站点上的模》引理
\ref{sites-modules-lemma-sections-over-quasi-compact} 结合即可。
\end{proof}

\begin{lemma}
\label{sites-cohomology-lemma-perfect-is-compact}
设 $(\mathcal{C}, \mathcal{O})$ 是一个环化站点。假设
$\mathcal{C}$ 具有下列性质：
\begin{enumerate}
\item $\mathcal{C}$ 有一个拟紧终对象 $X$；
\item $\mathcal{C}$ 的每个拟紧对象都有一个由有限覆盖组成的
共尾覆盖系，而且这些覆盖均由拟紧对象组成；
\item 对于 $U$ 和 $U_i$ 均拟紧的有限覆盖
$\{U_i \to U\}_{i \in I}$，纤维积 $U_i \times_U U_j$ 均拟紧。
\end{enumerate}
设 $K$ 是 $D(\mathcal{O})$ 的一个完美对象。则：
\begin{enumerate}
\item[(a)] $K$ 在下述意义下是 $D^+(\mathcal{O})$ 的紧对象：
如果 $M = \bigoplus_{i \in I} M_i$ 下有界，那么
$\Hom(K, M) = \bigoplus_{i \in I} \Hom(K, M_i)$。
\item[(b)] 如果 $(\mathcal{C}, \mathcal{O})$ 具有有限上同调维数，
也就是说，存在 $d$，使得对任意 $\mathcal{O}$-模 $\mathcal{F}$
以及 $i > d$ 都有 $H^i(X, \mathcal{F}) = 0$，那么
$K$ 是 $D(\mathcal{O})$ 的紧对象。
\end{enumerate}
\end{lemma}

\begin{proof}
设 $K^\vee$ 为 $K$ 的对偶；见引理
\ref{sites-cohomology-lemma-dual-perfect-complex}。则有
$$
\Hom_{D(\mathcal{O})}(K, M) =
H^0(X, K^\vee \otimes_\mathcal{O}^\mathbf{L} M)
$$
，并且这关于 $D(\mathcal{O})$ 中的 $M$ 是函子性的。
由于 $K^\vee \otimes_\mathcal{O}^\mathbf{L} -$ 与直和可交换，
只需证明 $R\Gamma(X, -)$ 与有关的直和可交换。

\medskip\noindent
证明 (a)。按上述方式重新表述后，这是引理
\ref{sites-cohomology-lemma-when-jshriek-lower-compact-worked-out} 中 $U = X$
的一个特殊情形。

\medskip\noindent
证明 (b)。由于 $R\Gamma(X, K) = R\Hom(\mathcal{O}, K)$，
并且由引理 \ref{sites-cohomology-lemma-colim-works-over-collection}，
$H^p(X, -)$ 与直和可交换，这是引理
\ref{sites-cohomology-lemma-when-jshriek-compact} 的一个特殊情形。
\end{proof}






\section{上同调层局部常值的复形}
\label{sites-cohomology-section-locally-constant}

\noindent
局部常值层在《站点上的模》第
\ref{sites-modules-section-locally-constant} 节中引入。
设 $\mathcal{C}$ 是一个站点，$\Lambda$ 是一个环。
用 $D(\mathcal{C}, \Lambda)$ 表示 $\mathcal{C}$ 上
$\underline{\Lambda}$-模的阿贝尔范畴的导出范畴。

\begin{lemma}
\label{sites-cohomology-lemma-locally-constant}
设 $\mathcal{C}$ 是一个具有终对象 $X$ 的站点。
设 $\Lambda$ 是一个诺特环。设 $K \in D^b(\mathcal{C}, \Lambda)$，
且 $H^i(K)$ 是有限型 $\Lambda$-模的局部常值层。
那么存在覆盖 $\{U_i \to X\}$，使得每个 $K|_{U_i}$
都由一个有限型 $\Lambda$-模的局部常值层复形表示。
\end{lemma}

\begin{proof}
选取 $a \leq b$，使得对 $i \not \in [a, b]$ 有 $H^i(K) = 0$。
对 $b - a$ 作归纳，证明存在覆盖 $\{U_i \to X\}$，使得
$K|_{U_i}$ 可以由复形 $\underline{M^\bullet}_{U_i}$ 表示，
其中 $M^p$ 是有限型 $\Lambda$-模，且对
$p \not \in [a, b]$ 有 $M^p = 0$。当 $b = a$ 时，这是显然的。
一般地，可以用一个覆盖的成员替换 $X$，并假设 $H^b(K)$
为常值层，记为 $H^b(K) = \underline{M}$。由《站点上的模》引理
\ref{sites-modules-lemma-locally-constant-finite-type}，
模 $M$ 是有限 $\Lambda$-模。选取由 $M$ 的生成元
$x_1, \ldots, x_r$ 给出的满射 $\Lambda^{\oplus r} \to M$。

\medskip\noindent
对引理 \ref{sites-cohomology-lemma-kill-cohomology-class-on-covering} 稍作推广
（细节略去），可得一个覆盖 $\{U_i \to X\}$，使得
$x_i \in H^0(X, H^b(K))$ 可以提升为 $H^b(U_i, K)$ 的一个元素。
因此，用 $U_i$ 替换 $X$ 后，便得到存在映射
$\underline{\Lambda^{\oplus r}}[-b] \to K$
，它在次数 $b$ 的上同调层上诱导满射。选取一个有区别三角
$$
\underline{\Lambda^{\oplus r}}[-b] \to K \to L \to
\underline{\Lambda^{\oplus r}}[-b + 1]
$$
。此时，$L$ 的上同调层只在区间 $[a, b - 1]$ 中非零，
并在区间 $[a, b - 2]$ 中与 $K$ 的上同调层相同；在次数
$b - 1$ 上有短正合列
$$
0 \to H^{b - 1}(K) \to H^{b - 1}(L) \to
\underline{\Ker(\Lambda^{\oplus r} \to M)} \to 0
$$
。由《站点上的模》引理
\ref{sites-modules-lemma-kernel-finite-locally-constant}，
$H^{b - 1}(L)$ 是有限型的局部常值层。由归纳假设，得到
$D(\mathcal{C}, \underline{\Lambda})$ 中的同构
$\underline{M^\bullet} \to L$，其中 $M^p$
是有限 $\Lambda$-模，且对 $p \not \in [a, b - 1]$ 有 $M^p = 0$。
映射 $L \to \Lambda^{\oplus r}[-b + 1]$ 给出映射
$\underline{M^{b - 1}} \to \underline{\Lambda^{\oplus r}}$，
它局部为常值（《站点上的模》引理
\ref{sites-modules-lemma-morphism-locally-constant}）。
因此，可以假设它由映射 $M^{b - 1} \to \Lambda^{\oplus r}$ 给出。
该有区别三角表明复合
$M^{b - 2} \to M^{b - 1} \to \Lambda^{\oplus r}$ 为零，
而三角范畴的公理产生同构
$$
\underline{M^a \to \ldots \to M^{b - 1} \to \Lambda^{\oplus r}}
\longrightarrow K
$$
，它位于 $D(\mathcal{C}, \Lambda)$ 中。
\end{proof}

\noindent
设 $\mathcal{C}$ 是一个站点，$\Lambda$ 是一个环。利用态射
$\Sh(\mathcal{C}) \to \Sh(pt)$，可见存在函子
$D(\Lambda) \to D(\mathcal{C}, \Lambda)$，
$K \mapsto \underline{K}$。

\begin{lemma}
\label{sites-cohomology-lemma-map-out-of-locally-constant}
设 $\mathcal{C}$ 是一个具有终对象 $X$ 的站点。设 $\Lambda$ 是一个环。
给定：
\begin{enumerate}
\item $D(\Lambda)$ 的一个完美对象 $K$；
\item 一个表示 $K$ 的、由有限投射 $\Lambda$-模组成的有限复形
$K^\bullet$；
\item 一个 $\Lambda$-模层复形 $\mathcal{L}^\bullet$；以及
\item $D(\mathcal{C}, \Lambda)$ 中的一个映射
$\varphi : \underline{K} \to \mathcal{L}^\bullet$。
\end{enumerate}
那么存在覆盖 $\{U_i \to X\}$ 和复形映射
$\alpha_i : \underline{K}^\bullet|_{U_i} \to \mathcal{L}^\bullet|_{U_i}$
，后者表示 $\varphi|_{U_i}$。
\end{lemma}

\begin{proof}
立即由引理 \ref{sites-cohomology-lemma-local-actual} 得出。
\end{proof}

\begin{lemma}
\label{sites-cohomology-lemma-locally-constant-map}
设 $\mathcal{C}$ 是一个具有终对象 $X$ 的站点。
设 $\Lambda$ 是一个环。设 $K, L$ 是 $D(\Lambda)$ 的对象，
且 $K$ 完美。设 $\varphi : \underline{K} \to \underline{L}$
是 $D(\mathcal{C}, \Lambda)$ 中的映射。存在覆盖
$\{U_i \to X\}$，使得对某个 $D(\Lambda)$ 中的映射
$\alpha_i : K \to L$，有
$\varphi|_{U_i} = \underline{\alpha_i}$。
\end{lemma}

\begin{proof}
这由引理 \ref{sites-cohomology-lemma-map-out-of-locally-constant} 和《站点上的模》引理
\ref{sites-modules-lemma-morphism-locally-constant} 得出。
\end{proof}

\begin{lemma}
\label{sites-cohomology-lemma-locally-constant-tensor-product}
设 $\mathcal{C}$ 是一个站点。设 $\Lambda$ 是一个诺特环。
设 $K, L \in D^-(\mathcal{C}, \Lambda)$。如果 $K$ 和 $L$
的上同调层都是有限型 $\Lambda$-模的局部常值层，那么
$K \otimes_\Lambda^\mathbf{L} L$
的上同调层也是有限型 $\Lambda$-模的局部常值层。
\end{lemma}

\begin{proof}
把这作为引理 \ref{sites-cohomology-lemma-locally-constant} 的一个应用来证明。
注意，$H^i(K \otimes_\Lambda^\mathbf{L} L)$ 与
$H^i(\tau_{\geq i - 1}K \otimes_\Lambda^\mathbf{L} \tau_{\geq i - 1}L)$
相同。因此，可以假设 $K$ 和 $L$ 有界。由引理
\ref{sites-cohomology-lemma-locally-constant}，可以假设 $K$ 和 $L$
由有限型 $\Lambda$-模的局部常值层复形表示。
然后，可以用有限自由 $\Lambda$-模的上有界复形替换这些复形。
在这种情形下，结论是显然的。
\end{proof}

\begin{lemma}
\label{sites-cohomology-lemma-locally-constant-bounded}
设 $\mathcal{C}$ 是一个站点。设 $\Lambda$ 是一个诺特环。
设 $I \subset \Lambda$ 是一个理想。设
$K \in D^-(\mathcal{C}, \Lambda)$。如果
$K \otimes_\Lambda^\mathbf{L} \underline{\Lambda/I}$ 的上同调层
都是有限型 $\Lambda/I$-模的局部常值层，那么
$K \otimes_\Lambda^\mathbf{L} \underline{\Lambda/I^n}$
的上同调层对所有 $n \geq 1$ 都是有限型
$\Lambda/I^n$-模的局部常值层。
\end{lemma}

\begin{proof}
回忆有限型 $\Lambda$-模的局部常值层构成所有
$\underline{\Lambda}$-模所成范畴的一个弱 Serre 子范畴；见《站点上的模》
引理 \ref{sites-modules-lemma-kernel-finite-locally-constant}。
因此，由上同调层均为有限型 $\Lambda$-模的局部常值层的那些复形
组成的 $D(\mathcal{C}, \Lambda)$ 的子范畴，构成
$D(\mathcal{C}, \Lambda)$ 的一个严格满、饱和三角子范畴；
见《导出范畴》引理
\ref{derived-lemma-cohomology-in-serre-subcategory}。
接着，考虑有区别三角
$$
K \otimes_\Lambda^\mathbf{L} \underline{I^n/I^{n + 1}} \to
K \otimes_\Lambda^\mathbf{L} \underline{\Lambda/I^{n + 1}} \to
K \otimes_\Lambda^\mathbf{L} \underline{\Lambda/I^n} \to
K \otimes_\Lambda^\mathbf{L} \underline{I^n/I^{n + 1}}[1]
$$
以及同构
$$
K \otimes_\Lambda^\mathbf{L} \underline{I^n/I^{n + 1}}
=
\left(K \otimes_\Lambda^\mathbf{L} \underline{\Lambda/I}\right)
\otimes_{\Lambda/I}^\mathbf{L} \underline{I^n/I^{n + 1}}
$$
结合引理 \ref{sites-cohomology-lemma-locally-constant-tensor-product}，便得到结论。
\end{proof}
